\subsection{The sign object, scalar compression, finite levels and descent}


\begin{statement}{A.1}{Proved}{The relative cohomology object of the two signs}
Assume the two normalized frame characters of supplied R4,
$a_3=\epsilon_3T$ and $a_2=\epsilon_2W:\HH\to\Ftwo$.
Let $I=[0,1]$, with its endpoints labelled $3,2$, and put
\[
 O=H^1(I,\partial I;\Ftwo)
   =\Ftwo^{\{3,2\}}/\Ftwo(1,1).
\]
The continuous map
$o_{\HH}(h)=[(a_3(h),a_2(h))]$ is a homomorphism.
It is the obstruction to extending the specified endpoint signs to a
locally constant sign on $I$. This definition does not assume that it
vanishes, and the R2 element has nonzero obstruction.
\end{statement}
\begin{proof}
The long exact sequence of the pair, or the cellular complex of an
interval, gives
\[
0\longrightarrow\Ftwo\xrightarrow{t\mapsto(t,t)}
\Ftwo^2\xrightarrow{(u,v)\mapsto u+v}\Ftwo\longrightarrow0.
\]
A locally constant function on the connected interval is constant;
its boundary values therefore agree. Conversely equal values extend
uniquely as a function, after the values have been specified. The
connecting homomorphism is their class modulo the diagonal. Composing
with the continuous homomorphism $(a_3,a_2)$ proves the assertions.
The interval is a precise algebraic model for the two-sign gluing
problem, not an asserted geometric path in a moduli stack.
\end{proof}

\begin{statement}{A.2}{Proved}{Descent to every multiplier and the cancelling sign}
Write $G=\GZ$, $\Gamma=G_{\Q}$, and $K=\Gamma\cap\HH$.
For $g\in G$, choose $\sigma\in\Gamma$ with
$\chi(\sigma)=\lambda(g)$ and define
\[
 o(g)=[(a_3(g\sigma^{-1}),a_2(g\sigma^{-1}))]\in O.
\]
This is well defined, continuous and a homomorphism. Identifying $O$
with $\{\pm1\}$ by $[(u,v)]\mapsto(-1)^{u+v}$, its inverse is the
unique sign that cancels it. Its additive exponent is exactly R4's
$\Delta(g)$.
\end{statement}
\begin{proof}
The cyclotomic character is onto, so a choice exists. Replacing
$\sigma$ by $\sigma'$ changes $h=g\sigma^{-1}$ to
$h(\sigma\sigma'^{-1})$, where the second factor belongs to $K$.
R4 and the common-field sign comparison say that the two characters
agree on $K$; the change is therefore diagonal and dies in $O$.
For $g_i=h_i\sigma_i$ the product is
$h_1(\sigma_1h_2\sigma_1^{-1})\sigma_1\sigma_2$.
The conjugation invariance of $a_3+a_2$ established in R4 gives the
homomorphism law, in this order of composition; reversing the order
gives the same additive value. The map $\HH\times\Gamma\to G$ is a
continuous compact surjection. The continuous function just defined
is constant on its fibres and hence descends continuously.
Finally $\{\pm1\}$ has exponent two, so the inverse sign is the same
sign. The definition agrees with
$\Delta(g)=a_3(g\sigma^{-1})+a_2(g\sigma^{-1})$, without a claim
that either sign is $+1$. The cancellation identity
$(-1)^{\Delta(g)}o(g)=1$ is a definition; the assertion $o(g)=1$
on $\GT$ is an additional theorem still required.
\end{proof}

\begin{statement}{A.3}{Refuted}{A vanishing class does not select a gauge-invariant torsor trivialization}
The gauge-invariant two-sign datum does not canonically select one
of its two lifts, even when $\Delta=0$. On the normalized subgroup
$\HH$, the ordered pair of signs is a finer continuous invariant.
It has no extension to an abelian-valued continuous homomorphism on
$G$ retaining both components.
\end{statement}
\begin{proof}
For $\omega\in O$ its inverse image
$P_\omega\subset\Ftwo^2$ is a torsor under the diagonal $\Ftwo$.
The simultaneous change $(u,v)\mapsto(u+1,v+1)$ preserves the
class and interchanges its two lifts. A choice natural under that
gauge transformation would be a fixed point, which does not exist.
In particular $P_0=\{(0,0),(1,1)\}$ has two trivializations as a
torsor. Declaring $(0,0)$ preferred uses a chosen signed coordinate;
it is not invariant under the allowed overall-sign change.
An ordinary torsor over a point is always trivializable; calling
$P_1$ a nontrivial torsor over a point would be incorrect. The
obstruction is the class in the quotient, or equivalently the
equivariant sign representation with character $\Delta$.

There exists $k\in K$ with $(a_3(k),a_2(k))=(1,1)$: the extension
$\Q^{\rm ab}((3/4)^{1/4})/\Q^{\rm ab}$ is quadratic, as shown in
A.12 below, and a nontrivial automorphism extends to an element of
$K$. If either local component extended to a character on $G$, its
restriction to $\Gamma$ would be an abelian character nonzero on
$K=\overline{[\Gamma,\Gamma]}$, contrary to Kronecker--Weber.
Thus the common bit is meaningful with the multiplier-one
normalization but cannot be promoted in the suggested way to all
multipliers. The R2 element has lifts $(1,0),(0,1)$ and class one.
For an involution $(-1,f)$, R7 gives only
$\Delta((-1,f))=\delta((1,f))$; no value for an unclassified
involution follows. Conjugates of complex conjugation have value
zero.
\end{proof}

\begin{statement}{A.4}{Proved}{The phase generalization and its limit}
For a finite abelian phase group $M$, the quotient
$M^2/\operatorname{diag}M$ is canonically $M$ by
$(u,v)\mapsto uv^{-1}$. A pair of continuous characters on $\HH$
therefore defines the same gluing obstruction with values in $M$.
It descends to $G=\HH\Gamma$ if its ratio is conjugation invariant
and trivial on $K$. The supplied data do not specify a higher phase
character or a degree-two gerbe class.
\end{statement}
\begin{proof}
The kernel and surjectivity of the ratio map give the quotient.
The descent proof is the proof of A.2 with multiplication in $M$.
An equivariant $M$-torsor is governed by degree-one cohomology;
a gerbe would require a degree-two cocycle or a specified central
extension. Neither follows from a pair of signs. For roots of unity
with nontrivial cyclotomic action the coefficient module must also
be specified; ordinary characters then have to be replaced by the
appropriate crossed homomorphisms. We introduce no such extra data.
\end{proof}

\begin{statement}{A.5}{Proved}{What the proposed product formula actually supplies}
Set $X=\mathbb P^1_{\Q}\setminus\{0,1728,\infty\}$ and
$u=1-1728/j\in\mathcal O(X)^\times$. Its Kummer class has geometric
boundary residues $(1,1,0)$ in $\Ftwo$, ordered as
$(0,1728,\infty)$. For a rational point $t\in X(\Q)$ its pullback
is the square class of $1-1728/t$. This class is not by itself a
comparison of the two CM frames.
\end{statement}
\begin{proof}
The divisor on the coarse projective line is
$[1728]-[0]$, and $u(\infty)=1$. The residue of the Kummer double
cover $v^2=u$ at a puncture is its order modulo two, giving the
displayed vector and its zero sum. Functoriality of Kummer theory
gives the value at a rational section. For an arbitrary continuous
section of $\pi_1(X/\Q)\to G_{\Q}$, restriction still gives an
element of $H^1(G_{\Q},\mu_2)$, but identifying it with a rational
evaluation requires additional information about that section.
These residues are fixed properties of $u$, independent of $g$.
Calling $(-1)^{\Delta(g)}$ the third local sign gives a true
identity by A.2; identifying it with an independently defined
cusp datum is an additional comparison, not Weil reciprocity
applied to the displayed divisor. The section-reciprocity claims
ANABELIAN D.27--D.28 are not used as proofs of that comparison.
\end{proof}

\begin{statement}{A.6}{Proved}{The CRT encoding is an encoding of parities}
There is a continuous homomorphism on $\HH$
\[
 u(h)\equiv(-1)^{a_3(h)}\pmod3,\qquad
 u(h)\equiv(-1)^{a_2(h)}\pmod4.
\]
It takes values in $\{1,5,7,11\}\subset(\Z/12)^\times$ and
has image modulo the global sign equal to $\delta$.
The complete table is
\[
\begin{array}{cc|ccc}
a_3&a_2&u&e&\delta\\\hline
0&0&1&0&0\\0&1&7&9&1\\1&0&5&4&1\\1&1&11&1&0
\end{array}
\]
where $e$ is the unique idempotent satisfying $2e=1-u$ in $\Z/12$.
The notation $(1-u)/2$ denotes this unique idempotent solution,
not division by an invertible element. For a general multiplier
only $[u(g\sigma^{-1})]$ is choice independent.
\end{statement}
\begin{proof}
CRT identifies the units with $\{\pm1\}\times\{\pm1\}$.
Multiplication of signs gives the homomorphism law. The idempotents
of $\Z/12\simeq\Z/3\times\Z/4$ are precisely the four pairs of
zeros and ones, and direct substitution gives the table. The two
solutions of $2e=1-u$ without the idempotence requirement need not
coincide. Changing $\sigma$ multiplies $u$ by $1$ or $11$, by
A.3. The ring $\Z/3\times\Z/4$ encodes the two signs as one
involution, but does not identify the two original frame matrices.
\end{proof}

\begin{statement}{A.7}{Refuted}{The uniform centre assertion at prime powers}
The assertion that
$\Ob(N)=Z(\SL_2(\Z/N))/\{\pm I\}$ is trivial at every prime
power is false. In particular $\Ob(8)\simeq C_2$.
For $N=2^a\prod_{i=1}^r p_i^{e_i}$ with $p_i$ odd,
the centre is the group of scalar square roots of one. Its rank
over $\Ftwo$ is $r+c(a)$, where
$c(0)=c(1)=0$, $c(2)=1$, and $c(a)=2$ for $a\ge3$.
For $N>2$, quotienting by the global sign reduces that rank by one.
\end{statement}
\begin{proof}
A matrix commuting with
$\left(\begin{smallmatrix}1&1\\0&1\end{smallmatrix}\right)$
has lower-left entry zero and equal diagonal entries; commuting
also with its transpose forces the upper-right entry zero. A
central matrix is therefore $aI$ with $a^2=1$. Conversely these
matrices are central. At an odd prime power,
$(a-1)(a+1)=0$ and the factors are coprime to each other at $p$,
so $a=\pm1$. At $2^a$, the roots are one for $a=1$, two for
$a=2$, and $\{1,-1,1+2^{a-1},-1+2^{a-1}\}$ for $a\ge3$,
as follows by comparing the two consecutive even valuations.
CRT finishes the proof. Modulo eight, $3I$ is central with
determinant one and is neither $I$ nor $-I$. The independent
enumeration in \texttt{check\_centres\_joint.py} checks the full
elementary-generator centralizers for $N=8,12$, and the roots for
$N=24,60$. The latter obstruction groups both have order four.
\end{proof}

\begin{statement}{A.8}{Refuted}{The proposed all-level formula is not a defined character}
The formula in Lead 8 cannot be used as a definition on all of
$G$ without new data: actual frames centralize individual torsion
elements, not the whole congruence group; moreover an arbitrary
multiplier is not a scalar element of $\SL_2(\Z/N)$.
\end{statement}
\begin{proof}
Even when an automorphism induces the identity on a finite
quotient, its frame satisfies $q(A)\in C_Q(q(r))$, which does
not imply $q(A)\in Z(Q)$. The noncentral Gaussian frame supplied
by A.25 below is an actual Galois example. We can strengthen it
to an example whose normalized congruence action is the identity.
The full Galois frame image in B.14 supplies an element with
$T=-I$ and $W=S$. Every positive pure $r$-frame belongs to $F'$
(as proved in B.15), and in $\SL_2(\mathbf F_3)^{\rm ab}=C_3$
the identity $T=q_3(A)^{-1}U_0^{8\rho}$ gives
$[T]=2\rho[U_0]$. Since $[-I]=0$ and $[U_0]$ generates $C_3$,
we obtain $\rho=0\pmod3$. Thus the R3 shear is the identity
modulo twelve, and both finite torsion images are fixed.
Nevertheless $W=S$ is noncentral projectively: it fails to commute
with the order-three generator. The finite matrix certificate
in \texttt{check\_pointed\_theta.py} checks that failure.
This refutes even the conditional centrality assertion when the
induced congruence action is literally the identity.
At $N=5$ the scalar
$2I$ has determinant four, not one. The diagonal matrix
$\operatorname{diag}(2,3)$ instead has determinant one but is
not central. Thus neither interpretation supplies the factor
$\langle\lambda\rangle$ in the asserted centre quotient. Also
R2 realizes one bit, and cannot by itself prove surjectivity to
groups of arbitrary rank. A.6 is the strongest all-multiplier
encoding established here from the supplied two parities.
\end{proof}

\begin{statement}{A.9}{Refuted}{The order-6912 candidate does not determine $\Delta$}
Let $U=F'F^3$ and let
\[
 J=\ker(q_{12})\cap(F'F^{12})\triangleleft\B,
 \qquad Q_{\rm joint}=\B/J.
\]
Here the second factor is the actual Fermat quotient of the
level-two subgroup, retaining its extension by $S_3$.
The R2 element and the identity induce the same automorphism of
$Q_{\rm joint}$ but have different $\Delta$ values. Consequently
$\Delta$ does not factor through this finite action, even though
$|Q_{\rm joint}|=6912$.
\end{statement}
\begin{proof}
Under the signed CRT identification
$q_{12}(F)=V_4\times\SL_2(\mathbf F_3)$, the image of $U$ is
$V_4\times Q_8$, a group of order 32. Indeed
$\SL_2(\mathbf F_3)'=Q_8$, cubes kill its quotient of order three,
and cubes generate the $V_4$ factor. In the Fermat quotient
$F/F'F^{12}\simeq(\Z/12)^2$, the image of $U$ is
$3(\Z/12)^2\simeq(\Z/4)^2$. Its image in the joint quotient
is therefore a subgroup of a product of two finite 2-groups,
and is nilpotent. By the definition of $R_{\rm nil}(U)$,
\[
 R_{\rm nil}(U)\subseteq J.
\]
The R2 automorphism is the identity on $\B/R_{\rm nil}(U)$,
so it is the identity on $\B/J$. It has $\Delta=1$, whereas
the identity has $\Delta=0$. This proves nonfactorization with
the exact comparison map in the required direction.

The independent program realizes the Fermat quotient by the
Bass--Serre graph cocycle in $(\Z/12)[S_3]\rtimes S_3$, with
$P(s)=0$, $P(r)=e_1-e_r$. It takes its joint image with $q_{12}$,
enumerates all 6912 elements, and normally generates the image
of $U$ using $x^3,y^3,[x,y]$. The image has order 128; its
congruence and Fermat projections have orders 32 and 16.
These numbers verify the finite input; no finite shadow is
asserted to be a full GT element by the computation.
\end{proof}

\begin{statement}{A.10}{Proved}{Existence of a detecting characteristic level}
There is a finite characteristic quotient of $\B$ whose
induced outer action determines $\Delta$. Every such quotient
must have nonnilpotent image of $U=F'F^3$. Neither continuity
nor this existence theorem identifies a smallest such quotient.
\end{statement}
\begin{proof}
A finitely generated profinite group has finitely many open
subgroups of any bounded index: their coset actions are among
finitely many maps to finite symmetric groups. Intersecting all
open normal subgroups of index at most $n$ gives a cofinal
sequence of open characteristic subgroups $K_n$. The maps
$\Out(\B)\to\Out(\B/K_n)$ separate points. Indeed, if an
automorphism is inner at every level, the compact sets of
possible conjugators at successively finer levels have the
finite intersection property and give a conjugator in $\B$.
The same compactness argument shows these kernels are a
neighbourhood basis of the identity in $\Out(\B)$.
The specified embedding of $G$ into this outer automorphism
group is a homeomorphism onto its compact image. Since
$\ker\Delta$ is open, one of these kernels meets $G$ inside it.
If a finite quotient has nilpotent image of $U$, its kernel
contains $R_{\rm nil}(U)$ and the R2 element acts identically;
it therefore cannot determine $\Delta$ by A.9's argument.
\end{proof}

\begin{statement}{A.11}{Proved}{The compactness certificate and its correct quantifiers}
Let $P:G\to\widehat\Gamma_{0,5}^{\rm pure}$ be the continuously
evaluated pentagon word, and let $L_n$ be cofinal open normal
subgroups of its target. Then
\[
 \Delta|_{\GT}=0
 \quad\Longleftrightarrow\quad
 \text{for some }n,\quad
 \{g\in G:\Delta(g)=1,\ P(g)\in L_n\}=\varnothing.
\]
The corresponding finite test must retain actual $G$ residues
and the same evaluation maps; arbitrary separately constructed
finite shadows do not satisfy this criterion automatically.
\end{statement}
\begin{proof}
The sets on the right are closed and nested in the compact
space $G$. Their intersection is
$\Delta^{-1}(1)\cap P^{-1}(1)$, which is empty precisely
when the restriction vanishes. If all the closed sets were
nonempty, compactness would give a point in their intersection.
Conversely an empty set at one level implies the intersection
is empty. Refine a detecting level from A.10 together with
the finitely many word evaluations to obtain a finite test.
A family of candidate residues which has not been shown to
arise as projections of a common compact space of gentle
automorphisms can contain extra points. Surviving such points
do not prove the existence of a compatible profinite point.
\end{proof}

\begin{statement}{A.12}{Proved}{The maximal abelian closures have exactly one common extra bit}
Inside a fixed algebraic closure of $\Q$, put
$K_4=\Q(i)$ and $K_3=\Q(\sqrt{-3})$. Then
\[
 K_4^{\rm ab}\cap K_3^{\rm ab}
       =\Q^{\rm ab}(3^{1/4}),\qquad
 [\Q^{\rm ab}(3^{1/4}):\Q^{\rm ab}]=2.
\]
This proves the field-intersection assertion of Lead 7. It
does not prove its claimed surjectivity for the adelic frames.
\end{statement}
\begin{proof}
Both maximal abelian closures are Galois over $\Q$, since
their quadratic base fields are Galois and the maximal
abelian extension is characteristic. Their intersection $L$
contains $\Q^{\rm ab}$. In $G_L=\Gal(L/\Q)$ the two
distinct subgroups $A_i=\Gal(L/K_i)$ are abelian of index
two. Their intersection $Z$ is central: it commutes with
both subgroups, and those subgroups generate $G_L$.
Moreover $G_L/Z\simeq C_2^2$. Choose coset representatives
$a,b$ which generate this quotient. Every commutator is
a power of $[a,b]$, which is central, and
$[a,b]^2=[a^2,b]=1$. Thus the closed derived subgroup
has order at most two. Its fixed field is exactly
$\Q^{\rm ab}$, so $[L:\Q^{\rm ab}]\le2$.

Let $t=3^{1/4}$ and $M=\Q(i,t)$. The polynomial $X^4-3$
is Eisenstein, and its splitting group is the dihedral group
of order eight: take $r(t)=it$, $r(i)=i$, and
$c(t)=t$, $c(i)=-i$. The subgroup fixing $K_4$ is
$\langle r\rangle\simeq C_4$; the subgroup fixing
$it^2=\sqrt{-3}$ is $\langle r^2,rc\rangle\simeq C_2^2$.
Hence $M/K_4$ and $M/K_3$ are abelian. Composing with
$\Q^{\rm ab}$ keeps them abelian over both quadratic fields,
so $\Q^{\rm ab}(t)\subseteq L$. But $t\notin\Q^{\rm ab}$:
otherwise the nonabelian Galois extension $M/\Q$ would be
contained in an abelian extension. Since $t^2=\sqrt3$ is
already abelian, the degree is two and the upper bound is
attained. Finally $(3/4)^{1/4}=t/\sqrt2$, and $\sqrt2$
is abelian, giving the same bit as in R6.
\end{proof}

\begin{statement}{A.13}{Proved}{The actual finite CM fields in the stated models}
For $E_3:y^2=x^3+1$ and $E_4:y^2=x^3-x$, with
$r_0=(3/4)^{1/4}$, one has
\[
 \Q(E_3[4])=\Q(i,r_0)=\Q(x(E_4[3])).
\]
The common bit on $G_{\Q^{\rm ab}}$ is the action on $r_0$.
\end{statement}
\begin{proof}
For a nonzero two-torsion point of $E_3$ whose coordinate is
$e\in\{-1,-\zeta_3,-\zeta_3^2\}$, the four halves have
$x=e(1\pm\sqrt3)$. For $e=-1$ their ordinate squares
are $-9\pm6\sqrt3$; a square root of the plus value is
$\sqrt3(\sqrt3-1)r_0$. The product of the two ordinate
roots is $3\sqrt{-3}$, up to sign. These expressions lie
in $\Q(i,r_0)$ and generate it: $i$ is obtained by the
Weil pairing on four-torsion, $\sqrt3$ by the abscissae,
and $r_0$ from the displayed ordinate.
The three-division polynomial for $E_4$ is
$3X^4-6X^2-1$. Choose roots $a,b$ with
$a^2=1+2/\sqrt3$, $b^2=1-2/\sqrt3$, $ab=i/\sqrt3$.
Then $w=a+ib$ satisfies $w^4=4/3$, and
\[
 a=(w+2/w)/2,\qquad b=(w-2/w)/(2i).
\]
The splitting field is exactly $\Q(i,1/w)=\Q(i,r_0)$.
These algebraic identities are independently checked in
\texttt{check\_root\_algebra.py}. This is a direct verification
of the unrefereed field statement, separate from the supplied
frame-to-torsion comparison in R4.
\end{proof}

\begin{statement}{A.14}{Proved}{The outer descent endomorphism exists}
Distinguish the abstract free profinite group $F_0=\langle X,Y\rangle$
from the level-two subgroup $F\subset\B$. Fix the identifications
of the prompt and the epimorphism
\[
 \varphi:F_0\twoheadrightarrow\B,
 \qquad X\mapsto s,\quad Y\mapsto sr.
\]
There is a continuous group endomorphism $D:G\to G$ with the
same multiplier, obtained by descending an outer action along
$\varphi$, then restricting to $F$ with its $S_3$ action
normalized. The outer compatibility group is exactly
$C_P=\operatorname{Fix}(D_P)$. Here $D_P=D$ uses the literal
current marking. It must not be identified with the original
manuscript's compatibility subgroup without changing markings.
\end{statement}
\begin{proof}
Put $Z=(XY)^{-1}$. The kernel of $\varphi$ is the closed
normal closure of $X^2,Z^3$. A gentle automorphism carries
each of $X,Y,Z$ to a conjugate of its common unit power.
It therefore preserves this kernel, in both directions, and
descends to a peripheral outer automorphism of $\B$.
Changing the representative by an inner automorphism of
$F_0$ changes the descendant by an inner automorphism of
$\B$, so descent is well defined and respects composition.

The subgroup $F=\ker(\B\to S_3)$ is characteristic:
the epimorphisms to $S_3$ form a single $\Aut(S_3)$ orbit.
An automorphism of $\B$ induces an inner automorphism on
$S_3$. Choose a lift $b\in\B$ of that conjugating
permutation and compose with $\operatorname{Ad}(b^{-1})$.
The normalized action induces the identity on $S_3$.
Its restriction to $F$ fixes the three boundary classes
up to the common unit power and commutes with the deck
$S_3$ in $\Out(F)$: the discrepancy between the images
of a deck lift is an element of $F$ and hence acts inner.
Different choices of $b$ differ by $F$, so define the
same outer restriction. Normalized representatives compose,
proving the homomorphism law for restriction.

This restriction is inverse to the given extension
$G\to\Out^{\rm per}(\B)$. To see injectivity explicitly,
an automorphism restricting inner to $F$ can be adjusted
to fix $F$ pointwise; then $b^{-1}\alpha(b)$ centralizes
$F$ for every $b\in\B$. The centralizer is trivial:
its $S_3$ image acts trivially on the faithful standard
action on $F^{\rm ab}$ and is therefore trivial, after
which the centre of the free profinite group is trivial.
Thus $\alpha$ is the identity. Continuity follows from
descent and restriction in the finite-quotient topology.
Composing restriction with descent gives $D$.

Compatibility of the two outer actions on $\B$ says
exactly that extension of $g$ equals descent of $g$;
applying the inverse restriction proves the fixed-point
criterion. The original manuscript uses a different free basis.
The required conversion, and a Galois counterexample to the
literal-current fixed-point assertion, are proved in A.60--A.64.
Thus the arithmetic functoriality sentence in checkpoint 01's
version of this statement was incorrect in the current marking
and is superseded here. The relevant manuscript passage is
\S47, cote 140--3, folios 415--420, scans 229--240.
No unproved equivalence with the harmonic equations is built
into this algebraic construction.
\end{proof}

\begin{statement}{A.15}{Proved}{The elementary factorization of the modular map}
With the Legendre coordinate $t$,
\[
 w=4t(1-t),\qquad j=64(4-w)^3/w^2
\]
has composite $256(1-t+t^2)^3/[t^2(1-t)^2]$.
This algebraic identity alone does not identify the
fixed-point equation of A.14 with relations (III$'$),
(NT1) and (NT2).
\end{statement}
\begin{proof}
Substitute $4-w=4(1-t+t^2)$ and
$w^2=16t^2(1-t)^2$. The factorization is over $\Q$.
An equality of coarse rational functions does not supply
the choices of tangential basepoints, inertia generators,
and conjugating paths needed to equate the respective
automorphism equations. Those are additional data in
the comparison requested by A.4(a).
\end{proof}

\begin{statement}{A.16}{Refuted}{The naive integral-specialization reading}
The stated Weierstrass models do not give two smooth
elliptic sections over $\Z$ whose torsion frames can simply
be read at both 2 and 3. The coarse $j$ coordinates do
coincide at those primes; this is weaker than the proposed
comparison of profinite frames.
\end{statement}
\begin{proof}
The discriminants are $-432$ for $E_3$ and $64$ for $E_4$.
Modulo three, $E_3$ has equation $y^2=(x+1)^3$ and is
singular at $(-1,0)$, whereas $E_4$ is smooth. Modulo
two both displayed cubics are singular. These assertions
follow by differentiating their homogeneous equations
and are checked in \texttt{check\_root\_algebra.py}.
After potential good reduction the common $j=0$ is
supersingular at these primes, but such a statement
requires a change of model and a base extension.
Moreover prime-to-residue-characteristic specialization
does not recover the wild-primary torsion signs used here.
The identity $1728=2^6 3^3$ proves the collision on the
coarse $j$ line, and supplies no map identifying the
two framed representations. Over $\Z[1/6]$ these two
coarse sections are disjoint.
\end{proof}

\begin{statement}{A.17}{Refuted}{The standard genus-two matrix action is not an action of the five-point sphere group}
The usual symplectic representation of $B_5$ from the genus-two
double cover does not descend even projectively to
$\Gamma_{0,[5]}$, including at level twelve. It therefore
cannot be substituted without a new comparison for the
five-point frame equations.
\end{statement}
\begin{proof}
Take a chain basis $e_1,\ldots,e_4$ of homology with
$J_{i,i+1}=1$, $J_{i+1,i}=-1$, and other entries zero.
It is unimodular. Its transvections
$T_i=I+e_i e_i^{\mathsf T}J$ satisfy the Artin braid
and distant-commutation relations. Direct multiplication gives
\[
 (T_1T_2T_3T_4)^5=-I,
\]
but the sphere relator is
\[
 T_4T_3T_2T_1^2T_2T_3T_4
 =-I+2(e_1+e_3)e_4^{\mathsf T}.
\]
It is not scalar modulo three, four or twelve. Killing
the central full twist thus does not kill the sphere
relator, and passing to the projective symplectic group
does not fix the failure. All matrix identities are
checked in \texttt{check\_genus\_two\_descent.py}.
Geometrically the curve $y^2=x^5-1$ has six branch
points, the five roots and infinity. A double cover
cannot have exactly five simple branch points, by
Riemann--Hurwitz. The extra point must be retained in
any proposed functorial comparison with five-point data.
This refutes the immediate substitution, not the
possibility of a different comparison theorem.
\end{proof}

\subsection{The three arithmetic CM fields}


\begin{statement}{A.18}{Proved}{A quadratic-by-quartic intersection bound}
Let $K/\Q$ be quadratic and $L/\Q$ cyclic quartic, with $K\cap L=\Q$.
Inside a fixed algebraic closure, put $M=K^{\mathrm{ab}}\cap L^{\mathrm{ab}}$.
Then $M/\Q^{\mathrm{ab}}$ has degree at most two.
\end{statement}
\begin{proof}
Both maximal abelian extensions are Galois over $\Q$, and contain
$\Q^{\mathrm{ab}}$. In $G=\Gal(M/\Q)$ put
$A=\Gal(M/K)$ and $B=\Gal(M/L)$. These are closed abelian normal
subgroups of indices two and four. Since $K\cap L=\Q$, we have
$AB=G$. The subgroup $Z=A\cap B$ is central: its elements commute
with $A$ and $B$, and these generate $G$. The quotient $G/Z$ is
$\Gal(KL/\Q)=C_2\times C_4$. Choose $a\in A$ generating its
order-four factor and $b\in B$ generating its order-two factor.
Then $G$ is topologically generated by $a,b,Z$, and all commutators
are central and generated by $[a,b]$. Because $b^2\in Z$,
$[a,b]^2=[a,b^2]=1$. Thus $G'$ has order at most two. The maximal
abelian subextension of $M/\Q$ is exactly $\Q^{\mathrm{ab}}$,
since that field is already contained in $M$. This proves the bound,
including for the infinite profinite extensions.
\end{proof}

\begin{statement}{A.19}{Proved}{The Gaussian--quintic intersection}
For $K_4=\Q(i)$ and $K_5=\Q(\zeta_5)$,
\[
 K_4^{\mathrm{ab}}\cap K_5^{\mathrm{ab}}
       =\Q^{\mathrm{ab}}(\sqrt[4]{5}).
\]
\end{statement}
\begin{proof}
Let $D_5=\Q(i,\sqrt[4]{5})$. The polynomial $X^4-5$ is Eisenstein
at five, and its real root generates a nonnormal real quartic field.
Thus $D_5/\Q$ has dihedral group of order eight. Its three quadratic
subfields are $\Q(i)$, $\Q(\sqrt5)$ and $\Q(\sqrt{-5})$.
Its extension over $K_4$ is cyclic quartic. The only quadratic
subfield of $K_5$ is $\Q(\sqrt5)$, and $K_5$ cannot be a quartic
Galois subfield of $D_5$, because the abelianization of the dihedral
group has exponent two. Consequently
$D_5\cap K_5=\Q(\sqrt5)$; the group over this intersection is the
Klein four group. Hence $D_5K_5/K_5$ is abelian too.
It follows that $\Q^{\mathrm{ab}}(\sqrt[4]5)$ lies in the claimed
intersection. It is quadratic over $\Q^{\mathrm{ab}}$: its square
is already there, whereas the fourth root itself is not, since
$D_5/\Q$ is nonabelian. A.18 supplies the opposite bound.
\end{proof}

\begin{statement}{A.20}{Proved}{The Eisenstein--quintic intersection and a finite defining field}
Let $K_3=\Q(\sqrt{-3})$ and $E_3:y^2=x^3+1$. The group of
$D_{35}=\Q(E_3[5])$ over $\Q$ is $C_{24}\rtimes_5 C_2$.
Let $M_{35}$ be its subfield fixed by the unique subgroup of
order three in the Cartan $C_{24}$. Then
\[
 \Gal(M_{35}/\Q)=C_8\rtimes_5 C_2,\qquad
 K_3^{\mathrm{ab}}\cap K_5^{\mathrm{ab}}
                    =\Q^{\mathrm{ab}}M_{35}.
\]
The extension on the right has degree two over $\Q^{\mathrm{ab}}$.
This description specifies $M_{35}$ unambiguously as a finite
subfield of the division field, and does not depend on a choice
of its primitive generator.
\end{statement}
\begin{proof}
The endomorphism $(x,y)\mapsto(\zeta_3x,y)$ is defined over $K_3$.
The $\mathbf F_5$-space $E_3[5]$ is one-dimensional over
$\mathcal O_{K_3}/5=\mathbf F_{25}$. Therefore the image of
$G_{K_3}$ is contained in its cyclic unit group of order 24.
The field $D_{35}$ contains $K_3$. Indeed, if a Galois element
trivial on $E_3[5]$ conjugated this endomorphism to its conjugate,
their difference would annihilate $E_3[5]$. That difference has
CM norm three, so it is invertible on five-torsion, a contradiction.

We now prove the full Cartan image, rather than assume a general
surjectivity theorem. Counting gives $\#E_3(\mathbf F_7)=12$.
Frobenius at either prime of $K_3$ above seven has trace $-4$
and norm seven, so it is multiplication by $-2+\sqrt{-3}$ or its
conjugate. In $\mathbf F_{25}=\mathbf F_5[t]/(t^2-2)$ these two
units $(3,1)$ and $(3,-1)$ have order 24. The attached exact
program checks the point count and both orders. Hence the image
is all of $C_{24}$. Complex conjugation has order two, acts on
this Cartan by the nontrivial field automorphism, and thus acts
by the exponent five. This proves the asserted order-48 group.
For a polynomial description of the finite division field, its
$x$-coordinates are the roots of
\[
       5X^{12}+380X^9-240X^6-1600X^3-256;
\]
one also adjoins their $y$-coordinates satisfying $y^2=x^3+1$.

Write the order-16 quotient as
$\langle b,c\mid b^8=c^2=1,\ cbc^{-1}=b^5\rangle$.
Its derived subgroup is $\langle b^4\rangle$ of order two.
The extension over $K_3$ is cyclic, generated by $b$.
The Weil pairing puts $K_5$ in $D_{35}$, and the odd subgroup
of order three fixes it. Choose the generator so that its
determinant is the chosen generator of $\mathbf F_5^\times$;
complex conjugation has determinant $-1$. The determinant kernel
in the order-16 group is
\[
       \{1,b^4,b^2c,b^6c\}=\langle b^2c\rangle\simeq C_4.
\]
Thus $M_{35}$ is abelian over both $K_3$ and $K_5$, but not over
$\Q$. Its compositum with $\Q^{\mathrm{ab}}$ is a nontrivial
quadratic extension, because its commutator subgroup has order
two. A.18 proves equality with the entire intersection.
\end{proof}

\begin{statement}{A.21}{Known}{The full elliptic torsion fields and the CM upper bound for the Jacobian}
Let $J_5=\operatorname{Jac}(y^2=x^5-1)$ and define Galois extensions
\[
 L_3=\Q^{\mathrm{ab}}(E_3[\mathrm{tors}]),\quad
 L_4=\Q^{\mathrm{ab}}(E_4[\mathrm{tors}]),\quad
 L_5=\Q^{\mathrm{ab}}(J_5[\mathrm{tors}]),
 \qquad E_4:y^2=x^3-x.
\]
Then $L_3=K_3^{\mathrm{ab}}$, $L_4=K_4^{\mathrm{ab}}$, and
$L_5\subseteq K_5^{\mathrm{ab}}$.
\end{statement}
\begin{proof}[Derivation from the stated CM theorem]
We use Weber's ray-class-field theorem in the precise form of
\citet[Theorem 8]{ClarkCookStankewicz2009}: for an elliptic curve
with the maximal CM order, adjoining its $j$-invariant and the
values of a Weber function on its $N$-torsion gives the ray class
field of modulus $N$. Here the Weber functions can be taken to
be $x^3$ for $E_3$ and $x^2$ for $E_4$. Both CM fields have class
number one, and the displayed curves and all their endomorphisms
are defined over those fields. Their full torsion fields are
abelian over the CM field: each Tate module is a rank-one module
over the corresponding completed CM order and Galois acts by
units. Every ray class field therefore lies in the torsion field,
and every finite abelian extension lies in a ray class field.
This proves the two equalities.

For the genus-two Jacobian the automorphism
$(x,y)\mapsto(\zeta_5x,y)$ induces an action of
$\mathcal O_{K_5}=\Z[\zeta_5]$ defined over $K_5$. On its
four-dimensional rational first homology this action has minimal
polynomial $\Phi_5$: the holomorphic differentials $dx/y$ and
$x\,dx/y$ have eigenvalues $\zeta_5,\zeta_5^2$, and their
conjugates supply the other two. Each Tate module is consequently
rank one over $\mathcal O_{K_5}\otimes\Z_\ell$. More explicitly,
a stable full lattice is torsion-free of rank one over each
local discrete valuation factor, hence free there, including at
$\ell=5$. Galois over $K_5$ acts by units of this commutative
ring. Every torsion extension over $K_5$ is abelian. Since
$\Q^{\mathrm{ab}}\subset K_5^{\mathrm{ab}}$, the asserted
containment follows. No equality $L_5=K_5^{\mathrm{ab}}$ is used.
\end{proof}

\begin{statement}{A.22}{Proved}{The Jacobian contains both predicted extra quadratic intersections}
The field $L_5$ contains $\sqrt[4]5$ and $\Q^{\mathrm{ab}}M_{35}$.
\end{statement}
\begin{proof}
For the first assertion we use the specialization statement,
not a generic-curve assertion: \citet[Corollary 1.2(a)]{Yelton2019}
places in the eight-division field the fourth roots of
$\prod_{j\ne i}(\alpha_j-\alpha_i)$, after the stated cyclotomic
adjoining, for the distinct roots $\alpha_i$ of an odd-degree
hyperelliptic polynomial. The necessary cyclotomic constants
are in $\Q^{\mathrm{ab}}$. At the root $\alpha_i=1$ of
$x^5-1$ this product is $f'(1)=5$. Thus
$\sqrt[4]5\in L_5$.

For the second assertion, put $D_J=\Q(J_5[3])$.
It contains $K_5$: a Galois element fixing the three-torsion
cannot change the endomorphism $\zeta_5$ to $\zeta_5^a$ for
$a\ne1$, since their difference has norm five and is invertible
on three-torsion. Let $U=\Gal(D_J/K_5)$, viewed through the
faithful rank-one representation as a subgroup of
$(\mathcal O_{K_5}/3)^\times=\mathbf F_{81}^\times$.
The Rosati involution for the canonical principal polarization
is complex conjugation, which is $u\mapsto u^9$ on this field.
The Weil pairing therefore gives
\[
                 u^{10}=u\overline u=\chi_3(u)\in\mathbf F_3^\times.
\]
So $U$ lies in the cyclic group of order 20 cut out by $u^{20}=1$.
The prime eleven splits in $K_5$, and is a prime of good reduction
of this curve. Frobenius at a prime above eleven belongs to
$G_{K_5}$ and has multiplier $11=-1\pmod3$. Its image satisfies
$u^{10}=-1$, forcing its two-primary part to have order four.
Thus the two-primary part of $U$ is exactly $C_4$.

Let $O$ be the unique odd-order subgroup of $U$. It is
characteristic in $U$, hence normal in $\Gal(D_J/\Q)$; it has
order one or five. Set $M_J=D_J^O$. There is an exact sequence
\[
  1\longrightarrow\langle u\rangle\simeq C_4
   \longrightarrow\Gal(M_J/\Q)\longrightarrow C_4
   \longrightarrow1.
\]
A generator of $\Gal(K_5/\Q)$ sends $\zeta_5$ to $\zeta_5^2$.
On $\mathbf F_{81}$ this is the third power of Frobenius, since
$3^3=2\pmod5$; on its subgroup $C_4$ it acts by exponent
$27=-1\pmod4$. Consequently the displayed group is nonabelian,
its derived subgroup is exactly $\langle u^2\rangle$, and that
subgroup is central. These conclusions do not require the exact
sequence to split.

The Weil pairing puts $\mu_3$ in $D_J$, and the odd group $O$
fixes it, so $K_3\subset M_J$. The multiplier character on
$\langle u\rangle$ sends its generator to $-1$. Choose a lift
$v$ of a generator of the quotient $C_4$ and, if necessary,
replace it by $uv$ to make its multiplier one. The multiplier
kernel is generated by $u^2$ and $v$: it has index two, since
these generate four quotient cosets and the central order-two
subgroup. The two generators commute, so this kernel is abelian.
It is exactly $\Gal(M_J/K_3)$. The group over $K_5$ is already
cyclic. Thus $M_J$ lies in $K_3^{\mathrm{ab}}\cap K_5^{\mathrm{ab}}$,
and its compositum with $\Q^{\mathrm{ab}}$ is the nontrivial
quadratic extension described in A.20. Since $M_J\subset D_J$,
that whole compositum lies in $L_5$.

The finite-field computations, both possible order-16 extension
presentations, their commutator and multiplier kernels, and the
Frobenius certificate in A.20 are checked in
\texttt{code/check\_cm\_fields.py}. They support the exact group
arguments; they do not substitute for the CM and specialization
comparisons proved above.
\end{proof}

\begin{statement}{A.23}{Proved}{All three pairwise torsion-field intersections and the triple intersection}
With the definitions in A.21, the complete arithmetic answer is
\[
\begin{aligned}
 L_3\cap L_4&=\Q^{\mathrm{ab}}(\sqrt[4]3),\\
 L_4\cap L_5&=\Q^{\mathrm{ab}}(\sqrt[4]5),\\
 L_3\cap L_5&=\Q^{\mathrm{ab}}M_{35},\\
 L_3\cap L_4\cap L_5&=\Q^{\mathrm{ab}}.
\end{aligned}
\]
The three nontrivial quadratic extensions in this display are
linearly disjoint over $\Q^{\mathrm{ab}}$ and have compositum
of degree eight. These are arithmetic field statements, not a
construction of three characters of $\GT$.
\end{statement}
\begin{proof}
The first equality is A.12; its finite radical may be written
$\sqrt[4]3$ because $\sqrt2\in\Q^{\mathrm{ab}}$.
The second and third follow from A.19--A.22: the maximal
abelian intersections give upper bounds, and A.22 supplies the
full quadratic subextensions inside the actual Jacobian field.

For the triple intersection it suffices to show that
$\sqrt[4]3\notin K_5^{\mathrm{ab}}$. Let
$D_3=\Q(i,\sqrt[4]3)$, whose dihedral Galois group has order
eight and whose quadratic subfields are $\Q(i)$,
$\Q(\sqrt3)$ and $\Q(\sqrt{-3})$.
Its intersection with $K_5$ is $\Q$: it contains no
$\Q(\sqrt5)$ and has no cyclic quartic quotient.
Thus $D_3K_5/K_5$ is still nonabelian. If the fourth root
belonged to $K_5^{\mathrm{ab}}$, then so would $i$, all of
$D_3$, and their compositum with $K_5$, a contradiction.
The triple intersection is therefore trivial over
$\Q^{\mathrm{ab}}$.

Finally the first and second quadratic fields are distinct,
since the first is not contained in $K_5^{\mathrm{ab}}$ whereas
the second is. Both lie in $L_4$, so their biquadratic
compositum lies in $L_4$. The third does not lie in $L_4$,
by the triple-intersection assertion, because it already lies
in $L_3\cap L_5$. It is therefore independent of the first
two, proving the degree-eight assertion.
\end{proof}

\begin{statement}{A.24}{Open}{The missing passage from these fields to order-five frame characters}
The equalities in A.23 give relations among the actual Galois
representations of the three fibres. They do not provide a
corrected homomorphism from the multiplier-one subgroup of
$\GT$ to an order-five norm-one target, or a comparison of that
homomorphism with the Jacobian representation. We have not
proved either of those maps. Consequently we have not extended
the three arithmetic quadratic relations to characters of
$\GT$, nor evaluated such proposed characters away from Galois.
\end{statement}
The smallest missing comparison is already the three-torsion
coordinate of $J_5$ used in A.22: we would need a genuine
five-point action, a specified path and normalization of its
order-five frame, and a proof that its finite coordinate has
the arithmetic quotient $M_J$ above. The usual genus-two braid
representation fails the sphere relation (A.17), so it cannot
be used without a different group and descent argument.

\subsection{The categorical and theta tests}


\begin{statement}{A.25}{Refuted}{the frames are central}
\end{statement}
Write $\mathcal H=\ker(\lambda:\widehat{\mathrm{GT}}_0\to
\widehat{\mathbf Z}^{\times})$, and use precisely the signed lifts and
cusp correction of LM7\_PUSH E.65.  The proved maps in that statement are
\[
 T:\mathcal H\longrightarrow
 C_{\mathrm{SL}_2(\mathbf F_3)}(R_0)
 =\langle-I\rangle\times\langle R_0\rangle,
 \qquad
 W:\mathcal H\longrightarrow
 C_{\mathrm{SL}_2(\mathbf F_3)}(S)=\langle S\rangle,
\]
with $S=\left(\begin{smallmatrix}0&1\\-1&0\end{smallmatrix}\right)$.
These are centralizers, not the scalar centre.  In particular the assertion
that both corrected frames are scalar modulo twelve is false even on
$G_{\mathbf Q}\cap\mathcal H$.

Here is an actual Galois counterexample, specified by a nonempty compact
coset of the absolute Galois group.  Put $r_0=(3/4)^{1/4}>0$.
The polynomial $4X^4-3$ is Eisenstein at $3$, so
$[\mathbf Q(r_0):\mathbf Q]=4$.  This real field is not normal over
$\mathbf Q$: it does not contain the root $ir_0$ of that irreducible
polynomial.  Every finite subextension of $\mathbf Q^{\rm ab}/\mathbf Q$
is normal, and consequently $r_0\notin\mathbf Q^{\rm ab}$.
On the other hand $r_0^2=\sqrt3/2$ belongs to $\mathbf Q^{\rm ab}$.
Thus $\mathbf Q^{\rm ab}(r_0)/\mathbf Q^{\rm ab}$ is quadratic.
The restriction map
\[
 G_{\mathbf Q^{\rm ab}}\longrightarrow
 \operatorname{Gal}(\mathbf Q^{\rm ab}(r_0)/\mathbf Q^{\rm ab})
\]
is a continuous surjection.  Take any element $\sigma$ in the inverse
image of its nonidentity element.  This inverse image is nonempty and
compact, and $\sigma(r_0)=-r_0$ and $\chi(\sigma)=1$.

The signed Galois comparison in the referee's repaired proof of
LM7\_PUSH E.65, using the two specific models
$E_3:y^2=x^3+1$ and $E_4:y^2=x^3-x$, gives
\[
 \epsilon_3(T_\sigma)=\epsilon_2(W_\sigma)=1.
\]
Therefore $W_\sigma\bmod3$ is $S$ or $-S$.  It is not scalar, even
projectively.  Indeed, for $U=\left(\begin{smallmatrix}1&1\\0&1\end{smallmatrix}\right)$,
\[
 SU=\begin{pmatrix}0&1\\2&2\end{pmatrix},\qquad
 US=\begin{pmatrix}2&1\\2&0\end{pmatrix}\quad\text{over }\mathbf F_3,
\]
and these matrices are neither equal nor negatives of each other.
Reduction modulo three of a central scalar class modulo twelve would be
scalar.  This proves the claimed counterexample.  It refutes the literal
central-frame assertion, not the parity compression that follows.

\begin{statement}{A.26}{Proved}{CRT compression and its precise ambiguity}
\end{statement}
Let $a(\Phi)=\epsilon_3(T_\Phi)$ and
$b(\Phi)=\epsilon_2(W_\Phi)$ for $\Phi\in\mathcal H$.
There is a unique continuous homomorphism
\[
 u:\mathcal H\longrightarrow\mu_2(\mathbf Z/12)
 =\{1,5,7,11\}
\]
whose reductions modulo three and four are respectively $(-1)^a$ and
$(-1)^b$.  Explicitly,
\[
\begin{array}{c|c|c|c}
(a,b)&u&e\text{ with }e^2=e,\;2e=1-u&\delta=a+b\\\hline
(0,0)&1&0&0\\
(1,0)&5&4&1\\
(0,1)&7&9&1\\
(1,1)&11&1&0
\end{array}
\]
Continuity and multiplicativity follow from LM7\_PUSH E.65 and CRT.
The idempotents of $\mathbf Z/12\simeq\mathbf F_3\times\mathbf Z/4$
are the four pairs with both coordinates zero or one; the table follows.
The expression $(1-u)/2$ is not ordinary division in $\mathbf Z/12$.
Each displayed equation $2e=1-u$ has two solutions in that ring, and
exactly one is idempotent.  This condition is essential.

The element of R2 has $u=5$.  Any independently established lift with
parities $(0,1)$ has $u=7$; the existence of that lift must be proved
separately and is not a consequence of this table.  On
$G_{\mathbf Q}\cap\mathcal H$ the values lie in $\{1,11\}$, and the
preceding counterexample shows that $11$ really occurs.

The ring $\mathbf F_3\times\mathbf Z/4$ combines the two \emph{signs}.
It does not combine the two actual frames as one scalar matrix.  Their
natural adelic centralizer algebras are different quadratic algebras,
the Eisenstein and Gaussian algebras.  Passing from their units to the
two indicated characters loses data.

For an arbitrary $\Phi\in\widehat{\mathrm{GT}}_0$, choose
$\sigma\in G_{\mathbf Q}$ with $\chi(\sigma)=\lambda(\Phi)$.
The well-defined output is
\[
 \overline u(\Phi)=
 [u(\Phi\sigma^{-1})]\ \in\
 \{1,5,7,11\}/\{1,11\}.
\]
If $\sigma$ is changed, the two arguments differ by an element of
$G_{\mathbf Q}\cap\mathcal H$ after conjugation, so the corresponding
sign-pairs differ by the diagonal sign.  The parity sum and its
conjugacy invariance, proved in LM7\_PUSH E.65, show that this formula
is independent of the choice, continuous, and multiplicative.
It corresponds to $\Delta(\Phi)$ under the unique isomorphism of its
target with $\mathbf F_2$.  The full lift $u(\Phi)$ has not been
canonically defined for general multiplier by this construction.

\begin{statement}{A.27}{Proved}{the theta identity and an induced action}
\citep{LitAtiyah1987}
\end{statement}
For $r\in\mathbf Z/12$ define
\[
 \theta_{6,r}(\tau,z)=
 \sum_{n\equiv r\;(12)}q^{n^2/24}e^{2\pi inz},
 \qquad q=e^{2\pi i\tau},\quad\operatorname{Im}\tau>0.
\]
Let $P_u$ permute the labels by $r\mapsto ur$.
All four units $u\in\{1,5,7,11\}$ satisfy $u^2=1\pmod{24}$;
thus they preserve the finite quadratic module
$A=\mathbf Z/12$, $q_A(r)=r^2/24\pmod{\mathbf Z}$.
In particular $P_u$ commutes with its diagonal $T$ matrix and its
finite Fourier $S$ matrix.  Euler's pentagonal identity gives
\[
 \eta(\tau)=q^{1/24}\prod_{n\ge1}(1-q^n)
 =\sum_{n\ge1}\chi_{12}(n)q^{n^2/24}
 =\theta_{6,1}(\tau,0)-\theta_{6,5}(\tau,0),
\]
where $\chi_{12}$ is zero on nonunits, positive at $1,11$, and
negative at $5,7$.  For example, multiplying Euler's exponent
$k(3k-1)/2$ by $24$ and adding $1$ gives $(6k-1)^2$;
sorting the positive representatives modulo twelve gives the stated
Dirichlet character.  Absolute convergence permits the rearrangement.
Also $\theta_{6,-r}(\tau,0)=\theta_{6,r}(\tau,0)$, so direct
substitution proves
\[
 P_u\eta=\chi_{12}(u)\eta,
 \qquad
 \chi_{12}(u(\Phi))=(-1)^{\delta(\Phi)}.
\]
At $z=0$ these are twelve labelled functions with linear dependencies;
on the formal twelve-dimensional label space the corresponding
eigenvector is $e_1+e_{11}-e_5-e_7$.

It follows that $\Phi\mapsto P_{u(\Phi)}$ is a continuous
representation of $\mathcal H$ on this label space and that its eta
line carries $(-1)^\delta$.  In this \emph{defined} representation
the R2 element acts by $P_5$ and gives eta eigenvalue $-1$; a lift
with parities $(0,1)$ acts by $P_7$ and has the same eta eigenvalue.
No comparison with an independently defined modular-functor action
has been proved.  This construction packages the known character;
it does not derive its restriction to $\widehat{\mathrm{GT}}$.

\begin{statement}{A.28}{Proved}{a pointed category kills the $f$-coordinate}
\citep[Section 4]{LitDrinfeld1991}
\end{statement}
We can make the issue completely explicit.  Work over
$K=\mathbf Q(\zeta_{24})$.  Use representatives $0,\ldots,11$
for $A=\mathbf Z/12$, and put
\[
 [a+b]=a+b-12\kappa(a,b),\qquad
 \kappa(a,b)=\left\lfloor\frac{a+b}{12}\right\rfloor.
\]
Consider finite-dimensional $A$-graded $K$-vector spaces with the
usual graded tensor product, associator and braiding on homogeneous
simple objects given by
\[
 \omega(a,b,c)=(-1)^{a\kappa(b,c)},\qquad
 c(a,b)=\zeta_{24}^{ab}.
\]
The pentagon follows from the carry identity
\[
 \kappa(b,c)+\kappa([b+c],d)
 =\kappa(c,d)+\kappa(b,[c+d]).
\]
More explicitly the exponent difference between its two sides is
$12\kappa(a,b)\kappa(c,d)$ modulo two.  For the first hexagon,
$c(a,[b+c])/(c(a,b)c(a,c))=(-1)^{a\kappa(b,c)}$,
and the associator corrections give the same value; the factors
containing $\kappa(a,c)=\kappa(c,a)$ cancel.  For the second,
$c([a+b],c)/(c(a,c)c(b,c))=(-1)^{c\kappa(a,b)}$,
and the same cancellation proves the equality.  This is a braided
pointed category with trace quadratic form
$c(a,a)=\zeta_{24}^{a^2}$ and nondegenerate double braiding
$\zeta_{12}^{ab}$.

For every pair of its double-braiding operators, evaluation of a
word $f\in\widehat F_2'$ is the identity.  Indeed the operators
are simultaneously diagonal on the homogeneous components of every
tensor product, with eigenvalues in the finite abelian group
$\mu_{12}$.  The unique continuous evaluation homomorphism from
$\widehat F_2$ to that finite diagonal group therefore kills
$\widehat F_2'$.  This statement holds on direct sums as well as
on individual simple objects.

Drinfeld's formula for the change of associator multiplies it by
such an evaluation of $f$ (\S4, equations (4.5), (4.10), and (4.12)
in the English edition).  Hence this correction is exactly one.
His braiding change, for $m=(\lambda-1)/2$, multiplies $c(a,b)$ by
the $m$th power of the double braiding, giving
\[
 c_\lambda(a,b)=c(a,b)^{\lambda}
 =\zeta_{24}^{\lambda ab},\qquad \omega_\lambda=\omega.
\]
These expressions make sense for every profinite unit $\lambda$
by reduction modulo twenty-four.  Since $\lambda$ is odd,
$\omega^\lambda=\omega$.  They also follow by applying the
field automorphism $\zeta_{24}\mapsto\zeta_{24}^{\lambda}$ to
the original coherence equations.  Therefore they satisfy both
hexagons and the pentagon for \emph{every} $f\in\widehat F_2'$,
including every gentle pair.  At $\lambda=1$ the resulting category
and its chosen structural maps are literally unchanged.

With the ribbon convention in which the trace quadratic form is
the twist, its modular data are
\[
 T_{ab}=\delta_{ab}\zeta_{24}^{a^2},\qquad
 S_{ab}=12^{-1/2}\zeta_{12}^{ab}.
\]
The opposite Fourier convention replaces the exponent by its
negative and makes no difference to the conclusion.  Cyclotomic
conjugation transforms $T$ by entrywise $\lambda$th powers.  If the
normalizing square root $\sqrt{12}=2\sqrt3$ is conjugated too,
then
\[
 \sigma_\lambda(S)_{ab}
 =\chi_{12}(\lambda)\,12^{-1/2}\zeta_{12}^{\lambda ab}
 =\chi_{12}(\lambda)S_{\lambda a,b}.
\]
The sign is obtained directly from
$\sqrt3=\zeta_{12}+\zeta_{12}^{-1}$.
In particular the pointed-category modular data depend on
$\lambda$, not on $f$.

\begin{statement}{A.29}{Refuted}{Refuted as a test; Open as a comparison}
\end{statement}
The assertion that the R2 element must fail the categorical pentagon
in the indicated pointed category is false.  Its multiplier is one,
so the previous paragraph proves that its proposed modification is
the identity modification and satisfies that pentagon.
This does \emph{not} prove that R2 satisfies the group pentagon.
The evaluation functor used here kills every $f$-word in the
commutator subgroup, and thus also kills the information that this
test was supposed to detect.

The distinction in Drinfeld's result is the quantifier: imposing
the group pentagon makes the construction coherent on all suitable
braided categories, naturally in braided functors.  A single
pointed category is a nonfaithful test of that condition.  The
unproved implication needed for the proposed route is an
identification, on the eta line or on the full label space, of
the actual frame representation with the independent modular
functor operation of a GT pair.  In our notation it would identify
$P_{u(\Phi)}$ with that operation after the specified cusp and
Galois normalizations.  For multiplier one the independent operation
is the identity whereas $P_{u(\Phi)}$ can be nonidentity.  Proving
that identification on all of $\widehat{\mathrm{GT}}\cap\mathcal H$
would already imply $\Delta=0$.  It is not supplied by the
category's pentagon or by the equality of modular-group
representations.  Consequently this route gives no decision of
LM7\_PUSH E.3.

\begin{statement}{A.30}{Proved}{abelianized phases cannot detect the bit}
\end{statement}
The projection to $B^{\rm ab}\simeq\mathbf Z/6$ cannot recover the
R2 parity from its pure frame.  Indeed the image of the level-two
free subgroup $F$ in $B^{\rm ab}$ has order three: with
$\tau=sr$ and $x=\tau^{-2}$, $y=sxs$, both generators map to
$-2\tau$.  Thus $F'F^3$ maps to zero.  The identity frame and
R2's frame in $F'F^3$ have the same abelianized image but different
parities.  The following additional argument establishes the
metaplectic version with an explicit, choice-independent convention.

In fact every multiplier-one positive $r$-frame belongs to $F'$.
The identities $\Phi(r)=f^{-1}r=ArA^{-1}$ and $f\in F'$
give $(I-R)[A]=0$ in $F^{\rm ab}=\widehat{\mathbf Z}^2$, where
conjugation by $r$ sends $(x,y,z)$ to $(y,z,x)$ and hence
\[
 R=\begin{pmatrix}0&-1\\1&-1\end{pmatrix},\qquad
 I-R=\begin{pmatrix}1&1\\-1&2\end{pmatrix},\qquad\det(I-R)=3.
\]
Multiplication by three is injective on $\widehat{\mathbf Z}$,
so the adjugate identity proves $[A]=0$. For the R2 element also
$\beta\in R_{\rm nil}(F'F^3)\subseteq(F'F^3)'\subseteq F'$,
because its retained order-two conjugator is one. Its shear is
therefore trivial, so its two corrected frames and their inverses
belong to $F'$.

Let $p:E\twoheadrightarrow\widehat B$ be any profinite central
extension, in particular the completion of
$\operatorname{Mp}_2(\mathbf Z)\to\operatorname{PSL}_2(\mathbf Z)$.
Choose lifts in $E$ of the two free generators of $F$. Freeness
gives a continuous homomorphism $\iota:F\to E$ lifting the
inclusion into $\widehat B$. For two such homomorphic lifts,
their pointwise ratio is a homomorphism from $F$ to the central
kernel, and therefore vanishes on $F'$. Thus $\iota$ restricted
to $F'$ is independent of these choices. In particular both
R2 frames have a well-defined lifted value, and their images
in $E^{\rm ab}$ are zero because $\iota$ is a homomorphism.
The identity has the same two abelianized images and bit zero,
whereas R2 has bit one. No function of these abelianized
metaplectic frame values can equal $\Delta$.

One may realize the usual metaplectic extension explicitly as
$B_3/\langle z_3^4\rangle$, mapping to
$B_3/\langle z_3\rangle=B$, with central kernel of order four;
its abelianization is $\mathbf Z/24$. The argument above does
not require this presentation. It also shows why arbitrary
pointwise lifts, which can be changed by a central element,
must not be confused with the homomorphic lift just specified.
The scalar eta multiplier and the label permutation $P_u$
are different operations even after this ambiguity is removed.

\begin{statement}{A.31}{Known}{the scalar eta phase and its central table}
\citep{LitAtiyah1987}
\end{statement}
For $M=\left(\begin{smallmatrix}a&b\\c&d\end{smallmatrix}\right)
\in\mathrm{SL}_2(\mathbf Z)$ with $c>0$, the classical formula is
\[
 \eta(M\tau)=
 \exp\!\left(\pi i\left(\frac{a+d}{12c}-s(d,c)-\frac14\right)\right)
 (c\tau+d)^{1/2}\eta(\tau),
\]
with the usual holomorphic square-root branch and
$s(d,c)=\sum_{k=1}^{c-1}((k/c))((dk/c))$.
This is the Dedekind-sum scalar multiplier; the values at the
generators are
\[
 \eta(\tau+1)=\zeta_{24}\eta(\tau),\qquad
 \eta(-1/\tau)=(-i\tau)^{1/2}\eta(\tau).
\]
See Atiyah (1987), equations (1.3)--(1.4), for the logarithmic
version, and Theorem 5.60 for its relation to signature defects,
eta invariants, and the Meyer cocycle. Those comparisons concern
the scalar automorphy factor on the modular group.

If a metaplectic element is written $(M,\varphi)$ with
$\varphi(\tau)^2=c\tau+d$, then
$\eta(M\tau)/(\varphi(\tau)\eta(\tau))$ is its scalar
character $\nu_\eta$. The cocycle law for $\varphi$ proves
multiplicativity directly. Its square descends to a character
$\nu_{\eta^2}:\mathrm{SL}_2(\mathbf Z)\to\mu_{12}$,
with $T\mapsto\zeta_{12}$ and
$S=\left(\begin{smallmatrix}0&-1\\1&0\end{smallmatrix}\right)
\mapsto\zeta_{12}^{-3}$.
The attached exact enumeration constructs a homomorphism
\[
 e:\mathrm{SL}_2(\mathbf Z/12)\longrightarrow\mathbf Z/12,
 \qquad e(T)=1,\quad e(S)=-3.
\]
It enumerates all $1152$ determinant-one matrices and verifies all
$4608$ right-generator edges, proving well-definedness and
multiplicativity. Equality on $S,T$ proves
$\nu_{\eta^2}(M)=\zeta_{12}^{e(M\bmod12)}$. Thus on the scalar
centre the exact table is
\[
\begin{array}{c|rrrr}
 u&1&5&7&11\\\hline
 \nu_{\eta^2}(uI)&1&1&-1&-1\\
 \chi_{12}(u)&1&-1&-1&1
\end{array}
\]
and the first row is $\chi_{-4}(u)$. This proves the proposed
eta-squared central-character calculation, including the distinction
between $5$ and $7$. It does not identify the actual frames with
these central matrices: A.25 refutes that identification. On the
homomorphically lifted R2 frames every scalar character is trivial
by A.30, while the separately defined label permutation has eta
eigenvalue $-1$. The remaining genus-zero/genus-one comparison is
therefore still the one in A.29.

\subsection{Algebraically ample subgroups}


\begin{statement}{A.32}{Known}{AA subgroups lie in character kernels}
\citep[Definition 5.12 and Theorem 5.17(ii)]{LitHoshiMochizukiTsujimura2025}
\end{statement}
In Hoshi--Mochizuki--Tsujimura (2025), ``AA'' means
\emph{algebraically ample}, with Definition 5.12 including the
COF, RGC, and ArBC prerequisites.  The proof of Theorem 5.17(ii)
in the author version dated 14 March 2025, printed page 78,
states that every closed AA subgroup of
$G=\widehat{\mathrm{GT}}$ is contained in a $G$-conjugate of
$G_{\mathbf Q}$; the maximal such subgroups are exactly those
conjugates.  The theorem is not an equality $G=G_{\mathbf Q}$.

Let $\eta:G\to C$ be any continuous homomorphism to an abelian
topological group which is trivial on $G_{\mathbf Q}$.  Then
\[
 \eta(g\sigma g^{-1})
 =\eta(g)\eta(\sigma)\eta(g)^{-1}=1.
\]
Thus every closed AA subgroup lies in $\ker\eta$.  This proves
the requested assertion for $\Delta|_G$ and for any additional
characters of A.3(h) if those characters are actually constructed.
No assertion about the existence of the latter is required.

The marked assertion in R10 also follows.  By
Lochak--Schneps (1997), Proposition 4(ii),
$C_G(c)=\langle c\rangle$ for $c=(-1,1)$.
Suppose $gG_{\mathbf Q}g^{-1}$ contains $c$.  All involutions of
$G_{\mathbf Q}$ are conjugate to complex conjugation, so there
is $h\in G_{\mathbf Q}$ with
$g^{-1}cg=hch^{-1}$.  Hence $gh\in C_G(c)=\langle c\rangle
\subseteq G_{\mathbf Q}$, so $g\in G_{\mathbf Q}$ and
$gG_{\mathbf Q}g^{-1}=G_{\mathbf Q}$.  Consequently
$G_{\mathbf Q}$ is the unique maximal closed AA subgroup
containing the marked $c$.  This proof uses no classification of
involutions of $G$ outside its Galois subgroup.

\subsection{Five-point frames and the harmonic cover}


\begin{statement}{A.33}{Refuted}{arbitrary-multiplier pure frames}
\end{statement}
The assertion that every $\Phi=(\lambda,f)\in\widehat{GT}$ admits
$\alpha_t\in P_5$ with $\Phi(t)=\alpha_t t^\lambda\alpha_t^{-1}$, for
$t=a,b,d$, is false in this normalization.
\emph{Proof.} The standard braid action sends every half twist to a conjugate by a
pure element of its $\lambda$-th power. We use throughout the cocycle
convention of LM7\_PUSH E.65.8, in which the first two generators map to
$\sigma_1^\lambda$ and
$f(\sigma_1^2,\sigma_2^2)^{-1}\sigma_2^\lambda
f(\sigma_1^2,\sigma_2^2)$. Since $\lambda$ is odd,
$\pi_5\Phi=\pi_5$. A pure conjugator would imply
$\pi_5(t)=\pi_5(t)^\lambda$. Since the permutations have orders
$5,4,3$, respectively, this requires $\lambda\equiv1$ modulo the
respective order. Complex conjugation $(-1,1)$ belongs to
$G_{\mathbb Q}\subseteq\widehat{GT}$ and fails all three congruences.
The finite permutation checks are in \texttt{check\_five\_points.py}.
This conclusion concerns the stated genuine action; changing an outer
representative requires specifying its conjugator and changes the displayed
permutation equation.

\begin{statement}{A.34}{Known}{the pentagon has two cocycle classes}
\citep[Theorem 2 and Proposition 3(iii)]{LitLochakSchneps1997}
\end{statement}
Let $\rho$ act on $P_5$ by conjugation by
$(\sigma_4\sigma_3\sigma_2\sigma_1)^2$ and put
$s=x_{34}x_{51}^{-1}x_{45}x_{12}^{-1}$. Lochak--Schneps (1997),
Theorem 2 and Proposition 3(iii), give
\[
 f(x_{12},x_{23})=
 \begin{cases}
 \rho(k)^{-1}k,&\lambda\equiv\pm1\pmod5,\\
 \rho(k)^{-1}s k,&\lambda\equiv\pm2\pmod5.
 \end{cases}
\]
The two rows are distinct classes in $H^1(\langle\rho\rangle,P_5)$.
Consequently A.3(b), if asserted without a multiplier restriction, is false:
choose a Galois element whose cyclotomic character is $2$ modulo $5$.
Existence follows from surjectivity of the cyclotomic character on
$\operatorname{Gal}(\mathbb Q(\zeta_5)/\mathbb Q)$.
For multiplier one the first row applies. The primitive is unique because
$P_5^{\rho}=1$, as shown next. This primitive is a path coordinate; its
identification with a particular conjugator of $a$ includes the base-path
conventions and is not an equality of undefined coordinates.

\begin{statement}{A.35}{Proved}{the three pure centralizers are trivial}
\citep[Theorem B\(^{\prime}\)]{HerfortRibes1985}
\end{statement}
We have
\[
 C_{P_5}(a)=C_{P_5}(b)=C_{P_5}(d)=1.
\]
\emph{Proof.} For $a$, Lochak--Schneps, Proposition 8(iv) and
Proposition 5(c), applied to $P_5\rtimes\langle a\rangle$, imply that
$\langle a\rangle$ is malnormal in that semidirect product; hence its
centralizer intersects $P_5$ trivially. Replacing $a$ by its square makes no
difference because the two elements generate the same order-five group.

For $b,d$, forget a marked point $p$ fixed by the rotation. This is a map
on the subgroup fixing $p$, not on all of $G_5$. The resulting profinite
Birman sequence has kernel
\[
 1\longrightarrow\widehat F_3\longrightarrow
 (G_5)_p\longrightarrow G_4\longrightarrow1.
\]
The pure version is the completion of the familiar split sequence
$1\to F_3\to\Gamma_{0,5}\to\Gamma_{0,4}\to1$.
Exactness follows directly from characteristic finite-index subgroups of the
finitely generated free kernel: they are invariant under the action of the
free quotient, and the induced finite action factors through a finite
quotient. Passing to the finite extension fixing $p$ retains the same
induced topology on the kernel.

The group $G_4$ has the finite normal Klein group $V_4$ and quotient
$\widehat{\operatorname{PSL}_2(\mathbb Z)}=C_2\amalg C_3$.
The pure subgroup maps injectively to the level-two subgroup of this
quotient. The images of $b,d$ in it have orders $2,3$, respectively.
Centralizers of torsion in this free profinite product meet that level-two
free group trivially. Hence an element of $C_{P_5}(b)$ or $C_{P_5}(d)$ lies
in the kernel $\widehat F_3$.

For $b$, realize the remaining four punctures as $\mu_4$ on
$\mathbb P^1$, with base point $p=\infty$, and $b:z\mapsto iz$.
The orbifold quotient of $\mathbb P^1\setminus\mu_4$ by $C_4$ has one
puncture and two cone points of order four. Its discrete orbifold group is
$C_4*C_4$. The split finite covering sequence therefore identifies
$\widehat F_3\rtimes C_4$ with $C_4\amalg C_4$; the splitting at
$\infty$ is a free factor. The centralizer of that factor is the factor
itself, and its intersection with the free kernel is trivial.

For $d$, take the remaining punctures to be $\{0\}\cup\mu_3$, the
base point $p=\infty$, and $d:z\mapsto\zeta_3z$. The quotient has two
punctures and one cone point of order three. Its group is $C_3*\mathbb Z$,
so its profinite completion is $C_3\amalg\widehat{\mathbb Z}$.
The same factor-centralizer argument applies. These free-product
centralizer facts are the standard torsion-centralizer theorem for free
profinite products: Herfort--Ribes (1985), Theorem B$'$ in Section 4,
printed p.~159, gives $C_G(a)=C_A(a)$ for $1\ne a\in A$ in
$G=A\amalg B$. The opened primary text states in its introductory
Theorem B that the centralizer ``is contained in A''; the displayed
equality is its strengthened formulation. Their Theorem A$'$ and
Theorem 1 also provide the torsion and finite-subgroup conjugacy used
in A.37 and B.35.

\begin{statement}{A.36}{Proved}{full centralizers}
\citep[Theorem 3.3]{Collas2012GenusZero}
\end{statement}
The full centralizers are $C_{G_5}(a)=\langle a\rangle$,
$C_{G_5}(b)=\langle b\rangle$ and
$C_{G_5}(d)=\langle d\rangle$.
\emph{Proof.} A.35 makes each centralizer inject into the centralizer of
its permutation. The latter is $C_5$ for $a$ and $C_4$ for $b$, proving
those two assertions. For $d$ it is $C_3\times C_2$, where the extra factor
swaps the two fixed labels. An extra element would therefore give an
involution in $G_5$ with permutation a single transposition. Collas (2012),
Theorem A / Theorem 3.3, makes every profinite involution geometric.
A nontrivial orientation-preserving involution of the sphere has two fixed
points, so on five marked points it must induce two transpositions and one
fixed label. A single transposition with three fixed labels is impossible.

\begin{statement}{A.37}{Proved}{geometric torsion of order four in $G_5$}
\end{statement}
Every element of order four in $G_5$ is conjugate to a discrete element.
\emph{Proof.} The pure group $P_5$ is torsion free, from its free-by-free
profinite Birman sequence. Thus the permutation of an order-four element
has order four and fixes a unique label $p$. Forget $p$. The free kernel is
torsion free, so the image still has order four in $G_4$.
Every finite-order element of $G_4$ is geometric: its image in
$C_2\amalg C_3$ is conjugate into a factor, and the preimage of either
factor under $G_4\to C_2\amalg C_3$ is the corresponding finite discrete
extension by $V_4$ (orders eight and twelve). The kernel itself is the
geometric $V_4$.
After a conjugation and lift, our element therefore maps to a geometric
order-four rotation of four points. Its preimage is the semidirect product
$\widehat F_3\rtimes C_4=C_4\amalg C_4$ from A.35. Every finite subgroup
of this free profinite product is conjugate into a discrete factor. Thus
the original element is conjugate to a discrete element of $\Gamma_{0,[5]}$.
Together with Collas's prime-order theorem this also handles all torsion:
an element of order six would have a permutation of type $(3)(2)$, whose
cube is a single transposition, contradicting the prime-order theorem as
in A.36. The injection of a finite cyclic subgroup into $\mathfrak S_5$
leaves no other orders.

\begin{statement}{A.38}{Proved}{corrected frame theorem}
\end{statement}
For the stated genuine action, a pure frame at $t=a,b,d$ exists exactly
when $\lambda\equiv1\pmod{\operatorname{ord}(t)}$; if it exists it is
unique and continuous in $\Phi$. In particular all three frames exist
uniquely on the multiplier-one subgroup of $\widehat{GT}$.
\emph{Proof.} Necessity is A.33. For $a$, Collas's power-conjugacy theorem
supplies $g\in G_5$ with $\Phi(a)=ga^\lambda g^{-1}$. When
$\lambda\equiv1\pmod5$, the permutation of $g$ centralizes the five-cycle,
so multiplication by a power of $a$ makes $g$ pure.

For $b$, A.37 gives $\Phi(b)=g b^k g^{-1}$, $k=1$ or $3$.
The unique fixed label is the same, so $g$ belongs to the subgroup
$L=(G_5)_p=\widehat{B_4/Z(B_4)}$. Its abelianization is
$\mathbb Z/12$, with every half twist mapping to $1$ and $b$ mapping to $3$.
The GT action on this abelianization is multiplication by $\lambda$.
Consequently $3k=3\lambda\pmod{12}$, so $k\equiv\lambda\pmod4$.
When $\lambda\equiv1\pmod4$, the permutation of $g$ centralizes the
four-cycle and a power of $b$ removes it.

For $d$, use Collas's theorem again:
$\Phi(d)=g d^\lambda g^{-1}$. Its conjugator can permute the two fixed
labels only by swapping them. If it swaps them, multiply it by the
geometric involution $n:z\mapsto1/z$ on
$\{0,\infty\}\cup\mu_3$, which swaps them and inverts $d$.
Then $gn\in L$ and $\Phi(d)=(gn)d^{-\lambda}(gn)^{-1}$.
But $d$ has abelianization $4$ in $L^{\mathrm{ab}}=\mathbb Z/12$.
The last equation would force $4\lambda=-4\lambda\pmod{12}$, which
contradicts invertibility of $\lambda$ modulo three. Thus $g$ fixes both
poles. If $\lambda\equiv1\pmod3$, its permutation is a power of the
three-cycle; multiplication by a power of $d$ makes it pure.
Uniqueness is A.35. The orbit map $P_5\to G_5$,
$u\mapsto utu^{-1}$, is a continuous injection from a compact space into a
Hausdorff space and hence a homeomorphism onto its image. This proves
continuity of the unique frame.

\begin{statement}{A.39}{Proved}{crossed composition law; homomorphisms need more}
\end{statement}
On the subgroup where the relevant pure frame exists, with composition
$\Phi\Psi=\Phi\circ\Psi$, one has
\[
 \alpha_t(\Phi\Psi)=\Phi(\alpha_t(\Psi))\alpha_t(\Phi).
\]
Indeed, apply $\Phi$ to the equality defining $\alpha_t(\Psi)$ and use
uniqueness. Reversing the product gives
$\alpha_t(\Psi\Phi)=\Psi(\alpha_t(\Phi))\alpha_t(\Psi)$.
This is a crossed law, not an assertion that a finite image is a
homomorphism. To obtain an ordinary homomorphism into an abelian target one
must prove both the required invariance of the quotient and triviality of
the induced action after the specified correction. A crossed cocycle from
$C_2$ to $C_3$, with the nontrivial element acting by inversion and cocycle
value $1$, already demonstrates that the crossed law alone is insufficient.

\begin{statement}{A.40}{Proved}{a two-generator presentation}
\end{statement}
Set $h=b^{-1}a$ and $q_i=a^{i-4}ha^{4-i}$ for $1\leq i\leq4$.
A presentation of the discrete group $\Gamma_{0,[5]}$ uses generators
$a,b$ and the following ten relators:
\begin{align*}
 q_iq_{i+1}q_i&=q_{i+1}q_iq_{i+1}&&(1\leq i\leq3),\\
 [q_i,q_j]&=1&&(1\leq i<j\leq4,\ j-i\geq2),\\
 q_4q_3q_2q_1^2q_2q_3q_4&=1,\\
 (q_1q_2q_3q_4)^5&=1,\\
 a&=q_1q_2q_3q_4,\qquad b=q_1q_2q_3.
\end{align*}
This also gives a profinite presentation on taking the closed normal
closure of the relators. To prove it, start from the Artin presentation
with the sphere and full-twist relators. The braid identity
$a\sigma_i a^{-1}=\sigma_{i+1}$ for $i=1,2,3$ and
$b^{-1}a=\sigma_4$ expresses every old generator as $q_i$.
The two last displayed equalities make the maps in both directions
inverse, so this is an explicit Tietze transformation. Also
$d=q_1^2q_2q_3$.
For multiplier one the exact frame constraints are obtained by replacing
$a,b$ in each of these ten relators by
$\alpha_a a\alpha_a^{-1},\alpha_b b\alpha_b^{-1}$ and adjoining
\[
 q_1(\Phi(a),\Phi(b))^2q_2(\Phi(a),\Phi(b))q_3(\Phi(a),\Phi(b))
 =\alpha_d d\alpha_d^{-1}.
\]
These are genuine group equations. Their reduction to a finite quotient
that reads the CM parities is still required before any parity conclusion
can be drawn.

\begin{statement}{A.41}{Proved}{abstract harmonic orbifold cover}
\end{statement}
The map $w=4z(1-z)$ is a degree-two orbifold covering from the sphere with
weights $(3,3,2)$ at $(0,1,\infty)$ to the sphere with weights $(3,2,4)$ at
$(0,1,\infty)$. Above $w=0$ lie $z=0,1$ unramified; above $w=1$ lies
$z=1/2$ with index two; and above infinity lies infinity with index two.
The condition $m_{\rm target}=e\,m_{\rm source}$ holds at every point,
and the orbifold Euler characteristics are $1/6$ and $1/12$.
Thus the abstract covering asserted in Lead 6 is correct.
It must not be identified without a comparison with a map of the
principal modular covers $X(3),X(4)$ over the same $j$-line.
Neither principal subgroup contains the other: $T^3$ belongs to
$\Gamma(3)$ but not $\Gamma(4)$, and $T^4$ belongs to
$\Gamma(4)$ but not $\Gamma(3)$. Both are normal in the modular group,
so conjugating them in that group cannot repair the inclusion.
The subgroup $A_4\subset S_4$ gives the abstract spherical-orbifold map;
that fact alone does not compare the two congruence representations or
prove equality of the CM parities.

\begin{statement}{A.42}{Known}{$X_0(6)$ and $V_4$}
\citep{Maier2009,Colo2022}
\end{statement}
A rational Hauptmodul is
\[
 h(\tau)=\frac{\eta(\tau)^5\eta(3\tau)}
 {\eta(2\tau)\eta(6\tau)^5},\qquad u=h+9.
\]
Its four cusp values in the $u$ coordinate are $0,1,9,\infty$.
Thus $Y_0(6)\simeq\mathbb P^1\setminus\{0,1,9,\infty\}$ over
$\mathbb Q$. The Atkin--Lehner involutions are
\[
 w_2(u)=\frac{u-9}{u-1},\qquad
 w_3(u)=\frac9u,\qquad
 w_6(u)=\frac{9(u-1)}{u-9}.
\]
They square to the identity, $w_2w_3=w_3w_2=w_6$, and act on the ordered
cusp list $(0,1,9,\infty)$ as
\[
 (0\ 9)(1\ \infty),\quad (0\ \infty)(1\ 9),\quad
 (0\ 1)(9\ \infty),
\]
respectively. All these identities are checked exactly in the program.
For the modular identification, Maier (2009), Table 2, gives the four
cusp values $0,-8,-9,\infty$ for its reciprocal-scaled Hauptmodul;
Fricke exchanges that coordinate with $h$. Col\`o (2022), Appendix A,
printed p.~245, gives the displayed $h$ and all three involutions directly.
Applying $u=h+9$ gives the formulas above.

\begin{statement}{A.43}{Proved}{Proved correction; comparison remains open}
\end{statement}
The geometric fibre of
$\mathcal M_{0,5}\to\mathcal M_{0,4}$ over the ordered point $9$ is
$\mathbb P^1\setminus\{0,1,9,\infty\}$ and its completed fundamental
group is $\widehat F_3$. This yields the abstract identification with the
completion of $\Gamma_0(6)/\{\pm I\}$. Both descriptions are over
$\mathbb Q$, so the arithmetic fundamental-group identifications respect
the Galois outer action after the chosen geometric base points and paths
are transported. There is no forgetting map from the full unordered
five-point moduli stack to the full unordered four-point stack without
first choosing the point to forget. The group map is on a point stabilizer,
or on the pure group.

The pentagon gives the GT action on the genus-zero tower and compatibility
with forgetting. To identify its action on this fibre with its restriction
from the modular group, one must additionally compare the path or
arithmetic section at $9$. An abstract isomorphism of the free fibre groups
and the identities among the three rational involutions do not establish
that comparison. The unresolved equality is equality, in
$\operatorname{Out}(\widehat F_3)$, of the modular restriction and the
transported fibre action for every $\Phi\in\widehat{GT}$, with the
transport path/section specified. No vanishing of $\Delta$ follows from
A.42 alone.

\subsection{Atkin--Lehner defects and finite five-point targets}


\begin{statement}{A.44}{Proved}{the level-six Atkin--Lehner defect is well typed}
\end{statement}
Let $H$ be the closure in $\widehat B$ of
$\Gamma_0(6)/\{\pm I\}$. Subject to the supplied R3 and the standard
arithmetic action on the modular cover $Y_0(6)$, every
$\Phi\in\widehat{GT}_0$ preserves the conjugacy class of $H$.
For multiplier one this follows by reducing the covariance in R3 modulo
six. For arbitrary multiplier choose a Galois element with that multiplier
and write $\Phi$ as a multiplier-one element times it; the Galois element
preserves the subgroup's conjugacy class since $Y_0(6)$ and its map to the
modular orbifold are defined over $\mathbb Q$.

The subgroup $H$ is self-normalizing. Indeed,
$\operatorname{PSL}_2(\mathbb Z/6)\simeq S_3\times A_4$, and the image
of $H$ is $C_2\times C_3$, the product of upper-triangular subgroups.
In $S_3$ the order-two subgroup is self-normalizing; in $A_4$ the
order-three subgroup is self-normalizing. Pullback through the finite
quotient proves the claim. The exact finite enumeration is recorded in
\texttt{check\_normalizers.py}.
Therefore correcting a representative of the outer action to preserve
$H$ produces a well-defined continuous homomorphism
\[
 \rho_6:\widehat{GT}_0\longrightarrow\operatorname{Out}(H).
\]
Two corrections differ by an element of $N_{\widehat B}(H)=H$, so give
the same outer automorphism of $H$. Composition proves that this is a
homomorphism.
For the rational involution $w_e$ of $Y_0(6)$, put
\[
 D_e(\Phi)=\rho_6(\Phi)w_e\rho_6(\Phi)^{-1}w_e^{-1}
 \quad\text{in }\operatorname{Out}(H),\qquad e=2,3,6.
\]
This is a continuous map, trivial on $G_{\mathbb Q}$ because $w_e$ is
rational. Its composition law is
\[
 D_e(\Phi\Psi)=
 \rho_6(\Phi)D_e(\Psi)\rho_6(\Phi)^{-1}D_e(\Phi).
\]
In particular it is a crossed defect, and calling it a finite-order
abelian character requires an additional theorem. Its trivial locus is
the inverse image of the centralizer of $w_e$. Neither
$D_3=1\Longleftrightarrow\Delta=0$ nor $D_6=1$ for all
$\widehat{GT}_0$ is proved by this typing and well-definedness argument.

\begin{statement}{A.45}{Proved}{cusps alone cannot detect the proposed defect}
\end{statement}
Every $\rho_6(\Phi)$ fixes each of the four cusps of $Y_0(6)$.
The cusp width associated to a double coset
$H g\overline{\langle\tau\rangle}$ is
\[
 [\overline{\langle\tau\rangle}:
  g^{-1}Hg\cap\overline{\langle\tau\rangle}].
\]
It is unchanged by an automorphism preserving the conjugacy class of $H$
and the conjugacy class of the procyclic cusp group, with a unit power on
that group. The four widths are $1,2,3,6$, all distinct; hence no two
cusps can be exchanged. As a result every $D_e(\Phi)$ induces the identity
permutation of the four cusps. Even if the element of R2 has
$D_3\ne1$, that nontriviality must be checked on the free fundamental
group or a stronger quotient; the permutation of the four cusps cannot
witness it.

\begin{statement}{A.46}{Proved}{an order-five norm-one target}
\end{statement}
Let $R_N=(\mathbb Z/N)[z]/(z^4+z^3+z^2+z+1)$, with involution
$\bar z=z^{-1}$, and put
\[
 A=\begin{pmatrix}0&0&0&-1\\1&0&0&-1\\0&1&0&-1\\0&0&1&-1\end{pmatrix},
 \qquad
 J=\begin{pmatrix}0&1&0&0\\-1&0&1&0\\0&-1&0&1\\0&0&-1&0\end{pmatrix}.
\]
Then $A^5=I$, $\det J=1$ and $A^{\mathsf T}JA=J$ over the integers.
For every $N$ its centralizer in the group preserving $J$ is canonically
\[
 T_N=\{u\in R_N^{\times}:u\bar u=1\}.
\]
Indeed $e_1,Ae_1,A^2e_1,A^3e_1$ is a basis over $\mathbb Z/N$.
A matrix commuting with $A$ is therefore uniquely a polynomial in $A$ of
degree at most three. Its adjoint for $J$ is obtained by
$A\mapsto A^{-1}$, so the symplectic equation is exactly $u\bar u=1$;
this equation already gives the inverse and no separate unit test is
needed. This argument also works when $N$ is divisible by five.

At $N=12$, exact enumeration of $12^4$ polynomial residues gives
$|T_{12}|=200$, $|T_{12}[2]|=8$, and
$|T_{12}/\{\pm1\}|=100$ with two-torsion of order four. One can also
read these orders from CRT: $R_3=\mathbb F_{81}$ and its norm-one group
has order ten; $R_4$ is the unramified degree-four extension of
$\mathbb Z/4$, whose norm-one group has order twenty. The latter has
principal-unit kernel $\mathbb F_4$ under reduction modulo two.
The four scalar elements $1,5,7,11$ all occur in $T_{12}$.
This proves the advertised form of a possible finite linear target,
but supplies no functorial homomorphism from the five-point mapping class
group to it and no comparison with the actual frame. In particular it
cannot turn a bare five-point primitive into a homomorphism to $T_{12}$.
The enumeration is \texttt{check\_order5\_torus.py} and its output.

\begin{statement}{A.47}{Proved}{the exact forgotten-point frame map}
\end{statement}
Use the genuine braid action in LM7\_PUSH E.65.8's convention. Let
$L=(G_5)_5$, and let
\[
 q:L\longrightarrow G_4\longrightarrow\widehat B
\]
forget label five and then divide by the geometric Klein four-group.
Writing $u,v$ for the standard braid generators of $\widehat B_3$ modulo
its central full twist, this map is given by
\[
 q(\sigma_1)=u,\qquad q(\sigma_2)=v,\qquad q(\sigma_3)=u.
\]
It follows that
\[
 q(b)=uvu=s,\qquad q(d)=u^2vu=u s=r^{-1},\qquad
 r=uv,\quad sr=u^{-1}.
\]
In particular $x=(sr)^{-2}=u^2$ and $y=sxs=v^2$, exactly the positive
free generators of R3 and R4. The restriction $q:P_5\to F$ is a
continuous surjection, since the pure elements $\sigma_1^2,\sigma_2^2$
map to $x,y$.

The GT tower action respects forgetting. Its action on the quotient
$\widehat B$ is determined by $u,v$ and is
\[
 u\longmapsto u^\lambda,\qquad
 v\longmapsto f(u^2,v^2)^{-1}v^\lambda f(u^2,v^2).
\]
This is the precise action used by E.65.8; at multiplier one it fixes
$sr=u^{-1}$ literally. The Klein subgroup is preserved: the GT action
preserves the permutation map; a torsion element of order two with
permutation a double transposition has trivial image in
$\widehat B=C_2\amalg C_3$, since any nontrivial order-two element there
has nontrivial image in $S_3$. Thus its image remains in $V_4$.
Applying $q$ to the unique frame equations and using the uniqueness of
pure $s,r$ frames in $\widehat B$ gives the exact identities
\[
 \boxed{\ \beta_\Phi=q(\alpha_b(\Phi)),\qquad
 A_\Phi=q(\alpha_d(\Phi))\ }
 \quad(\lambda(\Phi)=1).
\]
The second identity holds because a conjugator of $r^{-1}$ is the same
conjugator of $r$. Hence the CM parities are genuinely read from the
five-point frames after the already specified shear:
\[
 W_\Phi=q_{12}(q(\alpha_b(\Phi)))^{-1}U^{8\rho_\Phi},\qquad
 T_\Phi=q_{12}(q(\alpha_d(\Phi)))^{-1}U^{8\rho_\Phi},
\]
where $(q(\alpha_b(\Phi))^{-1})^{\mathrm{ab}}
=(\rho_\Phi,\rho_\Phi)$. No section at $9$ is needed for these
forgetful identities. The section-comparison gap in A.43 belongs to
the Atkin--Lehner realization of the fibre.

\begin{statement}{A.48}{Refuted}{descent of the usual genus-two linear representation}
\end{statement}
The usual genus-two symplectic representation of the five-strand braid
group does not descend to $G_5$ modulo twelve, even after quotienting by
the global scalar sign. Let $J_0$ be the block-diagonal alternating form
with two blocks $\left(\begin{smallmatrix}0&1\\-1&0\end{smallmatrix}\right)$.
For
\[
 v_1=(1,0,0,0)^t,\quad v_2=(0,1,0,0)^t,\quad
 v_3=(1,0,1,0)^t,\quad v_4=(0,0,0,1)^t
\]
put $T_i=I+v_i v_i^tJ_0$. Adjacent transvections satisfy the braid
relation and nonadjacent ones commute, so $\sigma_i\mapsto T_i$
defines a representation of $B_5$ in $\operatorname{Sp}_4(\mathbb Z)$.
The full twist maps to $-I$, but the sphere relator has image
\[
 T_4T_3T_2T_1^2T_2T_3T_4
 =\begin{pmatrix}
 -1&0&0&0\\0&-1&0&0\\0&0&-1&2\\0&0&0&-1
 \end{pmatrix}.
\]
This is not scalar modulo three or twelve. Thus it is not the identity
in either the quotient by $\{\pm I\}$ or by the full scalar center.
The relation required for the five-point sphere is exactly where this
candidate linear comparison fails. It is a representation of the
appropriate six-point hyperelliptic setting with one distinguished point,
and forgetting that point cannot be done silently. This counterexample
refutes this particular proposed genus-two comparison, not the existence
of every conceivable representation of $G_5$ in a symplectic group.

\begin{statement}{A.49}{Proved}{finite frame ambiguity in the natural five-block target}
\end{statement}
The signed homomorphism $L\to\operatorname{SL}_2(\mathbb Z)$ sending
$\sigma_1,\sigma_3$ to $U$ and $\sigma_2$ to $L_0=
\left(\begin{smallmatrix}1&0\\-1&1\end{smallmatrix}\right)$ is well defined,
since the central full twist of $B_4$ maps to $(UL_0U)^4=I$.
Conjugate the point stabilizer to obtain the analogous representation for
each of the five forgotten labels. On $P_5$, reduction modulo three gives
five homomorphisms $\rho_i:P_5\to\operatorname{SL}_2(\mathbb F_3)$.
Each is onto. Let $H$ be their joint finite image. Then
\[
 H=\{(g_i)\in\operatorname{SL}_2(\mathbb F_3)^5:
          \textstyle\sum_i\epsilon(g_i)=0\},\qquad
 H'=Q_8^5,
\]
where $\epsilon:\operatorname{SL}_2(\mathbb F_3)\to C_3$ sends
$U^2$ to $1$ and has quaternion kernel. In particular
$|H|=8^5\cdot3^4=2654208$.

Here is a proof that avoids enumerating pentagon shadows. For a positive
pure half-twist square $x_{jk}$, the $i$th coordinate of its abelian image
is zero when $i\in\{j,k\}$ and one otherwise. The former follows because
forgetting one of the colliding labels makes its projective image trivial,
so its signed image is central. The latter follows because the twist is
conjugate inside the $i$th stabilizer to $\sigma_1^2$.
The ten vectors with three ones span exactly
$C=\{(c_i)\in\mathbb F_3^5:\sum c_i=0\}$: their pairwise differences
include every $e_i-e_j$.

Modulo the five scalar centers, each quaternion factor becomes the
irreducible two-dimensional $\mathbb F_2$ module $V_i$ for its coordinate
$C_3$. The characters $c\mapsto c_i$ on $C$ are pairwise distinct even
up to inversion, so these five modules are pairwise nonisomorphic.
The image of $H'$ in $\bigoplus V_i$ is a $C$-submodule and projects onto
every $V_i$, since each projection of $H$ is the full
$\operatorname{SL}_2(\mathbb F_3)$. Maschke's theorem applies because
$|C|$ is odd. Multiplicity one then makes that image all of
$\bigoplus V_i$. Choose two elements of $H'$ whose images are a basis
of a single $V_i$ and zero in the other four factors. Their commutator is
$-I$ in factor $i$ and $I$ elsewhere. Thus $H'$ contains all five scalar
centers as well, proving $H'=Q_8^5$ and the asserted formula for $H$.

Induce the signed stabilizer representation to the full group $G_5$.
This gives a genuine finite five-block representation modulo twelve.
An element of $P_5$ acts block diagonally, and $b$ fixes the block with
forgotten label five, acting there by $S=UL_0U$.
The image of $P_5'$ is trivial modulo four in every block: a pure element
maps into the kernel of
$\operatorname{SL}_2(\mathbb Z/4)\to\operatorname{SL}_2(\mathbb F_2)$,
which is abelian. Combining this with $H'=Q_8^5$ and CRT shows that the
image of $P_5'$ contains the element supported only in the fifth block
with value
\[
 \widetilde S=\begin{pmatrix}9&4\\8&9\end{pmatrix}\pmod{12}.
\]
It is $I$ modulo four and $S$ modulo three, and it commutes with the
finite image of $b$. Multiplying a finite $b$-frame by this centralizer
element changes $\epsilon_2(W)$ by one without changing the frame
conjugacy equation. Similarly the element supported in the fifth block
with value $5I$ centralizes the finite image of $d$ and changes
$\epsilon_3(T)$ by one. Both toggles have lifts in $P_5'$, so their
forgotten-point abelian coordinates and hence the shear coordinate are
zero.

Thus even the identity automorphism of this quotient admits solutions of
the finite frame equations with either apparent CM parity. The
presentation equations of A.40 cannot disambiguate these finite primitive
choices, because the images of $a,b,d$ under the represented automorphism
have not changed. This is a precise failure of this proposed finite
reading of frames. It is not a mixed-sign element of $\widehat{GT}$, not
an assertion about extendability of the finite primitive choices, and not
a disproof that $\Delta$ factors through some action quotient with the
correct extendability condition. A finer finite quotient in which the
images of actual global primitives are separated is still needed.
The exact small group and module checks, including the two CRT matrices,
are recorded in \texttt{check\_five\_block\_modules.py} and its output.

\subsection{The restricted pentagon lifting problem}


\begin{statement}{A.50}{Proved}{relative products preserve a specified pentagon equation}
\end{statement}
Let $P_4=\widehat\pi_1(M_{0,4})=F$ and
$P_5=\widehat\pi_1(M_{0,5})$. Fix the five continuous
substitution maps $e_j:F\to P_5$ occurring in the group
pentagon, with all paths and product order fixed. Let
$N_5\triangleleft P_5$ be open, and let $N_4\triangleleft F$
be an open subgroup satisfying $e_j(N_4)\subseteq N_5$ for
all five maps. Such $N_4$ exists: intersect the five inverse
images and take a characteristic open subgroup contained in
that intersection. Write
\[
 D_i=P_i/N_i,\quad P_{i,p}=P_i/O^p(N_i),\quad
 P_{i,\rm nil}=P_i/R_{\rm nil}(N_i).
\]
Every $e_j$ induces continuous maps in these three categories,
and the actual quotient isomorphisms
\[
 P_{i,\rm nil}\simeq\prod_{p,D_i}P_{i,p}
\]
intertwine them. Consequently a parameter in the relative
pronilpotent source satisfies the pentagon there if and only if
all its relative pro-$p$ components satisfy the corresponding
pentagon, with their common finite-base coordinates retained.

\emph{Proof.}
For any homomorphism $e:N_4\to N_5$, composition with a finite
$p$-group quotient of $N_5$ proves
$e(O^p(N_4))\subseteq O^p(N_5)$; the same argument for finite
nilpotent quotients gives the analogous residual inclusion.
This proves the induced maps. The product comparison for any
profinite group with an open normal subgroup follows from
\[
 N/R_{\rm nil}(N)\simeq\prod_pN^{(p)}:
\]
finite nilpotent groups are the products of their Sylow groups,
and coprime Goursat together with compactness gives surjectivity.
For the extension over $D_i$, lift the common $D_i$ coordinate
once, then correct its tuple by an element of the displayed
kernel product. This proves the stated isomorphism without a
chosen splitting.

Word evaluation commutes with every continuous homomorphism.
If $\mathcal P(f)$ denotes the ordered pentagon defect, its
projection to each component is therefore exactly
$\mathcal P_p(f_p)$. The product is injective, so
$\mathcal P(f)=1$ in $P_{5,\rm nil}$ if and only if all these
components are one. The same observation applies to additional
two-cycle and three-cycle equations, provided their substitution
maps are included when choosing compatible finite bases.
Automorphisms of a whole quotient diagram also glue
componentwise when they induce one common finite-base diagram
automorphism. This is a statement about the diagram quotient;
it does not assert that arbitrary lists of four-point
automorphisms already define a five-point diagram automorphism.
\hfill$\square$

\begin{statement}{A.51}{Open}{lifting the restricted solution space}
\end{statement}
The identity in A.50 does not supply the epimorphism
\[
\begin{gathered}
 \{(1,f)\in\widehat{\mathrm{GT}}:\text{specified retained data}\}\\
 \longrightarrow
 \{\text{compatible relative pentagon solutions}\}.
\end{gathered}
\]
Neither R1 nor R2 states this epimorphism. Their proven
finite transport theorem retains a triple of peripheral
conjugators and the one peripheral product relation in a free
pants or modular source. Adding the pentagon requires a new
finite transport theorem for the five correlated substitutions
of the same parameter $f$, or an equivalent theorem for
automorphisms of the compatible four- and five-point diagram.
No such theorem follows from the supplied hypotheses alone.

\emph{Explanation with the exact logical boundary.}
The count of D.50.11 treats independent boundary values in
specified classes and extracts the coefficient of their product
being one. Its surface extension D.50.12 does the same for
\[
 \prod_{i=1}^g[a_i,b_i]\,c_1\cdots c_b=1.
\]
The pentagon instead requires an ordered relation between
$e_1(f),\ldots,e_5(f)$, all evaluations of one unknown $f$.
These five values are correlated by the substitution maps.
The independent-boundary character sum cannot silently be
restricted to that correlated subset; the Schur scalarization
and subsequent Mobius inversion count different sets.

There is also a distinct source issue. A five-punctured sphere
has free fundamental group of rank four. The group
$\pi_1(M_{0,5})$ appearing in the pentagon is the fundamental
group of a configuration space, with the split forgetting
sequence $F_3\rtimes F_2$ and abelianization of rank five.
Thus applying the genus-zero, five-boundary formula is not
applying a presentation of $P_5$. One cannot replace the
configuration group by the fiber curve merely because both
involve five marked points.

Even lifting the two free groups separately in the forgetting
sequence leaves the essential compatibility condition. If
$\rho:F_2\to\operatorname{Aut}(F_3)$ is its monodromy action,
candidate lifts $\beta,\gamma$ must satisfy
\[
 \beta\rho(b)\beta^{-1}
   =\operatorname{Inn}(h_b)\rho(\gamma(b))
       \quad(b\in F_2),
\]
with the appropriate continuous cocycle law for $h_b$.
R1 retains peripheral conjugators, but imposes no such
monodromy compatibility. B.21 provides an independent example
of precisely why an extra commutation requirement can obstruct
a lift even when a quotient automorphism is perfectly valid.

P.11 gives actual full gentle elements failing the pentagon,
so satisfying all the peripheral and generation conditions
used by R1 does not imply that additional relation. This is
an explicit obstruction to that inference, not a counterexample
to the still stronger surjectivity assertion displayed above.
In particular nothing here proves that the R2 mixed-sign
element itself fails the pentagon. Its absolute pro-prime
images are identity, so the degree-seven and degree-nine
detectors do not decide that element.

\begin{statement}{A.52}{Proved}{the exact compactness criterion}
\end{statement}
Let $\mathcal C=\{\Phi\in\widehat{\mathrm{GT}}_0:
\Delta(\Phi)=1\}$, and let
$\mathcal P:\widehat{\mathrm{GT}}_0\to P_5$ be the continuous
pentagon-defect map in the fixed convention. For each open
normal $J\triangleleft P_5$, set
\[
 \mathcal C_J=
 \{\Phi\in\mathcal C:\mathcal P(\Phi)\in J\}.
\]
Then
\[
 \widehat{\mathrm{GT}}\cap\mathcal C\ne\varnothing
 \quad\Longleftrightarrow\quad
 \mathcal C_J\ne\varnothing\ \hbox{for every open normal }J.
\]
Equivalently, vanishing of $\Delta$ on the full group is
witnessed by the emptiness of one exact finite pentagon test
on actual gentle residues. Arbitrary finite candidates that
have not been shown to lift to the gentle group are not such
residues.

\emph{Proof.}
The sets $\mathcal C_J$ are closed in the compact group
$\widehat{\mathrm{GT}}_0$. They are directed under finite
intersection because
$\mathcal C_{J_1\cap J_2}=\mathcal C_{J_1}\cap\mathcal C_{J_2}$.
If each is nonempty, compactness supplies an element of their
intersection. Its defect lies in every open normal subgroup
of $P_5$, hence is one. The reverse direction is immediate.
For a fixed $J$, the map $\mathcal P\bmod J$ and $\Delta$
are continuous finite-valued maps, so both factor through
one finite quotient of the compact profinite parameter group.
This proves the finite formulation. If one instead uses raw
finite pairs, the compactness argument must also impose all
finite automorphism and gentle relations simultaneously and
prove nonemptiness after every finite common refinement.
\hfill$\square$

\begin{statement}{A.53}{Open}{Scope of the present result}
\end{statement}
The underlying relative group product holds at every fixed
compatible finite base, including mixed two- and three-primary
bases; A.50 proves the coordinatewise assertion with the
pentagon. A smallest five-point level detecting $\Delta$ has
not been identified here, and no surviving shadow proves a
full mixed-sign element. The unresolved mathematical statement
is the restricted lifting theorem of A.51, or an independent
finite obstruction as in A.52. Calling the product decomposition
itself false would overstate what has been established.

\subsection{Finite quotients that detect actual frames}


\begin{statement}{A.54}{Proved}{a constructive finite quotient that separates frames}
\end{statement}
Let $B^\wedge=C_2\amalg C_3$, with generators $s,r$, and let
$q:B^\wedge\twoheadrightarrow Q$ be a continuous finite quotient in which
$s,r$ retain their respective orders. Write $H_2=\langle q(s)\rangle$,
$H_3=\langle q(r)\rangle$, and form the finite left permutation modules
\[
 M_2=\mathbb F_2[Q/H_2],\qquad M_3=\mathbb F_3[Q/H_3].
\]
Put $E=(M_2\oplus M_3)\rtimes Q$, with its ordinary semidirect-product
multiplication. If $e_2,e_3$ denote the basis vectors at the identity
cosets, define
\[
 S=(e_2,0,q(s)),\qquad R=(0,e_3,q(r)).
\]
Since the relevant cyclic subgroup fixes its identity coset, $S^2=R^3=1$.
The universal property of the free profinite product gives a continuous
homomorphism $\theta:B^\wedge\to E$ lifting $q$. Its finite image
$E_0=\langle S,R\rangle$ satisfies
\[
 \operatorname{pr}_Q C_{E_0}(S)=H_2,\qquad
 \operatorname{pr}_Q C_{E_0}(R)=H_3.
\]

We prove the first identity; the second is identical with three in place
of two. If $(m_2,m_3,g)$ centralizes $S$, then $g$ centralizes $q(s)$ and
\[
 (1-q(s))m_2=(1-g)e_2.
\]
Because $g$ centralizes $q(s)$, both cosets $H_2$ and $gH_2$ are fixed by
$q(s)$. In the coinvariants $M_2/(1-q(s))M_2$, basis vectors are indexed
by orbits of $\langle q(s)\rangle$ on $Q/H_2$. These two fixed cosets
therefore give distinct, independent basis vectors unless $gH_2=H_2$.
The displayed equation forces $g\in H_2$. Conversely $S^k$ centralizes
$S$ and projects to every element of $H_2$.

Take now $Q=\operatorname{SL}_2(\mathbb Z/12)/\{\pm I\}$, of order
$576$, and let $F$ be the level-two subgroup. If two elements of $E$ with
$Q$-projections in $q(F)$ conjugate $S$ to the same element, their
$Q$-projections are equal. Indeed their ratio centralizes $S$, its
projection lies in $H_2$, and
$H_2\cap q(F)=1$, since the map $Q\to S_3$ is injective on $H_2$.
The analogous assertion holds for $R$. Thus these two lifted finite
conjugacy equations determine the full $Q$-residues of the actual pure
frames, not merely their conjugacy classes in $Q$.

This construction has ambient order
$576\cdot2^{288}\cdot3^{192}$. It is not claimed minimal.
The exact program checks all $576$ elements of $Q$, the cyclic cosets,
and the fixed-coset obstructions. In $Q$ the centralizers of $q(s),q(r)$
have orders $16,18$; the construction excludes respectively $14,15$
unwanted centralizer residues. The large module group need not be
explicitly enumerated to verify this proof.

\begin{statement}{A.55}{Proved}{characteristic detector and evaluation on the multiplier-one group}
\end{statement}
There is an explicitly defined finite characteristic quotient of
$B^\wedge$ through which the genuine multiplier-one action determines
$\delta$ of R4. One construction is
\[
 K_* =\bigcap_{h\in\operatorname{Hom}_{\rm cont}(B^\wedge,E)}\ker h,
 \qquad Q_*=B^\wedge/K_*.
\]
The indexing set is finite: its elements are exactly pairs of elements
of $E$ whose squares and cubes, respectively, are one. Thus $K_*$ is
open, normal and characteristic; precomposition by an automorphism
permutes the finite set of homomorphisms. It lies in $\ker\theta$,
so $Q_*$ maps to $E_0$. Equivalently one may use the characteristic core
of $\ker\theta$, which is open because a finitely generated profinite
group has only finitely many subgroups of a specified finite index.

We use the \emph{genuine} action selected by E.65.8,
$u\mapsto u^\lambda$ and
$v\mapsto f(u^2,v^2)^{-1}v^\lambda f(u^2,v^2)$, with
$s=uvu$, $r=uv$, $sr=u^{-1}$.
This is an explicit representative, not an unspecified outer class.
Let $\rho:\widehat{GT}_0\to\operatorname{Aut}(Q_*)$ be the resulting
continuous homomorphism. For $\Phi\in\mathcal H$, project
$\rho(\Phi)(s),\rho(\Phi)(r)$ into $E_0$. Solve the two finite
conjugacy equations there, restricting the conjugators to have
$Q$-projection in $q(F)$. Actual pure frames guarantee existence.
A.54 guarantees that the two $Q$-projections obtained are unique and
are $q(\beta_\Phi),q(A_\Phi)$.

Their signed lifts in the specified image of $F$ in
$\operatorname{SL}_2(\mathbb Z/12)$ are unique: modulo four that image
is the displayed $V_4$ of R3, which does not contain $-I$.
For the actual multiplier-one frames the lift is $I$ modulo four.
Write $\epsilon:\operatorname{SL}_2(\mathbb F_3)\to\mathbb F_3$
for the abelianization with $\epsilon(U^2)=\epsilon(L_0^2)=1$.
Since $\beta_\Phi^{\rm ab}=(-\rho_\Phi,-\rho_\Phi)$,
\[
 \epsilon(q_3(\beta_\Phi))=-2\rho_\Phi=\rho_\Phi\pmod3.
\]
Consequently its residue already determines the shear $U^{8\rho_\Phi}$
modulo twelve. We can compute $T,W$ and their two parities by R3--R4.
This defines a well-defined character
$d:\rho(\mathcal H)\to\mathbb F_2$ with
$\delta=d\circ\rho|_{\mathcal H}$.
It is not asserted that every automorphism of $Q_*$ lies in
$\rho(\mathcal H)$; the detector evaluates actual action residues,
and finite admissibility is not a lifting theorem.

\begin{statement}{A.56}{Proved}{reduction for arbitrary multipliers; arithmetic table not computed}
\end{statement}
Retain $\rho$ and $d$ from A.55, and put
\[
 J=\chi\bigl(G_{\mathbb Q}\cap\ker\rho\bigr)
       \subset\widehat{\mathbb Z}^{\times}.
\]
This is open: $G_{\mathbb Q}\cap\ker\rho$ is open, and the surjective
continuous cyclotomic homomorphism of profinite groups is open.
Choose a positive integer $M$ such that
$\ker(\widehat{\mathbb Z}^{\times}\to(\mathbb Z/M)^\times)\subset J$.
For each $a\in(\mathbb Z/M)^\times$, choose a Galois element $\sigma_a$
whose multiplier has residue $a$; retain only the finite table
$a\mapsto\rho(\sigma_a)$. Then for every
$\Phi\in\widehat{GT}_0$,
\[
 \Delta(\Phi)=d\bigl(\rho(\Phi)\rho(\sigma_a)^{-1}\bigr),
 \qquad a=\lambda(\Phi)\bmod M.
\]
The argument of $d$ lies in $\rho(\mathcal H)$ even though
$\Phi\sigma_a^{-1}$ need not have multiplier exactly one.
Indeed its multiplier lies in $J$, so choose
$\sigma_0\in G_{\mathbb Q}\cap\ker\rho$ with that exact multiplier.
Then $h=\Phi\sigma_a^{-1}\sigma_0^{-1}$ lies in $\mathcal H$ and
$\rho(h)=\rho(\Phi)\rho(\sigma_a)^{-1}$. R4 gives
$\Delta(\Phi)=\delta(h)$. This also proves independence of the choices.

Adjoin a characteristic finite quotient covering the congruence quotient
modulo $M$. Since $sr$ has order $M$ in that congruence quotient and the
genuine action sends it to its $\lambda$-th power, the enlarged action
quotient records $\lambda\bmod M$ as well. Thus it detects the full
character. This is an exact reduction; neither $M$, the finite Galois
table, the smallest detector, nor all pentagon residues at that detector
have been computed here. Equivalently $M$ may be chosen as a cyclotomic
conductor for the abelian intersection of the finite Galois extension
cut out by $\rho|_{G_{\mathbb Q}}$.

\begin{statement}{A.57}{Proved}{the explicit order-four and order-three cocycles}
\end{statement}
Return to the global notation of A.33, and write
$u_i=\sigma_i$ for $1\leq i\leq4$ and
$z_i=(u_1\cdots u_{i-1})^i$ for the full twist on the first $i$
strands. For $\Phi=(1,f)\in\widehat{GT}$ put
\[
 F_i=f(z_i,u_i^2)\quad(2\leq i\leq4),\qquad F_1=1.
\]
All these words are read in the specified sphere quotient of the braid
group. The formula with $z_i$ agrees with the usual Drinfeld first
argument: their ratio is $z_{i-1}$, which commutes with that argument
and with $u_i^2$, so it can be inserted in a commutator word $f$.
In the convention of E.65.8,
\[
 \begin{split}
 E_b&=u_1F_2^{-1}u_2F_2F_3^{-1}u_3F_3
                         u_3^{-1}u_2^{-1}u_1^{-1},\\
 E_d&=u_1^2F_2^{-1}u_2F_2F_3^{-1}u_3F_3
                         u_3^{-1}u_2^{-1}u_1^{-2},\\
 E_a&=\left(\prod_{i=1}^{4}F_i^{-1}u_iF_i\right)a^{-1}.
 \end{split}
 \tag{FP.27.1}
\]
These are exactly $E_t=\Phi(t)t^{-1}\in P_5$.
Let $\vartheta_t$ be conjugation by $t$, of order $m=5,4,3$ respectively.
The defining torsion relation gives the nonabelian cocycle condition
\[
 E_t\,\vartheta_t(E_t)\cdots\vartheta_t^{m-1}(E_t)=1.
 \tag{FP.27.2}
\]
The unique pure frames of A.38 are the unique primitives, in the
explicit convention
\[
 E_t=\alpha_t\vartheta_t(\alpha_t^{-1}).
 \tag{FP.27.3}
\]
Indeed $(E_tt)^m=\Phi(t)^m=1$ proves (FP.27.2), and multiplying the
frame equation by $t^{-1}$ proves (FP.27.3). Conversely (FP.27.3)
recovers that frame equation. This is an explicit answer for the
order-four and order-three primitive cocycles; it does not identify an
arbitrary-multiplier Lochak--Schneps path primitive with an unspecified
torsion frame. The latter assertion was refuted in A.33--A.34.

There is also an exact relation between the two local frames which is
already visible before the fifth generator is used. Since $d=u_1b$ and
$\Phi(u_1)=u_1$,
\[
 \alpha_d d\alpha_d^{-1}
       =u_1\alpha_b b\alpha_b^{-1}.
 \tag{FP.27.4}
\]
If $h=\alpha_b^{-1}\alpha_d$, this is equivalently
\[
 b h b^{-1}=\alpha_b^{-1}u_1^{-1}\alpha_b h u_1.
 \tag{FP.27.5}
\]
Under the exact forgetful map of A.47 it becomes the corresponding
relation between $A,\beta$ in $C_2\amalg C_3$, with $q(b)=s$,
$q(d)=r^{-1}$ and $q(u_1)=(sr)^{-1}$. Every element of
$\mathcal H\subset\widehat{GT}_0$ already satisfies that four-point
relation. The verified R2 mixed-sign element therefore rules out
deducing equality of the two parities from (FP.27.4) or (FP.27.5)
alone. This pinpoints why the remaining relations involving the fifth
point are essential.

\begin{statement}{A.58}{Proved}{an actual finite five-point level separates the two frame residues}
\end{statement}
The detector of A.54 can be embedded in one explicitly specified finite
characteristic quotient of $G_5$. This makes precise the finite level
needed for A.3(e), without claiming a parity-vanishing enumeration.
Let $L=(G_5)_5$, let $q:L\to C_2\amalg C_3$ be the forgetful map
of A.47, and let $\theta:C_2\amalg C_3\to E_0$ be A.54.
Set
\[
 N=P_5\cap\ker(\theta q),\qquad m=[G_5:N],\qquad
 N_*=\bigcap_{h\in\operatorname{Hom}_{\rm cont}
                         (G_5,\mathfrak S_m)}\ker h.
 \tag{FP.28.1}
\]
The subgroup $N$ is open in $L$, and hence in $G_5$, because $L$
has index five. The final intersection is finite: a homomorphism is
determined by the images of the two generators $a,b$. It is invariant
under every automorphism of $G_5$. The coset action on $G_5/N$ is one
of these homomorphisms, so $N_*$ lies in its kernel, the normal core
of $N$. Consequently $N_*$ is open and characteristic, is contained
in $N$, and
\[
 \mathcal Q_5=G_5/N_*
\]
is a finite characteristic quotient. Since $N_*\subset P_5$, it
retains the permutation map to $\mathfrak S_5$. The image $L_*$ of
$L$ admits a specified homomorphism $L_*\to E_0$, since
$N_*\subset\ker(\theta q)$.

For an actual $\Phi\in\widehat{GT}\cap\mathcal H$, its genuine
action on $\mathcal Q_5$ determines the residues of both frames in
$Q=\operatorname{SL}_2(\mathbb Z/12)/\{\pm I\}$.
To see this, solve
\[
 \Phi(b_*)=h_b b_*h_b^{-1},\qquad
 \Phi(d_*)=h_d d_*h_d^{-1}
 \tag{FP.28.2}
\]
with $h_b,h_d$ in the image of $P_5$ in $\mathcal Q_5$.
Actual frames supply solutions. Two choices differ by a centralizer.
Projecting the equations into $E_0$ gives the conjugacy equations
for $S$ and $R^{-1}$, respectively. The uniqueness result of A.54
then forces their $Q$-projections to be
$q_{12}(\beta_\Phi)$ and $q_{12}(A_\Phi)$. Their signed lifts and
the shear are recovered exactly as in A.55. Thus the two individual
parities and $\Delta(\Phi)$ are functions of this genuine finite
five-point action residue.

The construction is finite and effective from the presentations:
enumerate the finite group $E$, the subgroup generated by $S,R$,
the coset action of $G_5$ on $G_5/N$, and then all pairs of
permutations in $\mathfrak S_m$ satisfying the ten relators of A.40.
The product of their homomorphisms has kernel $N_*$.
It is a very large bound, not the smallest level. The representation
relations of A.40 can now be tested at a level which really determines
the residues of actual frames; the ambiguity exhibited in A.49 is
absent after projection to $Q$. No enumeration of all pentagon shadows
in $\mathcal Q_5$ is claimed here. The unresolved step is to show
that all admissible pentagon action residues in such a detector have
equal parities, or to construct a compatible sequence of mixed-sign
residues giving an actual element. An isolated finite mixed-sign
solution would prove neither alternative.

\begin{statement}{A.59}{Proved}{the explicit detector also removes the outer-action gauge}
\end{statement}
The characteristic quotient $Q_*$ of A.55 detects $\delta$ from the
\emph{outer} action on the multiplier-one subgroup. The extra fact is
that the same module extension separates the cusp centralizer. Put
$\tau=sr$, $\bar\tau=q(\tau)$ and $T_\tau=SR=\theta(\tau)$.
Then
\[
 \operatorname{pr}_Q C_{E_0}(T_\tau)=\langle\bar\tau\rangle,
 \qquad |\bar\tau|=12.
 \tag{FP.29.1}
\]
To prove this, the $M_2$ component of $T_\tau$ is $e_2$.
If $(m,g)$ centralizes $T_\tau$, its projection $g$ centralizes
$\bar\tau$ and
\[
 (1-\bar\tau)m_2=(1-g)e_2.
\]
Passing to $\langle\bar\tau\rangle$-coinvariants of $M_2$ says that
the cosets $H_2$ and $gH_2$ belong to the same $\bar\tau$-orbit.
Thus $g\in\langle\bar\tau\rangle H_2$. But $H_2$ intersects
$C_Q(\bar\tau)$ trivially: the nonidentity element $q(s)$ does not
commute with $q(sr)$, as its displayed matrices verify. It follows
that $g\in\langle\bar\tau\rangle$. Powers of $T_\tau$ give the
reverse inclusion, proving (FP.29.1).

Now suppose $\Phi\in\mathcal H$ acts by an inner automorphism on
$Q_*$. Choose its finite inner conjugator $g_*$. Since the genuine
action fixes $\tau$ literally, the image of $g_*$ in $E_0$ centralizes
$T_\tau$. Thus its $Q$-projection is $\bar\tau^k$ for some
$k\in\mathbb Z/12$. The genuine action also fixes the permutation
quotient $S_3$ pointwise; hence that projection is pure, since
$Z(S_3)=1$. The two actual pure frames and $g_*$ conjugate $S,R$ in
the same way. A.54 then gives
\[
 q(\beta_\Phi)=q(A_\Phi)=\bar\tau^k.
\]
In the signed matrix normalization, $\bar\tau$ is represented by
$U^{-1}$. The actual frame lifts are $I$ modulo four. Of the two
possible signed lifts $\pm U^{-k}$ only the positive one can be $I$
modulo four, and it does so exactly when $k\equiv0\pmod4$.
Consequently
\[
 q_{12}(\beta_\Phi)=q_{12}(A_\Phi)=U^{-k},\qquad
 \rho_\Phi\equiv-k/2\pmod3.
\]
The second assertion follows by applying the abelianization
$\epsilon(U^2)=1$ as in A.55. CRT now gives
$8\rho_\Phi\equiv-k\pmod{12}$. Therefore
\[
 T_\Phi=W_\Phi=I,\qquad\delta(\Phi)=0.
 \tag{FP.29.2}
\]
Taking the quotient of the genuine action homomorphism by inner
automorphisms, write
$\rho_{\mathrm{out}}:\widehat{GT}_0\to\operatorname{Out}(Q_*)$.
We have proved that $\delta$ vanishes on
$\mathcal H\cap\ker\rho_{\mathrm{out}}$. Thus it factors through
$\rho_{\mathrm{out}}(\mathcal H)$, with no choice of a finite inner
conjugator or an unproved gauge recovery.

For arbitrary multipliers, apply A.56 with $\rho_{\mathrm{out}}$
in place of $\rho$. This gives a modulus $M$ and a finite Galois table
with the same exact evaluation formula. The multiplier residue can
also be retained in an outer finite action, as follows. The level-two
free subgroup $F\subset C_2\amalg C_3$ is characteristic: every
epimorphism onto $S_3$ takes $s$ to a transposition and $r$ to a
three-cycle, and all six such generating pairs are conjugate in $S_3$.
Thus they have the same kernel. The subgroup
\[
 F_M=\ker\bigl(F\longrightarrow F^{\mathrm{ab}}/M\bigr)
\]
is characteristic in $F$ and hence in $C_2\amalg C_3$.
In the finite quotient $(C_2\amalg C_3)/F_M$, the genuine GT action
on $F/F_M\cong(\mathbb Z/M)^2$ is the scalar $\lambda$.
Inner automorphisms act there through the six integral matrices of
the $S_3$ action on the three peripheral classes. For $M\geq2$, its
only scalar matrix is $I$; this follows directly from the transposition
$\left(\begin{smallmatrix}0&1\\1&0\end{smallmatrix}\right)$ and
the three-cycle
$\left(\begin{smallmatrix}0&-1\\1&-1\end{smallmatrix}\right)$.
Hence two genuine GT actions with equal outer residues have the same
$\lambda\bmod M$.
Adjoining this quotient, that is replacing $K_*$ by $K_*\cap F_M$,
therefore makes the full $\Delta$ a function of one finite
\emph{outer} action residue. The arithmetic modulus and Galois table
remain uncomputed; no claim of minimality is made.

The exact check of all $576$ elements of $Q$ finds
$|C_Q(\bar\tau)|=24$. The $M_2$ coinvariants exclude the twelve
noncusp centralizer residues and retain exactly the twelve cusp powers.
The three normalization cases $k=0,4,8$ give respectively
$\rho=0,1,2$ modulo three and $T=W=I$.
These checks are included in
\texttt{check\_frame\_separating\_module.py} and its output.

\subsection{The marking correction for the descent equation}


\begin{statement}{A.60}{Proved}{the literal map in the current marking is a different map}
\end{statement}
Let $\widehat B=\langle s,r\mid s^2=r^3=1\rangle^{\wedge}$,
$\tau=sr$, and let $F$ be its level-two subgroup.  Distinguish the
two free generating pairs
\[
 x_P=\tau^{-2},\quad y_P=sx_Ps=(rs)^{-2},\qquad
 x_L=(rs)^2=y_P^{-1},\quad y_L=\tau^2=x_P^{-1}.
\]
The subscripts $P,L$ mean the current prompt and LM2, respectively.
The original manuscript/LM2 quotient, expressed on the physical subgroup
$F$, is
\[
 \varphi_L(x_L)=s,\quad\varphi_L(y_L)=\tau,
 \qquad\text{hence}\qquad
 \varphi_L(x_P)=\tau^{-1},\quad\varphi_L(y_P)=s.
\]
The literal quotient specified in the current prompt is instead
\[
 \varphi_P(x_P)=s,\qquad\varphi_P(y_P)=\tau.
\]
Let $\kappa\in\operatorname{Aut}(F)$ send $x_P$ to $y_P^{-1}$ and
$y_P$ to $x_P^{-1}$. Then $\kappa^2=1$ and
\[
                 \varphi_L=\varphi_P\circ\kappa.                 \tag{M1}
\]
In particular, retaining the written values $x\mapsto s,y\mapsto\tau$
while changing from the LM2 basis to the current basis does not retain
the same physical homomorphism.

\emph{Proof.} All assertions follow by substituting the four displayed
words in $s,r$. The automorphism $\kappa$ is $\sigma\circ c$, where
$\sigma$ exchanges $x_P,y_P$ and $c$ inverts both generators.
The two factors commute.  Equation (M1) can be checked on the two
free generators.\hfill$\square$

\begin{statement}{A.61}{Proved}{the two descent endomorphisms differ by conjugation by $c$}
\end{statement}
Let $G=\widehat{\mathrm{GT}}_0$ be the outer centralizer of the deck
$S_3$ acting on $F$. Denote by $D_P,D_L$ the endomorphisms obtained by
outer descent along $\varphi_P,\varphi_L$, followed by the normalized
restriction to $F$. Then
\[
 D_L(g)=D_P(cgc)=cD_P(g)c,\qquad
 D_P(c)=c.                                                       \tag{M2}
\]
Thus the two compatibility subgroups are
\[
 C_L=\operatorname{Fix}(D_L)
     =\{g:D_P(g)=cgc\},\qquad
 C_P=\operatorname{Fix}(D_P).
\]
No equality of these two subgroups is asserted.

\emph{Proof.} For any quotient $\varphi$ and an automorphism $a$
preserving its kernel, write $\operatorname{Desc}_{\varphi}(a)$ for
the induced outer automorphism of its target. Equation (M1) gives
\[
 \operatorname{Desc}_{\varphi_L}(g)
 =\operatorname{Desc}_{\varphi_P}(\kappa g\kappa^{-1}).
\]
In the outer automorphism group, $g$ commutes with $\sigma$, while
$\kappa=\sigma c=c\sigma$. Therefore
$\kappa g\kappa^{-1}=cgc$. Applying normalized restriction proves
the first equality in (M2).

For $c(x_P)=x_P^{-1},c(y_P)=y_P^{-1}$, descent along $\varphi_P$
is the genuine automorphism
\[
 \iota(s)=s,\qquad\iota(\tau)=\tau^{-1},\qquad
 \iota(r)=sr^{-1}s.
\]
It induces the identity on $B/F=S_3$, since conjugation by the
transposition $s$ inverts the three-cycle. Its restriction sends
$x_P$ and $y_P$ to their inverses. Hence $D_P(c)=c$.
The homomorphism law for $D_P$ now proves the second equality in (M2).
All statements are about outer maps; inner choices disappear before
normalized restriction.\hfill$\square$

\begin{statement}{A.62}{Proved}{the original arithmetic compatibility and a counterexample to the literal-current fixed-point assertion}
\citep{LitLochakSchneps1997}
\end{statement}
In the original LM2 marking, arithmetic functoriality gives
$D_L(\sigma)=\sigma$ for every $\sigma\in G_{\mathbb Q}$.
Consequently the literal-current endomorphism satisfies
\[
                   D_P(\sigma)=c\sigma c
                   \quad(\sigma\in G_{\mathbb Q}).              \tag{M3}
\]
It does not fix every Galois element. In fact,
\[
                  C_P\cap G_{\mathbb Q}=\{1,c\}.
\]

\emph{Proof of the arithmetic comparison.} Let $\mathcal X$ be the
orbifold with coarse curve $\mathbb P^1_{\mathbb Q}\setminus\{0\}$,
with inertia orders two at $1$ and three at $\infty$. It is the
rigidified elliptic modular stack in coordinate $T=1728/j$.
Its level-two cover and its puncture-opening map are, respectively,
\[
 U\longrightarrow\mathcal X,\quad
 T=\frac{27t^2(1-t)^2}{4(1-t+t^2)^3},
 \qquad
 U=\mathbb P^1\setminus\{0,1,\infty\}\hookrightarrow\mathcal X.
\]
Here is a simultaneous choice of loops, including their orientation.
Choose a small rational $t_0>0$, for instance $t_0=1/100$, and base
the orbifold at $T_0=q(t_0)$, where $q$ is the displayed rational
function. Let $\tau$ be the positive, counterclockwise meridian
about $T=0$, and let $s$ be the positive meridian about $T=1$,
using stems along the real interval, slightly displaced when needed.
With the peripheral product ordered as $\ell_1\ell_0\ell_\infty=1$,
the positive infinity meridian is $(s\tau)^{-1}$. Thus setting
$r=s\tau$ gives $s^2=r^3=1$ and $\tau=sr$.

The cover has the deck involution $t\mapsto1-t$ and local degree
two at $t=0,1,\infty$. Its only other ramifications have degree
two over $T=1$ and degree three over $T=\infty$; explicitly
\[
 1-q(t)=\frac{(t+1)^2(2t-1)^2(t-2)^2}
                  {4(t^2-t+1)^3}.
\]
Together with the six anharmonic deck transformations this shows
that it is the degree-six $S_3$ orbifold cover.
Let $y_L$ be the positive meridian about $t=0$ based at $t_0$.
Since $q(t)=(27/4)t^2+O(t^3)$, its image is $\tau^2$, with the
positive exponent. The lift $\gamma$ of the positive loop $s$
starting at $t_0$ ends at $1-t_0$: it passes around the ramification
point $1/2$ and implements the deck involution $t\mapsto1-t$.
Define $x_L$ by transporting the positive meridian about $t=1$
back along $\gamma$. The latter meridian is the image of $y_L$
under the holomorphic involution, so it has the same positive
orientation. Hence
\[
 q_*(y_L)=\tau^2=(sr)^2,\qquad
 q_*(x_L)=s\tau^2s^{-1}=(rs)^2.
\]
The stems can be chosen with disjoint interiors, so $x_L,y_L$ are
a free peripheral basis of $\pi_1(U,t_0)$. These formulas therefore
identify $q_*$ with the physical level-two inclusion in precisely
the LM2 marking, rather than merely an unspecified conjugate or
orientation reversal of that inclusion.

For the open map use the real interval from $T_0$ to $t_0$ to
transport its basepoint. In the target orbifold the path $\gamma$
is homotopic to the interval from $t_0$ to $1-t_0$: its detour is
around $1/2$, which is an ordinary interior point of the target.
Thus the same two domain loops have images $\tau$ and $s$ under
the open map, respectively. This proves simultaneously
\[
 (x_L,y_L)\longmapsto((rs)^2,(sr)^2)\quad\text{under }q_*;
 \qquad (x_L,y_L)\longmapsto(s,\tau)\quad\text{under }j_*.
\]
Both algebraic maps are defined over $\mathbb Q$. The interval
used to transport basepoints need not be Galois invariant: changing
its Galois transform changes the induced comparison by a geometric
inner automorphism, exactly the ambiguity already removed in the
outer descent and restriction. Functoriality of the arithmetic
fundamental-group extensions therefore makes these two outer
actions equivariant. Consequently $D_L$ fixes the Galois image.
This reproves the typed content of LM2 E.1 with the same physical
marking on both maps; it does not prove the same assertion with
the unconverted formula for $\varphi_P$.
There is one final identification of arithmetic images to check.
Changing the positive LM2 basis to the prompt basis conjugates the
outer Galois image by $\kappa$. Since $\kappa=\sigma c$, the
interchange $\sigma$ commutes with this image in the outer group,
and $c$ belongs to the image, one has
$\kappa G_{\mathbb Q}\kappa^{-1}=G_{\mathbb Q}$ as subgroups.
Thus the pointwise fixed subgroup statement for $D_L$ applies to
every element of the current Galois subgroup, without identifying
the old and new labels of individual arithmetic automorphisms.
Combining it with (M2) proves (M3).

For completeness the needed Galois centralizer fact has a direct
field-theoretic proof. An automorphism commuting with $c$ preserves
its fixed field $\mathbb R_{\rm alg}=\overline{\mathbb Q}\cap\mathbb R$.
The positive elements of this field are exactly its nonzero squares,
so every field automorphism preserves its order. It fixes the dense
subfield $\mathbb Q$, hence fixes every real algebraic number: two
different real numbers are separated by a rational number. Finally
$\overline{\mathbb Q}=\mathbb R_{\rm alg}(i)$, so an automorphism
fixing $\mathbb R_{\rm alg}$ is either $1$ or $c$. Therefore
$C_{G_{\mathbb Q}}(c)=\{1,c\}$, proving the displayed intersection
without a GT centralizer theorem.\hfill$\square$

Here is an explicit compact family of witnesses that does not rely
only on the centralizer theorem. Choose compatible positive real
roots $2^{1/n}$ and standard roots of unity $\zeta_n$. Define the
Kummer cocycle by
\[
 \sigma(2^{1/n})/2^{1/n}=\zeta_n^{\rho_2(\sigma)}.
\]
The number $2^{1/4}$ does not belong to $\mathbb Q^{\mathrm{ab}}$:
$X^4-2$ is Eisenstein, and its real root generates a nonnormal real
quartic field, whereas every finite subextension of
$\mathbb Q^{\mathrm{ab}}/\mathbb Q$ is Galois. Its square does belong
to $\mathbb Q^{\mathrm{ab}}$. Thus
\[
 \mathcal C=
 \{\sigma\in G_{\mathbb Q^{\mathrm{ab}}}:
                    \sigma(2^{1/4})=-2^{1/4}\}
\]
is a nonempty compact coset, and each of its elements has
$\rho_2(\sigma)\equiv2\pmod4$.
Complex conjugation fixes every chosen positive root and inverts
every $\zeta_n$, so
\[
                 \rho_2(c\sigma c)=-\rho_2(\sigma).
\]
In $\widehat{\mathbb Z}$, a nonzero element cannot equal its negative:
multiplication by two is injective. More concretely, a residue equal
to two modulo four is two or six modulo eight, and changes under
negation. Every $\sigma\in\mathcal C$ therefore satisfies
$D_P(\sigma)=c\sigma c\ne\sigma$. The claimed literal-current
Galois fixed-point assertion is refuted by this entire compact coset.
The same witnesses refute both
$\widehat{\mathrm{GT}}\subseteq C_P$ and $C_P=\ker\Delta$:
they belong to $G_{\mathbb Q}\subseteq\widehat{\mathrm{GT}}$ and
have $\Delta=0$ by R4, but do not belong to $C_P$. These
counterexamples do not refute the corresponding proposed inclusion
for the original group $C_L$.

\begin{statement}{A.63}{Proved}{the exact multiplier-one formula in the current marking}
\end{statement}
Write $x=x_P$, $y=y_P$, $z=(xy)^{-1}$ and
$\sigma=\operatorname{Ad}(s)|_F$, $\rho=\operatorname{Ad}(r)|_F$.
For a normalized pair $(1,f)$, set
\[
 C=\varphi_P(f)=f(s,\tau),\qquad
 H=\varphi_P(f(z,x))=f(r^{-1},s).
\]
Then $C\in F$. Write $C^{\mathrm{ab}}=(p,q)$ in the basis $[x],[y]$,
put $a=q-p$, and write $H\bmod F=r^l$ with $l\in\mathbb Z/3$.
The normalized pair and the two pure frames of $D_P(1,f)$ are
\[
 \begin{split}
 f_{D_P}&=y^a\sigma(C)C^{-1}x^{-a},\\
 \beta_{D_P}&=x^a C,\\
 A_{D_P}&=x^a CHr^{-l}.
 \end{split}                                                    \tag{M4}
\]
In particular, the current literal fixed-point condition is exactly
\[
             f=y^{q-p}\sigma(\varphi_P(f))\varphi_P(f)^{-1}
                    x^{p-q}.                                   \tag{M5}
\]
The corresponding original compatibility condition is
$D_P(1,f)=c(1,f)c$, equivalently
$f_{D_P}=f(x^{-1},y^{-1})$.

\emph{Proof.} The standard multiplier-one representative fixes $x$
and sends $y$ to $f^{-1}yf$. Its extension to $B$ is
$\Phi(s)=sf$, $\Phi(r)=f^{-1}r$, $\Phi(\tau)=\tau$.
The three-cycle relation gives
\[
 \Phi(z)=f^{-1}\rho(f)^{-1}z\rho(f)f
        =f(z,x)z f(z,x)^{-1}.
\]
Thus its descendant $\psi$ satisfies
\[
 \psi(s)=s,\qquad\psi(\tau)=C^{-1}\tau C,
 \qquad\psi(r)=HrH^{-1}.
\]
Because $f\in F'$, both $C,H$ map into the derived subgroup $A_3$
of $S_3$. The last formula therefore fixes $r$ modulo $F$, and the
first fixes $s$. Hence $\psi$ induces the identity on $S_3$.
The middle formula now shows that $C\bmod F$ centralizes the
transposition $\tau$, and $A_3\cap C_{S_3}(\tau)=1$, proving $C\in F$.

The map $\operatorname{Ad}(C)\psi$ fixes $\tau$ and has pure frames
$C$ and $CHr^{-l}$. Its pair coordinate is
$\sigma(C)C^{-1}$, with abelianization $(q-p,p-q)=(a,-a)$.
Composing by $\operatorname{Ad}(x^a)$ fixes the cusp and sends this
coordinate to $y^a\sigma(C)C^{-1}x^{-a}\in F'$. This proves (M4),
including the normalization exponent. Since both normalized pairs
fix $x$ and have coordinate in $F'$, they represent the same outer
class only when their coordinates are equal; their possible inner
conjugator centralizes $x$ and its abelianization forces its exponent
to vanish. This proves (M5).\hfill$\square$

\begin{statement}{A.64}{Proved}{an independent check from the NT formulas}
\citep[Theorems 1 and 2]{LitTsunogaiNakamura2002}
\end{statement}
The primary report of Tsunogai and Nakamura (2002, Theorems 1 and 2)
uses braid generators $t_1,t_2$ and $\eta=t_1t_2t_1$.
The exact quotient to the current $B$ marking is
\[
 t_1=\tau^{-1}=r^{-1}s,\qquad t_2=s\tau^{-1}s=sr^{-1},
 \qquad t_1^2=x,\ t_2^2=y,\ \eta=s,\ t_1t_2=r.
\]
Let $\gamma\in G_{\mathbb Q^{\mathrm{ab}}}$, and write its harmonic
primitive as $g=\beta^{-1}$. The source gives
\[
 g=f(x,s)t_1^{4\rho_2},\qquad
 f(t_1,s)=s^{\rho_3-2\rho_2}f(x,s)t_1^{6\rho_2-3\rho_3}.
\]
Apply these formulas to $c\gamma c$. Its pair coordinate is
$f^c(X,Y)=f(X^{-1},Y^{-1})$, its Kummer exponents are
$-\rho_2,-\rho_3$, and its pure harmonic frame is $\iota(\beta)$,
where $\iota$ is the reflection used in A.61. As both Kummer
exponents of $\gamma$ are even, the order-two factor disappears.
Taking the inverse yields
\[
 C=f(s,\tau)
    =t_1^{2\rho_2-3\rho_3}\,\iota(\beta).
\]
Since $\beta^{\mathrm{ab}}=(-\rho_2,-\rho_2)$, the second
coordinate of $C^{\mathrm{ab}}$ is $q=\rho_2$.
Formula (M4) gives $\rho_2(D_P\gamma)=-q=-\rho_2(\gamma)$,
in agreement with (M3), and its normalization makes
$\beta_{D_P\gamma}=\iota(\beta_\gamma)$ exactly.
This independently detects the conjugation twist using the opened
primary equations, rather than a presumed identification with them.

\paragraph{Consequences for the unresolved comparison.}
The manuscript equation (30) remains correct in its original marking:
with $\xi=y_L^{(\lambda-1)/2}f_L$ and an allowed gauge,
\[
 \xi\rho(\xi)\rho^2(\xi)=1,\qquad
 \xi=[\varphi_L(\xi),s]x_L^{(\lambda-1)/2}.
\]
Its literal unconverted use with $x_P,y_P,\varphi_P$ changes the
compatibility subgroup. This correction does not by itself prove
$\widehat{\mathrm{GT}}\subseteq C_L$ or $C_L\subseteq\ker\Delta$.
For the finite-character route, however, the distinction is harmless
at the level of the composed character: because $\Delta$ is invariant
under conjugation by $c$,
\[
                         \Delta\circ D_L=\Delta\circ D_P.
\]
A proof that this common character vanishes would put both fixed
subgroups in $\ker\Delta$. No such vanishing is presumed here.

\begin{statement}{A.65}{Proved}{an exact NT primitive comparison and two conditional inclusions}
\citep[Theorems 1 and 2]{LitTsunogaiNakamura2002}
\end{statement}
Let $\Phi=(1,f)\in\mathcal H$, and put
\[
 a=\tau^{-1},\quad b=s\tau^{-1}s,\quad a^2=x,\quad b^2=y,
 \quad aba=s,\quad ab=r.
\]
Write $g=\beta^{-1}$, $h=A^{-1}$, so that
\[
 f=\sigma(g)^{-1}g=\rho(h)^{-1}h,
 \qquad g^{\rm ab}=(\rho_2,\rho_2),\quad h^{\rm ab}=0.
\]
The last equality follows from the order-three primitive equation:
the matrix $1-\rho$ on $F^{\rm ab}$ has determinant three and is
injective over $\widehat{\mathbb Z}$. The coordinate $\rho_2$ is even
on the whole multiplier-one group, as established in the supplied
R4 derivation.

Set $P=f(a,s)$. Then $P\in F$. Write $P^{\rm ab}=(p,q)$ and
$b_0=p-q$. Descent along the \emph{original physical} quotient
$\varphi_L(x)=a,\varphi_L(y)=s$ gives
\[
 \beta_{D_L}=x^{b_0}P^{-1}.                                     \tag{M6}
\]
Consequently
\[
 \Phi\in C_L\quad\Longleftrightarrow\quad
       \beta^{-1}=f(a,s)x^{-b_0}.                              \tag{M7}
\]
In particular, the current fixed-point dictionary identifies the
original compatibility with the \emph{GF0 primitive form}. It does
not directly identify it with the harmonic GF1 form.

More precisely, fix any even $\kappa\in\widehat{\mathbb Z}$ and consider
the four equations in $\widehat B$
\[
\begin{aligned}
 (\mathrm{GF}_0)_\kappa:\quad
       g&=f(a,s)a^{-2\rho_2+3\kappa},\\
 (\mathrm{GF}_1):\quad
       g&=f(x,s)a^{4\rho_2},\\
 (\mathrm{HF}_0)_\kappa:\quad
       h&=r^{\rho_2-3\kappa/2}f(a,r)a^{3\kappa-2\rho_2},\\
 (\mathrm{HF}_1)_\kappa:\quad
       h&=f(x,r)a^{3\kappa}.
\end{aligned}                                                   \tag{M8}
\]
These are the central-quotient forms of NT Theorem 2 at multiplier
one, with $\kappa$ in place of its Kummer parameter $\rho_3$.
Then
\[
 \begin{split}
 \Phi\in C_L
 &\Longleftrightarrow
   (\mathrm{GF}_0)_\kappa\text{ for some even }\kappa\\
 &\Longleftrightarrow
   (\mathrm{HF}_0)_\kappa\text{ for some even }\kappa.
 \end{split}                                                    \tag{M9}
\]
The parameter is unique when either equation holds. Moreover the
following two sufficient conditions for $C_L$ are proved:
\begin{align}
 &(\mathrm{GF}_1)\quad\text{and}\quad
 f(a,s)=f(x,s)a^{6\rho_2-3\kappa};                              \tag{M10}\\
 &(\mathrm{HF}_1)_\kappa\quad\text{and}\quad
 f(a,r)=r^{-\rho_2}f(x,r)a^{2\rho_2}.                           \tag{M11}
\end{align}
The second equations in (M10) and (M11) are respectively the
central-quotient forms of NT (V) and $(\mathrm{IV}'_{\rm bis})$.
NT's harmonic equation $(\mathrm I')$ is equivalent to
$(\mathrm{GF}_1)$, and its equianharmonic equation $(\mathrm{II}')$
is equivalent to $(\mathrm{HF}_1)$ with the stated parameters.
Thus (harmonic $+$ V) and (equianharmonic $+$ IV-bis) each give a
precisely typed inclusion in the original compatibility group.
This is an inclusion between explicitly defined equation loci;
it is not a claim that all of GT satisfies those additional equations.

\emph{Proof.} The descendant along $\varphi_L$ fixes $a$ and sends
$s$ to $P^{-1}sP$. Its image of $r$ is $QrQ^{-1}$, where
$Q=f(a,r)^{-1}$. Since $P,Q$ are commutators in $B$, their reductions
lie in $A_3$; the descendant fixes $r$ and the transposition $a$
modulo $F$, hence fixes all of $S_3$. This proves $P\in F$.
The initial $s$-frame is $P^{-1}$ and has abelianization $(-p,-q)$.
The unique normalization making the pair coordinate a commutator
is conjugation by $x^{p-q}$. This proves (M6). The normalized map
fixes $\tau$, so its $s$-frame determines the whole automorphism;
this proves (M7).

Suppose first that $\Phi\in C_L$. Equation (M7) gives $q=\rho_2$.
As $P\in B'$, its image in $B^{\rm ab}$ is zero. The images of
$x,y$ in the order-three factor of $B^{\rm ab}$ coincide, so
$p+q\equiv0\pmod3$. Hence
\[
      \kappa=\frac{2(\rho_2-b_0)}3
\]
is well defined and even, and (M7) is exactly
$(\mathrm{GF}_0)_\kappa$. Conversely this primitive equation makes
$P^{\rm ab}=(2\rho_2-3\kappa/2,\rho_2)$, so
$b_0=\rho_2-3\kappa/2$ and (M7) follows. The cusp subgroup is
$\widehat{\mathbb Z}$, on which multiplication by three is injective,
so $\kappa$ is unique.

For the order-three frame write $Q\bmod F=r^j$.
The normalized descendant has
\[
 A_{D_L}=x^{b_0}Qr^{-j}.
\]
Its abelianization in $F$ is zero. Passing further to $B^{\rm ab}$
therefore gives $j\equiv b_0\pmod3$. Inverting this formula yields
$h_{D_L}=r^{b_0}f(a,r)x^{-b_0}$. It is
$(\mathrm{HF}_0)_\kappa$ with the same $\kappa$ when $\Phi\in C_L$.
Conversely, if $(\mathrm{HF}_0)_\kappa$ holds, its inverse is
$A=x^dQr^{-d}$ with $d=\rho_2-3\kappa/2$. Purity gives
$d\equiv j\pmod3$, and comparison with $A_{D_L}$ makes their
quotient $x^{d-b_0}$. Both frames have zero abelianization, so
$d=b_0$. The two normalized maps fix $\tau$ and have the same
$r$-frame, hence are equal. This proves (M9).

Combining the two equations in (M10) gives
$g=f(a,s)a^{-2\rho_2+3\kappa}$ directly. Combining those in (M11)
gives
\[
 r^{\rho_2-3\kappa/2}f(a,r)a^{3\kappa-2\rho_2}
 =r^{-3\kappa/2}f(x,r)a^{3\kappa}
 =f(x,r)a^{3\kappa}=h,
\]
where the middle equality uses $r^3=1$ and evenness of $\kappa$.
Finally, inserting the GF1 or HF1 expression in the corresponding
primitive equation produces NT $(\mathrm I')$ or $(\mathrm{II}')$;
the reverse implication follows from uniqueness of the pure
primitive. All substitutions use $a^2=x,b^2=y,aba=s,ab=r$, so
there is no unrecorded change of braid generators.\hfill$\square$

\paragraph{What this says about (III$'$).}
Nakamura--Schneps (2000, Theorem 1.1) puts its relation (III$'$) in
the five-point mapping class group, with
$x_{45}=(t_1t_2)^3$ and $x_{51}=(t_2t_3)^3$; Proposition 5.4
gives the central correction when lifting it to $B_4$.
Tsunogai--Nakamura (2002, p.200) separately displays
$(\mathrm I'),(\mathrm{II}'),(\mathrm{III}'),(\mathrm{IV}')$ and (V).
The opened texts do not state that (III$'$) alone supplies (V), or
that all GT elements satisfy GF0. Equations (M8)--(M11) specify
the remaining comparison exactly: to pass from the harmonic
primitive to the original $C_L$ condition one needs the displayed
V comparison (or an independently proved substitute for it).

\subsection{Independent descent calculation and arithmetic witness}


\begin{statement}{A.66}{Proved}{actual normalized descent frames at multiplier one}
\end{statement}
Put $\tau=sr$, $x=\tau^{-2}$, $y=sxs$, and let
$\varphi(x)=s$, $\varphi(y)=\tau$. For $\Phi=(1,f)\in\mathcal H$ set
\[
 C=\varphi(f),\qquad H=\varphi\bigl(f(z,x)\bigr).
\]
Then $C\in F$. Write $C^{\rm ab}=(p,q)$ in the basis $x,y$, put
$a=q-p$, and let $l\in\{0,1,2\}$ be defined by
$H\bmod F=r^l$. The descent in the prompt's stated marking has
the normalized pair and frames
\begin{align*}
 f_D&=y^a\theta(C)C^{-1}x^{-a},\qquad \theta(w)=sws^{-1},\\
 \beta_D&=x^aC,\qquad A_D=x^aCHr^{-l},\qquad \rho_D=-q.
\end{align*}
In particular these are actual pure frames, and the exponent $a$ is
uniquely specified; an unrecorded inner change is not being absorbed
into a signed coordinate.

\emph{Proof.}
Conjugation by $r$ cyclically permutes $x,y,z$. The normalized gentle
extension has $\Phi(r)=f^{-1}r$. Applying this twice, and using
$f(z,x)f(y,z)f(x,y)=1$, gives
\[
 \Phi(z)=f(z,x)z f(z,x)^{-1}.
\]
The kernel of $\varphi$ is the closed normal subgroup generated by
$x^2,z^3$. It is preserved by $\Phi$ and its inverse. Thus the descent
$\psi$ is an actual automorphism of $\widehat B$, with
\[
 \psi(s)=s,\qquad \psi(\tau)=C^{-1}\tau C,
 \qquad \psi(r)=HrH^{-1}.
\]
Since $f\in F'$, both $C,H$ lie in $\widehat B'$. Their images in
$S_3$ lie in $A_3=\langle r\rangle$. In particular $H$ centralizes
$r$ there, so $\psi$ induces the identity on $S_3$. The cusp equation
then says that the image of $C$ centralizes the transposition $\tau$.
Since $A_3\cap C_{S_3}(\tau)=1$, this proves $C\in F$.

The automorphism $\psi_0=\operatorname{Ad}_C\psi$ fixes $\tau$ and
has raw frames $C,CH$. Right multiplication of the latter by $r^{-l}$
makes it pure and does not change its conjugation of $r$. Now set
$\Phi_D=\operatorname{Ad}_{x^a}\psi_0$. This still fixes $\tau$.
Its pure $s$ frame is $x^aC$, with abelianization $(p+a,q)$.
The element
\[
 s\Phi_D(s)=y^a\theta(C)C^{-1}x^{-a}
\]
has abelianization $(q-p-a,p-q+a)$, so it belongs to $F'$ precisely
when $a=q-p$. The normalized representative then has
$\Phi_D(x)=x$, $\Phi_D(y)=f_D^{-1}yf_D$,
$\Phi_D(s)=sf_D$, and $\Phi_D(r)=f_D^{-1}r$.
These actual identities prove that it is the normalized gentle pair
representing the descended outer action. They also give the two
displayed frames and $\rho_D=-q$.

We record two integral consequences needed below. The image of
$C\in\widehat B'$ in $\widehat B^{\rm ab}=C_6$ gives
$p+q\equiv0\pmod3$. Further, $A_D\in F'$: the equality
$f_D=rA_Dr^{-1}A_D^{-1}$ implies
$(\operatorname{Ad}_r-1)A_D^{\rm ab}=0$. On
$F^{\rm ab}=\widehat{\mathbb Z}^2$ the determinant of this linear
operator is $3$. Its kernel is zero because multiplication by $3$
is injective on $\widehat{\mathbb Z}$.

\begin{statement}{A.67}{Proved}{the signed finite formula for $\Delta\circ D$}
\end{statement}
Use the matrices of R4,
\[
 U=\begin{pmatrix}1&1\\0&1\end{pmatrix},\quad
 S=\begin{pmatrix}0&1\\-1&0\end{pmatrix},\quad
 R_0=\begin{pmatrix}0&-1\\1&-1\end{pmatrix}.
\]
Let
\[
 \mathsf C=f(S,U^{-1}),\qquad
 \mathsf H=f(R_0^{-1},S)
       \quad\hbox{in }\mathrm{SL}_2(\mathbb Z/12).
\]
These are the commutator lifts of $C,H$; they must not be substituted
for the signed pure lifts without a sign correction. With $a,l$ as
in A.66,
\[
 a\equiv l\pmod3,\qquad
 q_{12}(\beta_D)=(-I)^aU^{2a}\mathsf C,\qquad
 q_{12}(A_D)=(-I)^aU^{2a}\mathsf C\mathsf H R_0^{-l}.
\]
Modulo three write
$\mathsf C=S^j$ and $\mathsf C\mathsf H=(-I)^e$; such
$j\in\mathbb Z/4$ and $e\in\mathbb Z/2$ exist for every actual
gentle pair. Then
\[
 \boxed{\ \Delta(D\Phi)=a+e+j\pmod2.\ }
\]
The parity of $a$ is itself finite: modulo four,
$\mathsf C=(-U^2)^a$. Thus the displayed formula is an exact
finite evaluation formula on actual gentle pairs.

\emph{Proof.}
The central extension $\widehat{\mathrm{SL}_2(\mathbb Z)}\to
\widehat B$ has kernel $\{\pm I\}$. Its abelianization is $C_{12}$,
in which $U,L$ have class $1$, $S=ULU$ has class $3$,
$-I$ has class $6$, and $R_0=-UL$ has class $8$.
Its derived subgroup maps isomorphically onto $\widehat B'$:
the two possible lifts of an element of $\widehat B'$ have
abelianizations $0,6$, and exactly one is in the derived subgroup.
The words defining $\mathsf C,\mathsf H$ choose this latter lift.
For $H$, the lift of $\varphi(z)$ differs from $R_0^{-1}$ by $-I$,
which does not change evaluation of the commutator word $f$.

The signed pure lift of $C$ instead has abelianization $2(p+q)$.
Since $p+q\equiv0\pmod3$ and
$p+q\equiv q-p=a\pmod2$, it differs from the commutator lift by
$(-I)^a$. This proves the formula for $\beta_D$. The word
$U^{2a}\mathsf C\mathsf H R_0^{-l}$ has abelianization
$2a-8l$, while the signed lift of $A_D\in F'$ has abelianization
zero. The two lifts differ by a sign. Consequently
$a\equiv l\pmod3$, and that sign is again $(-I)^a$, since
$(a-4l)/3\equiv a\pmod2$. This proves the second formula.

Inserting $\rho_D=-q$ into the R4 cusp correction gives, modulo
three,
\[
 W_D=(-I)^a\mathsf C^{-1},\qquad
 T_D=(-I)^aR_0^l(\mathsf C\mathsf H)^{-1}.
\]
Indeed the remaining upper-unipotent exponent is
$-2a-8q=2p-10q\equiv2(p+q)=0\pmod3$.
The R4 centralizer condition implies $\mathsf C\in\langle S\rangle$.
Both $\mathsf C,\mathsf H$ modulo three lie in
$\mathrm{SL}_2(\mathbb F_3)'=Q_8$. Since
$C_{\mathrm{SL}_2(\mathbb F_3)}(R_0)\cap Q_8=\{\pm I\}$,
their product is $(-I)^e$. The sign character of $T_D$ is $a+e$,
while that of $W_D$ is $j$, proving the boxed formula. Finally
$q_4(\beta_D)=I$ in R4 and $U^{-2}=U^2$ modulo four give
$\mathsf C=(-U^2)^a$ modulo four.

\begin{statement}{A.68}{Refuted}{the claim that the stated $D$ fixes the Galois image}
\end{statement}
For the precise current marking in A.66 one has
\[
 \boxed{\ D(\sigma)=c\sigma c
       \quad(\sigma\in G_{\mathbb Q^{\rm ab}}),\qquad
       \rho(D\sigma)=-\rho(\sigma).\ }
\]
There exist such $\sigma$ with $D(\sigma)\ne\sigma$.
Consequently the literal inclusion
$\widehat{\mathrm{GT}}\subseteq\operatorname{Fix}(D)$ is false
with this marking. Identifying this fixed subgroup with the manuscript's
arithmetically compatible subgroup without a marking change is also
false.

\emph{Proof.}
Let $u=\tau^{-1}$ and $s=uvu$. Write $g=\beta_\sigma^{-1}$.
With positive real radical roots, Nakamura--Schneps prove
\[
 g^{\rm ab}=(\rho_2,\rho_2),\qquad \rho=\rho_2
\]
\cite[Proposition~2.4 and Definition~5.1]{LitNakamuraSchneps2000}.
The two actual braid identities in the authors' opened report are
\begin{align*}
 g(u^2,v^2)&=f(u^2,s)u^{4\rho_2},\\
 f(u,s)&=s^{\rho_3-2\rho_2}f(u^2,s)u^{6\rho_2-3\rho_3}
\end{align*}
\cite[Theorem~2, $(\mathrm{GF}_1)$; Theorem~1, $(\mathrm V)$]
{LitTsunogaiNakamura2002}.
For $\sigma\in G_{\mathbb Q^{\rm ab}}$ the Kummer exponents
$\rho_2,\rho_3$ are even: $\sqrt2,\sqrt3\in\mathbb Q^{\rm ab}$.
Passing to $\widehat B$, where $s^2=1$, therefore gives
\[
 f(u,s)=g\,u^{2\rho_2-3\rho_3}.
\]
Let $c_B$ denote the actual complex-conjugation automorphism of
$\widehat B$: $c_B(u)=u^{-1}$, $c_B(v)=v^{-1}$, $c_B(s)=s$.
The two-cycle relation and the preceding equality yield
\begin{align*}
 C=f(s,u^{-1})
   &=c_B(f(u,s))^{-1}\\
   &=u^{2\rho_2-3\rho_3}c_B(\beta_\sigma)
     =\tau^{-2\rho_2+3\rho_3}c_B(\beta_\sigma).
\end{align*}
Since $c_B$ negates $F^{\rm ab}$ and $\beta_\sigma^{\rm ab}
=(-\rho_2,-\rho_2)$, it follows that
\[
 C^{\rm ab}=(2\rho_2-3\rho_3/2,\rho_2),\qquad
 a=-\rho_2+3\rho_3/2.
\]
The correction of A.66 now cancels exactly:
\[
 \beta_D=x^aC=c_B(\beta_\sigma).
\]
Both $\Phi_D$ and $c_B\Phi_\sigma c_B$ fix $\tau$ literally;
the latter has the pure $s$ frame $c_B(\beta_\sigma)$.
They agree on $s,\tau$, which generate $\widehat B$, and hence are
equal as actual automorphisms. This proves the boxed identities.

For an actual counterexample put $K=\mathbb Q(\zeta_8)$ and
$L=K(\alpha)$ with $\alpha=2^{1/8}>0$. The field
$\mathbb Q(2^{1/4})$ is not Galois over $\mathbb Q$:
$T^4-2$ is Eisenstein and its real root field does not contain
the nonreal root $i2^{1/4}$. It therefore is not contained in
the abelian extension $K/\mathbb Q$. Since $\mu_4\subset K$
and $\alpha^4=\sqrt2\in K$, the extension $K(\alpha)/K$
is cyclic of degree dividing four; degree at most two would
put $\alpha^2=2^{1/4}$ in $K$. Thus $[L:K]=4$ and
$[L:\mathbb Q]=16$. It is the splitting field of $T^8-2$.

Write its automorphisms as $(v,b)$, meaning
$\zeta_8\mapsto\zeta_8^v$ and
$\alpha\mapsto\zeta_8^b\alpha$. Their exact set and law are
\[
 \begin{gathered}
 v\in\{1,3,5,7\},\quad b\in\mathbb Z/8,\qquad
 b\equiv\begin{cases}0&v=1,7,\\1&v=3,5\end{cases}\pmod2,\\
 (v,b)(v',b')=(vv',b+vb').
 \end{gathered}
\]
The parity condition is exactly the image of
$\alpha^4=\sqrt2=\zeta_8+\zeta_8^{-1}$. Exact enumeration
gives the commutator subgroup $\{(1,0),(1,2),(1,4),(1,6)\}$.
Choose an actual $\sigma\in G_{\mathbb Q^{\rm ab}}$ restricting
to $(1,2)$; this is possible since restriction maps the closed
commutator subgroup of $G_{\mathbb Q}$ onto the commutator
subgroup of every finite Galois quotient. Then
$\rho_2(\sigma)\equiv2\pmod8$, whereas
$\rho_2(D\sigma)\equiv6\pmod8$. Complex conjugation is
$(7,0)$, and conjugation by it indeed sends $(1,2)$ to $(1,6)$.
This proves $D\sigma\ne\sigma$ in an actual degree-sixteen
Galois quotient. All group laws, the commutator subgroup, an
explicit commutator for $(1,2)$, and the distinction modulo eight
are checked in \texttt{check\_descent\_marking.py}.

\begin{statement}{A.69}{Proved}{the exact conversion from the LM2 marking}
\end{statement}
The old LM2 generators are
\[
 x_L=(rs)^2=y^{-1},\qquad y_L=(sr)^2=x^{-1}.
\]
Consequently its map $\varphi_L(x_L)=s$, $\varphi_L(y_L)=\tau$,
when written on the current generators, is
\[
 \varphi_L(x)=\tau^{-1},\qquad \varphi_L(y)=s.
\]
Put $k(x)=y^{-1}$, $k(y)=x^{-1}$. Then
$\varphi_L=\varphi\circ k$ and $k=\theta\circ c$ on $F$.
Every gentle outer class commutes with the geometric interchange
$\theta$. Writing $D_L$ for the old descent followed by restriction
to the same embedded pure subgroup therefore gives
\[
 D_L(\Phi)=D(c\Phi c)=cD(\Phi)c.
\]
Here $D(c)=c$ follows directly from
$\psi(s)=s$, $\psi(\tau)=\tau^{-1}$. Thus the corrected
endomorphism for the manuscript fixed-point convention is
$D^*=\operatorname{Ad}_c\circ D$. The two functions
$\Delta\circ D^*$ and $\Delta\circ D$ agree by R4's
conjugation invariance, but their fixed subgroups are different.
This change is required before transporting equations (30), NT
compatibilities, or the assertion that the Galois image is fixed.

\emph{Proof.}
The displayed word substitutions prove the marking formulas
without a choice of path. Descent along $\varphi\circ k$ is descent
along $\varphi$ after conjugating the source automorphism by $k$.
The geometric interchange commutes with gentle outer classes, so
this conjugation is the same as conjugation by $c$ at the outer
level. Descent and restriction are homomorphisms on these outer
groups. Together with $D(c)=c$ this gives the equality of
endomorphisms. The final equality of characters is precisely the
conjugation invariance already proved in R4. A.68 independently
checks the resulting nontrivial twist on the actual Galois action,
including all cusp and frame normalizations.

\begin{statement}{A.70}{Open}{Boundary of the calculation, not a profinite conclusion}
\end{statement}
The exact finite formula A.67 makes a direct test possible. The
six substitutions of the same $f$ into the pairs
\[
 (S,U^{-1}),\ (U^{-1},Z),\ (Z,S),\
 (U^{-1},S),\ (Z,U^{-1}),\ (S,Z),\qquad Z=(SU^{-1})^{-1},
\]
give a finite group of order $8153726976$, with derived subgroup of
order $56623104$. Exact membership, signed-frame conditions, and
the two-cycle and three-cycle equations retain twelve candidates
of each value of the right-hand side of A.67. Including the original
inclusion evaluation $(U^2,L^2)$ and its required R3 value for a fixed
point still does not force zero. The opposite source sign required by
the corrected locus $D(\Phi)=c\Phi c$ retains the same counts.
These computations do not establish
that any retained candidate has a compatible profinite lift. Their
only purpose is to identify why this finite-evaluation attempt has
not proved $\Delta\circ D=0$ or the asserted fixed-subgroup inclusion.
The actual remaining question requires further profinite or pentagon
constraints. The exact implementation and its output are retained
as \texttt{check\_descent\_evaluations.py} and its records.

\subsection{Adelic covariance and the full-image assertion}


\begin{statement}{A.71}{Proved}{The signed covariance at every level}
Let $\widehat B_3=\langle u,v\mid uvu=vuv\rangle^{\wedge}$,
let $q:\widehat B_3\to\mathrm{SL}_2(\widehat\Z)$ send
$u$ to $U=\left(\begin{smallmatrix}1&1\\0&1\end{smallmatrix}\right)$
and $v$ to $L=\left(\begin{smallmatrix}1&0\\-1&1\end{smallmatrix}\right)$,
and use the genuine gentle action of R4, E.65.8. There is a unique
continuous homomorphism
\[
 \mathcal V:G\longrightarrow
 \left\{\left(\begin{smallmatrix}\lambda&t\\0&1\end{smallmatrix}\right):
       \lambda\in\widehat\Z^{\times},\ t\in\widehat\Z\right\}
\]
such that $q\Phi_g=\operatorname{Ad}(\mathcal V_g)q$ and
the upper-left entry of $\mathcal V_g$ is $\lambda(g)$.
On $\HH$ we have $\mathcal V_h=U^{t_h}$ with
$t:\HH\to\widehat\Z$ an additive continuous character and
$t_h\equiv8\rho_h\pmod{12}$.
\end{statement}
\begin{proof}
We give the integral calculation behind
\citet[Section 4, Lemmas 4.1--4.3]{LitMarin2008} to include
the primes two and three. Fix a prime $\ell$ and put
$A=U^{\lambda}$, $B=q_\ell\Phi_g(v)$.
The second matrix is conjugate to $L^{\lambda}$, so its trace
is two and its determinant is one. Write
$B=\left(\begin{smallmatrix}a&b\\c&2-a\end{smallmatrix}\right)$.
The lower-left entry of $ABA=BAB$ is
$c(1+\lambda c)=0$. The coefficient ring $\Z_\ell$ is an
integral domain. If $c=0$, the determinant gives $(a-1)^2=0$
and the upper-right relation gives $b=\lambda$; thus $B=A$.
This is impossible: $\Phi_g$ is an automorphism, so
$q_\ell\Phi_g(\widehat B_3)=q_\ell(\widehat B_3)$,
which is nonabelian since $U,L$ do not commute.
Therefore $c=-\lambda^{-1}$. Put $w=1-a$. The determinant
then gives $b=\lambda w^2$, and direct multiplication yields
\[
 B=\begin{pmatrix}1-w&\lambda w^2\\
             -\lambda^{-1}&1+w\end{pmatrix}
 =\begin{pmatrix}\lambda&\lambda w\\0&1\end{pmatrix}
       L
  \begin{pmatrix}\lambda&\lambda w\\0&1\end{pmatrix}^{-1}.
\]
The same matrix conjugates $U$ to $U^{\lambda}$. It is unique:
the simultaneous centralizer of $U,L$ consists of scalar
matrices, and the lower-right entry has been fixed to one.
Its entries are continuous functions of $B$ and $\lambda$.

Taking the product over all primes gives $\mathcal V_g$.
For the composition convention $\Phi_{gh}=\Phi_g\Phi_h$,
both $\mathcal V_{gh}$ and $\mathcal V_g\mathcal V_h$
give the same intertwiner with the same normalization; hence
they agree. This also checks the order of composition.
At multiplier one the product law is addition of $t$.
Modulo twelve R3 supplies the intertwiner $U^{8\rho_h}$;
uniqueness in the finite ring follows directly from the same
simultaneous-centralizer computation, so the residues agree.
The script \texttt{check\_adelic\_covariance.py} verifies the
polynomial braid calculation and the centralizer identities.
No inference from solutions over rings with zero divisors is
used to obtain the domain calculation.
\end{proof}

\begin{statement}{A.72}{Proved}{Two actual signed adelic norm-one characters}
Use the signed free-subgroup lift $F\to\widehat B_3$ given by
$x=u^2,y=v^2$, and set $S=ULU$, $R_0=-UL$.
For the positive pure frames of $h\in\HH$, define
\[
 \mathcal T_h=q(A_h)^{-1}\mathcal V_h,\qquad
 \mathcal W_h=q(\beta_h)^{-1}\mathcal V_h.
\]
These are continuous homomorphisms with targets
\[
 \begin{split}
 \mathbb T_3&=\{aI+bR_0:a,b\in\widehat\Z,
                       \ a^2-ab+b^2=1\},\\
 \mathbb T_4&=\{aI+bS:a,b\in\widehat\Z,
                       \ a^2+b^2=1\}.
 \end{split}
\]
They are the norm-one groups of
$\Z[\zeta_3]\otimes\widehat\Z$ and
$\Z[i]\otimes\widehat\Z$ in the indicated integral bases.
Both characters have value $I$ modulo four.
Their reduction modulo twelve is exactly R4's pair $(T,W)$.
\end{statement}
\begin{proof}
We first check the signed comparison. The kernel of
$\widehat B_3\to\widehat B$ is the closure of the central
subgroup generated by $z_3=(uv)^3$. The exponent-sum map
to $\widehat\Z$ sends $z_3^e$ to $6e$, so this central
procyclic subgroup is isomorphic to $\widehat\Z$.
The frame equalities in $\widehat B$ lift to equalities in
$\widehat B_3$ up to $z_3^e$. Their two sides have equal
exponent sums (two for $uv$, three for $uvu$), since
$h$ has multiplier one. Thus $6e=0$ and $e=0$.
The comparison is consequently a signed matrix equality.
A.71 now shows that $\mathcal T_h$ centralizes $R_0$ and
$\mathcal W_h$ centralizes $S$.

Solving each two-by-two commutation equation gives respectively
$aI+bR_0$ and $aI+bS$ over any commutative coefficient ring;
their determinants are the two norm forms displayed above.
These descriptions are valid also at two and three, since
no division by either prime occurs. They are commutative groups.
The pure-frame multiplication law is
$A_{hk}=\Phi_h(A_k)A_h$ and
$\beta_{hk}=\Phi_h(\beta_k)\beta_h$.
Using A.71 and cancelling in the specified order gives
\[
 \begin{split}
 \mathcal T_{hk}
 &=q(A_h)^{-1}\mathcal V_h q(A_k)^{-1}\mathcal V_k
   =\mathcal T_h\mathcal T_k,\\
 \mathcal W_{hk}&=\mathcal W_h\mathcal W_k.
 \end{split}
\]
Continuity follows from frame continuity and matrix multiplication.
Finally R4's integral sign convention gives
$q_4(A_h)=q_4(\beta_h)=I$, while A.71 gives
$\mathcal V_h\equiv I\pmod4$. This proves the last assertions
without changing the global sign.
\end{proof}

\begin{statement}{A.73}{Refuted}{The literal surjectivity onto the full projective adelic product}
The pair $(\mathcal T,\mathcal W)$ is not onto
$(\mathbb T_3/\{\pm I\})\times(\mathbb T_4/\{\pm I\})$.
Its image is contained in the image of
\[
 \mathbb T_3(4)\times\mathbb T_4(4),\qquad
 \mathbb T_i(4)=\ker\bigl(\mathbb T_i\to
                           \mathrm{SL}_2(\Z/4)\bigr).
\]
Thus Lead 7's word ``whole'', with the full norm-one targets
as stated there, is false. The same counterexample works
if one quotients by all scalar matrices instead of only the
global sign.
\end{statement}
\begin{proof}
The pair $([R_0],[I])$ belongs to the displayed full target:
$R_0$ has norm one. Neither $R_0$ nor $-R_0$ is $I$ modulo
four. In fact $R_0$ is not scalar modulo four. A.72 therefore
excludes this pair from the image, even after the stated
projective identifications. The pair $([I],[S])$ is another
excluded value. The four residue comparisons are checked
in \texttt{check\_adelic\_covariance.py}.

This counterexample concerns the explicitly full target, not
a target redefined to incorporate the positive level-four
normalization. In that smaller target, the exact adelic
image and the assertion that equality of the two parities
is its only arithmetic constraint remain undecided here.
The proved image modulo twelve is $C_6\times C_4$, by B.14;
the Galois image there is the equal-parity cyclic subgroup
of order twelve. The maximal-abelian-field intersection in
A.12 explains that one finite arithmetic bit; it does not
prove surjectivity onto an unspecified adelic target.
\end{proof}

\subsection{Retaining the genuinely five-point part of the pentagon}


\begin{statement}{A.74}{Proved}{every induced forgetful detector annihilates the pentagon defect}
\end{statement}
There is an exact obstruction to obtaining the missing parity relation
by merely inducing finer four-point quotients to five points.
Write
\[
 (X_1,X_2,X_3,X_4,X_5)
   =(x_{12},x_{23},x_{34},x_{45},x_{51})
\]
for the standard cyclic boundary twists in $P_5$. In the label-five
stabilizer $B_4/\langle z_4\rangle$, with generators $a,b,c$ as in
B.35, they have the exact words
\[
 (X_1,X_2,X_3,X_4,X_5)
       =(a^2,b^2,c^2,(aba)^2,(cbc)^2).
 \tag{FP.30.1}
\]
For $X_4,X_5$, this is the equality of the positive twist around two
points with the positive twist around the complementary three points
on a punctured sphere. It keeps the central factor until passage to
$B_4/\langle z_4\rangle$. Under the literal fold of A.47,
\[
 q_5(X_1,X_2,X_3,X_4,X_5)=(x,y,x,1,1),
 \quad x=u^2,\quad y=v^2.
 \tag{FP.30.2}
\]
Indeed $q(a)=q(c)=u$, $q(b)=v$, and
$q((aba)^2)=q((cbc)^2)=(uvu)^2=1$ in $C_2\amalg C_3$.
The order-five automorphism $\rho$ of Lochak--Schneps (1997),
Section 1, is conjugation by $(\sigma_4\sigma_3\sigma_2\sigma_1)^2$
and sends $X_i$ to $X_{i+3}$, with indices modulo five.
Its permutation is $(1\,4\,2\,5\,3)$, so the maps
$q_5\rho^k$ run through the five forgotten labels. Their images of
the $X_i$ are exactly the five cyclic rotations of (FP.30.2), with
the specified base-path conjugations already included.

For any $f\in\widehat F_2'$ satisfying the 2-cycle relation, put
\[
 \mathcal P(f)=f(X_3,X_4)f(X_5,X_1)f(X_2,X_3)
                         f(X_4,X_5)f(X_1,X_2).
 \tag{FP.30.3}
\]
This is the actual five-cycle defect, before it is assumed trivial.
Since a commutator word is trivial when either argument is one,
\[
 q_5(\mathcal P(f))=f(y,x)f(x,y)=1.
\]
For every cyclic rotation there are again exactly two surviving
factors, one $f(x,y)$ and one $f(y,x)$; all intervening factors are
one. Hence
\[
 \mathcal P(f)\in\mathcal B_5:=
                \bigcap_{i=1}^{5}\ker(q_i|_{P_5}).
 \tag{FP.30.4}
\]
The subgroup $\mathcal B_5$ is the closed Brunnian subgroup: its
members become trivial after forgetting any marked point. The
projective quotient in $q_i$ introduces no extra pure kernel, because
$P_4$ injects into $G_4/V_4$.

Now let $\eta:L\to D$ be \emph{any} finite homomorphism factoring
through the forgetful map $q_5$, including
$\eta=\theta q_5$ from A.54. Induction from the index-five subgroup
$L$ gives a homomorphism
\[
 G_5\longrightarrow D\wr\mathfrak S_5.
\]
For a pure element its five diagonal blocks are the conjugate
forgetful images. Equation (FP.30.4) therefore makes
$\mathcal P(f)$ the identity in this induced quotient for every
2-cycle word, even if it is nontrivial in $P_5$. Changing coset
representatives only conjugates those five identity blocks.
Thus no enlargement of the modules \emph{inside the forgotten
four-point groups}, followed by this induction alone, makes its
pentagon equation a new condition. In particular the proposed
five-block linear calculation with $E_0\wr S_5$ has an identically
zero pentagon equation. Other consistency or lifting conditions are
not asserted to be automatic.

This obstruction does \emph{not} apply automatically to A.58:
the intersection of kernels of \emph{all} homomorphisms
$G_5\to\mathfrak S_m$ is smaller than the kernel of the one induced
representation, and need not contain $\mathcal B_5$. A genuine
five-point detector must retain an appropriate part of this Brunnian
subgroup as well as the frame data. The exact substitutions for all
five labels, the cycle permutation, and a nontrivial element of
$\mathcal B_5$ are checked in
\texttt{check\_brunnian\_detector.py}.

\begin{statement}{A.75}{Proved}{a small affine quotient retaining genuine Brunnian data}
\end{statement}
There is an explicit next quotient beyond the forgetful family.
In the literal Ferrari splitting of B.35, reduce the kernel's
abelianization modulo three and use the signed modular action:
\[
 \eta_3:L=B_4/\langle z_4\rangle
       \longrightarrow\mathbb F_3^2\rtimes\operatorname{SL}_2(\mathbb F_3),
 \qquad
 a\longmapsto(0,U),\quad b\longmapsto(0,L_0),\quad
 c\longmapsto(e_1,U).
 \tag{FP.31.1}
\]
The target has order $9\cdot24=216$ and this map is onto.
The three Artin relations hold by direct multiplication.
Its kernel image contains both basis translations: $ca^{-1}$ maps
to $e_1$, and conjugation by $b$ supplies the second direction.
The projection of $a,b$ generates $\operatorname{SL}_2(\mathbb F_3)$,
which proves surjectivity.
The element $d=abc$ maps to $(-e_2,S)$, so
\[
 \eta_3(d)^4=\bigl(-(I+S+S^2+S^3)e_2,I\bigr)=1.
\]
Thus the map really descends through the four-strand centre and is
a representation of the five-point stabilizer, not merely of $B_4$.

For the commutator convention $[g,h]=ghg^{-1}h^{-1}$, the element
\[
 w=[[X_1,X_2],[X_3,X_4]]
\]
lies in $\mathcal B_5$: under each pattern in FP.30.2 one of the two
inner commutators is one. Nevertheless exact multiplication gives
\[
 \eta_3(w)=
 \begin{pmatrix}1&0&0\\0&1&2\\0&0&1\end{pmatrix}\ne I.
 \tag{FP.31.2}
\]
This proves both that the Brunnian subgroup is nontrivial and that
the quotient (FP.31.1) retains data absent from all five forgetful
quotients. The same affine construction modulo two sends every
$X_i$ to the identity. Hence three is the smallest coefficient
modulus in \emph{this affine family} with nontrivial pure or Brunnian
data. No assertion that order $216$ is the smallest possible finite
Brunnian quotient is made.

For a diagnostic of the equation itself, the 2-cycle word
$f_0=[x,y]$ has
\[
 \eta_3(\mathcal P(f_0))=
 \begin{pmatrix}1&0&1\\0&1&1\\0&0&1\end{pmatrix}\ne I.
 \tag{FP.31.3}
\]
Thus the pentagon is no longer an identically invisible word in this
quotient. The pair $(1,f_0)$ is \emph{not} asserted to satisfy the
3-cycle relation or define an element of $\widehat{GT}_0$; this is
a check of the comparison map, not a mixed-sign shadow or a lifting
claim.

To retain both types of data, induce the product
\[
 L\longrightarrow E_0\times
          (\mathbb F_3^2\rtimes\operatorname{SL}_2(\mathbb F_3)),
 \qquad h\longmapsto(\theta q_5(h),\eta_3(h))
\]
to $G_5$ and, if a characteristic quotient is required, apply the
all-homomorphisms construction of A.58 to its finite coset
representation. The resulting quotient retains the actual frame
residues and the nontrivial Brunnian component (FP.31.2).
It is an explicit candidate for the next necessary pentagon equation.
The unresolved comparison is whether its equation, on actual
multiplier-one GT action residues, links the two frame parities;
neither (FP.31.2) nor (FP.31.3) proves such a link.

\subsection{An integral five-point parity coordinate and an affine pentagon equation}


\begin{statement}{A.76}{Proved}{an integral cubic frame character and the signed-lift correction}
\end{statement}
Let $\mathcal H_5=\widehat{GT}\cap\mathcal H$. On this group there
are canonical continuous additive characters
$\kappa_a,\kappa_b,\kappa_d:\mathcal H_5\to\widehat{\mathbb Z}$
such that, in the ordered cyclic basis $(X_1,\ldots,X_5)$ of
$P_5^{\mathrm{ab}}\cong\widehat{\mathbb Z}^{5}$,
\[
 \begin{split}
 \alpha_a^{\mathrm{ab}}&=\kappa_a(1,1,1,1,1),\\
 \alpha_b^{\mathrm{ab}}&=\kappa_b(-1,-2,-1,1,1),\\
 \alpha_d^{\mathrm{ab}}&=\kappa_d(0,0,0,0,1),\qquad
 \kappa_b=\rho/2.
 \end{split}
 \tag{FP.32.1}
\]
Moreover the cubic CM parity has the exact five-point expression
\[
 \boxed{\ \epsilon_3(T_\Phi)=\kappa_d(\Phi)\bmod2\ },
 \qquad
 \Delta(\Phi)=\kappa_d(\Phi)+\epsilon_2(W_\Phi)\pmod2.
 \tag{FP.32.2}
\]

We give the lattice and sign proofs. The five cyclic boundary classes
form a basis because the presentation of $P_5$ on them has only
commutator relators in abelianization. The other five boundary classes
are
\[
 [x_{i,i+2}]=-[X_i]-[X_{i+1}]+[X_{i+3}],
 \quad i\pmod5.
\]
These identities follow from the five sphere relations
$\sum_{j\ne i}[x_{ij}]=0$ together with the full-twist relation;
they also verify those relations directly. Conjugation on this lattice
is the ordinary label-permutation action. For $b=(1234)$ and
$d=(234)$ its matrices are
\[
 B=\begin{pmatrix}
 0&0&0&0&-1\\1&0&1&0&0\\0&1&0&0&1\\
 0&0&-1&0&0\\0&0&-1&1&-1
 \end{pmatrix},\qquad
 D=\begin{pmatrix}
 -1&0&0&-1&0\\-1&0&-1&0&0\\0&1&-1&1&0\\
 1&0&0&0&0\\0&0&1&-1&1
 \end{pmatrix}.
 \tag{FP.32.3}
\]
The action of $a$ is the cyclic permutation of the five basis vectors.
The kernels of $B-I,D-I,a-I$ are the three rank-one lattices in
(FP.32.1). This is an integral assertion: the nonzero Smith invariants
are respectively $(1,1,1,2)$, $(1,1,1,3)$ and $(1,1,1,1)$.
The same kernels over $\widehat{\mathbb Z}$ follow because every
nonzero integer is a non-zero-divisor there.

Each conjugating word $F_i$ in A.57 lies in $P_5'$. Hence a
multiplier-one genuine action is the identity on $G_5/P_5'$ and
$E_t^{\mathrm{ab}}=0$. Its frame must therefore lie in the indicated
fixed lattice. On $P_5^{\mathrm{ab}}$ the action is the identity,
so the crossed frame law becomes addition; uniqueness and continuity
of the frames give the three continuous characters.
The forgetful abelian map is
\[
 \begin{pmatrix}1&0&1&0&0\\0&1&0&0&0\end{pmatrix}.
\]
Its value on the $b$-fixed primitive vector is $(-2,-2)$, whereas
$\beta^{\mathrm{ab}}=(-\rho,-\rho)$ by R3. This proves
$2\kappa_b=\rho$ and the displayed normalization.

The signed projection of $\eta_3$ onto
$\operatorname{SL}_2(\mathbb F_3)$ is \emph{not} the signed lift
of the forgotten free-group frame. Write $\zeta:P_5\to\mathbb F_2$
for their difference:
\[
 Q_L(h)=(-I)^{\zeta(h)}q_3(q_5(h)),\qquad h\in P_5.
 \tag{FP.32.4}
\]
Both sides have the same projective image; their central ratio makes
$\zeta$ a continuous homomorphism. On the five generators its values
are $(0,0,0,1,1)$, because the signed images of $X_4,X_5$ are
$S^2=-I$, whereas their forgotten images are one.
Consequently $\zeta(\alpha_d)=\kappa_d\bmod2$.

By B.35 and R3, the actual action on the affine quotient satisfies
\[
 \eta_3\Phi_4=\operatorname{Inn}(0,V_\Phi)\eta_3,
 \qquad V_\Phi=U^{8\rho_\Phi}.
\]
Here the homology matrix is exactly $V_\Phi$: it fixes $e_1$,
commutes with $U$, and has the normalized covariance with $L_0$.
There is no remaining scalar choice in this step.
The element $d=a^2bc$ has affine image
\[
 \eta_3(d)=\left((-1,-1)^t,R_0^{-1}\right),\qquad
 R_0^{-1}=\begin{pmatrix}-1&1\\-1&0\end{pmatrix}\pmod3.
\]
Its centralizer projects exactly to $\langle R_0\rangle$, the
order-three subgroup of its order-six linear centralizer. Indeed
the translation has nonzero class in
$\mathbb F_3^2/(I-R_0^{-1})\mathbb F_3^2$; the order-three subgroup
fixes that class, while multiplication by $-I$ negates it.
The affine frame normalized by $(0,V_\Phi)$ centralizes this element.
Its linear projection is
$V_\Phi^{-1}(-I)^{\zeta(\alpha_d)}q_3(A_\Phi)$.
It lies in the order-three subgroup, so applying the unique sign
character of the order-six centralizer gives
$\epsilon_3(T_\Phi)=\zeta(\alpha_d)$, which proves (FP.32.2).
The exact centralizer has order nine and its linear projection has
order three; both are independently enumerated in the attached check.

The integral lattice alone supplies no equality with the harmonic
parity. In $G_5/P_5'$ the identity automorphism admits independently
every one of the three invariant frame vectors in (FP.32.1). Their
conjugations fix $a,b,d$ in that quotient, so all presentation relations
hold for every choice of coefficients. In particular the abelianization
of (FP.27.5) reduces only to $D\alpha_d^{\mathrm{ab}}
=\alpha_d^{\mathrm{ab}}$; it imposes no relation between
$\kappa_b$ and $\kappa_d$. These are abelianized primitive choices,
not actual GT frames or a lifting assertion. The remaining genuine
comparison is precisely
\[
 \epsilon_2(W_\Phi)\stackrel{?}{=}\kappa_d(\Phi)\pmod2
 \quad(\Phi\in\mathcal H_5).
\]
On Galois this follows from R4 and equals the specified common
quadratic character. Identifying the entire integral $\kappa_d$ with
a normalized half-Kummer character would require its own path and
normalization proof; only the parity identity is asserted here.
The lattice, signed correction, and both affine centralizers are
checked in \texttt{check\_five\_point\_lattice.py} and its output.

\begin{statement}{A.77}{Proved}{the exact affine pentagon equation and its limit}
\end{statement}
For $\Phi=(1,f)\in\mathcal H$, let
\[
 D_6=\langle p,t\mid p^3=t^2=(pt)^2=1\rangle\cong S_3,
 \qquad f(p,t)=p^{a_f},\quad a_f\in\mathbb F_3.
 \tag{FP.33.1}
\]
The exponent exists uniquely because $f\in\widehat F_2'$ and
$D_6'=\langle p\rangle$. Then the exact value of the cyclic
pentagon defect in the affine quotient of A.75 is
\[
 \boxed{\ \eta_3(\mathcal P(f))
             =\bigl((a_f-\rho_\Phi)e_1,I\bigr)\ }.
 \tag{FP.33.2}
\]
Consequently every genuine multiplier-one GT element satisfies the
necessary equation
\[
 a_f=\rho_\Phi\pmod3.
 \tag{FP.33.3}
\]
The difference $\eta_{\mathrm{aff}}(\Phi)=a_f-\rho_\Phi$ is a
continuous additive character $\mathcal H\to\mathbb F_3$.
This is an actual equation for globally defined pairs; finite
admissibility or an unspecified lift is not used.

Here is a complete calculation of (FP.33.2).
In $V\rtimes\operatorname{SL}_2(\mathbb F_3)$, $V=\mathbb F_3^2$,
the five pure generators have images
\[
 (0,U^2),\quad(0,L_0^2),\quad(-e_1,U^2),\quad
 (0,-I),\quad(e_2,-I).
 \tag{FP.33.4}
\]
Set $p=(e_1,I)$, $t=(0,-I)$ and write the same letters $X_i$
for these five images. The element $h_0=(0,U^2)$ commutes with
$p,t$, and $X_3=p^{-1}h_0$. Inserting a central argument in a
commutator word changes nothing. Conjugation by $t$ inverts $p$,
so
\[
 f(X_3,X_4)=f(p^{-1},t)=p^{-a_f}.
\]
For $t'=X_5$, one has $t'pt'^{-1}=p^{-1}$ and
$t'X_1t'^{-1}=pX_1$. Thus $h=p^2X_1$ commutes with $p,t'$,
and $X_1=ph$. A translation in the $e_2$ direction conjugates
$t$ to $t'$ while fixing $p$ and every element of $\langle p\rangle$.
The 2-cycle relation gives
\[
 f(X_5,X_1)=f(t',p)=f(p,t')^{-1}=p^{-a_f}.
\]
The pair $(X_2,X_3)$ is the conjugate, by the translation $-e_2$,
of the linear pair $(L_0^2,U^2)$. If
$F=f(U^2,L_0^2)$, its contribution is therefore
\[
 f(X_2,X_3)=\bigl((F^{-1}-I)e_2,F^{-1}\bigr).
\]
For the fourth factor put $q=(e_2,I)$, so $X_4=t$ and $X_5=qt$.
The multiplier-one 3-cycle relation at $(x,y,z)=(t,qt,q)$ is
\[
 f(q,t)f(qt,q)f(t,qt)=1.
\]
Conjugating $(t,q)$ by $q^2$ gives $(qt,q)$; the value of a
commutator word is in $\langle q\rangle$ and is unaffected by this
conjugation. Hence $f(qt,q)=f(t,q)=f(q,t)^{-1}$, and
$f(X_4,X_5)=1$. Finally $f(X_1,X_2)=(0,F)$.
Multiplying in the order of (FP.30.3) gives
\[
 \eta_3(\mathcal P(f))
       =\bigl(-2a_f e_1+(F^{-1}-I)e_2,I\bigr).
\]
R3 supplies
\[
 F=\begin{pmatrix}1&\rho\\\rho&1+\rho^2\end{pmatrix}\pmod3,
 \qquad (F^{-1}-I)e_2=-\rho e_1,
\]
which is exactly (FP.33.2).

To verify the character assertion, the standard pure automorphism of
the pair descends under $x\mapsto p,y\mapsto t$ to
\[
 p\longmapsto p,\qquad
 t\longmapsto p^{-a_f}tp^{a_f}.
\]
It is an inner automorphism of $D_6$ and is the identity on its
derived group. The GT composition formula for $f$ therefore adds
the two $a_f$ exponents. The harmonic coordinate $\rho$ is additive
on $\mathcal H$ by the abelianized frame law, proving the claim.

At affine modulus four, every matrix of a generator $X_i$ is
congruent to $I$ modulo two. The group
$1+2\operatorname{Mat}_3(\mathbb Z/4)$ is abelian. Hence every
$f(X_i,X_j)$ is already one for $f\in\widehat F_2'$, and the
pentagon has no further component there. CRT shows that the entire
affine modulus-twelve equation is precisely (FP.33.3).

This necessary equation cannot force the CM parity equality, even on
actual globally liftable elements of $\widehat{GT}_0$.
Let $g\in\mathcal H$ be the verified R2 element, so
$\Delta(g)=1$. Its \emph{GT-group cube} $g^3$ satisfies
\[
 \eta_{\mathrm{aff}}(g^3)=3\eta_{\mathrm{aff}}(g)=0,
 \qquad \Delta(g^3)=3\Delta(g)=1.
 \tag{FP.33.5}
\]
Thus its pair satisfies the full affine modulus-twelve pentagon
equation and has mixed CM parity. This is the cube in the group of
automorphisms, not the literal substitution $f\mapsto f^3$.
It is not claimed to satisfy the actual five-point pentagon, nor to
define an automorphism of every finite five-point quotient. What it
proves is the exact failure of this necessary finite equation to
imply $\Delta=0$, with an already global $\widehat{GT}_0$ witness.
The next comparison must retain additional Brunnian information
which is not reduced to this order-three character.
All finite affine identities used in the derivation are checked in
\texttt{check\_brunnian\_detector.py}; no pentagon-shadow census is
used as a substitute for a lifting theorem.

\subsection{A nontrivial kernel of descent and an actual cubic pentagon obstruction}


\begin{statement}{A.78}{Refuted}{the proposed nilpotent-residual exclusion at four}
\end{statement}
Let $U_0=F'F^3$, and evaluate $F$ in
$\mathrm{SL}_2(\mathbb Z/4)$ by $x\mapsto S$, $y\mapsto U^{-1}$.
The image of $U_0$ is the whole group of order $48$, and its
nilpotent residual is its derived subgroup of order $12$.
In particular it contains $-U^2$. There is even an explicit
actual profinite word $w\in R_{\rm nil}(U_0)$ with this value.
Thus membership $f\in R_{\rm nil}(U_0)$ alone cannot force
the first descent character to vanish.

\emph{Proof.}
The finite group has abelianization $C_4$, so commutators together
with cubes generate it. Its lower central series has successive
orders $48,12,12$, as the complete finite calculation verifies.
For an explicit source word let $e_p\in\widehat{\mathbb Z}$ be
the idempotent with $p$-adic component one and every other prime
component zero. Put
\[
 A=(x^6y^6)^{e_2},\qquad B=(x^3y^{-3})^{e_3},\qquad
 w=[BAB^{-1},B].
\]
Both $A,B$ lie in $U_0$. The first is a pro-two element and the
second is pro-three. Their indicated commutator maps to one in
every finite nilpotent quotient of $U_0$, whose different Sylow
subgroups commute. Hence $w\in R_{\rm nil}(U_0)$.
At the displayed modulo-four evaluation, $A=-U^2$ and $B=R_0$,
with respective orders two and three, and
\[
 [R_0(-U^2)R_0^{-1},R_0]=-U^2.
\]
This is an actual source counterexample to the residual exclusion.
It is not claimed to be a gentle parameter. The finite computations
are in \texttt{check\_iterated\_lifting.py}.

\begin{statement}{A.79}{Proved}{a simultaneous full-kernel word with nontrivial frame residue}
\end{statement}
Let $\varphi:F\to\widehat B$ be the current map of A.66, and let
\[
 K=\operatorname{core}_{\widehat B}(\ker\varphi)
   =\bigcap_{b\in\widehat B} b(\ker\varphi)b^{-1}
       \quad\text{as a closed subgroup of }F.
\]
Then the explicit word
\[
                   A_0=[[x^4,y^4],z^4]^2
\]
belongs to $F'\cap K$ and has $q_{12}(A_0)=[5I]$.

\emph{Proof.}
The kernel of $\varphi$ is normal in $F$, so the displayed core
requires only representatives for the six cosets $\widehat B/F$.
Their conjugations permute the three peripheral classes of $F$.
Under each resulting evaluation into $\widehat B$, one of the
three entries $x,y,z$ is conjugate to $s$ and has order two.
Its fourth power is therefore one. The iterated commutator
$[[x^4,y^4],z^4]$ vanishes if any of its three entries is one.
Thus $A_0$ lies in all six kernels in the full profinite target,
not just at a congruence level. It is a commutator word, so belongs
to $F'$ as well. For the original inclusion use
$x=U^2$, $y=L^2$. Direct multiplication modulo twelve gives
$q_{12}(A_0)=[5I]$; the accompanying program gives all seven
matrix values, the six other values being $I$.

\begin{statement}{A.80}{Proved}{the quadratic character is nontrivial on the kernel of descent}
\end{statement}
There exists an actual $k\in\mathcal H$ such that
\[
 \boxed{\ D(k)=1,\qquad\Delta(k)=1.\ }
\]
It can simultaneously be chosen with
\[
 f_k\in R_{\rm nil}(U_0),\qquad
 q_{12}(A_k)=[5I],\quad q_{12}(\beta_k)=1,\quad \rho_k=0,
 \qquad U_0=F'F^3.
\]
In particular it is invisible on every absolute pro-prime pure action.
This is a full profinite existence theorem, not a finite survival result.

\emph{Proof.}
Let $J\triangleleft\widehat B$ be any open normal subgroup, and set
\[
 N_J=U_0\cap\operatorname{core}_{\widehat B}
                              \bigl(\varphi^{-1}(J)\bigr).
\]
This is open and normal in $\widehat B$, contained in $F$.
The word $A_0$ of A.79 belongs to every $N_J$.
Retain $Q=\mathrm{SL}_2(\mathbb Z/12)/\{\pm I\}$ and
$q=q_{12}:\widehat B\to Q$. R2's finite calculation gives
$q(U_0)=V_4\times Q_8$, a two-group. Consequently $q(N_J)$
is a two-group containing $q(A_0)=[5I]$, and $q$ factors through
$\Gamma_2=\widehat B/O^2(N_J)$.

Put $D_J=\widehat B/N_J$ and retain the actual joint image
\[
 E_J=\operatorname{im}(\widehat B\to Q\times D_J).
\]
Its kernel over $D_J$ is exactly $q(N_J)$, and the kernel of
$\Gamma_2\twoheadrightarrow E_J$ is pro-two. The element
$([5I],1)$ belongs to $E_J$, because it is the image of $A_0$,
and it centralizes the image of $r$. The order of that image is
three, detected by the retained $S_3$ quotient. Coprime
centralizer lifting therefore produces
\[
 c_2\in C_{\Gamma_2}(r),\qquad
 c_2\longmapsto([5I],1)\in E_J.
\]
Here the elementary lifting argument is the one used in R2:
two complements to a pro-two kernel for the same order-three
subgroup are conjugate by the kernel, with the same generator
image; finite complement conjugacy and compactness give the
profinite assertion.

Use the identity automorphism in every relative pro-$p$ component
of $\widehat B/R_{\rm nil}(N_J)$. Give $s,\tau$ conjugator one,
and give $r$ conjugator $c_2$ at two and one at all other primes.
All these coordinates have the same value one in $D_J$.
The relative product theorem and R2 lift them to an actual full
peripheral automorphism. Normalize its cusp by the inverse of
the lifted cusp conjugator, which lies in $R_{\rm nil}(N_J)$.
The resulting normalized gentle pair $k_J=(1,f_J)$ has
\[
 f_J,\beta_J\in R_{\rm nil}(N_J),\qquad
 q(A_J)=[5I],\quad q(\beta_J)=1.
\]
Since $R_{\rm nil}(N_J)\subseteq N_J'\subseteq F'$, its
harmonic frame has abelianization zero, hence $\rho_J=0$.
R4 now gives $\Delta(k_J)=1$. Moreover
$\varphi(f_J)\in J$.

These are witnesses for a compactness argument on the actual
compact group $\mathcal H$. Impose the closed conditions
\[
 f\in R_{\rm nil}(U_0),\quad\beta\in R_{\rm nil}(U_0),
 \quad q(A)=[5I],\quad q(\beta)=1,
 \quad\varphi(f)\in J.
\]
Every finite collection is met by the construction for the
intersection of its open normal subgroups $J$.
Indeed $N_J\subseteq U_0$ implies
$R_{\rm nil}(N_J)\subseteq R_{\rm nil}(U_0)$.
The normalized parameter and the unique pure frames depend
continuously on the gentle element, so all these sets are closed.
Compactness gives an actual $k$ satisfying every condition and
$\varphi(f_k)=1$.

In A.66 this says $C=1$, $p=q=a=0$, and $\beta_D=1$.
The normalized descendant fixes both $s$ and $\tau$, hence is
the identity. Thus $D(k)=1$, while the retained frame values
give $\Delta(k)=1$. The inclusion
$R_{\rm nil}(U_0)\subseteq R_{\rm nil}(F)$ proves the final
absolute pronilpotent invisibility assertion.

\begin{statement}{A.81}{Proved}{consequences for fixed subgroups and iterated characters}
\end{statement}
The descent endomorphism $D$ is not injective. Both
$\operatorname{Fix}(D)$ and the correctly marked
$\operatorname{Fix}(\operatorname{Ad}_c\circ D)$ are proper
closed subgroups of $\widehat{\mathrm{GT}}_0$. For every $n\geq1$,
\[
 \Delta\circ D^n\ne\Delta,
\]
and $\Delta$ does not belong to the finite linear span of these
positive iterates. More strongly,
\[
 (D,\Delta):\widehat{\mathrm{GT}}_0
       \longrightarrow\operatorname{im}(D)\times\mathbb F_2
\]
is onto; the corresponding multiplier-one statement is also true.
Using the independently proved arithmetic marking comparison
$D(\sigma)=c\sigma c$ on $G_{\mathbb Q}$, every
$\chi_n=\Delta\circ D^n$ vanishes on the Galois image.
Consequently $\chi_1$ is either zero or linearly independent of
$\Delta$ as a relative quadratic character. Its nonvanishing
is not established by the present construction.
Moreover, neither fixed subgroup contains $\ker\Delta$.
In particular the proposed equality
$\operatorname{Fix}(\operatorname{Ad}_c\circ D)=\ker\Delta$
is false even in the corrected original marking. More precisely,
$\ker D\cap\mathcal H$ contains a closed subgroup
$Z\simeq\mathbb Z_2$ on which $\Delta$ is reduction modulo two.
Every nonidentity element of $2Z$ lies in $\ker\Delta$ and in
neither fixed subgroup. No pentagon assertion about these
elements is made.

\emph{Proof.}
Use the actual element $k$ of A.80. It is nonidentity since
$\Delta(k)=1$, but both descent endomorphisms send it to one.
Thus it is in neither fixed subgroup. Every positive iterate
$D^n(k)$ is one. Evaluating the proposed equality of characters,
or a finite linear combination of positive iterates, at $k$
proves the stated distinctions. Finally, for any chosen preimage
$g$ of a point of $\operatorname{im}(D)$, the elements $g,gk$
have the same descendant and opposite $\Delta$ values.
Both stay in $\mathcal H$ when $g\in\mathcal H$.
For the arithmetic assertion use the simultaneous positive-loop
comparison proved in A.62: $D$ preserves the Galois image and
acts there by conjugation by $c$. R4 therefore makes every
$\chi_n$ vanish on that image. Since $\chi_1\ne\Delta$,
the only remaining possibilities in a vector space of quadratic
characters are $\chi_1=0$ or linear independence.

For the fixed-subgroup comparison, R7 implies $k^2\ne1$,
because $k\in\mathcal H$ is nonidentity and $\mathcal H$ is
torsion-free. But $D(k^2)=1$ and $\Delta(k^2)=0$.
Both $D$ and $\operatorname{Ad}_c\circ D$ send $k^2$ to one,
so neither fixes it. For the stronger closed-subgroup statement
put $k_2=k^{e_2}$, using the prime idempotent of A.78.
Continuity gives $D(k_2)=1$ and $\Delta(k_2)=1$.
The closed subgroup $Z=\overline{\langle k_2\rangle}$ is
procyclic and pro-two, nontrivial by its quadratic value, and
torsion-free by R7. Hence $Z\simeq\mathbb Z_2$.
Its nonzero homomorphism to $\mathbb F_2$ is reduction modulo
two, with kernel $2Z$. The two descent endomorphisms are
trivial on $Z$, proving all remaining assertions.

\begin{statement}{A.82}{Proved}{an actual independent affine-pentagon obstruction}
\end{statement}
Let $p,t$ generate $S_3$ with $p^3=t^2=(pt)^2=1$, and define
$a_f\in\mathbb F_3$ by $f(p,t)=p^{a_f}$.
There exists an actual multiplier-one gentle pair with
\[
                  a_f=0,\qquad \rho=-1\pmod3.
\]
It is identity on every absolute pro-prime pure action. Combined
with the independently verified five-point affine pentagon formula
$(a_f-\rho)e_1$, this is an actual profinite counterexample to the
pentagon, detected by that genuine five-point quotient.

\emph{Proof.}
Let $N$ be the $\widehat B$-normal core of the kernel of
$F\to S_3$, $x\mapsto p$, $y\mapsto t$. Every geometric
permutation of the boundary triple gives a generating $S_3$
triple with orders $3,2,2$. The two explicit words
\[
 b_x=x^2yx^2y^{-1},\qquad b_y=y^2xy^2x^{-1}
\]
vanish under every such evaluation. For $b_x$, either the image
of $x$ has order two, or it has order three and the image of $y$
is an involution inverting it. These are exactly the two cases
and both give one. The proof for $b_y$ is identical. Their
abelianizations are $(4,0)$ and $(0,4)$, so $N$ maps onto
$F/F'F^3=\mathbb F_3^2$, and $b_0=b_xb_y$ has image $(1,1)$.

Use the actual finite quotient $Q_3=\widehat B/(F'F^3)$ of
order $54$. Its pure kernel is $\mathbb F_3^2$, and $s$
interchanges the two coordinates. The image of $N$ is this
three-group, and the image of $b_0$ is the nonzero diagonal
translation, which centralizes $s$.
In the actual joint quotient of $Q_3$ and $\widehat B/N$,
the element $((1,1),1)$ therefore lifts to a centralizer of $s$
in $\widehat B/O^3(N)$, because $s$ has order two and the
kernel is pro-three.

Choose the identity relative automorphism, this conjugator
for $s$ at three, and conjugator one for $s$ at every other
prime and for $r,\tau$ at every prime. R2 and cusp normalization
give an actual full gentle pair with
\[
 f\in R_{\rm nil}(N)\subseteq N,\qquad
 \beta^{\rm ab}\equiv(1,1)\pmod3.
\]
The correction at the cusp lies in $R_{\rm nil}(N)$ and thus
has zero abelianization, so the retained frame congruence is
literal. It gives $\rho=-1\pmod3$, whereas $f\in N$ gives
$f(p,t)=1$ and $a_f=0$. Finally
$R_{\rm nil}(N)\subseteq R_{\rm nil}(F)$ proves the claimed
absolute pronilpotent invisibility. The finite seed identities
are checked on all eighteen generating pairs of $S_3$ by the
accompanying program. Existence of the full pair is supplied
by the displayed relative lift, not inferred from that enumeration.

\begin{statement}{A.83}{Proved}{all iterated quadratic tests miss an actual nonpentagonal element}
\end{statement}
There is an actual pro-three element $h\in\mathcal H$ which is
identity on every absolute pro-prime pure action, fails the pentagon,
and satisfies
\[
                   \Delta(D^n h)=0\qquad(n\geq0).
\]
Its closed procyclic subgroup is isomorphic to $\mathbb Z_3$.
Thus the intersection of all these quadratic kernels, even together
with all absolute pronilpotent tests, is strictly larger than its
intersection with $\widehat{\mathrm{GT}}$.

\emph{Proof.}
Take the actual element $h_0$ from A.82 and put
$h=h_0^{e_3}$, where $e_3$ is the prime idempotent of A.78.
The two maps $a_f,\rho\bmod3$ are additive on $\mathcal H$.
For $a_f$, evaluate the normalized action in the preserved quotient
$S_3=\langle p,t\rangle$: it fixes $p$ and conjugates $t$
by a power of $p$, so its action on $\langle p\rangle$ is trivial.
The pair multiplication law therefore adds the exponent $a_f$.
Additivity of $\rho$ is R4's frame multiplication law.
It follows that $(a_f-\rho)(h)=1$, since $e_3\equiv1\pmod3$.
The affine pentagon detector consequently remains nonzero.
Absolute pro-prime invisibility is preserved by taking a profinite
power. The closed subgroup generated by $h$ is pro-three, so every
continuous homomorphism from it to $\mathbb F_2$ is zero.
Apply this to each $\Delta\circ D^n$.
Finally $h\ne1$, and R7 excludes nontrivial finite-order elements
in $\mathcal H$. A nontrivial torsion-free procyclic pro-three
group is $\mathbb Z_3$, which proves the last assertion.
This conclusion only identifies an actual nonpentagonal element
in the joint zero locus. It assumes no vanishing of any
$\chi_n$ on all of $\widehat{\mathrm{GT}}$.

\subsection{The exact remaining decision and comparisons}


\begin{statement}{A.84}{Open}{The main quadratic decision on the full profinite pentagon group}
We have not proved or refuted
\[
                    \Delta|_{\GT}=0.
\]
With $\mathcal H_5=\GT\cap\HH$, the exact remaining identity
can now be written
\[
 \epsilon_2(W_h)=\kappa_d(h)\pmod2
                      \qquad(h\in\mathcal H_5),
\]
where $\kappa_d$ is the canonical integral character of A.76.
Equivalently, by the genuine quadratic pentagon consequence
A.112, the remaining identity is
\[
                 \nu_4(h)=\kappa_d(h)\pmod2
                    \qquad(h\in\mathcal H_5).
\]
Here $\nu_4$ is the explicit continuous $S_4$ evaluation
character, not a choice of a finite frame primitive.
No assertion in this return assigns a pentagon value to the R2
element, its cube, or the element of $\ker D$ in A.80.
\end{statement}
The equivalence with the full-multiplier question is exact.
For $g\in\GT$, choose $\sigma\in G_{\Q}\subseteq\GT$ with
$\chi(\sigma)=\lambda(g)$. Then $h=g\sigma^{-1}\in\mathcal H_5$
and $\Delta(g)=\delta(h)$. Conversely a counterexample in
$\mathcal H_5$ is already one in $\GT$. A.76 proves
$\delta(h)=\kappa_d(h)+\epsilon_2(W_h)$, giving the displayed
form. On the Galois subgroup both bits equal the quadratic
character of $\Q^{\rm ab}(\sqrt[4]3)/\Q^{\rm ab}$; on every
closed AA subgroup vanishing follows from A.32.

The progress and the exact failures of the tested routes are as
follows. The pointed category kills every commutator parameter,
so its coherence cannot distinguish the R2 witness. The explicit
finite characteristic quotient of A.54--A.59 does retain the
actual frame signs and the outer-action ambiguity. However,
A.74 proves that quotients induced only from forgetting points
annihilate every two-cycle pentagon defect. The genuine affine
Brunnian quotient of A.75 imposes the proved equation
$a_f=\rho\pmod3$, but the actual R2 cube satisfies that equation
at modulus twelve while retaining $\Delta=1$ (A.77).
Thus that equation does not force the displayed quadratic equality.
A.112 now adds the actual quadratic consequence
$\epsilon_2(W)=\nu_4$; its finite cubic-frame ambiguity explains
exactly why this does not yet identify $\nu_4$ with $\kappa_d$.
The next required comparison must retain additional five-point
information, or identify these last two characters by an
independent theorem.

The mixed beta reduction A.103 expresses the same problem as
$\Delta=d_2+e_3$, so the required conclusion is $d_2=e_3$;
proving both zero is sufficient but is a stronger assertion.
A.107 proves that either descent fixed subgroup kills $d_2$.
A.109 supplies an actual gentle element with every cyclic-cover
Kummer coefficient zero, beta series one, and nevertheless
$\Delta=1$. Thus even the entire coefficient package leaves
the cubic frame comparison essential.

A.94 gives a single actual $k\in\mathcal H$ with
$\Delta(k)=1$, $D(k)=1$, and
$f_k\in R_{\rm nil}(F'')\subseteq F'''$.
A.96 and A.113 prove that it passes all pentagon evaluations
in the specified pure-solvable-length-three quotients, the
entire two-adic Burau family, and the first relative
characteristic-three homology extension of the order-$216$
five-point quotient. These are exact failures of those tests,
not hypothetical finite survivors. A new test must escape
their nilpotent second-derived evaluation images.
A.114 strengthens the Burau conclusion to the entire adelic
congruence family at a primitive cube root of unity, using
one actual witness and exact integral identities.
A.128 strengthens that construction simultaneously to the
whole quotient $L/R_{\rm nil}(C')$, where
$C=\ker(L\to\widehat B)=C_2\amalg C_2\amalg C_2$.
It retains the full profinite forgetful base. A.129 also proves
that the corresponding braid formulas define an actual
automorphism of this quotient, fixing its Ferrari kernel
pointwise. Thus every finite five-point test with nilpotent
image of $C'$ is insufficient for this implication.
The next cover uses $D=C''(C')^2$, free of rank $129$, and
the actual action on $H_1(D,\mathbb F_3)$. A.131 finds a
nonnilpotent harmonic second-derived image here, independently
checked by a second program. Nevertheless A.140 proves
$\mathcal P(f)\in D$ for every actual multiplier-one gentle
element. Its conjugation action on every $H_1(D,R)$ is
therefore identity. Thus even this nonnilpotent linear test
is universally insufficient. The next precise target is
$L/(D'D^3)$, where the residual class in $D/(D'D^3)$
survives as an affine translation. Linear homology discards
exactly that translation; retaining it is a distinct test.
A.147 now constructs an actual mixed-sign gentle element
passing this affine test as well, simultaneously with the
earlier complete families. Its cyclic seed has nilpotent
normal image in every evaluation and every peripheral core;
this is a proved input to R2, not a finite-shadow lift.
A.148 upgrades it to an actual automorphism of
$L/(D'D^3)$ fixing $C/(D'D^3)$ pointwise.
A.153 supplies that separate higher-level argument: exact
transfer isolates the rank-five factor, and congruence
kernels control its complementary affine factor. One
actual witness therefore satisfies the pentagon in the
whole quotient $L/D'$, and A.154 proves its actual
automorphism there, fixing $C/D'$ pointwise. This retains
every coefficient level of $D^{\rm ab}$; it does not
retain the nonabelian pro-prime completions of $D$.
The next specified cover is $D_4=C''(C')^4$. Its full
Fox chain representations retain the augmentation line
as well as the homology, so their translation information
is not silently discarded.
The canonical homology detector A.98 and its full-multiplier
bound A.110 do detect $\Delta$ as an action character, but no
pentagon vanishing theorem for that action image is proved.

The mixed-sign restricted lifting problem in A.51 is still
unproved: R1 and R2 do not lift the five correlated evaluations
of one pentagon parameter. We have no compatible system giving
an actual mixed-sign element of $\GT$. Nor have we computed a
finite detector at which the pentagon excludes every actual
gentle residue of sign one. The smallest such level is unknown.
These are the precise two missing alternatives, not an inference
from the survival of any finite shadow.

If vanishing is eventually proved, P.13 has already established
that $\GT\subsetneq\ker\Delta\subsetneq\GZ$.
If it fails, the restriction is onto $\Ftwo$, and its kernel is
an index-two closed subgroup of $\GT$ containing $G_{\Q}$.
The latter alternative would prove $G_{\Q}\ne\GT$; no element
establishing that alternative is claimed here.

\begin{statement}{A.85}{Open}{The corrected compatibility subgroup and the first iterated character}
Put $C_L=\operatorname{Fix}(\operatorname{Ad}_c\circ D_P)$.
The assertions $\GT\subseteq C_L$, $C_L\subseteq\ker\Delta$,
and the proposed characterization of $C_L$ or $\ker\Delta$
by the harmonic and equianharmonic equations alone remain
undecided here. The literal-current assertions
$\GT\subseteq\operatorname{Fix}(D_P)$ and
$\operatorname{Fix}(D_P)=\ker\Delta$ are already false.
\end{statement}
A.60--A.65 prove the marking conversion and the exact
GF0/HF0 primitive comparisons. A.80 proves $C_L$ proper,
without assuming any parity theorem for its elements.
The equality $C_L=\ker\Delta$ is also false: the element
$k^2$ from A.80 is nonidentity by R7, has $\Delta(k^2)=0$,
and has descendant one, so lies outside either fixed subgroup.
The character $\chi_1=\Delta\circ D_P$ vanishes on Galois
and differs from $\Delta$. It is either zero or independent
of $\Delta$, by A.81. If it is zero, then both fixed subgroups
lie in $\ker\Delta$. We have proved neither its zero value
nor a nonzero value on any actual gentle element.
The common-word finite calculation A.70 retains both bits
but supplies neither missing conclusion. The stronger
$D_P(k)=1$, $\Delta(k)=1$ theorem proves noninjectivity of
descent, not vanishing of the character on its image.
A.107 does prove that both fixed subgroups annihilate the mixed
duplication defect $d_2$. Their remaining $\Delta$ comparison
is precisely the cubic defect, not the already settled
harmonic half-Kummer identity.

\begin{statement}{A.86}{Open}{The well-defined level-six Atkin--Lehner comparison}
For the crossed defect $AL_{6,3}$ defined in A.44, the
equivalence $AL_{6,3}(g)=1\Longleftrightarrow\Delta(g)=0$
on its specified gentle domain has not been proved or refuted.
The actual free-group value of this defect on the R2 element
has not been computed. The analogous full-GT compatibility
of the two actions at the section with coordinate nine
also remains undecided.
\end{statement}
The rational four-cusp model and all three Atkin--Lehner maps
are explicit in A.42. Arithmetic functoriality proves the
Galois comparison. A.45 proves that the defect always fixes
the four cusps, so the proposed cusp permutation test cannot
detect its value. A calculation on a genuinely marked free
basis, with the section and conjugating paths retained,
is still required. The general $\Delta_N$ construction of
Lead 8 was refuted in A.7--A.8; it cannot fill this gap.

\subsection{The harmonic sign as a mixed beta phase}


\begin{statement}{A.87}{Proved}{the harmonic parity is an integral diagonal beta phase}
\citep[Proposition 4.4]{LitNakamuraTsunogaiYasuda2010}
\end{statement}
There is an exact finite mixed-coefficient formula for the other
parity. Set
\[
 R=(\mathbb Z/4)[\omega]/(\omega^2+\omega+1),\qquad
 c=(1-\omega)^{-1}=2+3\omega.
\]
In the finite affine group $R_{\rm add}\rtimes\langle-\omega\rangle$
define
\[
 s\longmapsto(c,-1),\qquad r\longmapsto(c,\omega).
 \tag{FP.34.1}
\]
Their orders are two and three. Thus this is a continuous
homomorphism from $\widehat B=C_2\amalg C_3$; its image has order
$96$. The exact peripheral images are
\[
 \tau=sr\longmapsto(0,-\omega),\qquad
 x=(sr)^{-2}\longmapsto(0,\omega),\qquad
 y=sxs\longmapsto(1,\omega).
\]
For $\Phi=(1,f)\in\mathcal H$, write the image of $f$ as
$(k_f,1)$ and set
\[
 \mathcal B_f=1-(1-\omega)k_f\in R.
\]
Then
\[
 \boxed{\ \mathcal B_f
       =\omega^{\rho_\Phi}(-1)^{\epsilon_2(W_\Phi)}\quad\text{in }R\ }.
 \tag{FP.34.2}
\]
This formula retains the coefficient-two component together with
the order-three eigenvalue, and does not pass to a pro-two quotient.

To prove it, write the image of the actual harmonic frame as
$\beta=(b,\omega^\rho)$. Its linear component is as stated because
$\beta^{\mathrm{ab}}=(-\rho,-\rho)$, so it is
$\omega^{-2\rho}=\omega^\rho$. The exact equality
$\beta s\beta^{-1}=sf$ becomes
\[
 2b+\omega^\rho c=c-k_f.
 \tag{FP.34.3}
\]
The image of the shear $V=U^{8\rho}$ is $(0,\omega^\rho)$,
because $U=\tau^{-1}$ has zero translation and linear component
$-\omega^{-1}$. Hence
\[
 W=\beta^{-1}V\longmapsto(w,1),\qquad
 w=-\omega^{-\rho}b.
\]
Modulo two, (FP.34.1) is the projective level-three quotient
$A_4=\mathbb F_4\rtimes C_3$ with the prescribed images of
$x,y,s,r$. The corrected harmonic frame belongs to the cyclic
centralizer of the signed $S$; its projective image is therefore
either $1$ or $s$. Accordingly
\[
 w/c\bmod2=\epsilon_2(W)\in\{0,1\}\subset\mathbb F_4.
\]
Multiplying (FP.34.3) by $c^{-1}$ yields
$2b/c=\mathcal B_f-\omega^\rho$, and hence
\[
 \mathcal B_f=\omega^\rho(1-2w/c)
             =\omega^\rho(-1)^{\epsilon_2(W)}\pmod4.
\]
This proves (FP.34.2), including the division by two which would
be lost by an initial reduction modulo two.

In the usual metabelian notation
$f\equiv A_f(x,y)[x,y]\pmod{F''}$, with
$[x,y]=xyx^{-1}y^{-1}$, the commutator image in (FP.34.1) is
$\omega-1$. Therefore
\[
 k_f=(\omega-1)A_f(\omega,\omega),\qquad
 \mathcal B_f=1+(\omega-1)^2A_f(\omega,\omega)
            =1-3\omega A_f(\omega,\omega).
 \tag{FP.34.4}
\]
The plus sign in this intrinsic multiplier-one beta normalization
agrees with the opened primary proof of
Nakamura--Tsunogai--Yasuda (2010), Proposition 4.4, printed p.~443.
Identifying its exact labelled value with Ihara's
$B_\Phi(1/3,1/3)$ additionally requires matching the free-group
marking and the action convention. Simultaneous inversion of the
roots of unity, or inversion of the resulting beta value, leaves
its sign unchanged but can change the displayed power of $\omega$.
No distribution formula with its cyclotomic exponent is imported
without that comparison. The intrinsic formula (FP.34.2) is
independent of this naming issue.

For genuine GT elements, combining (FP.32.2) and (FP.34.2) reduces
the main sign equality to the precise finite identity
\[
 \mathcal B_f\stackrel{?}{=}
            \omega^{\rho_\Phi}(-1)^{\kappa_d(\Phi)}
 \quad\text{in }(\mathbb Z/4)[\omega].
 \tag{FP.34.5}
\]
Full profinite beta/gamma factorization supplies
$B(s,t)=\Gamma(s)\Gamma(t)/\Gamma(s+t)$, but that formula alone
does not identify the finite phase in (FP.34.5) with a Kummer or
cubic-frame coordinate. Ihara's further multiplication conditions
must be distinguished from the proved factorization; their use
would require a separate theorem valid on all $\widehat{GT}$.
The exact ring calculations, group order, peripheral identities,
and all six $(\rho\bmod3,\epsilon_2)$ phase cases are checked in
\texttt{check\_harmonic\_beta\_phase.py} and its output.

\subsection{An integral Kummer half-character and the limit of factorization}


\begin{statement}{A.88}{Proved}{a full profinite definition with no logarithm}
\citep[Theorem 1 and Section 2]{LitIhara2000BetaGammaII}
\end{statement}
Let $F=\widehat F(x,y)$, $N_n=\ker(F\to C_n)$, where $x$ maps
to a chosen generator and $y$ to one.  A free profinite basis is
\[
 x^n,\quad y_j=x^jyx^{-j}\quad(0\leq j<n).
\]
For $f\in F'$, write its image in $N_n^{\rm ab}$ as
\[
 0\cdot[x^n]+\sum_{j=0}^{n-1}m_j[y_j].
\]
The coefficient of $x^n$ is zero and $\sum_jm_j=0$, because
$f$ has zero exponent sums in $F^{\rm ab}$.
Define $\rho_n(f)=m_0$.
For normalized automorphisms $\Phi_g(x)=x^{\lambda_g}$,
$\Phi_g(y)=f_g^{-1}y^{\lambda_g}f_g$ with $f_g\in F'$, this is
a continuous crossed homomorphism:
\[
 \rho_n(gh)=\rho_n(g)+\lambda_g\rho_n(h).
\]
Indeed, $N_n$ is preserved, and its abelianization sends
$[y_j]$ to $\lambda_g[y_{\lambda_gj}]$ and $[x^n]$ to
$\lambda_g[x^n]$.  Conjugation by a power of $x^n$ does not
affect this statement.  Since multiplication of normalized pairs
gives $f_{gh}=f_g\Phi_g(f_h)$, the displayed formula follows.
Thus the restriction to the multiplier-one group is an additive
character with values in $\widehat{\mathbb Z}$.

This is the $\tau_n^{(0)}$ definition in Ihara's primary
RIMS-1224 manuscript, pp.3 and 6 (the preprint of
\emph{On beta and gamma functions associated with the
Grothendieck--Teichm\"uller group. II}, J. reine angew. Math.
527 (2000), 1--11). Theorem 1, manuscript p.4, writes
$(n^{-r})^{\sigma-1}=e[r\tau_n^{(0)}(\sigma)]$.
Thus, in that paper's monodromy convention, it is the negative
of the exponent $k_n$ defined by
$\sigma(n^{1/m})/n^{1/m}=\zeta_m^{k_n(\sigma)}$ for the positive roots.
Replacing both free generators by their inverses also changes the
Schreier coefficient $m_0$ to $-m_0$.
The even half modulo two used below is invariant under either change;
no identification of the integral signs in two different markings
is being made.

\begin{statement}{A.89}{Proved}{an order-48 formula for $\rho_3$ modulo four}
\citep[p.~657]{LitSchneps2006}
\end{statement}
Let
\[
 R=\mathbb Z_2[\omega],\qquad \omega^2+\omega+1=0,
 \qquad (t,u)(t',u')=(t+ut',uu')
\]
in $R\rtimes\langle\omega\rangle$.  Evaluate at
$X=(0,\omega)$ and $Y=(1,1)$, and write
$f(X,Y)=(t_f,1)$ for $f\in F'$.  Then
\[
 \rho_{3,2}(f)=\frac{\operatorname{Tr}_{R/\mathbb Z_2}(t_f)}3.
 \tag{A.89}
\]
In fact $y_j$ maps to $(\omega^j,1)$, so
$t_f=m_0+m_1\omega+m_2\omega^2$.
Using $m_0+m_1+m_2=0$ and
$\operatorname{Tr}(1)=2$, $\operatorname{Tr}(\omega)
=\operatorname{Tr}(\omega^2)=-1$ gives
$\operatorname{Tr}(t_f)=3m_0$.
Modulo four the target is the explicitly finite group
\[
 ((\mathbb Z/4)[\omega],+)\rtimes C_3
\]
of order $48$.  If $t_f=a+b\omega$, the formula reads
$\rho_{3,2}\equiv3(2a-b)\pmod4$.

Equivalently, with $[x,y]=xyx^{-1}y^{-1}$ and
$f\equiv A_f(x,y)[x,y]\pmod{F''}$, one has
\[
 t_f=(\omega-1)A_f(\omega,1),\qquad
 \rho_{3,2}(f)=\frac13\operatorname{Tr}
       ((\omega-1)A_f(\omega,1)).
\]
This is the coefficient-two component of the displayed formula
in Schneps (2006), p.657.  The Schreier proof fixes the additive
normalization independently of the evident typographical errors
in the target and root-of-unity subscript printed beside that formula.
The finite target retains order three; it is not a pro-two
quotient of $F$.

\begin{statement}{A.90}{Proved}{the half-character is integral}
\end{statement}
Suppose $\Phi=(1,f)$ is an actual gentle element, with the standard
two-cycle and three-cycle equations, so that
\[
 \Phi(z)=f(z,x)z f(z,x)^{-1},\qquad z=(xy)^{-1}.
 \tag{KM.3a}
\]
Then
\[
 t_f\in\mathbb Z_2\subset R,\qquad
 \rho_{3,2}(f)\in2\mathbb Z_2,\qquad
 \frac{\rho_{3,2}(f)}2=\frac{t_f}3.
 \tag{KM.3b}
\]
In particular its half modulo two is computed in the order-twelve
group $\mathbb F_4\rtimes C_3\simeq A_4$: the translation of
$f(X,Y)$ is either $0$ or $1$, and is exactly this bit.
Since two is a unit in every odd-adic factor, the full
$\rho_3$ consequently lies in $2\widehat{\mathbb Z}$.
Its unique half is a continuous additive character on the
multiplier-one gentle group, by A.88 and injectivity of
multiplication by two on $\widehat{\mathbb Z}$.

For completeness, (KM.3a) follows directly from the two three-cycle
equations available for an actual gentle element.
Put $g=f(z,x)$ and $h=f(y,z)$. The multiplier-one equation gives
$ghf=1$. Right composition with $c=(-1,1)$ gives the normalized
pair $(-1,f)$, whose three-cycle, with $m=-1$, is
$gz^{-1}hy^{-1}fx^{-1}=1$. Substituting $h=g^{-1}f^{-1}$ yields
$gz^{-1}g^{-1}=xf^{-1}yf=\Phi(z^{-1})$, proving (KM.3a).
This uses the stated group structure of the actual gentle group,
not the pentagon or a lifting assertion.

Here are the remaining details, to separate this elementary evenness from
any extra gamma multiplication relation. First evaluate at
$X_0=(0,\omega)$, $Y_0=(1,\omega^{-1})$.  The product
$X_0Y_0$ is a translation.  The conjugator in (KM.3a) is also a
translation, since $f\in F'$, and therefore acts trivially on
$X_0Y_0$.  If $f(X_0,Y_0)=(v,1)$, the equality
$\Phi(xy)=xy$ in this target gives
$(1-\omega)v=0$.  Since $1-\omega$ is a unit in $R$, $v=0$.
Since $[X_0,Y_0]=(\omega-1,1)$, this proves
\[
 A_f(\omega,\omega^{-1})=0.
 \tag{KM.3c}
\]
The two-cycle implies $A_f(U,V)=A_f(V,U)$ in the free
metabelian module.  Now return to $X=(0,\omega)$,
$Y=(1,1)$ and $Z=(XY)^{-1}=(-1,\omega^{-1})$.
The three-cycle is
$f(Z,X)f(Y,Z)f(X,Y)=1$.  Its first factor is one by
(KM.3c); its other translations are, respectively,
\[
 (1-\omega^{-1})A_f(1,\omega^{-1})=-\overline{t_f},
 \qquad t_f.
\]
Consequently $t_f=\overline{t_f}$, where the bar exchanges
$\omega$ and $\omega^{-1}$.  Its fixed ring is $\mathbb Z_2$.
Formula (A.89) now proves (KM.3b).  This proof uses no pentagon
and does not identify either modular frame sign with this character.

\begin{statement}{A.91}{Proved}{the exact remaining type of comparison}
\citep[Proposition 4.4]{LitNakamuraTsunogaiYasuda2010} \citep[Introduction and Section 2]{LitIhara2000BetaGammaII}
\end{statement}
For multiplier one the intrinsic metabelian action on the
commutator module is multiplication by
\[
 B_f(U,V)=1+(U-1)(V-1)A_f(U,V).
\]
To check the sign, use affine coordinates with linear parts $U,V$:
$\Phi(y)=f^{-1}yf$ changes its translation by $(V-1)f$,
and the translation of $[x,\Phi(y)]$ changes by
$(U-1)(V-1)A_f$ times that of $[x,y]$.
The same plus convention is used in the proof of NTY (2010),
Proposition 4.4, p.443.  At the diagonal order-three component,
\[
 B_f(\omega,\omega)
   =1+ (\omega-1)^2A_f(\omega,\omega)
   =1-3\omega A_f(\omega,\omega).
\]
The five-point audit's affine comparison of this value with the
Gaussian frame concerns this \emph{diagonal} component, whereas
(A.89) concerns $(\omega,1)$.  Equating the resulting bit with
$\rho_{3,2}/2$ is a further assertion, not a consequence of the
definition of $\rho_3$ or of its evenness.

Ihara II, introduction pp.1--2 and Section 2, makes the relevant
source boundary explicit: full profinite pentagon gives
$B(s,t)=\Gamma(s)\Gamma(t)/\Gamma(s+t)$, but Anderson's
multiplication formula is imposed as an additional condition
defining $\mathrm{GT}_A$.  Its logarithmic-derivative version
likewise defines $\mathrm{GT}_{A'}$.  Proposition 2 proves
$\mathrm{GT}_K\subseteq\mathrm{GT}_{A'}$, and Proposition 3
identifies $(A')$ with $(D\log A)$ and yields
$\mathrm{GT}_A\subseteq\mathrm{GT}_{A'}$.
The final Remark explicitly warns that the derivation has a
nonconstant kernel from mixed prime components.  Thus these
opened results do not license replacing a coefficient-two,
order-three unit or its sign by logarithmic Kummer data.

\begin{statement}{A.92}{Proved}{factorization and logarithmic data do not determine the sign}
\end{statement}
This is a countermodel to a proposed deduction between coefficient-ring
identities, not a constructed element of GT.
Let $\Lambda=\mathbb Z_2[[\widehat{\mathbb Z}]]$ and let
$\Lambda_3=\mathbb Z_2[[\mathbb Z_3]]$ be its subalgebra from
the pro-three direct factor of the exponent group.  The invariant
derivation $D(x^a)=a_2x^a$, where $a_2$ is the two-adic component
of $a$, vanishes identically on $\Lambda_3$.

There is an integral idempotent $e\in\Lambda_3$ supported on
characters of exact order three.  Explicitly, in the quotient
group ring $\mathbb Z_2[C_{3^n}]$ with generator $v$, for $n\geq1$
its image is
\[
 e_n=\frac1{3^{n-1}}\sum_{j=0}^{3^{n-1}-1}v^{3j}
       -\frac1{3^n}\sum_{j=0}^{3^n-1}v^j.
\]
These are compatible idempotents: the first summand projects
onto the characters of order dividing three and the second onto
the trivial character.  Put $\Gamma=1-2e$.  Then
\[
 \Gamma^2=1,\quad D\Gamma=0,\quad \Gamma(1)=1,
 \quad \Gamma(\omega)=\Gamma(\omega^{-1})=-1.
\]
Moreover $\Gamma(U^{-1})=\Gamma(U)^{-1}$, since inversion fixes
$e$.  Define in the completed two-variable coefficient ring
\[
 B(U,V)=\Gamma(U)\Gamma(V)\Gamma(UV)^{-1}.
\]
It is a normalized symmetric multiplicative two-cocycle and
satisfies $B(U,U^{-1})=1$.  Nevertheless
\[
 B(\omega,\omega)=-1,
 \qquad D_VB(\omega,1)=0.
\]
The first assertion follows from $\omega^2=\omega^{-1}$;
the second follows from $D\Gamma=0$.  Thus the diagonal unit
sign cannot be recovered from the off-diagonal logarithmic
derivative using only these identities.  To apply a stronger
identity to the actual GT image, that additional identity must
be proved or cited in its full profinite, mixed-component domain.

The usual integrality constraint on the metabelian coefficient
does not remove this particular coefficient-ring example.
The odd-character decomposition of the completed group algebra
is integral over $\mathbb Z_2$.  On any component where the
pro-three character of $U$ or of $V$ is trivial, $B(U,V)=1$.
On every remaining component, both $U-1$ and $V-1$ are units:
their nontrivial roots of odd order remain nontrivial in residue
characteristic two.  Thus an integral $A$ exists with
$B=1+(U-1)(V-1)A$, taking $A=0$ on the former components.
It has $A(\omega,1)=0$ and
$A(\omega,\omega)=2/(3\omega)$.
No profinite free-group lift satisfying the pentagon is claimed.

\subsection{Descent discrepancy and a sharp solvable-source boundary}


\begin{statement}{A.93}{Proved}{the cubic pentagon character measures a descent discrepancy}
\end{statement}
On $\mathcal H$ let $\psi(\Phi)=a_f-\rho\pmod3$, with
$a_f$ as in A.82. In the prompt's marking the following are
identities of continuous homomorphisms on $\mathcal H$:
\[
 \rho\circ D=-a_f\pmod3,
 \qquad
 \boxed{\ \psi=-\rho\circ D-\rho\pmod3.\ }
\]
Consequently $\rho\circ D=-\rho\pmod3$ is a necessary
pentagon equation, and A.82 gives an actual gentle pair violating
it. These identities concern $\mathcal H$; they do not assert
a full-multiplier extension of $\psi$.

\emph{Proof.}
Use the actual notation of A.66. Reducing
$H=\varphi(f(z,x))$ modulo the pure subgroup gives
$f(r^{-1},s)=r^{-a_f}$ in $S_3$. Here simultaneous conjugation
shows that the defining exponent $a_f$ is independent of the
chosen generating pair $(p,t)$ of orders three and two, and
replacing $p$ by $p^{-1}$ changes the displayed power in the
specified way. Thus $l=-a_f\pmod3$.
A.67 proves $a=l\pmod3$, whereas A.66 gives
$a=q-p$, $p+q=0\pmod3$ and $\rho_D=-q$.
It follows that $a=-q=-a_f\pmod3$, hence
$\rho_D=-a_f$. Subtracting $\rho$ gives the boxed identity.
The genuine five-point affine pentagon calculation makes
$\psi=0$ necessary for the pentagon, and A.82 proves that this
necessary equation is not a consequence of the gentle relations.

\begin{statement}{A.94}{Proved}{the character survives beyond the third derived action}
\end{statement}
The element $k$ of A.80 can be chosen so that, in addition,
\[
       f_k,\beta_k\in R_{\rm nil}(F'')\subseteq F''',
       \qquad A_k\in F'',
       \qquad D(k)=1,\qquad\Delta(k)=1.
\]
Here $F''$ is the closed second derived subgroup of the full
free profinite group, and $F'''=(F'')'$ is its closed third
derived subgroup. The genuine normalized action of $k$ is
identity on $F/F'''$, hence on $F/F''$, on the completed Alexander module
$F'/F''$, and on $K/K'$ for every open subgroup $K$ containing
$F''$. In particular the canonical homology of finite abelian
covers, including the Fermat covers, does not contain a line
affording $(-1)^\Delta$ over a field of characteristic different
from two. The same is true of the completed Alexander module
after extension to such coefficients, and of the homology of
any finite metabelian cover preserved by the gentle action.
This assertion concerns these canonical modules; it does not
assert the analogous statement for arbitrary nonabelian covers
such as the full level-twelve cover.
It is also identity on $\widehat B/R_{\rm nil}(F'')$, hence
on $\widehat B/\widehat B'''$. In particular it is invisible
in every finite modular-source quotient in which the image
of $F''$ is nilpotent.
Consequently no action on any finite solvable quotient of
$\widehat B$ of derived length at most three can detect
$\Delta$, even if all such actions are retained jointly.
Together with the independently proved finite characteristic
$\widehat B$ detector of derived length four in A.95, this
makes four the exact minimum derived length for such a
finite modular-source action detector. This is a statement
about quotients of $\widehat B$, not about the derived length
of arbitrary finite representations of the gentle group.

\emph{Proof.}
Replace the seed of A.79 by the word
\[
                    A_1=[[x^4,y^4],[y^4,z^4]].
\]
Both inner commutators belong to $F'$, so $A_1\in F''$.
Under every $\widehat B/F$ conjugate of $\varphi$, one of
$x^4,y^4,z^4$ is one. If it is $x^4$ or $z^4$, the respective
inner commutator is one; if it is $y^4$, both are one.
Consequently $A_1\in\operatorname{core}_{\widehat B}\ker\varphi$.
Direct multiplication in the original signed level-twelve
representation gives $q_{12}(A_1)=5I$. The complete six
orbifold values and the original value are checked in
\texttt{check\_iterated\_lifting.py}.

Let $J\triangleleft\widehat B$ be open normal, and let
$M\triangleleft\widehat B$ be open normal with $F''\subseteq M$.
Set
\[
 N_{J,M}=U_0\cap M\cap
       \operatorname{core}_{\widehat B}(\varphi^{-1}J).
\]
The seed $A_1$ belongs to every $N_{J,M}$. Hence its level-twelve
image is a two-group containing $[5I]$, exactly as in A.80.
The same coprime lift at $r$, followed by R2 and cusp
normalization, gives an actual $k_{J,M}\in\mathcal H$ with
\[
 \begin{gathered}
 f_{J,M},\beta_{J,M}\in R_{\rm nil}(N_{J,M}),
 \qquad A_{J,M}\in N_{J,M},\\
 q_{12}(A_{J,M})=[5I],\quad q_{12}(\beta_{J,M})=1,
 \quad\rho_{J,M}=0,\quad\varphi(f_{J,M})\in J.
 \end{gathered}
\]
The additional statement about $A_{J,M}$ follows from its
prescribed base coordinate one in $\widehat B/N_{J,M}$.
The cusp correction lies in $R_{\rm nil}(N_{J,M})$ and
therefore preserves membership of all three entries in
$N_{J,M}$. It also preserves their retained finite values.
Thus $f_{J,M},\beta_{J,M},A_{J,M}\in M$.

In fact $R_{\rm nil}(N_{J,M})\subseteq R_{\rm nil}(M)$:
every finite nilpotent quotient of $M$ restricts to one of
$N_{J,M}$. Impose the closed conditions of A.80 together with
$f,\beta\in R_{\rm nil}(M)$ and $A\in M$, for every such $M$.
A finite collection is
met by taking the intersections of its $J$'s and $M$'s.
The continuity of the normalized parameter and pure frames
and compactness of $\mathcal H$ give a simultaneous limit.
For completeness the two exact closed-subgroup identities are
\[
 \bigcap_{\substack{M\triangleleft_o\widehat B\\F''\subseteq M}}M=F'',
 \qquad
 \bigcap_{\substack{M\triangleleft_o\widehat B\\F''\subseteq M}}
          R_{\rm nil}(M)=R_{\rm nil}(F'').
\]
The first is the separation of the profinite quotient
$\widehat B/F''$ by its finite quotients. For the second,
monotonicity gives containment from right to left. Conversely
let $q:\widehat B\to Q$ be any finite quotient killing
$R_{\rm nil}(F'')$. The finite group $q(F'')$ is nilpotent,
since $F''/R_{\rm nil}(F'')$ is pronilpotent.
Take $M=F''\ker q$, which is open normal and contains $F''$.
Then $q(M)=q(F'')$ is nilpotent, so $q(R_{\rm nil}(M))=1$.
All such $q$ separate the quotient by $R_{\rm nil}(F'')$,
proving the reverse containment. Normality here is legitimate:
$F''$ is characteristic in the normal subgroup $F$, and its
nilpotent residual is characteristic in $F''$.

The limit therefore has $f,\beta\in R_{\rm nil}(F'')\subseteq F'''$
and $A\in F''$, and retains
$\varphi(f)=1$, the two frame values, and $\rho=0$.
A.66 and R4 again give $D(k)=1$ and $\Delta(k)=1$.
It also retains the absolute pronilpotent invisibility of A.80.

The normalized genuine automorphism fixes $x$ and sends
$y$ to $f_k^{-1}yf_k$. Since $f_k\in F'''$, it induces identity
on $F/F'''$. Its restriction to $F'/F''$ is therefore identity.
If $K$ is open and contains $F''$, then $F'''\subseteq K'$.
Moreover $k$ preserves $K$ since it is identity on $F/F'''$.
For every $w\in K$ the element $k(w)w^{-1}$ lies in
$F'''\subseteq K'$, so the action on $K/K'$ is identity.
These are actual actions from the fixed normalized pair, so
no inner-automorphism ambiguity enters the assertion.
The normalized modular extension has
$k(s)=sf_k$, $k(r)=f_k^{-1}r$ and $k(\tau)=\tau$.
Since $f_k\in R_{\rm nil}(F'')$, it is identity on
$\widehat B/R_{\rm nil}(F'')$. Any finite quotient whose
image of $F''$ is nilpotent kills this residual and thus sees
the identity action. The inclusion
$R_{\rm nil}(F'')\subseteq F'''\subseteq\widehat B'''$
gives the asserted identity on $\widehat B/\widehat B'''$.
Every finite solvable quotient
of $\widehat B$ of derived length at most three kills
$\widehat B'''$; hence $k$ preserves each such quotient map
and acts identically on its target. Its nonzero $\Delta$
therefore rules out every detector from that joint collection.
A.95 supplies the matching upper bound by an actual finite
characteristic quotient with a faithful enough outer-action
detector; no statement about the full pentagon is used here.
When $K$ contains $F'$, it is preserved by the full gentle
group as well: its image in $F^{\rm ab}$ is an open additive subgroup,
preserved by multiplication by every unit
$\lambda\in\widehat{\mathbb Z}^{\times}$.
For general metabelian covers we impose this stability when
regarding their individual homology as a module for the full
group; characteristic covers in particular satisfy it.

Over a field of characteristic different from two, a line
affording $(-1)^\Delta$ would make this same element $k$
act by $-1$. Since its action on each indicated module is
identity, no such line exists. More generally, for any nonzero
one of these coefficient modules $V$, its twist
$V\otimes(-1)^\Delta$ is not isomorphic to $V$ as a
$\widehat{\mathrm{GT}}_0$ module: an intertwiner would satisfy
$T=-T$ by applying $k$ and hence be zero. Their restrictions
to $G_{\mathbb Q}$ are canonically identical because R4
gives $\Delta|_{G_{\mathbb Q}}=0$.

Finally let $V_*$ denote the actual joint image of
$\widehat{\mathrm{GT}}_0$ in its absolute pro-prime pure
actions, its normalized full action on $F/F'''$, and the
canonical homologies of all finite abelian covers and all
finite metabelian covers preserved by the full group; include
also the canonical outer action on
$\widehat B/R_{\rm nil}(F'')$ (and hence on
$\widehat B/\widehat B'''$).
The pure and homology actions are the continuous genuine
actions from normalized pairs, while the last action is
explicitly an outer action. Cover stability has been
specified above. The retained closed condition
$f_k\in R_{\rm nil}(U_0)$ gives the first invisibility,
and the preceding proof gives all the other ones. Thus
\[
 \widehat{\mathrm{GT}}_0\longrightarrow V_*\times\mathbb F_2,
 \qquad g\longmapsto(\operatorname{image}(g),\Delta(g))
\]
is onto: $g$ and $gk$ have identical first coordinate and
opposite second coordinates. In particular $\Delta$ is not
a function of this entire joint collection of action and
homology data. This conclusion does not concern arbitrary
solvable quotients.

\subsection{A sharp derived-length bound and a simultaneous pentagon-test obstruction}


\begin{statement}{A.95}{Proved}{a finite solvable detector of derived length at most four}
\end{statement}
The normalized finite detector of (A.55), and the outer detector
of (A.59), can be chosen as actions on a finite characteristic
quotient of $\widehat B$ whose derived length is at most four.
On $\mathcal H$ this detects $\delta$ directly. For all of
$\widehat{GT}_0$, use exactly the finite multiplier residue and
Galois normalization table in (A.56); neither an explicit value
of its modulus $M$ nor a smaller source quotient is asserted here.

Indeed, the finite group
$Q=\operatorname{SL}_2(\mathbb Z/12)/\{\pm I\}$ has derived
series of orders
\[
                     576,\quad96,\quad8,\quad1.
\]
For comparison the local series have orders
$48,12,4,1$ at four and $24,8,2,1$ at three. CRT gives the
first displayed series: the simultaneous central element
$(-I,-I)$ is not in the first derived product, because
$-I\notin\operatorname{SL}_2(\mathbb Z/4)'$.
These are independently enumerated in
\texttt{check\_solvable\_detector.py} and its output.

The group $E=(M_2\oplus M_3)\rtimes Q$ of (A.54) is an
extension of a group of derived length three by an abelian group.
Its fourth derived group is therefore trivial. The characteristic
quotient $Q^*$ of (A.55) embeds in the finite product of the
images of all homomorphisms $\widehat B\to E$. Subgroups and
finite products preserve the bound on derived length, so
$(Q^*)^{(4)}=1$.

For the full multiplier version of the outer detector, take the
joint characteristic quotient with $\widehat B/F_M$ as in
(A.59). That added factor is an extension of $S_3$ by the
abelian group $F/F_M\simeq(\mathbb Z/M)^2$, hence has derived
length at most three. The joint source quotient still has derived
length at most four. The conclusion concerns the finite
\emph{source group} on which the automorphisms act; it does not
assert the same bound for its automorphism group, nor assert that
every finite automorphism or pentagon assignment lifts globally.

\begin{statement}{A.96}{Proved}{a simultaneous obstruction to low-derived and two-adic Burau pentagon certificates}
\end{statement}
The lifting audit's strengthened A.94 supplies an \emph{actual}
$k=(1,f_k)\in\mathcal H$ satisfying
\[
 \Delta(k)=1,\qquad
 f_k,\beta_k\in R_{\rm nil}(F'')\subset F''',\qquad A_k\in F''.
 \tag{FP.36.1}
\]
Here the derived groups and the nilpotent residual are closed.
The strengthened residual assertion follows from its compactness
construction as follows. In the profinite quotient
$\widehat B/R_{\rm nil}(F'')$, let $q$ be any finite quotient
and put $M=F''\ker(q)$. This is an open normal subgroup of
$\widehat B$, and $q(M)=q(F'')$ is nilpotent. Consequently
$q$ kills $R_{\rm nil}(M)$. Taking all such $q$ proves
\[
 \bigcap_{M\mathrel{\triangleleft_o}\widehat B,\ F''\subset M}
             R_{\rm nil}(M)=R_{\rm nil}(F'').
\]
The reverse inclusion follows by restriction to $F''$ of each
finite nilpotent quotient of $M$. Finally,
$R_{\rm nil}(F'')\subset(F'')'=F'''$ because finite abelian
quotients are nilpotent.

\emph{Simultaneous pure-derived-length-three consequence.}
In every finite quotient of $G_5$ whose pure image has derived
length at most three, all five evaluations
$f_k(X_i,X_j)$ in the pentagon are one. Hence its pentagon
word is one in every such quotient, simultaneously, although
$\Delta(k)=1$. More generally, every standard Drinfeld
generator formula at every number of marked points has the form
of a half twist conjugated by an evaluation of $f_k$ at two pure
elements (with multiplier one). It evaluates to that same half
twist modulo $P_n'''$. Thus these formulas give the literal
identity automorphism on all $G_n/P_n'''$ simultaneously.
This does not assert an action on $G_n$ itself, or membership of
$k$ in $\widehat{GT}$. It proves that no certificate confined to
these quotients can force the CM sign equality. The witness is
the already lifted profinite element (FP.36.1), not a collection
of unrelated finite solutions.

\emph{A genuinely five-point quotient, with its central relation checked.}
Let $R_n=(\mathbb Z/2^n)[\omega]/(\omega^2+\omega+1)$ and set
\[
 A=\begin{pmatrix}1-\omega&\omega&0\\1&0&0\\0&0&1\end{pmatrix},
 \quad B=\begin{pmatrix}1&0&0\\0&1-\omega&\omega\\0&1&0\end{pmatrix},
 \quad C=\begin{pmatrix}1&0&-1\\0&1&-1\\0&0&-\omega\end{pmatrix}.
 \tag{FP.36.2}
\]
These reduced Burau matrices satisfy the three braid and distant
commutation relations. Before projectivization they satisfy
$(ABC)^4=\omega I$. Therefore their projective classes give an
actual continuous quotient of the label-five stabilizer
$L=B_4/\langle z_4\rangle$.
Inducing this homomorphism from the actual index-five subgroup
$L\subset G_5$ gives a finite quotient of $G_5$ in a wreath
product. In particular no unverified direct representation of
$B_5$ through the spherical relator is being used.

For $n=2$, the projective image has order $55\,296$; its
modulo-two image has order $216$ and its reduction kernel is
$(\mathbb F_2)^8$. This is genuinely finer than the earlier
affine quotient. Nevertheless it still cannot force $\Delta=0$.
In pentagon order, the five evaluation pair images have derived
series
\[
\begin{array}{c|c|c}
 \text{pair}&\text{modulus }4&\text{modulus }8\\ \hline
 (X_3,X_4),(X_5,X_1)&144,12,4,1&1152,96,32,2,1\\
 (X_2,X_3),(X_1,X_2)&192,64,4,1&1536,512,8,1\\
 (X_4,X_5)&48,12,4,1&384,96,32,2,1 .
\end{array}
 \tag{FP.36.3}
\]
Thus the modulus-four equation is killed already by
$f_k\in F'''$. The modulus-eight groups do reach derived length
four, but (FP.36.1) still kills every evaluation: the second
derived group of each pair image is a finite two-group, hence
nilpotent.

In fact this obstruction holds at \emph{every} $2^n$ in the
same family. At modulus two the pair images have series
$18,3,1$, $24,8,2,1$, and $6,3,1$, respectively.
Their second derived groups are trivial or $C_2$. The kernel
of reduction from $\operatorname{PGL}_3(R_n)$ to
$\operatorname{PGL}_3(\mathbb F_4)$ is a two-group: its
successive congruence quotients are additive matrix groups in
characteristic two, modulo their scalar subspaces. Hence the
second derived group of every lifted pair image is a two-group.
The image of $F''$ under that evaluation lies in this nilpotent
group, so $R_{\rm nil}(F'')$ maps to one. Therefore
\[
           \eta_{2^n}\bigl(\mathcal P(f_k)\bigr)=1
           \quad\text{for every }n\geq1,
           \qquad \Delta(k)=1.
 \tag{FP.36.4}
\]
The same conclusion holds after induction to $G_5$ in every
block. Choose the five coset representatives to be the powers
of the transitive rotation used in (A.74). Its conjugation
cyclically permutes the five $X_i$, so the pairs in each block
are again the five cyclically adjacent pairs listed in
(FP.36.3); each corresponding $f_k$ value is one.
The program \texttt{check\_burau\_pair\_filtration.py} and
the three outputs \texttt{check\_burau\_pair\_mod2.out},
\texttt{check\_burau\_pair\_mod4.out}, and
\texttt{check\_burau\_pair\_mod8.out} verify the literal
matrices, central word, and all finite series used here.
This failure leaves open quotients whose evaluation pair has a
non-nilpotent second derived image; it makes no assertion that
the full pentagon is satisfied by $k$.

\subsection{A canonical congruence-cover homology detector}


\begin{statement}{A.97}{Proved}{The level-twelve relation lattice and its first two cohomology groups}
Let $q=q_{12}:\widehat B\twoheadrightarrow Q$, where
$Q=\operatorname{SL}_2(\mathbb Z/12)/\{\pm I\}$ has order $576$,
and put $N=\ker q$ and $M=N^{\rm ab}$.  We use the genuine
normalized gentle action, with $\Phi(\tau)=\tau^{\lambda_\Phi}$.
It preserves $N$ by the integral covariance of A.71.  Thus it
gives a continuous linear representation
\[
 h_{12}:\widehat{GT}_0\longrightarrow
          \operatorname{Aut}_{\widehat{\mathbb Z}}(M),
 \qquad M\simeq\widehat{\mathbb Z}^{97}.
\]
For the deck action of $Q$ on this lattice,
\[
 H^1(Q,M)=0,\qquad H^2(Q,M)\simeq\mathbb Z/96.
\]
Here cohomology is continuous finite-group cohomology with the
displayed profinite coefficient module. The deck representation
on $M/576M$ is faithful.
\end{statement}
\begin{proof}
Write $H_2=\langle q(s)\rangle$ and $H_3=\langle q(r)\rangle$.
The kernel $N$ meets every conjugate of either finite factor
trivially.  The ordinary Bass--Serre tree of $C_2*C_3$, divided
by the discrete congruence kernel, is therefore a finite graph
with edge set $Q$ and vertex sets $Q/H_2$ and $Q/H_3$.
Its fundamental group is free of rank
$576-288-192+1=97$.  Completion identifies its homology lattice
with $M$, and gives the exact sequence of $Q$-modules
\[
 0\longrightarrow M\longrightarrow C_1\longrightarrow C_0
       \longrightarrow\widehat{\mathbb Z}\longrightarrow0,
 \quad
 C_1=\widehat{\mathbb Z}[Q],\quad
 C_0=\widehat{\mathbb Z}[Q/H_2]\oplus
                \widehat{\mathbb Z}[Q/H_3].
 \tag{RH.1.1}
\]
All terms and maps are already defined over $\mathbb Z$; their
profinite completions are exact because the underlying graph
homology groups are free.  Equivalently one may take the inverse
limit of their exact finite-coefficient graph complexes.

Put $L=\ker(C_0\to\widehat{\mathbb Z})$.  Invariants of $C_0$
are pairs of coefficients of the two orbit sums. Their
augmentation is $(a,b)\mapsto288a+192b$. Consequently
\[
 L^Q=\widehat{\mathbb Z}(2,-3).
\]
The norm edge sum maps to exactly $(2,-3)$, up to the harmless
choice of orientation of all edges.  Shapiro's lemma gives
$H^j(Q,C_1)=0$ for $j>0$.  The long exact sequence of
$0\to M\to C_1\to L\to0$ therefore gives $H^1(Q,M)=0$.
Also $H^1(Q,C_0)=0$: by Shapiro it is the sum of the groups
$\operatorname{Hom}(C_2,\widehat{\mathbb Z})$ and
$\operatorname{Hom}(C_3,\widehat{\mathbb Z})$, both zero.
The second short exact sequence in (RH.1.1) now gives
\[
 H^2(Q,M)\simeq H^1(Q,L)
  \simeq\operatorname{coker}
   \bigl((288,192):\widehat{\mathbb Z}^{,2}
                     \to\widehat{\mathbb Z}\bigr)
  \simeq\mathbb Z/96.
\]

For $g\ne1$ in $Q$, the ordinary integral character of $M$ is
\[
 \operatorname{tr}(g\mid M)
       =1-\#(Q/H_2)^g-\#(Q/H_3)^g.
 \tag{RH.1.2}
\]
A nonidentity element cannot fix a coset of both types, because
its order would divide both two and three. Hence
$96\leq97-\operatorname{tr}(g\mid M)\leq384<576$.
If its matrix were the identity modulo $576$, this strictly
positive integer would be divisible by $576$, a contradiction.
This proves the faithful finite deck action.

The program \texttt{check\_congruence\_homology.py} constructs
the actual graph, a spanning tree, its $97$ integral cycle
generators and the deck matrices of $s,r$.  It checks their
orders, the character formula for all $576$ deck transformations,
and the primitive cusp class used in A.99.  The cohomology
conclusions follow from the exact sequences just proved; they
are not inferred from a numerical rank computation.
\end{proof}

\begin{statement}{A.98}{Proved}{The quadratic character is detected by homology modulo 576 at multiplier one}
On $\mathcal H$ the character $\Delta$ factors through the
finite linear action
\[
 h_{12,576}:\mathcal H\longrightarrow
        \operatorname{GL}_{97}(\mathbb Z/576).
 \tag{RH.2.1}
\]
This is an explicit canonical homology level. It is not claimed
to be the smallest detecting quotient or modulus.  In
particular, the quotient of $\widehat B$ by $N'N^{576}$ has
order $576^{98}$ and suffices for this action detector; the
larger all-homomorphisms characteristic core of A.55 is not
required here.
\end{statement}
\begin{proof}
Suppose $\Phi\in\mathcal H$ is the identity on $M/576M$.
Equivariance of its action on $M$ with its induced action
$\Phi_Q$ on $Q$ gives, modulo $576$,
$T_{\Phi_Q(g)}=T_g$ for every $g\in Q$.
Faithfulness in A.97 implies $\Phi_Q=1$.
For a positive integer $n$, put
\[
 \Gamma_n=\widehat B/(N'N^n),\qquad
 1\longrightarrow M/nM\longrightarrow\Gamma_n
       \longrightarrow Q\longrightarrow1.
\]
The induced automorphism of $\Gamma_{576}$ is the identity
both on its kernel and on its quotient. For a lift $b$ of
$g\in Q$, the difference $\Phi(b)b^{-1}\in M/576M$
is therefore independent of the lift and is a one-cocycle.
Its class belongs to $H^1(Q,M/576M)$.

The exact sequence $0\to M\xrightarrow{n}M\to M/nM\to0$
and A.97 identify
\[
 H^1(Q,M/nM)\simeq H^2(Q,M)[n].
\]
Under these identifications reduction from modulus $576$ to
modulus $6$ is multiplication by $576/6=96$ on $H^2(Q,M)$.
This follows by the diagram of the two exact coefficient
sequences, whose left vertical map is multiplication by $96$
and middle vertical map is the identity.  A.97 shows that this
reduction of the cocycle class is zero. Thus the induced
automorphism of $\Gamma_6$ is conjugation by an element of
$M/6M$.

We now use the single finite representation from A.54, rather
than its all-homomorphisms core.  Its target is the actual image
$E_0$ of $\widehat B$ in
\[
 \bigl(\mathbb F_2[Q/H_2]\oplus
             \mathbb F_3[Q/H_3]\bigr)\rtimes Q.
\]
Its restriction to $N$ has abelian image of exponent dividing
six. Hence this representation factors through $\Gamma_6$.
Let $u$ denote the image in $E_0$ of the conjugating element;
its projection to $Q$ is one.  From
$\Phi(s)=\beta s\beta^{-1}$ and $\Phi(r)=ArA^{-1}$,
we find that $\beta^{-1}u$ centralizes the image of $s$ and
$A^{-1}u$ centralizes the image of $r$.
The two proven centralizer-projection statements of A.54 give
$q(\beta)\in H_2$ and $q(A)\in H_3$.
Both frames are pure, and neither $H_2$ nor $H_3$ meets $q(F)$
nontrivially. Consequently
\[
 q(\beta)=q(A)=1.
\]
The pure frame abelianization, or the explicit R3 formula,
then gives $\rho\equiv0\pmod3$. The signed normalization
is the one equal to $I$ modulo four, so the R4 shear is also
$I$ modulo twelve. We obtain $T=W=I$, and therefore
$\Delta(\Phi)=0$.
This proves containment of the kernel of (RH.2.1) in
$\ker\Delta$, which is exactly the asserted factorization.
The extension $\Gamma_{576}$ has kernel
$(\mathbb Z/576)^{97}$ and quotient of order $576$, proving
the displayed order.  Its invariance under the actual gentle
action is sufficient; no claim of characteristicity under
all automorphisms of $\widehat B$ is used.
\end{proof}

\begin{statement}{A.99}{Proved}{Full-multiplier homology factorization and its exact scope}
The full character $\Delta:\widehat{GT}_0\to\mathbb F_2$
factors continuously through the genuine representation
$h_{12}$ of A.97, and hence through $h_{12,n}$ for some positive
integer $n$.  The value $n=576$ has been proved sufficient on
$\mathcal H$ only.  This construction provides a sign line in
a finite permutation module obtained from this homology action;
it does not assert that the sign line occurs in $M$ itself.
\end{statement}
\begin{proof}
The element $\tau^{12}$ lies in $N$. Its homology class is a
primitive nonzero vector: in the finite graph of A.97 its
cusp cycle is an embedded cycle with coefficients $1$ and
$-1$. An edge outside any spanning tree on that cycle has
coefficient $1$ or $-1$ in its fundamental-cycle coordinates.
The program also checks this assertion in the displayed basis.
The normalized action sends this class to its multiple by
$\lambda_\Phi$. Thus $h_{12}(\Phi)=1$ forces
$\lambda_\Phi=1$; primitivity in the torsion-free profinite
lattice is essential here. A.98 now gives $\Delta(\Phi)=0$.

The continuous map from the compact group $\widehat{GT}_0$
onto its closed image $K\subseteq\operatorname{GL}_{97}
(\widehat{\mathbb Z})$ is a quotient map. The kernel
containment just proved therefore descends $\Delta$ to a
continuous character of $K$. Its kernel is open, so it
contains $K\cap\ker(\operatorname{GL}_{97}
(\widehat{\mathbb Z})\to\operatorname{GL}_{97}
(\mathbb Z/n))$ for some $n$.
This paragraph gives a topological existence argument.
A.110 below supplies the explicit full-multiplier bound
$n_*=2^{51942}3^{14136}$ by genuine tangential specialization
and a cyclotomic conductor estimate. Neither argument computes
the smallest modulus or a full arithmetic normalization table.

For completeness the promised module is explicit once such an
$n$ is fixed. Let $K_n$ be the actual image on $M/nM$ and let
$\mathcal O_n$ be the orbit of a chosen ordered basis of
$(\mathbb Z/n)^{97}$.  The action of $K_n$ on this orbit is
free, so the permutation module $k[\mathcal O_n]$ over any
field of characteristic different from two contains the
nonzero vector
\[
 v_\Delta=\sum_{g\in K_n}(-1)^{\Delta(g)}
                   e_{g\cdot\mathbf b}.
\]
Every group element acts on its span by $(-1)^\Delta$.
At multiplier one one may take $n=576$ throughout.
This module consists of functions on finite ordered homology
bases of the congruence cover; its construction uses the
specified tower action.  Determining the occurrence of this
character in the raw homology lattice, rather than in this
permutation functor of its finite reduction, remains a distinct
question.  A.94 proves its absence from the canonical Alexander
module of $F$ and from the homologies of its finite abelian covers.
\end{proof}

\subsection{The exact mixed duplication and cubic-frame defects}


\begin{statement}{A.100}{Proved}{conventions and the peripheral identity}
\citep[Theorem 1]{LitIhara2000BetaGammaII} \citep[p.~100]{LitFurusho2003GTA}
\end{statement}
In this module $\Phi=(1,f)$ acts by
$x\mapsto x$, $y\mapsto f^{-1}yf$, and
$[x,y]=xyx^{-1}y^{-1}$. We write
\[
 f\equiv A_f(U,V)[x,y]\pmod{F''},\qquad
 B_f(U,V)=1+(U-1)(V-1)A_f(U,V).
\]
The module action is left conjugation. We denote by $\Psi_n(\Phi)$
the coefficient $m_0$ in A.88, without assigning it an arithmetic
sign by the symbol $\rho_n$. The conventions needed below are
\[
\begin{array}{c|l}
\text{symbol}&\text{definition}\\\hline
\Psi_n&m_0\text{ in }\ker(F\to C_n)^{\rm ab}\text{ with }y_0=y,\\
\rho&\beta_\Phi^{\rm ab}=(-\rho,-\rho),\\
\kappa_d&\alpha_d(\Phi)^{\rm ab}=\kappa_d e_5\quad(\Phi\in\mathcal H_5),\\
k_n&\sigma(n^{1/m})/n^{1/m}=\zeta_m^{k_n(\sigma)}.
\end{array}
\]
Here $\mathcal H_5=\widehat{GT}\cap\mathcal H$.
Ihara II, Theorem 1, writes the arithmetic identity with $n^{-r}$,
so his $\tau_n^{(0)}$ is $-k_n$ in that monodromy convention.
Furusho (2003), p.100, explicitly records the corresponding minus
sign for the NT characters. The older NS (2000) formulation uses
its frame-defined $\rho_2$ and relates it to its own Alexander
coefficient in Proposition 5.3, on its specified subgroup.
These source symbols are not identified across markings here.

The simultaneous substitution $x\mapsto x^{-1}$,
$y\mapsto y^{-1}$ sends $y_j=x^jyx^{-j}$ to $y_{-j}^{-1}$,
up to the harmless conjugation by a power of $x^n$ in the Schreier
basis. It sends $\Psi_n$ to $-\Psi_n$; the half modulo two of an
even $\Psi_n$ is therefore unchanged.

We also give the direct justification for A.90's peripheral identity.
Put $g=f(z,x)$ and $h=f(y,z)$, where $z=(xy)^{-1}$. The
multiplier-one three-cycle is $ghf=1$. Right composition with the
actual gentle element $c=(-1,1)$ has normalized pair $(-1,f)$.
Its three-cycle, with $m=-1$, is
\[
 gz^{-1}hy^{-1}fx^{-1}=1.
\]
Substitution of $h=g^{-1}f^{-1}$ gives
$gz^{-1}g^{-1}=xf^{-1}yf=\Phi(z^{-1})$. Consequently
\[
 \Phi(z)=f(z,x)z f(z,x)^{-1}.
\]
This deduction uses the group of actual gentle automorphisms and
the two three-cycle equations; it uses no pentagon or finite lift.

\begin{statement}{A.101}{Proved}{the off-diagonal beta value is the half-Kummer sign}
\end{statement}
For every $\Phi\in\mathcal H$, in
$R_4=(\mathbb Z/4)[\omega]$, $\omega^2+\omega+1=0$, one has
\[
 \boxed{\quad B_f(-1,\omega)=(-1)^{\Psi_3(\Phi)/2}\quad}.
 \tag{A.101}
\]
The right side means the two-adic half reduced modulo two.
Its integrality was proved in A.90, not assumed from arithmetic.
Indeed, the two-cycle gives $A_f(U,V)=A_f(V,U)$.
Reduction modulo two identifies the continuous evaluation maps
$U\mapsto-1$ and $U\mapsto1$, even for the full completed
profinite group ring. Therefore
\[
\begin{split}
 B_f(-1,\omega)
 &=1-2(\omega-1)A_f(-1,\omega)\\
 &=1-2(\omega-1)A_f(1,\omega)\pmod4\\
 &=1-2t_f,
 \qquad t_f=(\omega-1)A_f(\omega,1).
\end{split}
\]
A.90 gives $t_f\in\mathbb Z_2$ and $\Psi_{3,2}/2=t_f/3$.
Since $3\equiv1\pmod2$, the last displayed unit is precisely
the sign in (A.101). No gamma factorization is needed.

\begin{statement}{A.102}{Proved}{the order-three part of duplication is already pentagonal}
\end{statement}
Use the affine realization of $S_3$ over $\mathbb F_3$ with
$p=(1,1)$ and $t=(0,-1)$. Let $a_f$ be defined by
$f(p,t)=p^{a_f}$. Then, for every multiplier-one gentle element,
\[
 a_f=\Psi_2(\Phi)\pmod3.
 \tag{MD.3a}
\]
To prove this directly, evaluate the index-two Schreier basis at
$x=t$, $y=p$. The two conjugates of $y$ have translations $1,-1$,
and their coefficients are $m_0,-m_0$. Thus the translation of
$f(t,p)$ is $2\Psi_2$. The two-cycle also gives
$f(t,p)=p^{-a_f}$. Hence $2\Psi_2=-a_f$ in $\mathbb F_3$,
which is (MD.3a).

For $\Phi\in\mathcal H_5$, the explicit genuine five-point
pentagon calculation A.77 gives $a_f=\rho\pmod3$. Consequently
\[
 \rho=\Psi_2\pmod3\qquad(\Phi\in\mathcal H_5).
 \tag{MD.3b}
\]
This is only a congruence. It does not identify the entire
frame-defined integer $\rho$ with $\Psi_2$.

\begin{statement}{A.103}{Proved}{Proved reduction; the vanishing assertions remain open}
\end{statement}
For $\Phi\in\mathcal H_5$, define the mixed duplication unit
\[
 D_2(\Phi)=
 \omega^{2\Psi_2(\Phi)}
       \frac{B_f(\omega,\omega)}{B_f(-1,\omega)}\in R_4^\times.
 \tag{MD.4a}
\]
Combining A.101, A.102 and the intrinsic phase calculation A.87 gives
\[
 \boxed{\quad
 D_2(\Phi)=(-1)^{d_2(\Phi)},\qquad
 d_2(\Phi)=\epsilon_2(W_\Phi)+\Psi_3(\Phi)/2\pmod2.
 \quad}
 \tag{MD.4b}
\]
In fact A.87 says
$B_f(\omega,\omega)=\omega^\rho(-1)^{\epsilon_2(W)}$,
and $2\Psi_2+\rho\equiv0\pmod3$ by A.102.
Both summands in $d_2$ are continuous additive characters, so
$d_2:\mathcal H_5\to\mathbb F_2$ is an actual continuous
character, independently of any unknown distribution theorem.
The same expression defines a character on all $\mathcal H$.
More precisely, write
$\eta_{\mathrm{aff}}=a_f-\rho=\Psi_2-\rho\pmod3$.
The formula for $D_2$ continues to make sense there and gives
\[
 D_2=\omega^{-\eta_{\mathrm{aff}}}(-1)^{d_2}
 \qquad(\Phi\in\mathcal H).
 \tag{MD.4b$'$}
\]
Indeed $2\Psi_2+\rho=-(\Psi_2-\rho)\pmod3$.
Thus A.77 is exactly the vanishing of the order-three component
of this specialized duplication unit; its order-two component is
the remaining character $d_2$. In particular $D_2^3=(-1)^{d_2}$
on all $\mathcal H$.

Define a second continuous character on $\mathcal H_5$ by
\[
 e_3(\Phi)=\kappa_d(\Phi)+\Psi_3(\Phi)/2\pmod2.
 \tag{MD.4c}
\]
The five-point frame identity A.76 then yields the exact decomposition
\[
 \boxed{\quad\Delta|_{\mathcal H_5}=d_2+e_3.\quad}
 \tag{MD.4d}
\]
In particular the main parity equality on $\mathcal H_5$ is
equivalent to $d_2=e_3$ there. Proving both zero would suffice,
but their simultaneous vanishing is stronger than that equivalence.
The equality $d_2=0$ is exactly the single finite specialization
\[
 B_f(\omega,\omega)
       =\omega^{-2\Psi_2(\Phi)}B_f(-1,\omega)\pmod4.
 \tag{MD.4e}
\]
The equality $e_3=0$ is exactly the identification of the
five-point cubic-frame parity with the index-three half-character.
Neither assertion has been deduced here from the full pentagon.
On the arithmetic subgroup $G_{\mathbb Q}\cap\mathcal H$,
Anderson duplication gives $d_2=0$ in the same intrinsic
convention. The supplied arithmetic vanishing of $\Delta$
then gives $e_3=0$ there as well. This uses no identification
of an unnormalized source Kummer exponent with the frame coordinate.

\begin{statement}{A.104}{Proved}{completed-ring, reflection and cyclotomic audit of A.92}
\end{statement}
We strengthen the coefficient-level scope of the example, while
keeping it separate from any profinite GT lifting claim.
Write $\widehat{\mathbb Z}=\mathbb Z_2\times H_{\rm odd}$,
where $H_{\rm odd}=\prod_{p\ne2}\mathbb Z_p$.
At a finite odd cyclic quotient, its group algebra over
$\mathbb Z_2$ is finite etale. After the corresponding finite
unramified coefficient extension it splits into its individual
character components. Passing to inverse limits identifies
the full group algebra componentwise with the two-adic
power-series rings in the pro-two exponent variable, with
compatible odd-character labels. An inverse on a character block
is integral whenever its reduction modulo two is a unit.

The idempotents $e_n$ of A.92 are integral because their
denominators are powers of three. Reduction
$C_{3^{n+1}}\to C_{3^n}$ sends $e_{n+1}$ to $e_n$:
each coset sum acquires the factor three canceled by its
denominator. They therefore define $e$ in
$\mathbb Z_2[[\mathbb Z_3]]$ and, through the direct-factor
inclusion, in $\mathbb Z_2[[\widehat{\mathbb Z}]]$.
On odd-character components $e$ has value one precisely when
the three-primary part of the character has exact order three.
It is zero otherwise, independently of the remaining exponent factors.

For $\Gamma=1-2e$, put
$B(U,V)=\Gamma(U)\Gamma(V)/\Gamma(UV)$.
If the three-primary character of $U$ is trivial then
$\Gamma(U)=1$ and $\Gamma(UV)=\Gamma(V)$, so $B=1$.
The same holds with $U,V$ exchanged. On every remaining
character block both $U-1$ and $V-1$ are integral units:
their nontrivial odd parts have nontrivial reduction modulo two.
Thus there is an integral continuous $A$ with
\[
 B=1+(U-1)(V-1)A,
\]
zero on the former blocks and the indicated quotient on the latter.
Its componentwise definitions are compatible under finite-quotient
transition maps, proving the claimed completed-ring integrality.
They imply $A(U,1)=A(1,V)=0$, and in particular all the
off-diagonal Kummer data of this example vanish.

Every $\lambda\in\widehat{\mathbb Z}^{\times}$ preserves the
exact order of each three-primary character. Therefore
\[
 \Gamma(U^\lambda)=\Gamma(U),\qquad
 B(U^\lambda,V^\lambda)=B(U,V).
\]
The coefficients are in $\mathbb Z_2$ and are unchanged by
unramified coefficient automorphisms. In addition, multiplication
of the exponent by two is an automorphism of the pro-three
factor, so $\Gamma(U^2)=\Gamma(U)$ and
$B(U^2,V^2)=B(U,V)$. Thus, on its odd-root evaluations, the
coefficient-two Frobenius substitution also fixes this example;
the values $\pm1\in\mathbb Z_2$ are fixed by Witt Frobenius.
This is Frobenius substitution, not a multiplication formula
relating shifted arguments. Inversion fixes $e$;
since $\Gamma^2=1$, the reflection identity
$\Gamma(U^{-1})=\Gamma(U)^{-1}$ holds exactly.
Accordingly ordinary cyclotomic substitution covariance and
reflection do not exclude the example. This is a statement about
these specified coefficient identities, not all possible
compatibilities of a GT action on a tower.

The example also obeys both coefficient forms of the
multiplier-one metabelian three-cycle. Indeed $A$ is symmetric,
and, for $W=(UV)^{-1}$, reflection gives
$B(U,V)=B(V,W)=B(W,U)$. Clearing the regular factors
$U-1,V-1,W-1$ gives
\[
\begin{split}
 V^{-1}A(W,U)+(UV)^{-1}A(V,W)+A(U,V)&=0,\\
 V^{-1}A(W,U)+A(V,W)+U A(U,V)&=0.
\end{split}
\]
For example, the first cleared numerator is
$(V-1)/V+(U-1)/(UV)+(W-1)=0$.
The second is $(V-1)/V+(U-1)+U(W-1)=0$.
The relevant factors are regular: on each odd-character block
they are units or nonzero power-series parameters over an
unramified discrete valuation ring.

Nevertheless
\[
 B(\omega,\omega)=-1,\qquad B(-1,\omega)=1,
 \qquad A(\omega,1)=0.
\]
Thus the mixed duplication specialization MD.4e fails when
$\Psi_2=\Psi_3=0$. The exact finite-ring program
\texttt{check\_mixed\_idempotent.py} independently verifies
idempotency, compatibility, integral divisibility, reflection,
all unit substitutions, and the displayed metabelian identities
in $(\mathbb Z/4)[C_{12}]$, $(\mathbb Z/4)[C_{36}]$ and
$(\mathbb Z/4)[C_{108}]$ and their two-variable rings.
It retains the order-four exponent factor through the CRT embedding.
Its output is PASS in each case. The inverse-limit proof above,
not these finitely many checks, proves the completed-ring statement.
No lift of this coefficient datum to a free profinite word satisfying
the full five-point pentagon is asserted.

\begin{statement}{A.105}{Known}{the opened duplication theorem}
\citep[Theorems 3.1 and 3.4]{LitFurusho2003GTA} \citep[Theorems 2--3 and Appendix]{LitLochakNakamuraSchneps2004Eigenloci}
\end{statement}
Furusho (2003), Theorem 3.1, proves that the relations
$(I),(II),(IV),(K_2)$ imply Anderson duplication $(A_2^0)$;
Theorem 3.4 gives $GT_K\cap I\Gamma\subseteq GT_{A,2^\infty}$.
The explicit group-ring form printed on p.102 is
\[
 B_\sigma\{X,X\}
   =X^{-2\Psi_2^{(0)}(\sigma)}B_\sigma\{-1,X\}.
\]
Here $\Psi_2^{(0)}$ is defined on p.100 by exactly the
index-two Schreier coefficient. In that same coefficient/action
convention, the specialization $X=\omega$ modulo four is
MD.4e. Thus this theorem proves $d_2=0$ for elements satisfying
its extra equations, once those equations are expressed in the
same marking. It supplies neither these extra hypotheses on
all full profinite GT nor the independent comparison $e_3=0$.
An indexed author correction dated 3 January 2008 retains the
same Theorem 3.1 statement; its full PDF could not be retrieved
in this pass, so no detailed claim from the correction is used.

The focused search checked the concrete multiplication and
special-locus leads, including Ihara II, Furusho's author version,
the LNS eigenlocus paper and its conventions appendix, and the
earlier opened NTY comparison. No opened theorem in that search
implies either missing parity comparison on all full profinite GT.
This is a bounded source finding; it is not a theorem of
nonexistence, a proof of independence from the full pentagon,
or a novelty claim.

\begin{statement}{A.106}{Proved}{two precise equation loci kill the duplication sign}
\end{statement}
Let $C_L$ be the fixed subgroup for the original physical quotient
of A.60--A.65, rather than the literal current-marking subgroup $C_P$.
Then
\[
 C_L\cap\mathcal H\ \subseteq\ \ker d_2.
 \tag{MD.7a}
\]
Also the multiplier-one locus satisfying the harmonic primitive
equation $(\mathrm{GF}_1)$ lies in $\ker d_2$:
\[
 \{\Phi\in\mathcal H:g=f(x,s)a^{4\rho}\}\subseteq\ker d_2,
 \qquad a=\tau^{-1},\quad g=\beta^{-1}.
 \tag{MD.7b}
\]
Statement A.65 identifies this second locus with its precisely marked
NT harmonic equation $(\mathrm I')$.

We prove both inclusions in the projective reduction modulo two
of the affine representation used in A.87. Put
$c_0=(1-\omega)^{-1}\in\mathbb F_4$. Its relevant values are
\[
 x=(0,\omega),\qquad a=(0,\omega^2),\qquad
 s=(c_0,1),\qquad V=(0,\omega^\rho).
\]
For any $f\in F'$, its evaluation at
$(0,\omega),(c_0,1)$ is the translation $c_0t_f$.
Its evaluation at $(0,\omega^2),(c_0,1)$ is the translation
$c_0\overline{t_f}$. By A.90, $\overline{t_f}=t_f$ and
$t_f\bmod2=\Psi_3/2\bmod2$.

For $\Phi\in C_L$, equation (M7) in A.65 gives
$g=f(a,s)x^{-b_0}$. Since the rightmost factor has zero
translation, the translation of $g$ is $c_0t_f$.
Under $(\mathrm{GF}_1)$ the same conclusion follows from
$g=f(x,s)a^{4\rho}$. Finally $W=gV$, and $V$ also has
zero translation. The proof of A.87 identifies
$\operatorname{trans}(W)/c_0$ with $\epsilon_2(W)$.
Therefore both hypotheses give
\[
 \epsilon_2(W)=t_f=\Psi_3/2\pmod2,
\]
which proves (MD.7a)--(MD.7b).

Neither argument uses the pentagon, an equality between the
existential GF0 parameter and $\Psi_3$, nor a full Anderson
multiplication law. In particular this proves only the harmonic
half-character comparison on $C_L$. On
$C_L\cap\mathcal H_5$ the remaining main obstruction is exactly
$e_3=\kappa_d+\Psi_3/2$; its vanishing is still open here.

\begin{statement}{A.107}{Proved}{the half-character is a descended harmonic sign}
\end{statement}
The half-character has a universal interpretation, stronger than
the fixed-locus statement A.106:
\[
 \boxed{\quad
 \Psi_3(\Phi)/2\bmod2
 =\epsilon_2(W_{D_L\Phi})
 =\epsilon_2(W_{D_P\Phi})
 \qquad(\Phi\in\mathcal H).
 \quad}
 \tag{MD.8a}
\]
Consequently, as characters on $\mathcal H$,
\[
 d_2=\epsilon_2\circ W+\epsilon_2\circ W\circ D_L
     =\epsilon_2\circ W+\epsilon_2\circ W\circ D_P.
 \tag{MD.8b}
\]
In particular both $C_L\cap\mathcal H$ and
$C_P\cap\mathcal H$ lie in $\ker d_2$. The two fixed groups
themselves remain distinct; this statement does not identify them.

For the proof, keep the same affine quotient and $c_0$ as in A.106.
The established pure normalizations of the two descents are
\[
 \beta_{D_L}=x^{b_0}P^{-1},\quad P=f(a,s),\qquad
 \beta_{D_P}=x^{a_0}C,\quad C=f(s,\tau).
\]
The image of $x$ has zero translation, while
\[
 \operatorname{trans}(P)=c_0t_f,\qquad
 \operatorname{trans}(C)=-c_0t_f.
\]
The first identity was proved in A.106. For the second, the
modulo-two images are $s=(c_0,1)$ and $\tau=(0,\omega)$;
the two-cycle turns the evaluation into the inverse of the
one defining $c_0t_f$.
Taking inverses of the displayed pure normalizations, the
rightmost powers of $x$ still have zero translation. It follows
that both $\beta_{D_L}^{-1}$ and $\beta_{D_P}^{-1}$ have
translation $c_0t_f$. The respective normalizing shears
$V_{D_L\Phi}$ and $V_{D_P\Phi}$ also have zero translation.
Thus both corrected harmonic frames have translation $c_0t_f$,
and A.87 gives (MD.8a). No comparison of the integer gauges
$a_0,b_0$ or of the individual old/new Galois labels is required.

\begin{statement}{A.108}{Proved}{the square-pair comparison is not a word identity}
\end{statement}
The following exact finite lemma isolates an additional issue in
the modulo-four Burau pentagon comparison. Let
\[
 p=(123),\qquad t=(1234)\quad\text{in }S_4.
\]
Then $(p,t)$ has peripheral orders $(3,4,4)$ and generates $S_4$,
whereas $(p,t^2)$ has peripheral orders $(3,2,3)$ and generates
$A_4$. The joint same-word evaluation
\[
 F\longrightarrow S_4\times A_4,
 \qquad x\longmapsto(p,p),\quad y\longmapsto(t,t^2)
 \tag{MD.9a}
\]
is onto. Consequently its restrictions have images
\[
 F'\longmapsto A_4\times V_4,
 \qquad F''\longmapsto V_4\times\{1\}.
 \tag{MD.9b}
\]

For the proof, both projections in (MD.9a) are onto. The nontrivial
quotients of $S_4$ are $S_4,S_3,C_2$, and those of $A_4$ are
$A_4,C_3$. The groups have no common nontrivial quotient.
Goursat's lemma therefore gives the full direct product.
Taking its first and second derived subgroups proves (MD.9b).
The discrete word
\[
 w=[[x,y],x[x,y]x^{-1}]
\]
provides a concrete check: $w(p,t)=(14)(23)$ but
$w(p,t^2)=1$. The independent exact permutation program
\texttt{check\_square\_pair\_goursat.py} verifies the pair orders,
the joint derived-series orders $288,48,4,1$, and this witness.

The square pair is relevant to the preceding finite Burau
calculation for a precise reason. After removing its central
order-three factor, that calculation gives a group generated by
$p,t$ with peripheral orders $(3,4,4)$. Its square-pair subgroup
is $A_4\times C_2$; if $e$ denotes its central involution, then
$p,et^2$ generate an $A_4$ complement with orders $(3,2,3)$.
For every $f\in F'$, the central factor cancels and
$f(p,t^2)=f(p,et^2)$. This identifies the square-pair value with
the half-Kummer $A_4$ evaluation, after the stated marked
isomorphism. The central decomposition and that marked
isomorphism are finite-Burau inputs, not consequences of
Goursat's lemma.

Thus an equality comparing the relevant coordinate of $f(p,t)$
with the half-Kummer value of $f(p,t^2)$ is a substantive
distribution statement; it cannot be inferred merely from
using a single commutator word in both places. The lemma does
not assert that $w$ satisfies the gentle equations, nor that
either independent finite assignment lifts to a genuine GT
element. It therefore does not prove independence from the
gentle equations or the full profinite pentagon.

\begin{statement}{A.109}{Proved}{duplication alone cannot settle the gentle cubic sign}
\end{statement}
Extend the second defect to the actual gentle group by
\[
 \widetilde e_3(\Phi)=\epsilon_3(T_\Phi)+\Psi_3(\Phi)/2\pmod2.
\]
Its restriction to $\mathcal H_5$ is $e_3$ by A.76.
The actual gentle element $k$ supplied by A.94 satisfies
\[
 B_k=1,\qquad \Psi_n(k)=0\ (n\geq1),\qquad
 \rho(k)=0,\qquad d_2(k)=0,\qquad
 \widetilde e_3(k)=\Delta(k)=1.
 \tag{A.109}
\]
Indeed A.94 gives $f_k,\beta_k\in R_{\rm nil}(F'')\subseteq F'''$
and $\Delta(k)=1$. In particular $f_k$ has zero image in
$F'/F''$, proving $A_{f_k}=0$ and $B_k=1$.
For every $n$, the Schreier subgroup $N_n$ contains $F'$, so
$F''\subseteq N_n'$ and every coefficient defining $\Psi_n(k)$
vanishes. The abelianization of $\beta_k$ is zero, giving
$\rho(k)=0$. Formula A.87 now gives $\epsilon_2(W_k)=0$,
and the other assertions follow.

Consequently this element satisfies the beta-factorization,
reflection and multiplication coefficient equations with the
choice $\Gamma_k=1$ and all Kummer exponents zero, while retaining
the nontrivial cubic defect. This does not put $k$ in $GT_A$,
whose definition already requires membership in full profinite GT.
The full pentagon for $k$ remains unproved. The conclusion is
instead the precise logical limitation that a beta/Kummer
comparison, including the complete vanishing of $d_2$, still
needs a cubic-frame comparison to establish the desired parity.

\subsection{An explicit homology modulus and the absence of a raw rational sign constituent}


\begin{statement}{A.110}{Proved}{a full-multiplier finite homology detector}
\end{statement}
Let $N=\ker q_{12}$, let $M=N^{\rm ab}\simeq
\widehat{\mathbb Z}^{97}$, and use the genuine normalized
homology action $h_{12}$ of A.97. Put
\[
                  \boxed{\ n_*=2^{51942}3^{14136}.\ }
\]
The full character $\Delta:\widehat{\mathrm{GT}}_0\to\mathbb F_2$
factors through the finite representation
\[
 h_{12,n_*}:\widehat{\mathrm{GT}}_0\longrightarrow
             \operatorname{GL}_{97}(\mathbb Z/n_*).
\]
Thus the ordered-homology-basis permutation module of A.99
can be taken at this explicit modulus for all multipliers.
The underlying gentle-stable modular-source quotient
\[
 \Gamma_{n_*}=\widehat B/(N'N^{n_*})
\]
has order
\[
       576n_*^{97}=2^{5038380}3^{1371194}.
\]
These are sufficient bounds. No optimality of this modulus
or quotient order is asserted.

\emph{Proof.}
We first prove the required genuine arithmetic ramification
statement. The reduction $Q/q_{12}(F)\simeq S_3$ implies
$N\subseteq F$. Consequently
\[
                 F/(N'N^{576})
\]
is an actual finite quotient of $F$, of order dividing
$|\Gamma_{576}|=576^{98}$; in particular its only prime
divisors are two and three. The standard tangential Galois
action on $F$ is the literal restriction of the genuine
action used here. Indeed R4, E.65.8 gives
$u\mapsto u^\lambda$ and
$v\mapsto f(u^2,v^2)^{-1}v^\lambda f(u^2,v^2)$; its
restriction to $x=u^2,y=v^2$ is exactly
$x\mapsto x^\lambda$, $y\mapsto f^{-1}y^\lambda f$.
There is no change of outer representative in this comparison.

For every rational prime $p\geq5$, use the smooth pair
$(\mathbb P^1_{\mathbb Z_p},\{0,1,\infty\})$ and the
unit tangential base point $\overrightarrow{01}$.
Petrov's genuine tangential specialization isomorphism
identifies its geometric prime-to-$p$ fundamental group with
that of the special fibre, compatibly with Galois actions,
the latter factoring through the residue Galois group
\cite[Lemma~8.7]{LiftPetrov2023}.
Thus inertia at $p$ acts identically on every finite
$\{2,3\}$ quotient of this tangential $F$. Apply this to
the displayed quotient of $F$. Its subgroup
$N/(N'N^{576})=M/576M$ is therefore fixed pointwise.
It follows that the genuine finite Galois representation
$h_{12,576}|_{G_{\mathbb Q}}$ is unramified outside two
and three. This argument requires neither a multiplier-one
assumption nor any comparison of outer actions on the
modular orbifold.

Let
\[
 K=\ker(h_{12,576}|_{G_{\mathbb Q}}),\qquad
 L=\overline{\mathbb Q}^{\,K},\qquad
 \Lambda=\lambda(K)\subseteq\widehat{\mathbb Z}^{\times},
 \qquad D=|\operatorname{GL}_{97}(\mathbb Z/576)|.
\]
The field $L/\mathbb Q$ is finite Galois, unramified outside
two and three, and $[L:\mathbb Q]$ divides $D$.
Moreover
\[
 [\widehat{\mathbb Z}^{\times}:\Lambda]
       =[L\cap\mathbb Q^{\rm cyc}:\mathbb Q]
       \quad\hbox{divides }D.
\]
The equality is the restriction Galois correspondence; the
divisibility also follows directly because this finite
quotient of the multiplier group is a quotient of
$G_{\mathbb Q}/K$.

For clarity we give the conductor bound in group form.
In the cyclotomic unit product
$\widehat{\mathbb Z}^{\times}=\prod_p\mathbb Z_p^{\times}$,
inertia at a prime $p$ has image the full $p$-factor and
is trivial on the other factors. Since $L$ is unramified
at every $p\geq5$, the closed subgroup $\Lambda$ contains
all these factors, and therefore their product by closedness.
This is also the usual cyclotomic
ramification description of its associated abelian field
\cite[Chapter~I, Theorem~4.16 and its proof]{LiftMilneCFT2020}.
The images of $1+4\mathbb Z_2$ and $1+3\mathbb Z_3$ in
the remaining finite quotient have orders dividing
$2^{v_2(D)}$ and $3^{v_3(D)}$, respectively. The elementary
principal-unit power identities give
\[
 \begin{split}
 (1+4\mathbb Z_2)^{2^{v_2(D)}}
       &=1+2^{v_2(D)+2}\mathbb Z_2\subseteq\Lambda,\\
 (1+3\mathbb Z_3)^{3^{v_3(D)}}
       &=1+3^{v_3(D)+1}\mathbb Z_3\subseteq\Lambda.
 \end{split}
\]
Consequently the principal congruence kernel of modulus
$2^{v_2(D)+2}3^{v_3(D)+1}$ is contained in $\Lambda$.

The standard exact order formula
\[
 |\operatorname{GL}_d(\mathbb Z/p^e)|
   =p^{(e-1)d^2+d(d-1)/2}\prod_{i=1}^d(p^i-1)
\]
and $576=2^6 3^2$ yield
\[
 \begin{split}
 v_2(D)&=5\cdot97^2+\frac{97\cdot96}{2}
               +\sum_{i=1}^{97}v_2(3^i-1)=51940,\\
 v_3(D)&=97^2+\frac{97\cdot96}{2}
               +\sum_{i=1}^{97}v_3(2^i-1)=14135.
 \end{split}
\]
The two cross sums are $239$ and $70$. Exact direct order
valuations and independent principal-power valuations are
checked by \texttt{check\_homology\_bound.py}.
Thus the displayed principal congruence kernel is that of
$n_*$, and $576$ divides $n_*$.

Now suppose $g\in\widehat{\mathrm{GT}}_0$ satisfies
$h_{12,n_*}(g)=1$. The primitive cusp class of $\tau^{12}$
in A.99 is multiplied by $\lambda(g)$. Its primitivity
therefore gives $\lambda(g)\equiv1\pmod{n_*}$.
By the preceding inclusion $\lambda(g)\in\Lambda$.
Choose an actual $\sigma\in K$ with
$\lambda(\sigma)=\lambda(g)$, as an equality in the full
group $\widehat{\mathbb Z}^{\times}$.
Then $g\sigma^{-1}\in\mathcal H$ and
$h_{12,576}(g\sigma^{-1})=1$.
A.98 gives $\Delta(g\sigma^{-1})=0$, and R4 gives
$\Delta(\sigma)=0$. Hence $\Delta(g)=0$.
This proves the required finite kernel containment for
every multiplier. The order of $\Gamma_{n_*}$ follows
from its kernel $(\mathbb Z/n_*)^{97}$ and quotient $Q$ of
order $576$.

\begin{statement}{A.111}{Proved}{the raw homology has no characteristic-zero $\Delta$ constituent}
\end{statement}
For a prime $\ell$, put
$M_\ell=M\otimes_{\widehat{\mathbb Z}}\mathbb Z_\ell$
and $M_{\mathbb Q_\ell}=M_\ell\otimes_{\mathbb Z_\ell}\mathbb Q_\ell$.
Then
\[
                 M_{\mathbb Q_\ell}^{G_{\mathbb Q}}=0.
\]
More strongly, no subquotient of this finite-dimensional
$G_{\mathbb Q}$ representation is the trivial representation.
Consequently the raw homology $M_{\mathbb Q_\ell}$ has no
$\widehat{\mathrm{GT}}_0$ subquotient affording the character
$(-1)^\Delta$, and in particular no such line. There is also
no free rank-one $\mathbb Z_\ell$ submodule of $M_\ell$
with this action. These statements concern the full gentle
group and its full Galois subgroup; they do not assert the
corresponding statement after restriction to $\mathcal H$.

\emph{Proof.}
Let $X=\mathbb P^1_{\mathbb Q}\setminus\{0,1,\infty\}$,
with its standard rational tangential base point. The genuine
action comparison in A.110 identifies its arithmetic
fundamental group with $F\rtimes G_{\mathbb Q}$, using
exactly the action under consideration. Since $N\subseteq F$
is stable under this action, the open subgroup
\[
                  N\rtimes G_{\mathbb Q}
                  \ \subseteq\ F\rtimes G_{\mathbb Q}
\]
corresponds to a finite \emph{geometrically connected} etale
cover $Y\to X$ over $\mathbb Q$, of degree $[F:N]=96$.
Indeed it projects onto $G_{\mathbb Q}$, and its geometric
intersection is exactly $N$. This is the usual finite-cover
correspondence, applied to the specified genuine section
\cite[\S3]{LiftMilneLEC2013}. Its geometric first homology
is $M_{\mathbb Q_\ell}$ with precisely the given Galois
action. Changing a geometric base point only conjugates
by $N$, which has no effect on its abelianization.

Let $\overline Y$ be the smooth projective compactification,
let $D=\overline Y\setminus Y$, and let
$J=\operatorname{Jac}(\overline Y)$. The Gysin sequence
for this smooth curve and its finite boundary, dualized,
gives the natural Galois-equivariant exact sequence
\[
 0\longrightarrow
 \mathbb Q_\ell(1)\otimes
 \left(\mathbb Q_\ell[D(\overline{\mathbb Q})]
       /\mathbb Q_\ell\!\cdot\!\sum_{d\in D(\overline{\mathbb Q})}d\right)
 \longrightarrow M_{\mathbb Q_\ell}
 \longrightarrow V_\ell J\longrightarrow0.
\]
Here $V_\ell J=T_\ell J\otimes\mathbb Q_\ell$; the compact
curve comparison follows from its Jacobian's universal
abelian covering property and the Tate-module description
of first homology
\cite[Corollary~16.2]{LiftMilneLEC2013}
\cite[Proposition~9.1 and Corollary~9.6]{LiftMilneJac2021}
\cite[Theorem~15.1 and Remark~15.5]{LiftMilneAV2022}.

Choose a rational prime $p\ne\ell$ of good reduction for
$\overline Y$ at which its finite boundary is unramified.
Such primes exist because these data spread out over
$\mathbb Z[1/S]$ for a finite set $S$. Let $\mathrm{Fr}_p$
be an arithmetic Frobenius lift. On the boundary term its
eigenvalues are $p$ times roots of unity: the finite point set
is permuted, and the Tate twist is multiplication by $p$.
Thus this term has no eigenvalue one.

For the compact term, good reduction identifies the
$\ell$-power torsion with that of the special-fibre Jacobian
$J_p/\mathbb F_p$, compatibly with arithmetic Frobenius.
One can see this directly from the finite etale group schemes
$J[\ell^m]$ over $\mathbb Z_p$
\cite[Proposition~20.7]{LiftMilneAV2022}.
The endomorphism $\operatorname{Frob}_p-1$ of $J_p$ is an
isogeny: its geometric kernel is the finite set
$J_p(\mathbb F_p)$, and its differential is $-1$, so this
kernel is a finite etale group scheme. An endomorphism of an
abelian variety with finite kernel is surjective and finite.
It therefore becomes invertible on $V_\ell J_p$.
Equivalently, arithmetic Frobenius has no eigenvalue one
on $V_\ell J$.

The characteristic polynomial on the middle term is the
product of those on the two end terms. Hence
$\det(\mathrm{Fr}_p-1\mid M_{\mathbb Q_\ell})\ne0$.
Every invariant subquotient has a characteristic polynomial
dividing this polynomial, so no trivial Galois subquotient
can occur. R4 says $\Delta|_{G_{\mathbb Q}}=0$; a full
gentle-group $\Delta$ constituent would therefore restrict
to a trivial Galois constituent, a contradiction. A free
$\mathbb Z_\ell$ line would yield such a line after tensoring
with $\mathbb Q_\ell$.

This does not exclude a constituent with finite coefficients:
the nonzero determinant above need not be a unit modulo a
chosen integer. It is also consistent with A.110, whose
sign constituent lies in the permutation module of finite
ordered homology bases, rather than in the raw homology.

\subsection{A quadratic pentagon character and the next relative homology obstruction}


\begin{statement}{A.112}{Proved}{an exact quadratic necessary pentagon character}
\end{statement}
There is a continuous homomorphism
$\nu_4:\mathcal H\to\mathbb F_2$, defined by one literal
$S_4$ evaluation, such that
\[
 \boxed{\quad \epsilon_2(W_\Phi)=\nu_4(\Phi)
                 \qquad(\Phi\in\widehat{GT}\cap\mathcal H).\quad}
 \tag{FP.37.1}
\]
Thus $\epsilon_2(W)+\nu_4$ is a finite continuous character
of the actual gentle group which vanishes on its genuine GT
subgroup. Its nonvanishing on the gentle group is not asserted.
It is not identified here with $\Delta$, with the cubic-frame
character, or with the half-Kummer character.

\emph{Definition and homomorphism law.}
In $S_4$ put $p=(123)$ and $t=(1234)$, with products acting
right to left. The orders of $p,t,pt$ are $3,4,4$.
For $\Phi=(1,f)\in\mathcal H$, write $R=f(p,t)\in A_4$.
The actual peripheral automorphism of the free group preserves
the conjugacy class of $xy$, hence $pR^{-1}tR$ has order four.
Writing $A_4=V_4\rtimes\langle p\rangle$, the four choices
of $v\in V_4$ give order four for $pv^{-1}tv$ exactly when
$v\in\{1,t^2\}$; the other two choices give order two.
Consequently there are unique $\nu_4\in\mathbb F_2$ and
$a\in\mathbb F_3$ with
\[
                 f(p,t)=t^{2\nu_4}p^a.
 \tag{FP.37.2}
\]
The quotient $S_4/V_4=S_3$ identifies $a$ with the actual
$a_f$ of (A.77). Moreover the induced action on this $S_4$
quotient is exactly $\operatorname{Inn}(p^{-a_f})$:
it fixes $p$ and sends $t$ to $p^{-a_f}tp^{a_f}$.
In particular its kernel is preserved by every actual gentle
element; no invariance of a Burau kernel has been assumed.
For two such elements the GT multiplication rule gives
\[
 R_{\Phi\Psi}
   =R_\Phi p^{-a_\Phi}R_\Psi p^{a_\Phi}
   =t^{2(\nu_4(\Phi)+\nu_4(\Psi))}
                  p^{a_\Phi+a_\Psi}.
\]
This proves the homomorphism assertion and continuity follows
from the finite evaluation.

\emph{Reduction of the actual pentagon.}
Use the modulus-four matrices in (FP.36.2), with
$X_1=A^2,X_2=B^2,X_3=C^2,X_4=(ABA)^2,X_5=(CBC)^2$.
Let $D$ be the image of the reversing half twist
$abcaba$. Then $D^2=1$, and its conjugation interchanges
$X_1,X_3$ and $X_4,X_5$, while fixing $X_2$.
Put $F=f(X_1,X_2)$ and $R=f(X_3,X_4)$. The two-cycle
therefore gives the exact matrix pentagon
\[
 \mathcal P(f)=
 R(DR^{-1}D)(DF^{-1}D)\,f(X_4,X_5)\,F.
 \tag{FP.37.3}
\]
Inside $\langle X_3,X_4\rangle$ one can strip a central
element $h$ of order three from $X_3$, obtaining
$p_0=X_3h^{-1}$ and $t_0=X_4$ which generate a group of
order $48$. Its centre is generated by an involution $e_0$;
its quotient by that centre is the marked $S_4$ of
(FP.37.2). Its derived subgroup is an $A_4$ mapping
isomorphically onto the quotient's derived subgroup, so the
commutator evaluation is exactly
\[
 R=(e_0t_0^2)^{\nu_4}p_0^{a_f}.
 \tag{FP.37.4}
\]
Here $h$ disappears from every evaluation of $f\in F'$.

For completeness all conjugacies needed for the three-cycle
reduction are literal finite matrices. Set
$z=(X_4X_5)^{-1}$. There is an invertible matrix $C_0$ with
$C_0p_0C_0^{-1}=z$ and $C_0t_0C_0^{-1}=X_4$.
Furthermore $z^{-1}X_4z$ differs from $X_5$ by the central
involution of $\langle X_4,X_5\rangle$. Put
$H=C_0RC_0^{-1}=f(z,X_4)$. The multiplier-one three-cycle,
and the fact that central factors disappear from a commutator
word, give
\[
 f(X_5,z)=z^{-1}H^{-1}z,
 \qquad f(X_4,X_5)=z^{-1}HzH^{-1}.
 \tag{FP.37.5}
\]
Equations (FP.37.3)--(FP.37.5) reduce the test to the actual
two parameters $F$ and $R$ of one and the same $f$.

Here is an exact small certificate for the remaining matrix
identity. Encode $a+b\omega\in R_2$ by the integer $a+4b$
with $0\leq a,b<4$, and list matrices in row-major order.
The constants can be chosen as
\[
\begin{split}
 h&=(8,1,0;15,14,0;0,0,9),\\
 e_0&=(1,10,0;8,3,0;0,0,11),\\
 C_0&=(1,4,9;10,4,0;1,13,0).
\end{split}
\]
The group $H_{12}=\langle X_1,X_2\rangle$ has order $192$.
Its positive level-three map sends $X_1$ to $U^2$ and $X_2$
to $L_0^2$. Enumerate the images $\beta\in H_{12}$ of the
\emph{actual} harmonic frame, with
$q_3(\beta)^{-1}U^{8\rho}\in\langle S\rangle$.
The equation $F=(ABA)^{-1}\beta(ABA)\beta^{-1}$ gives
exactly four possible matrices $F$ for each $\rho\bmod3$,
two for each $\epsilon=\epsilon_2(W)$. They are
\[
                   F=F^0_{\rho,\epsilon}Z^c,
             \qquad Z=(ABA)^4,\quad c\in\mathbb F_2,
\]
where the six base matrices are
\[
\begin{array}{c|c|l}
 \rho&\epsilon&F^0_{\rho,\epsilon}\\ \hline
0&0&(1,0,0;0,1,0;0,0,1)\\
0&1&(1,10,0;8,9,10;0,8,3)\\
1&0&(1,7,0;10,0,14;3,10,15)\\
1&1&(1,13,0;0,0,14;9,0,5)\\
2&0&(8,1,0;4,9,12;7,12,10)\\
2&1&(0,1,0;14,3,4;15,6,0).
\end{array}
\]
For a genuine pentagon, the independently proved necessary
equation (A.77) gives $a_f=\rho$.
Substitution in (FP.37.3)--(FP.37.5)
now yields, for all the remaining $24$ cases, the exact formula
\[
                  \mathcal P(f)=v^{\epsilon+\nu_4}w_\rho^c,
 \tag{FP.37.6}
\]
where $v,w_\rho$ are independent commuting involutions, with
encoded matrices
\[
\begin{split}
 v&=(1,10,8;8,3,8;0,0,9),\\
 w_0&=(1,8,10;10,9,8;10,2,3),\\
 w_1&=(1,2,8;2,9,2;10,2,3),\\
 w_2&=(1,8,2;2,9,10;10,2,1).
\end{split}
\]
Thus $\mathcal P(f)=1$ forces $\epsilon=\nu_4$, proving
(FP.37.1). The program
\path{check_burau_quadratic_character.py} and its output
check every conjugacy, the marked $S_4$ quotient, the 36
coordinate multiplication laws, and all 24 instances of
(FP.37.6). This is a necessary equation for actual GT elements;
none of the enumerated parameter combinations is asserted to
lift to an actual gentle element merely because it passes the
finite equations.

\emph{Exact limitation of the cubic-frame comparison in this quotient.}
The same Burau target cannot identify $\nu_4$ with
$\kappa_d\bmod2$ by the cubic conjugator equation alone.
Its entire pure image $G$ is generated by $X_1,X_2,X_3$, all
of order three: $X_4\in\langle X_1,X_2\rangle$ and
$X_5\in\langle X_2,X_3\rangle$. Hence $G$ has no
nontrivial elementary-two quotient. Goursat's lemma then makes
the joint pure map
\[
            P_5\longrightarrow G\times(P_5^{\rm ab}/2)
                              =G\times\mathbb F_2^5
\]
surjective. Choose an actual source element $h\in P_5$ with
Burau image one and abelianization $e_5$ modulo two.
There is a literal homomorphism from the order-$216$ reduction
of the Burau image to $\operatorname{SL}_2(\mathbb F_3)$
sending $A,C$ to $U$ and $B$ to $L_0$. The certificate checks
all generator multiplications on all $216$ elements and
surjectivity onto its $24$-element target. This is precisely
the signed $Q_L$ representation in (A.76), since the three
generators have the prescribed images. Thus Burau$(h)=1$
implies $Q_L(h)=I$. The signed identity in (A.76) gives
$q_3(q_5(h))=-I$, because $\zeta(h)=1$;
the abelian forgetting map gives $q_4(q_5(h))=I$.
Thus $q_{12}(q_5(h))=[5I]$.
Its image centralizes $d$ in the joint quotient retaining
$G$, $P_5^{\rm ab}/2$, and $q_{12}q_5$: its three coordinates
are respectively $1$, the $d$-fixed vector $e_5$, and a central
scalar. Multiplication by this finite centralizer element
changes the candidate cubic-frame parity by one while preserving
all those finite conjugator equations. This proves a precise
failure of parity extraction from this target. It does not
produce a nontrivial element of the actual profinite centralizer
$C_{P_5}(d)$, which is trivial, and it does not refute an eventual
comparison using additional five-point quotients or a new theorem.

\begin{statement}{A.113}{Proved}{the first relative characteristic-three homology extension is still insufficient}
\end{statement}
There is a concrete further test which adds characteristic-three
homology to the genuinely five-point quotient of order $216$;
it also fails to force $\Delta=0$, for an exact structural reason.
Let $L=B_4/\langle z_4\rangle$ and let
\[
 q:L\twoheadrightarrow Q=\mathbb F_3^2\rtimes
                 \operatorname{SL}_2(\mathbb F_3),\qquad |Q|=216,
\]
send $a,b,c$ respectively to $(0,U),(0,L_0),(e_1,U)$.
Put $N=\ker(q)$. Use the presentation with three generators
and relators
\[
 aba(bab)^{-1},\quad bcb(cbc)^{-1},\quad aca^{-1}c^{-1},
                  \quad(abc)^4.
\]
The Fox derivatives of these four relators and all their
$Q$-translates form an $864\times648$ matrix over $\mathbb F_3$.
The exact rank is $422$. Therefore the quotient of
$\mathbb F_3[Q]^3$ by its row module is a $226$-dimensional
$Q$-module $W$. The literal assignments
\[
       a\longmapsto(e_a,q(a)),\quad
       b\longmapsto(e_b,q(b)),\quad
       c\longmapsto(e_c,q(c))
           \qquad\text{in } W\rtimes Q
 \tag{FP.38.1}
\]
kill every defining relator and give an actual finite quotient
of $L$. They realize the usual relative abelian extension
$L/[N,N]N^3$ as their image. The boundary map to
$\mathbb F_3[Q]$ has rank $215$, so
\[
                    \dim_{\mathbb F_3}H_1(N,\mathbb F_3)=11.
\]
Neither this construction nor the pentagon evaluation assumes
that an arbitrary gentle action preserves $N$.

Here is the decisive calculation on the evaluation subgroups.
For a pair image $P$, write $H=q(P)$ and
$U_P=P\cap W$. Schreier rewriting on the finite Cayley graph
of $H$ computes $U_P$ exactly as the span of
$v_h+h v_g-v_{hg}$, for the two generators $g$ and the chosen
lifts $(v_h,h)$. It gives
\[
\begin{array}{c|c|c|c|c}
 \text{pair}&(|H|,|H'|,|H''|)&\dim U_P&|P|
       &H''\text{ action on }U_P\\ \hline
 (X_1,X_2),(X_2,X_3)&(24,8,2)&4&1944&\text{trivial}\\
 (X_3,X_4),(X_5,X_1)&(18,3,1)&5&4374&\text{trivial}\\
 (X_4,X_5)&(6,3,1)&4&486&\text{trivial}.
\end{array}
 \tag{FP.38.2}
\]
The reversing half twist gives the paired equalities in the
first two rows. The action in the last column is checked on
an exact basis of the entire kernel $U_P$, not just on a
selected commutator subspace.

Each $U_P$ is elementary abelian of characteristic three.
If $H''=1$, then $P''\subset U_P$ is abelian. In the remaining
case, the preimage of $H''=C_2$ is a central extension of a
cyclic group by $U_P$, because the displayed action is trivial;
that preimage is abelian as well. Thus every pair $P$ has
derived length at most three. The actual element $k$ in
(FP.36.1) consequently has $f_k(X_i,X_j)=1$ in all five pair
images. Its whole pentagon is one in this quotient and in the
induced five-point quotient, while $\Delta(k)=1$.

The exact rank, Fox-boundary identities, all affine relators,
Schreier spans and triviality of the entire $C_2$ action are
checked by \texttt{check\_fox\_relative\_module.py} and its
output. The complete quotient map, the $216$ module actions,
and the finite base multiplication are stored in
\texttt{fox\_relative\_module.npz}. This resolves this specific
relative homology test negatively; it does not rule out a
different finite base or a further nonabelian relative extension.

\Needspace{10\baselineskip}
\subsection{A simultaneous actual obstruction to the entire adelic Burau family}


\begin{statement}{A.114}{Proved}{one actual witness passes the entire adelic primitive-cube-root Burau pentagon family}
\end{statement}
Let $\mathcal O=\mathbb Z[\omega]/(\omega^2+\omega+1)$,
and use the three literal Burau matrices $A,B,C$ of FP.36.2.
For every prime $p$ and $n\geq1$, reduce their projective
classes modulo $p^n$. This gives a genuine finite quotient
$\eta_{p^n}$ of $L=B_4/\langle z_4\rangle$, and hence an
induced genuine five-point quotient as in A.96.
There exists one actual $k=(1,f_k)\in\mathcal H$ such that
\[
 \Delta(k)=1,\qquad Dk=1,\qquad\rho(k)=0,\qquad
 f_k,\beta_k\in R_{\rm nil}(F''),\quad A_k\in F'',
\]
retaining all the invisibility assertions of A.94, and
\[
 f_k(X_3,X_4)=f_k(X_5,X_1)=f_k(X_2,X_3)
       =f_k(X_4,X_5)=f_k(X_1,X_2)=1
       \quad\text{under every }\eta_{p^n}.
\]
Thus its pentagon is one in this entire adelic congruence
family, including every induced five-point quotient, although
$\Delta(k)=1$. This does not assert the full profinite
pentagon for $k$.

\emph{The two harmonic pairs.}
Over $\mathcal O$, the matrices $A,B$ fix the primitive
column $v=(1,1,1)^{\mathsf T}$ and the row
$\ell=(1,\omega,\omega^2)$, with $\ell v=0$.
The basis with columns
\[
 v,\qquad(0,1,-\omega^2)^{\mathsf T},\qquad
                    (0,0,\omega)^{\mathsf T}
\]
has determinant $\omega$, a unit. Hence these matrices are
upper triangular in an integral basis, with diagonal entries
$(1,-\omega,1)$. The group generated by $X_1=A^2,X_2=B^2$
therefore has derived subgroup in the upper unitriangular
group. Over $\mathcal O/p^n$ this is a finite $p$-group.
In particular its second derived group is nilpotent.
The reversing half twist identifies the pair $(X_2,X_3)$
with the first pair, up to interchanging its entries.
Consequently $f\in R_{\rm nil}(F'')$ kills both harmonic
pair evaluations at every prime and every level.

\emph{The three remaining pairs and their uniform seed.}
The following are exact identities over $\mathcal O$:
\[
 (X_4-I)^2=(X_5-I)^2=0,
 \qquad X_i\equiv I\pmod{(1-\omega)}\quad(1\leq i\leq5).
\]
For example,
\[
 X_5=\begin{pmatrix}
 1&\omega-1&-2\omega-1\\0&1&0\\0&0&1
 \end{pmatrix}.
\]
Thus $X_4^{p^n}=X_5^{p^n}=1$ modulo $p^n$.
At $p=3$, the ring $\mathcal O/3^n$ is local with residue
field $\mathbb F_3$ and nilpotent maximal ideal generated
by $1-\omega$. The kernel of its projective linear
reduction to the residue field is a finite three-group.
All the $X_i$ lie in this kernel, so every pair image is a
three-group. Its evaluations are again killed by
$R_{\rm nil}(F'')$.

Now take any finite list of levels $(p,n)$ with $p\ne3$,
and choose a positive integer $m$ divisible by two and by
every listed $p^n$, with $3\nmid m$; their least common
multiple has these properties. Put
\[
             a_m=[[x^m,y^m],[y^m,z^m]]\in F''.
\]
For any group triple with product one, this word is one
if any one of the three displayed $m$th powers is one.
The triples for $(X_3,X_4),(X_5,X_1),(X_4,X_5)$ each
contain $X_4$ or $X_5$. Therefore $a_m$ vanishes at these
three pairs and at all six peripheral permutations, at
every listed level.

These six tests give the full $B$-normal core condition,
not merely a selection of substitutions. Indeed conjugation
by $r$ cyclically permutes $(x,y,z)$, while conjugation
by $s$ swaps $x,y$ and sends $z$ to $(yx)^{-1}$.
The latter three ordered pairs are simultaneously conjugate
to $(y,x),(x,z),(z,y)$, respectively. Thus the six ordered
peripheral pairs represent the six outer cosets in $B/F$;
their kernel intersection is exactly the required $B$-core.

The same word belongs to every full orbifold kernel:
each peripheral permutation of $\varphi$ sends one peripheral
to an element of order two, and $m$ is even. Moreover
\[
                       q_{12}(a_m)=[5I].
\]
This follows directly for the four possible classes
$m\equiv2,4,8,10\pmod{12}$ using
$x\mapsto U^2$, $y\mapsto L_0^2$ and $z=(xy)^{-1}$.
The certificate checks all four classes; equivalently the
modulo-four value is $I$ and the modulo-three value is
$-I$ in each case.

\emph{Actual lifting and compactness.}
Let $K$ be the intersection of the $B$-normal cores of the
three nonharmonic evaluation kernels at the finitely many
listed levels. For any open $B$-normal $M\supseteq F''$
and any finite orbifold target with kernel $J$, set
\[
 N=U_0\cap M\cap K\cap
                   \operatorname{core}_B(\varphi^{-1}J).
\]
This is open and $B$-normal, contains $a_m$, and has
two-group $q_{12}$ image containing $[5I]$.
Apply precisely the coprime $r$-centralizer lift and R2
normalization of A.94. It supplies an actual normalized
gentle element with the prescribed frame values and
$f,\beta\in R_{\rm nil}(N)$, $A\in N$.
The cusp correction remains in $R_{\rm nil}(N)$, so
it preserves membership in every added evaluation kernel.

Take the compact intersection over all finite lists of
levels, all such $M$, and all $J$. Finite intersections
are nonempty by the construction just given. The A.94
intersection argument still yields
$f,\beta\in R_{\rm nil}(F'')$, while every nonharmonic
evaluation is one at every level away from three.
The flag argument handles the two harmonic pairs at all
primes, and the residue argument handles every pair at
three. The retained finite frames give $\Delta(k)=1$,
and the full orbifold condition gives $Dk=1$.
The coset representatives used to induce to the full
five-point group may be taken to be powers of the cyclic
rotation of the five $X_i$; hence each block again uses
the same five pairs, all of whose evaluations are one.

For the first odd-field example, an entirely explicit word
which kills even the two harmonic pair maps is
\[
              a_7=\bigl([[x^{28},y^{28}],
                              [y^{28},z^{28}]]\bigr)^7.
\]
It has original value $[5I]$ and all thirty pair/permutation
values one over $\mathbb F_7$. The general proof uses
the uniform unipotent and flag identities, rather than an
enumeration of further finite pentagon assignments.

\subsection{The exact Ferrari kernel and its marked geometric cover}


\begin{statement}{A.115}{Proved}{the involution kernel and its exact derived subgroup}
\end{statement}
Use the literal generators of B.35:
\[
 L=\widehat B_4/\overline{\langle z_4\rangle},\quad
 z_4=(abc)^4,\quad z_3=(ab)^3,\quad
 x=ca^{-1},\quad y=aba c^{-1}b^{-1}a^{-1}.
\]
The fold sends $a,c$ to $a$ and $b$ to $b$. Its free kernel
before quotienting the centres is $K=\widehat F(x,y)$, with
the exact Nielsen actions
\[
 a:(x,y)\mapsto(x,yx),\qquad
 b:(x,y)\mapsto(xy^{-1},y).
 \tag{FC.1a}
\]
These are established directly from the Artin presentation in B.35;
no geometric identification is needed for the group calculation below.
The fold of $z_4$ is $z_3^2$. Since $z_3$ has infinite
profinite order, the closed central subgroup generated by $z_4$
meets $K$ trivially. Thus $K$ embeds in $L$, with quotient
$\widehat B_3/\overline{\langle z_3^2\rangle}$.

Put $d=abc$ and $j=d^2$. The identity $c=xa$ and
$(ab)x(ab)^{-1}=y^{-1}$ give $d=y^{-1}(aba)$.
Using (FC.1a) once more gives
\[
       j=x^{-1}y^{-1}z_3,\qquad
       j^2=1,\qquad jxj=x^{-1},\quad jyj=y^{-1}.
 \tag{FC.1b}
\]
Consequently the kernel $C$ of the further quotient
$L\to\widehat B_3/\overline{\langle z_3\rangle}$ is
$K\rtimes\langle j\rangle$. The three elements
\[
             t_0=j,\qquad t_1=xj,\qquad t_2=yj
\]
are involutions. The inverse substitutions
$x=t_1t_0$, $y=t_2t_0$, $j=t_0$ identify the presentations,
so
\[
              C=C_2\amalg C_2\amalg C_2.
 \tag{FC.1c}
\]
The substitutions are valid as continuous maps of the indicated
profinite presentations, as well as for the discrete groups.

There is the useful exact identity
\[
                    \boxed{\ C'=K'K^2.\ }
 \tag{FC.1d}
\]
Here and below all derived and power subgroups are closed.
Indeed $C/(K'K^2)$ is abelian: the action of $j$ on
$K/(K'K^2)\simeq(C_2)^2$ is trivial. Conversely $K'\subset C'$
and $[x,j]=x^2$, $[y,j]=y^2$, which prove the reverse inclusion.
The subgroup $C'$ is therefore an index-four subgroup of the
free group $K$ of rank two. It is free of rank five, and
\[
  C'/(C''(C')^3)\simeq\mathbb F_3^5,\qquad
  |C/(C''(C')^3)|=8\cdot3^5=1944.
 \tag{FC.1e}
\]

\begin{statement}{A.116}{Proved}{exact actions, cusp, and the rank-five decomposition}
\end{statement}
The actions of $a,b$ on $C$ are
\[
\begin{split}
 a:(t_0,t_1,t_2)&\mapsto(t_0t_1t_0,t_0,t_2),\\
 b:(t_0,t_1,t_2)&\mapsto(t_2,t_1,t_2t_0t_2).
\end{split}
 \tag{FC.2a}
\]
For example (FC.1b), the centrality of $z_3$ in the
three-strand section, and (FC.1a) give
$aja^{-1}=x^{-1}j$ and $bjb^{-1}=yj$; the other four
images follow immediately. The generator $c=xa$ acts as
$\operatorname{Inn}(t_1t_0)$ followed by $a$.
Both $a$ and $b$ fix the peripheral element
\[
 w=t_1t_0t_2=xyj=z_3^{-1},\qquad w^2=[x,y].
 \tag{FC.2b}
\]
The last equality is a direct free-product reduction. In particular
the word order in $w$ is fixed here; permuting its three letters
would alter the marked cusp.

Write $Q=C/C'=(C_2)^3$, with the ordered basis $(t_0,t_1,t_2)$,
and $M=C'^{\rm ab}/3$. The Cayley covering graph for $C'$
is the cube, with eight vertices and twelve unoriented edges.
Its cellular modules over $\mathbb F_3$ are
\[
 C_0=\mathbb F_3[Q],\qquad
 C_1=\bigoplus_{i=0}^2
    \operatorname{Ind}_{\langle t_i\rangle}^{Q}\operatorname{sgn}.
\]
For a nontrivial character $\chi:Q\to\{\pm1\}$ of weight
$k$, the edge multiplicity is $k$, the vertex multiplicity is one,
and $H_0$ has multiplicity zero. Thus its multiplicity in $M$
is $k-1$. The trivial character occurs in neither $H_1$ nor any
edge summand. Hence
\[
 M=B\oplus P,\quad
 B=\chi_{110}\oplus\chi_{101}\oplus\chi_{011},\quad
 P=\chi_{111}^{\oplus2}.
 \tag{FC.2c}
\]
The even-parity subgroup $Q_0=K/C'$ is the kernel of
$\chi_{111}$. Since its order four is invertible in
$\mathbb F_3$, averaging over $Q_0$ is an idempotent projector
with image $P$. Transfer identifies this image with
$H_1(K,\mathbb F_3)$. Because $K\triangleleft L$ and
$C'\triangleleft L$, the splitting $M=B\oplus P$ is
$L$-invariant and canonical.

The independent program \texttt{check\_ferrari\_cube.py}
constructs the cube and a spanning tree, obtains five free
Schreier loops, and calculates every deck and $a,b$ matrix by
substituting the literal words (FC.2a). It imports no saved finite
quotient or matrix. It verifies the eight character multiplicities,
the two inverse automorphisms, the braid relation, (FC.1a), and
(FC.2b). The boundary linear image of $\langle a,b\rangle$
has order six, the proper linear image order twenty-four, and
their joint image order one hundred and forty-four.

The cusp class $[w^2]$ is a nonzero element of $B$ and is not
$Q$-invariant. This is verified by its explicit cube cycle, and
also follows from the geometric description below. Therefore the
central element $z_3^2=w^{-2}$ of the three-strand section acts
on $C/(C''(C')^3)$ by a nontrivial inner automorphism even though
its action on $M$ is trivial. Retaining only the rank-five linear
matrices would lose this action. The affine models used in A.119
and A.121 retain it.

\begin{statement}{A.117}{Proved}{the characteristic elliptic cover and the scope of its geometric interpretation}
\citep[Sections 2.9 and 2.15]{LitNakamuraTsunogaiYasuda2010}
\end{statement}
The geometric Ferrari identification of Nakamura, Tsunogai, and
Yasuda (2010, Sections 2.9 and 2.15) identifies $K$ with the
geometric fundamental group of a once-punctured elliptic curve
$E\setminus\{O\}$. Under this identification, $K'K^2$ is the
kernel of the mod-two abelianization map. Its associated connected
degree-four cover is exactly
\[
       [2]:E\setminus E[2]\longrightarrow E\setminus\{O\}.
 \tag{FC.3a}
\]
Indeed the proper multiplication-by-two covering corresponds to
$2H_1(E,\mathbb Z)$, and filling the single puncture identifies
$H_1(E\setminus\{O\},\mathbb Z/2)$ with
$H_1(E,\mathbb Z/2)$. This proves the asserted equality of covering
subgroups. A path to a chosen lift of the basepoint changes only
the inner representative of the identification.

The resulting geometric exact sequence is
\[
 0\longrightarrow
 \mathbb F_3[E[2]]\Big/\mathbb F_3\sum_{P\in E[2]}P
 \longrightarrow H_1(E\setminus E[2],\mathbb F_3)
 \longrightarrow H_1(E,\mathbb F_3)\longrightarrow0.
 \tag{FC.3b}
\]
It has dimensions $3,5,2$. The deck group $Q_0=E[2]$
acts trivially on the proper quotient and regularly on the four
punctures, so averaging over $Q_0$ gives precisely the canonical
splitting (FC.2c). In this interpretation the involution with
class $(1,1,1)\in Q$ fixes all four punctures. The individual
$t_i$ represent that involution composed with the three nonzero
two-torsion translations; they do not each fix all punctures.
This accounts for the boundary weights being two rather than one.

This geometric identification is a statement about a characteristic
cover and its geometric homology. It supplies no unproved
all-$\widehat{GT}$ comparison of the special-fibre paths in
Nakamura--Tsunogai--Yasuda, Proposition 4.13. In particular it does
not identify the cubic frame character $\kappa_d$ with the
quadratic character $\nu_4$.

\subsection{The complete obstruction to odd relative refinements of a forgetful base}


\begin{statement}{A.118}{Proved}{odd relative homology over a forgetful base stays forgetful}
\end{statement}
There is a structural obstruction to adding a characteristic-three
Fox module over the forgetful base $Q_{12}$, and it holds for
every finite forgetful base and every odd relative characteristic.
Let $q_L:L=B_4/\langle z_4\rangle\to\widehat B$ be the
fold followed by quotienting the three-strand centre. Then
\[
                  C=\ker(q_L)\simeq C_2\amalg C_2\amalg C_2.
 \tag{FP.40.1}
\]
Indeed the Ferrari kernel is $K=\widehat F(\xi_1,\xi_2)$
by (B.35). The element $j=(abc)^2$ has order two in $L$,
and its conjugation inverts each $\xi_i$, by squaring
(FP.26.6). The kernel $C$ is $K\rtimes\langle j\rangle$:
upstairs $q((abc)^2)=z_3$, and quotienting $z_4$ makes
this remaining central quotient have order two. Its presentation
\[
 \langle\xi_1,\xi_2,j\mid j^2=1,
           j\xi_1j=\xi_1^{-1},\ j\xi_2j=\xi_2^{-1}\rangle
\]
becomes the free product of three involutions
$j,\xi_1j,\xi_2j$, proving (FP.40.1), discretely and after
completion.

Let $N_B\triangleleft_o\widehat B$ be any finite base kernel
and $N_L=q_L^{-1}(N_B)$. For any odd prime $p$, the three
involution generators of $C$ belong to $N_L^p$, since an
involution equals its own $p$th power. Hence
$C\subseteq N_L'N_L^p$. The natural quotient is therefore
exactly
\[
 \boxed{\quad L/(N_L'N_L^p)
           \simeq\widehat B/(N_B'N_B^p).\quad}
 \tag{FP.40.2}
\]
To see equality of the indicated kernels, quotient $N_L$ by
$C$: commutators and $p$th powers map onto those of $N_B$,
and $C$ is already killed. Thus the purported relative homology
extension contains no additional Brunnian information. Its
induced five-point quotient factors through the five forgetful
maps, where (A.74) makes the pentagon identically invisible.
The same proof applies to relative exponent $p^r$ for every
$r\geq1$. In particular the proposed characteristic-three
Fox refinement over the order-$576$ congruence base is excluded
without an unnecessary $2304\times1728$ matrix computation.

In fact the obstruction holds for the entire relative pro-$p$
completion. Write $O^p(N)$ for the intersection of all open
normal subgroups of $N$ with $p$-group quotient. Every homomorphism
from $N_L$ to a finite odd $p$-group kills the three involutions,
and hence kills $C$. The universal property gives
\[
 L/O^p(N_L)\simeq\widehat B/O^p(N_B).
 \tag{FP.40.3}
\]
Thus iterating odd relative homology over the same forgetful
base cannot evade this obstruction.

\subsection{Explicit mixed Ferrari quotients and their nonlinear coordinates}


\begin{statement}{A.119}{Proved}{the first mixed homology stage and its precise pentagon limitation}
\end{statement}
Put $V=\mathcal C^{\rm ab}\simeq(\mathbb F_2)^3$ with
basis the three $t_i$. Its commutator subgroup is free of
rank five: the eight-sheeted Cayley covering is the cube,
with twelve unoriented edges and eight vertices.
Over $\mathbb F_3$, its homology is
\[
 H=\mathcal C'^{\rm ab}/3
   =B_0\oplus P_0,
 \qquad\dim B_0=3,\quad\dim P_0=2,
\]
with respective $V$ character masks
$(3,5,6)$ and $(7,7)$. A mask acts by minus one precisely
when its dot product with the base element is one.
Equivalently the three distinct pair-sign characters occur
once, and the all-minus character occurs twice.
The absence of other characters follows from the cellular
character formula
$2\mathbb F_3[V]-\sum_i\mathbb F_3[V/\langle t_i\rangle]
 +\mathbb F_3$.

An explicit model of
$\mathcal C/(\mathcal C''\mathcal C'^3)$ is
$H\rtimes V$, of order $8\cdot3^5=1944$, generated by
\[
 t_0=(0,e_0),\quad
 t_1=((1,0,0,1,0),e_1),\quad
 t_2=((0,1,1,0,1),e_2).
\]
The exact conjugation actions induced by $a,b$ are
\[
 \begin{split}
 a:(t_0,t_1,t_2)&\longmapsto(t_0t_1t_0,t_0,t_2),\\
 b:(t_0,t_1,t_2)&\longmapsto(t_2,t_1,t_2t_0t_2),
 \end{split}
\]
and $c$ acts as $\operatorname{Inn}(t_1t_0)$ followed by
the action of $a$. For example
$j=\xi_1^{-1}\xi_2^{-1}z_3$, so the Nielsen formulas give
$aja^{-1}=\xi_1^{-1}j$ and $bjb^{-1}=\xi_2j$.
This derives the displayed actions without a choice of
extension of an outer automorphism.

The affine action on $H$ splits into its two displayed
blocks. Each of the five pure $X_i$ has identity linear
part on $B_0$, so those five images are translations
there. Every evaluation of $f\in F'$ is therefore one on
that block. On $P_0$, simultaneous conjugation by the
translation $(1,1)$ gives exactly
\[
           a\mapsto(0,U),\qquad b\mapsto(0,L_0),
                    \qquad c\mapsto(e_1,U),
\]
the already audited order-$216$ affine quotient.
Adjoining $q_{12}$ through the fold does not alter this
pentagon test: its five $X_i$ images are
$(x,y,x,1,1)$, and the resulting pentagon is
$f(y,x)f(x,y)=1$ by antisymmetry.
In particular A.94's actual witness passes this joint
test with $\Delta=1$. The exact affine matrices and all
source relators are checked in
\texttt{check\_mixed\_ferrari\_quotient.py}.

\begin{statement}{A.120}{Proved}{an actual class-two quotient retaining the six mixed commutators}
\end{statement}
Let
\[
 Z=B_0\otimes P_0,\qquad\dim Z=6,
 \qquad N=H\oplus Z.
\]
Give $N$ the class-two exponent-three group law
\[
 (v,z)(v',z')=
 \left(v+v',\ z+z'+\frac12[v,v']\right),
 \quad
 [(b,p),(b',p')]=b\otimes p'-b'\otimes p.
\]
Here $1/2=2$ in $\mathbb F_3$.
The $V$ action on $Z$ has masks $(4,4,2,2,1,1)$.
Set $E=N\rtimes V$, of order
\[
                         |E|=8\cdot3^{11}=1\,417\,176.
\]
Use the same three $t_i$ as above, with zero $Z$
coordinates. They are involutions, since their degree-one
coordinates lie in the minus eigenspaces of the respective
base elements. They generate $E$: their commutator subgroup
surjects onto $H$, as checked in A.119, and the Frattini
criterion for the finite three-group $N$ then gives
surjectivity onto $N$; equivalently their mixed brackets
span all of $Z$.

This is a canonical $L$-stable quotient of the full
class-two exponent-three quotient of $\mathcal C'$.
Indeed that full quotient has centre $\bigwedge^2H$.
The discarded subspace
$\bigwedge^2B_0\oplus\bigwedge^2P_0$ is stable under
the exact $a,b,c$ actions, because their degree-one
matrices preserve $B_0,P_0$ and the action on commutators
is their exterior square. The quotient retains exactly
the $B_0\otimes P_0$ component. Since the first-stage
group $H\rtimes V$ is centreless and $Z^V=0$, the group
$E$ is centreless. Its conjugation action therefore gives
an actual finite quotient of $L$, faithful on this
specified quotient of $\mathcal C$.

Here are complete small coordinates for that action.
Write $M_a=\operatorname{diag}(B_a,U)$,
$M_b=\operatorname{diag}(B_b,L_0)$,
$M_c=\operatorname{diag}(B_c,U)$ over $\mathbb F_3$, where
\[
 B_a=\begin{pmatrix}1&0&0\\0&0&1\\0&1&0\end{pmatrix},
 \quad
 B_b=\begin{pmatrix}0&0&2\\0&1&0\\2&0&0\end{pmatrix},
 \quad
 B_c=\begin{pmatrix}1&0&0\\0&0&2\\0&2&0\end{pmatrix}.
\]
The corresponding inner translations $n_g=(u_g,z_g)$
are
\[
\begin{array}{c|c|c}
g&u_g&z_g\\ \hline
a&(1,0,0,1,0)&(2,1,0,0,0,0)\\
b&(2,2,2,0,2)&(0,1,2,2,0,2)\\
c&(2,0,0,2,0)&(2,1,0,0,0,0).
\end{array}
\]
Thus $g$ acts as $\operatorname{Inn}(n_g)$ followed by
the graded action $M_g$ on $H$, $B_g\otimes P_g$ on
$Z$, and its induced permutation on $V$.
Equivalently these are elements $(n_g,M_g)$ of the finite
semidirect group with multiplication
\[
       (n,M)(n',M')=(n\,M(n'),MM').
\]
Their central coordinates are not guessed from an exterior
power representation. They are obtained by evaluating
the three exact involution images above with the BCH law,
then solving their central-coordinate equations. Each
central character is nontrivial, so the solution is unique;
the solution is checked on all three generators.

The code \texttt{check\_mixed\_ferrari\_class2.py} verifies
the four defining relations of $L$, as well as
\[
 (abc)^2=\operatorname{Inn}(t_0),\qquad
 ca^{-1}=\operatorname{Inn}(\xi_1),\qquad
 aba c^{-1}b^{-1}a^{-1}=\operatorname{Inn}(\xi_2)
\]
inside this exact model. It supplies all five $X_i$
coordinates and reusable group operations.
This paragraph constructs the genuine finite target.
It makes no assertion that its pentagon forces $\Delta=0$,
that a finite pentagon assignment lifts, or that its pair
second derived groups are nonnilpotent; those are separate
questions about this target.

\Needspace{10\baselineskip}
\subsection{The first homology and full class-two Ferrari tests are insufficient}


\begin{statement}{A.121}{Proved}{the first mixed two--three Ferrari cover adds no pentagon equation}
\end{statement}
The preceding obstruction motivates retaining $C^{\rm ab}$
before taking homology in characteristic three. Put
\[
 C=\langle t_0,t_1,t_2\mid t_0^2=t_1^2=t_2^2=1\rangle,
 \qquad
 \overline C=C/(C''(C')^3).
\]
The subgroup $C'$ is free of rank five and index eight, so
$|\overline C|=8\cdot3^5=1944$. In the notation of (A.118),
$t_0=j$, $t_1=\xi_1j$, $t_2=\xi_2j$. The exact conjugations
by the generators of $L$ are
\[
\begin{split}
 a:(t_0,t_1,t_2)&\longmapsto(t_0t_1t_0,t_0,t_2),\\
 b:(t_0,t_1,t_2)&\longmapsto(t_2,t_1,t_2t_0t_2),\\
 c&=\operatorname{Inn}(t_1t_0)\,a.
\end{split}
 \tag{FP.41.1}
\]
These follow directly from the Ferrari Nielsen formulas in
(B.35). In particular this constructs an actual homomorphism
$L\to\operatorname{Aut}(\overline C)$.

The character decomposition of $V=C'^{\rm ab}/3$ under
$C^{\rm ab}=(C_2)^3$ is
\[
 V=B\oplus P,
 \qquad
 B=\chi_{110}\oplus\chi_{101}\oplus\chi_{011},
 \qquad P=\chi_{111}^{\oplus2}.
 \tag{FP.41.2}
\]
For completeness, the Cayley graph of the three involutions
has edge module $\bigoplus_i\operatorname{Ind}_{\langle t_i\rangle}^{(C_2)^3}
\operatorname{sgn}$ and vertex module the regular representation.
A nontrivial character of weight $k$ therefore occurs in its
first homology with multiplicity $k-1$. This proves (FP.41.2).
The summand $P$ identifies canonically with the homology of the
index-two Ferrari subgroup $K$, by averaging over the even
subgroup $(C_2)^2$. The other three lines are its boundary
homology after passing to the index-four cover of $K$.

There is a precise collapse in the pentagon test. In a faithful
affine representation of the automorphism group on $V$, all
five $X_i$ act on $B$ by translations, with vectors
\[
 (2,0,0),\quad(0,1,0),\quad(1,0,0),\quad
 (2,1,1),\quad(1,1,2).
 \tag{FP.41.3}
\]
They preserve the direct sum with the affine $P$ block. Their
$P$ blocks, after one common conjugation by translation $(1,1)$,
are precisely the images of the same five words under
\[
 a\longmapsto(0,U),\qquad
 b\longmapsto(0,L_0),\qquad
 c\longmapsto(e_1,U)
 \quad\hbox{in }\mathbb F_3^2\rtimes\operatorname{SL}_2(\mathbb F_3).
 \tag{FP.41.4}
\]
This is exactly the quotient $\eta_3$ of (A.75), with its
marking preserved, not merely an abstract group isomorphism.

Every $f\in\widehat F'$ is one on an abelian target.
Consequently all five $f$-factors have identity $B$ block,
and their entire pentagon product is exactly the old
$\eta_3$ product. For actual gentle elements the necessary
equation is still $a_f=\rho$ from (A.77). Adjoining the
forgetful $q_{12}$ component changes nothing: there the
five $X_i$ have images $(x,y,x,1,1)$, and the pentagon
cancels by the two-cycle relation.

More strongly, every evaluation-pair group in this joint
quotient has derived length at most three. The boundary
projection is abelian; the compact pair series are
$24,8,2,1$ for the two harmonic pairs and
$18,3,1$ or $6,3,1$ for the others. The forgetful level-twelve
pair groups also have derived length at most three.
The actual element from A.94, with $f\in F'''$ and
$\Delta=1$, therefore kills every one of the five factors,
not just their product. The induced five-point quotient has
the same obstruction.

The construction and every automorphism identity on all $1944$
elements are independently checked in the lifting audit's
\texttt{check\_mixed\_ferrari\_quotient.py}. The exact common
conjugacy, all boundary blocks, and all compact pair series
are checked in \texttt{check\_mixed\_ferrari\_collapse.py} and
its output. This result leaves open the next nonabelian
stage, where commutators between $B$ and $P$ are retained.

\begin{statement}{A.122}{Proved}{the complete class-two Ferrari stage is also insufficient}
\end{statement}
Retain the mixed commutators suggested by (A.121). Over
$\mathbb F_3$ put $V=B\oplus P$, $W=B\otimes P$ and
\[
 N=V\oplus W,\qquad
 (v,z)(v',z')=(v+v',z+z'+\tfrac12[v,v']),
 \quad
 [v,v']=v_B\otimes v'_P-v'_B\otimes v_P.
 \tag{FP.42.1}
\]
The group $A=(C_2)^3$ acts with the characters in (FP.41.2)
on $V$ and their products on $W$. The six characters on $W$
are the three characters of weight one, each twice. Let
$E=N\rtimes A$. It has order $8\cdot3^{11}=1417176$.
In the coordinates of the certificate, the three elements
\[
 t_0=((0,0),e_0),\quad
 t_1=(((1,0,0,1,0),0),e_1),\quad
 t_2=(((0,1,1,0,1),0),e_2)
\]
are involutions. Their images generate $E/[N,N]$ by (A.121),
and hence generate $E$: their intersection with $N$ generates
$N/\Phi(N)$, so the finite-$3$-group basis theorem applies.
Thus $E$ is an actual quotient of $C$, not just a module
with a putative action. It is obtained from the full free
class-two exponent-three quotient of $C'$ by killing
$\bigwedge^2B\oplus\bigwedge^2P$.

The two sets of characters on $V$ and $W$ are disjoint.
Consequently $\operatorname{Hom}_A(V,W)=0$, and
$C_N(A)=1$. Schur--Zassenhaus, followed by this vanishing,
shows that every induced automorphism has the unique form
$\operatorname{Inn}(n)$ followed by its graded linear map.
This reduces the exact computation to six additional central
coordinates. The images (FP.41.1) determine those coordinates
uniquely, and give a homomorphism
$L\to N\rtimes\operatorname{GL}(V)$. Every braid,
distant-commutation and full-twist relator is checked in the
actual class-two multiplication, including all nonlinear terms.

There is an explicit structural explanation of the pair groups.
Let $h\in N$ have
\[
 h_V=(0,0,0,1,1),\qquad
 h_W=(0,2,2,0,2,0).
\]
Conjugation by $h$ sends the harmonic pair to
\[
 hX_1h^{-1}=(2e_1,U^2),\qquad
 hX_2h^{-1}=(e_2,L_0^2),
 \tag{FP.42.2}
\]
with zero compact and mixed translations. Hence that pair
embeds into the direct product of an abelian boundary group
and $\operatorname{SL}_2(\mathbb F_3)$. The reversing half
twist gives the other harmonic pair. In each remaining pair,
the two compact linear parts commute. Its first derived group
therefore lies in $P\oplus W$, an abelian subgroup of $N$.
The exact orders and derived series are
\[
\begin{array}{c|c}
 (X_1,X_2),(X_2,X_3)&72,8,2,1\\
 (X_3,X_4),(X_5,X_1),(X_4,X_5)&486,27,1.
\end{array}
 \tag{FP.42.3}
\]
Schreier rewriting and an independent enumeration of each
complete pair group agree. The harmonic pairs have no kernel
in the six new mixed coordinates; each other pair has a
two-dimensional such kernel. Thus the available mixed module
does not produce a coprime action in the harmonic second
derived subgroup.

The same conclusion holds for the \emph{full} class-two
quotient
\[
 E_{\rm full}=C/(\gamma_3(C')(C')^3),
       \qquad |E_{\rm full}|=8\cdot3^{15}.
 \tag{FP.42.4}
\]
Here all powers and subgroups are understood topologically.
The omitted kernel inside the nilpotent subgroup is
$Z_0=\bigwedge^2B\oplus\bigwedge^2P$, of dimension four.
Consider the kernel $\mathcal K$ of the induced map on the
automorphism groups preserving these canonical character
summands. An automorphism in $\mathcal K$ is the identity on
$V$, hence on $\bigwedge^2V$, and is described by a
$Z_0$-valued cocycle. These cocycles add under composition,
so $\mathcal K$ is elementary abelian of characteristic three.
Its two types of coordinates are
\[
 \operatorname{Inn}(\bigwedge^2B),\qquad
 \operatorname{Hom}_A(B,\bigwedge^2B),
 \tag{FP.42.5}
\]
each of dimension three. The invariant line
$\bigwedge^2P$ is central in $E_{\rm full}$ and gives no
additional inner automorphism. There are no other equivariant
maps from $V$ to $Z_0$ by their character decompositions.

One must distinguish this automorphism kernel from the
four-dimensional central kernel in the nilpotent subgroup.
A pure $X_i$ can shear the derivation coordinate in (FP.42.5)
by an inner correction, because its boundary translation need
not vanish. Every element of the second derived pair group,
however, has zero boundary translation. Its boundary linear
part is the identity, and its compact part is $\pm I$ in the
harmonic case, or the identity in the other cases.
It therefore acts trivially on both summands in (FP.42.5).
Explicitly, write a kernel element as
$\operatorname{Inn}(z)\,d_D$, where
$d_D(v,w)=(v,w+D(v_B))$. Conjugation by an element with
boundary linear part $I$ and boundary translation $u_B$
sends $(z,D)$ to $(z-D(u_B),D)$; the compact scalar $\pm I$
does not affect these coordinates. This proves the asserted
triviality when $u_B=0$.
The second derived subgroup in the full quotient is a central
extension of $C_2$ (harmonic case), or of the trivial group,
by a subgroup of $\mathcal K$. A central extension of a cyclic
group is abelian. Consequently all five full pair images still
have derived length at most three.

The actual A.94 element with $f\in F'''$ and $\Delta=1$
therefore satisfies every pentagon factor separately in both
class-two quotients, also after adjoining any of the previous
forgetful level-twelve factors. No assertion that this element
belongs to GT is made. The exact nonlinear construction is
\texttt{agent\_lifting/code/check\_mixed\_ferrari\_class2.py};
the independent pair proof and complete enumeration are
\texttt{check\_mixed\_class2\_pairs.py} and its output.

\subsection{An actual mixed-sign witness on the entire relative pronilpotent Ferrari quotient}


\begin{statement}{A.123}{Proved}{the compact target and its congruence kernel}
\end{statement}
Retain the notation of A.115--A.117 and A.118:
\[
 L=\widehat B_4/\overline{\langle z_4\rangle},\qquad
 C=\ker(L\longrightarrow\widehat B)
      =C_2\amalg C_2\amalg C_2.
\]
Let $O^3(C')$ be the intersection of the kernels of all continuous
homomorphisms from $C'$ to finite three-groups, and put
\[
 \Gamma=L/O^3(C'),\qquad
 C_\infty=C/O^3(C'),\qquad
 N=C'/O^3(C').
 \tag{FR.1a}
\]
The subgroup $C'$ is characteristic in $C$, which is normal in $L$;
$O^3(C')$ is characteristic in $C'$. Thus it is a closed normal
subgroup of $L$, and every group and quotient in (FR.1a) is a
genuine profinite group. There are exact sequences
\[
 1\longrightarrow N\longrightarrow C_\infty
 \longrightarrow Q=(C_2)^3\longrightarrow1,
 \qquad
 1\longrightarrow C_\infty\longrightarrow\Gamma
 \longrightarrow\widehat B\longrightarrow1.
 \tag{FR.1b}
\]
Here $N$ is free pro-three of rank five, since $C'$ is free
profinite of rank five. In particular $N$ is finitely generated,
and so are $C_\infty$ and $\Gamma$.

Write
\[
 M=N/\Phi(N)=N/(N'N^3)\simeq\mathbb F_3^5.
\]
The $Q$-action on $M$ has characters $110,101,011,111,111$
by A.116. These characters span $Q^*$, so that action is faithful.
The centralizer of $N$ in $C_\infty$ is trivial: an element
centralizing $N$ acts trivially on $M$, hence has trivial image
in $Q$; it then belongs to the centre of the nonabelian free
pro-three group $N$, which is trivial. Here is a direct
proof of the free-pro-$p$ centre assertion used at every prime.
For a free pro-$p$ group of rank $r\geq2$, the universal map
\[
 x_i\longmapsto1+T_i
 \quad\hbox{into}\quad
 1+(T_1,\ldots,T_r)\mathbb F_p\langle\!\langle
 T_1,\ldots,T_r\rangle\!\rangle
\]
is injective. Indeed the target is pro-$p$ by truncation in
degree. For every finite $p$-group quotient $H$, substitution
$T_i\mapsto\bar x_i-1$ is defined in $\mathbb F_p[H]$
because its augmentation ideal is nilpotent. It recovers
the original quotient map, so an element in the displayed
kernel is trivial in every finite quotient. Nilpotence of
that augmentation ideal follows by induction on $|H|$:
choose a central element $h$ of order $p$; the central ideal
$(h-1)$ has $p$th power zero, and the quotient is the group
algebra of $H/\langle h\rangle$.

A formal series commuting with $T_1$ contains only powers
of $T_1$. To see this coefficientwise, moving trailing
$T_1$'s to the beginning preserves the coefficient by
$T_1A=AT_1$; any remaining word ending in another letter
has coefficient zero by the same identity. Thus the
centralizer of $T_1$ is $\mathbb F_p[[T_1]]$.
Its intersection with the centralizer of $T_2$ is
$\mathbb F_p$. A central group element has augmentation
one, so its series is one. This proves the asserted
centrelessness and hence
\[
           C_{C_\infty}(N)=1,
           \qquad Z(C_\infty)=1.
 \tag{FR.1c}
\]

Conjugation defines a continuous homomorphism
\[
 \rho:\Gamma\longrightarrow\operatorname{Aut}(C_\infty),
 \qquad \rho_0:\Gamma\longrightarrow\operatorname{GL}(M).
\]
The finite generation ensures that these automorphism groups
carry their usual profinite topology: finite characteristic
quotients are cofinal. All images below are compact, hence
closed. The subgroup $N$ is characteristic in $C_\infty$,
being its unique pro-three Sylow subgroup, so restriction gives
$\operatorname{Aut}(C_\infty)\to\operatorname{Aut}(N)$.
It is injective by (FR.1c): if an automorphism fixes $N$
pointwise, the difference between each element and its image
centralizes $N$.

The kernel of
$\operatorname{Aut}(N)\to\operatorname{GL}(M)$ is pro-three.
Here is a filtration proof. Set
\[
 N_1=N,\qquad N_{r+1}=\overline{N_r^3[N_r,N]}.
\]
These characteristic open subgroups are cofinal, because the
corresponding series terminates in every finite three-group.
The quotients $N_r/N_{r+1}$ are central elementary abelian
sections and are generated by the commutators and powers
coming from $N/N_2$. An automorphism trivial on $N/N_2$
therefore acts trivially on every such section. If $U_r$
denotes the subgroup of automorphisms trivial on $N/N_r$,
then, for $r\geq2$, the map
\[
 U_r/U_{r+1}\longrightarrow
 \operatorname{Hom}(N/N_2,N_r/N_{r+1}),\qquad
 \alpha\longmapsto\bigl(g\mapsto\alpha(g)g^{-1}\bigr)
 \tag{FR.1d}
\]
is an injective group homomorphism. Its target is elementary
abelian of exponent three. It follows by induction that
$U_2/U_r$ is a finite three-group for every $r$, and hence
$U_2=\varprojlim_rU_2/U_r$ is pro-three.
Consequently the congruence kernel
\[
 \ker\bigl(\rho(\Gamma)\longrightarrow\rho_0(\Gamma)\bigr)
                       \quad\hbox{is pro-three.}
 \tag{FR.1e}
\]
One may avoid the centreless restriction argument in proving
(FR.1e): the additional kernel is a group of cocycles from
$Q$ into $Z(N)$, itself pro-three. The centreless statement
is needed later for the stronger conclusion inside $\Gamma$.

\begin{statement}{A.124}{Proved}{the exact braid centre controls both harmonic pairs}
\end{statement}
Let $G$ be the closed subgroup of $\operatorname{Aut}(C_\infty)$
generated by $\rho(a),\rho(b)$, and let $G_0$ be its image
on $M=B\oplus P$. On the boundary space $B$, its image
is $S_3$; on the proper space $P$, it is
$\operatorname{SL}_2(\mathbb F_3)$, with
\[
 a\longmapsto U=\begin{pmatrix}1&1\\0&1\end{pmatrix},
 \qquad
 b\longmapsto L_0=\begin{pmatrix}1&0\\-1&1\end{pmatrix}.
 \tag{FR.2a}
\]
The boundary assertion follows either from the exact actions
(FC.2a), or the cube certificate: the two boundary generators
are involutions satisfying the braid relation and generate
a group of order six. Thus $G_0$ embeds in
$S_3\times\operatorname{SL}_2(\mathbb F_3)$ and projects
onto both factors. Since
\[
 S_3''=1,\qquad
 \operatorname{SL}_2(\mathbb F_3)'=Q_8,\qquad
 \operatorname{SL}_2(\mathbb F_3)''=\{\pm I\},
\]
its second derived group is exactly
\[
              G_0''=\{(I_B,I_P),(I_B,-I_P)\}\simeq C_2.
 \tag{FR.2b}
\]
No assertion about independence of the two linear factors is
needed for this conclusion: the inclusion follows from the
two displayed derived groups, and surjectivity onto the proper
factor proves equality.

The source element $z_3=(ab)^3$ is central in the actual
three-strand section. Its image $z=\rho(z_3)$ is therefore
central in $G$. On $M$ it has the value $(I_B,-I_P)$.
Let $e_2\in\widehat{\mathbb Z}$ be one in the two-adic
factor and zero in all other prime factors. Since the kernel
of $G\to G_0$ is pro-three, the procyclic group generated
by $z$ has order dividing $2\cdot3^\infty$ and maps onto
the displayed $C_2$. Hence
\[
                         \iota=z^{e_2}
\]
is a central involution of $G$ with the same nontrivial
linear image. If $U_G=\ker(G\to G_0)$, then
\[
  (G\longrightarrow G_0)^{-1}(G_0'')
                 =U_G\times\langle\iota\rangle.
 \tag{FR.2c}
\]
The right-hand side is pronilpotent, being the direct product
of a pro-three group and a group of order two. It contains
$G''$. Every closed subgroup of a pronilpotent group is
pronilpotent. Thus the second derived group of
$\overline{\langle\rho(a^2),\rho(b^2)\rangle}$ is
pronilpotent. Reversing the braid generators gives the identical
conclusion for
$\overline{\langle\rho(b^2),\rho(c^2)\rangle}$.

This central lift is essential. An arbitrary extension of
$C_2$ by a pro-three group need not be pronilpotent. The
argument uses centrality of the literal braid $z_3$ in
the actual harmonic subgroup; it does not assume that
$-I$ acts trivially on every direction in the ambient
automorphism group.

\begin{statement}{A.125}{Proved}{all five pair second-derived images are pronilpotent}
\end{statement}
Use the five literal pentagon inputs of A.96 and A.121:
\[
 X_1=a^2,\quad X_2=b^2,\quad X_3=c^2,\quad
 X_4=(aba)^2=(ab)^3,\quad X_5=(cbc)^2=(cb)^3.
 \tag{FR.3a}
\]
Every $X_i$ has identity boundary linear action. On the
proper two-dimensional block their images are respectively
\[
                    U^2,\quad L_0^2,\quad U^2,\quad -I,\quad -I.
 \tag{FR.3b}
\]
The two harmonic pairs $(X_1,X_2)$ and $(X_2,X_3)$
are covered by A.124. Each of the other three pairs
$(X_3,X_4)$, $(X_5,X_1)$, $(X_4,X_5)$ has abelian
linear image, by (FR.3b). Its first derived action is
therefore contained in the pro-three congruence kernel
of (FR.1e). In particular its second derived action
is pro-three. All five action-pair groups consequently
have pronilpotent second derived subgroup.

Let $R_{\rm nil}(F'')$ be the kernel of the maximal
pronilpotent quotient of the second derived subgroup of
the free profinite group $F=\widehat F(x,y)$. For every
\[
                         f\in R_{\rm nil}(F''),
 \tag{FR.3c}
\]
the image of $f$ under each of the five pair evaluations
in $\operatorname{Aut}(C_\infty)$ is one. Indeed a pair
map sends $F''$ onto the closed second derived group
of its image. A continuous map from $F''$ to a
pronilpotent group kills $R_{\rm nil}(F'')$, by its
defining universal property.

This conclusion is simultaneous at every finite characteristic
three-group quotient of $C'$ and in their inverse limit.
No selection of finite residues or passage from finite
solutions to a genuine $\widehat{GT}$ element occurs.
Adjoining the forgetful $q_{12}$ action does not change it:
the five input images there are $(x,y,x,1,1)$; the two
nontrivial pair groups have nilpotent second derived
group. Thus the joint pair second-derived image is a
closed subgroup of a product of pronilpotent groups.

\begin{statement}{A.126}{Proved}{the pentagon holds in the actual quotient with full forgetful base}
\end{statement}
Suppose in addition to (FR.3c) that
$f(y,x)=f(x,y)^{-1}$. Let $\mathcal P(f)$ denote the
literal five-point pentagon product in $L$. Its image
under $L\to\widehat B$ is one: the five inputs become
$(x,y,x,1,1)$, so the only nontrivial factors cancel
by this two-cycle identity, as in A.74. Consequently
the image of $\mathcal P(f)$ in $\Gamma=L/O^3(C')$
belongs to $C_\infty$.

By A.125, each factor has trivial conjugation action on
$C_\infty$. Their product therefore centralizes
$C_\infty$. Equation (FR.1c) now gives
\[
             \boxed{\ \mathcal P(f)=1
                       \quad\hbox{in }L/O^3(C').\ }
 \tag{FR.4a}
\]
This is stronger than equality in an automorphism image:
the quotient retains the complete profinite forgetful
base $\widehat B$. Individual factors are asserted to
have trivial action; they need not themselves be one
in this larger quotient. It is their pentagon product
that belongs to the centreless kernel and hence is one.

A.94 supplies an actual gentle element $k$ with
\[
            f_k\in R_{\rm nil}(F''),\qquad\Delta(k)=1.
\]
Its two-cycle identity is part of the genuine gentle
equations. Thus this same actual element satisfies
(FR.4a) and all of its finite relative three-group
quotients, while retaining $\Delta(k)=1$. The entire
relative pro-three Ferrari family is therefore
insufficient to force the required parity identity.
This does not assert the full profinite pentagon,
membership of $k$ in $\widehat{GT}$, or an analogous
result for relative primes other than three.

\begin{statement}{A.127}{Proved}{the integral marking and the other relative primes}
\end{statement}
For every prime $\ell$, put
\[
 C_\ell=C/O^\ell(C'),\qquad
 \Gamma_\ell=L/O^\ell(C').
 \tag{FR.5a}
\]
For odd $\ell$, the cube idempotents of A.116 are defined over
$\mathbb Z[1/2]$. More intrinsically, averaging over the deck
group $K/C'$ of order four splits the integral homology after
inverting two into puncture homology and proper elliptic homology.
The proper block is identified with $H_1(K)$ by pushforward
and transfer, whose composite is multiplication by four.
The actions of $a,b,c$ on that block are respectively $U,L_0,U$
over $\mathbb Z$. On the boundary block each pure $X_i$ acts
trivially: $a,b,c$ permute the four two-torsion punctures,
their squares are trivial on them, and each of $(ab)^3,(cb)^3$
has trivial permutation. Thus for every odd $\ell$ the five
mod-$\ell$ linear actions are exactly
\[
        I_B\oplus U^2,\quad I_B\oplus L_0^2,\quad
        I_B\oplus U^2,\quad I_B\oplus(-I),\quad
        I_B\oplus(-I).
 \tag{FR.5b}
\]
This assertion is not obtained by extrapolating from characteristic
three. The integral part of \texttt{check\_ferrari\_cube.py}
constructs the integer homology matrices and verifies
$X_i(4I-\sum_{q\in K/C'}q)=4I-\sum_{q\in K/C'}q$ for
each $i$, together with the five exact pushforward identities
on the proper homology. These are integer equalities, and
division by four is legitimate at every odd prime.

The same proof as A.123, with three replaced by $\ell$, shows
that the kernel of the action on $H_1(C',\mathbb F_\ell)$
is pro-$\ell$. Consequently the procyclic action groups of
$X_1,X_2,X_3$ are pro-$\ell$, and those of $X_4,X_5$
are pro-$\{2,\ell\}$. This is only a statement about each
individual input; no assertion that their generated pair
group is pro-$\{2,\ell\}$ is needed.

At $\ell=2$ there is the exact identification
\[
                   C_2=C/O^2(C')=C^{(2)}.
 \tag{FR.5c}
\]
Indeed the left-hand side is an extension of a pro-two group
by $(C_2)^3$, hence is pro-two; every finite two-group quotient
of $C$ kills $O^2(C')$, proving its universal property.
Its abelianization is $(C_2)^3$, and its Frattini quotient is
the same group, since $C$ is generated by involutions. The
exact actions (FC.2a) show that each pure $X_i$ acts trivially
on this quotient. The Frattini automorphism kernel is pro-two
by the filtration argument of A.123, which applies to every
finitely generated pro-$p$ group. Thus all five input actions
are pro-two in this case.

Each $C_\ell$ is centreless. For odd $\ell$ this follows
as in A.123 from the faithful $(C_2)^3$ action on the rank-five
mod-$\ell$ homology and the centrelessness of the free pro-$\ell$
kernel. For two, the split presentation $C=K\rtimes\langle j\rangle$
gives
\[
                  C^{(2)}=K^{(2)}\rtimes\langle j\rangle,
\]
where $K^{(2)}$ is free pro-two of rank two. An element of
$K^{(2)}$ central in this extension is one, and an element
outside it induces $-I$ on its abelianization $\mathbb Z_2^2$,
so cannot be central. This proves the required centreless
assertion without reducing $-I$ modulo two.

\begin{statement}{A.128}{Proved}{one actual witness for the full relative pronilpotent quotient}
\end{statement}
Let $e\in\widehat{\mathbb Z}$ have $e_3=1$ and $e_p=0$
for every $p\ne3$, and set
\[
                 A_e=[[x^e,y^e],[y^e,z^e]],
                 \qquad z=(xy)^{-1}.
 \tag{FR.6a}
\]
For every $\ell\ne3$, (A.127) implies that each of the five
$X_i^e$ is one in $\operatorname{Aut}(C_\ell)$. Thus $A_e$
is in the kernel of every pentagon pair evaluation at that
prime. It is in each required peripheral conjugate of these
kernels as well: the peripheral $S_3$ permutes, up to inner
conjugacy, the three source peripheral elements $x,y,z$;
the original two input powers are one, and if any one of the
three powered entries in (FR.6a) is one, its double commutator
is one. Passing to the finite $\widehat B$-normal core therefore
preserves membership of $A_e$.

The same seed lies in $F''$, has $e\equiv4\pmod{12}$ and
$q_{12}(A_e)=5I$, and lies in the full original orbifold cores:
one source peripheral has order two there and $e$ is even.
These are precisely the seed conditions used in A.94--A.114.
For any finite list of finite action quotients at primes
$\ell\ne3$, intersect their peripheral normal cores with
the finite normal subgroups used to impose the A.94 conditions.
The single seed $A_e$ remains in this intersection and retains
its prescribed signed value. The relative R2 construction of
A.114 consequently supplies a genuine gentle element satisfying
that finite list, with $\Delta=1$ and the corresponding
relative pronilpotent conditions.

All these conditions are closed subsets of the compact group
of genuine gentle elements: finite evaluation equations are
clopen, $\Delta=1$ is clopen, and the relative pronilpotent
conditions are closed inverse-image conditions. The finite
intersection property, with all open normal subgroups needed
in A.94 included in the indexing family, therefore gives one
actual element $k_{\rm rel}$ such that
\[
 \Delta(k_{\rm rel})=1,\qquad
 f_{k_{\rm rel}}\in R_{\rm nil}(F''),
 \qquad
 f_{k_{\rm rel}}(X_i,X_j)=1
       \text{ in }\operatorname{Aut}(C_\ell)
       \quad(\ell\ne3)
 \tag{FR.6b}
\]
for every one of the five pairs. The prime-three assertion
then follows from A.125, using the same actual element.
This does not assert that the particular earlier A.94 witness
already satisfies all the new prime conditions; it constructs
a simultaneous strengthening by the same authorized R2 theorem.

The entire A.114 Eisenstein--Burau family may be imposed in
the same finite-intersection construction. At any finite list
of Burau congruence levels, the exponent $e$ reduces to an
integer $m$ divisible by every required prime-to-three
exponent and congruent to one modulo a sufficiently large
power of three. In particular $m$ is even, prime to three,
and congruent to four modulo twelve. These are the allowed
A.114 seed exponents. Its square-zero peripheral identities
therefore kill the same nonharmonic pair kernels away from
three; its harmonic and three-adic conditions follow from
$f\in R_{\rm nil}(F'')$ exactly as in A.114. Thus the fixed
seed $A_e$ retains all A.114 conditions simultaneously with
the new Ferrari conditions, rather than requiring two
unrelated witnesses.

For each prime, antisymmetry puts the pentagon defect inside
$C_\ell$, while (FR.6b) and A.125 make its inner action trivial.
The centreless result in A.127 proves
$\mathcal P(f_{k_{\rm rel}})=1$ in every $\Gamma_\ell$.
Finally
\[
       \bigcap_\ell O^\ell(C')=R_{\rm nil}(C'),
 \tag{FR.6c}
\]
because a finite nilpotent group is the direct product of its
Sylow subgroups. The natural map
$L/R_{\rm nil}(C')\to\prod_\ell L/O^\ell(C')$ is therefore
injective, as well as continuous with compact closed image.
Consequently
\[
 \boxed{\quad\mathcal P(f_{k_{\rm rel}})=1
           \text{ in }L/R_{\rm nil}(C'),\qquad
           \Delta(k_{\rm rel})=1.\quad}
 \tag{FR.6d}
\]
The quotient retains the full profinite forgetful base and
the full relative pronilpotent Ferrari kernel. This proves
the insufficiency of this entire family, conditional only
on the already specified R2 lifting input used in A.94--A.114.
It is not the full five-point pentagon, and it does not
produce an element of $\widehat{GT}$ with nonzero discrepancy.

\begin{statement}{A.129}{Proved}{an actual automorphism of the relative quotient}
\end{statement}
There is a further action-level conclusion that does not assume
kernel preservation merely from the pentagon product. Put
\[
 \Gamma_{\rm nil}=L/R_{\rm nil}(C'),\qquad
 C_{\rm nil}=C/R_{\rm nil}(C').
\]
The group $C_{\rm nil}$ is centreless. Indeed it embeds into
$\prod_\ell C_\ell$, and its projection to each $C_\ell$
is surjective. A central element therefore has trivial image
in every centreless $C_\ell$, and is one.

For the witness $k_{\rm rel}=(1,f)$ of A.128, define in
$\Gamma_{\rm nil}$
\[
 f_0=f(a^2,b^2)=f(X_1,X_2),\qquad
 e=f(z_3,c^2)=f(X_4,X_3),
\]
and assign
\[
       a\longmapsto a'=a,\qquad
       b\longmapsto b'=f_0^{-1}bf_0,\qquad
       c\longmapsto c'=e^{-1}ce.
 \tag{FR.7a}
\]
These are exactly the multiplier-one braid formulas in
the marking of FP.26.4. By A.128, $f_0$ acts trivially on
every $C_\ell$; the same is true of $e$, using antisymmetry
to reverse the pair $(X_3,X_4)$. They therefore centralize
$C_{\rm nil}$. Hence the three new generators have exactly
the same conjugation action on $C_{\rm nil}$ as the old
generators.

Under the fold $q:\Gamma_{\rm nil}\to\widehat B$, one has
$q(c)=q(a)$ and $q(z_3)=1$, so $q(e)=f(1,q(c)^2)=1$.
Consequently their folded images are
\[
 q(a')=\Phi_3(q(a)),\qquad
 q(b')=\Phi_3(q(b)),\qquad
 q(c')=\Phi_3(q(a)),
 \tag{FR.7b}
\]
where $\Phi_3$ is the existing genuine gentle automorphism
of the three-strand central quotient in this marking.
That source automorphism, not any five-point assumption,
is the input to (FR.7b).

Now evaluate any of the four defining relators
\[
 aba=bab,\qquad bcb=cbc,\qquad ac=ca,\qquad(abc)^4=1
\]
on $(a',b',c')$. Its fold is one by (FR.7b), and its
conjugation action on $C_{\rm nil}$ is one because the
three actions have not changed. The relator therefore
belongs to $Z(C_{\rm nil})=1$. The assignment (FR.7a)
thus defines a continuous homomorphism
$\widetilde\Phi:L\to\Gamma_{\rm nil}$ directly from
the profinite presentation.

For every $w\in C$, both $\widetilde\Phi(w)$ and its
original quotient image $\bar w$ belong to $C_{\rm nil}$,
and they induce the same inner automorphism there.
Centrelessness gives $\widetilde\Phi(w)=\bar w$.
In particular $R_{\rm nil}(C')$ is killed. Thus the map
descends to an endomorphism
\[
            \Phi_{\rm nil}:\Gamma_{\rm nil}
                         \longrightarrow\Gamma_{\rm nil}
\]
that fixes $C_{\rm nil}$ pointwise and induces the
automorphism $\Phi_3$ on the quotient $\widehat B$.
It is injective: an element of its kernel first maps
trivially to $\widehat B$, hence lies in the fixed kernel
$C_{\rm nil}$. It is surjective: lift a preimage under
$\Phi_3$ of the desired folded element, and then correct
the difference by an element of the fixed kernel.
A continuous bijection of compact Hausdorff groups is
a topological isomorphism. Therefore
\[
 \boxed{\ \Phi_{\rm nil}\in\operatorname{Aut}(\Gamma_{\rm nil}),
 \qquad\Phi_{\rm nil}|_{C_{\rm nil}}=\operatorname{id},
 \qquad\Delta(k_{\rm rel})=1.\ }
 \tag{FR.7c}
\]
This proves kernel stability and an actual relative action.
It does not promote $k_{\rm rel}$ to an element of the
full profinite Grothendieck--Teichm\"uller group.

\Needspace{10\baselineskip}
\subsection{A nonnilpotent mixed cover and its exact boundary pentagon relation}


\begin{statement}{A.130}{Proved}{an actual rank-129 mixed-prime homology quotient}
\end{statement}
Let $C=C_2\amalg C_2\amalg C_2$ as above and put
\[
 D=C''(C')^2,\qquad Q=C/D,
 \qquad V_D=D^{\rm ab}/3.
 \tag{FP.44.1}
\]
Since $C'$ is free of rank five, $C'/D=(C_2)^5$ and
$|Q|=8\cdot32=256$. The subgroup $D$ is free of rank
$1+32(5-1)=129$. All these subgroups are characteristic
in $C$ and hence invariant under the literal Ferrari action
of $L$. Conjugation therefore defines the actual continuous
homomorphism
\[
           \eta_D:L\longrightarrow
                    \operatorname{GL}(V_D)
                    =\operatorname{GL}_{129}(\mathbb F_3).
 \tag{FP.44.2}
\]
This target has an additional two-primary layer inside $C'$
before its relative three-primary step; it is not covered by
the relative-pronilpotent obstruction for $C'$ itself.

Here is an exact finite construction. The three-cube has
eight vertices and twelve unoriented edges. Over $\mathbb F_2$,
map the three involutions to $(e_i,t_i)$ in its edge-chain
module semidirect $(C_2)^3$. Their image consists of the
$256$ path chains with boundary equal to their endpoint minus
the basepoint. This realizes $Q$ with its prescribed marking.
Apply the exact involution-word automorphisms (FP.41.1) to
obtain the three automorphisms of $Q$.

The Schreier graph of $Q$ with its three involution generators
has $256$ vertices and $384$ unoriented edges. A spanning
tree leaves $129$ chords. Their based cycles give a literal
basis of $V_D$. For each automorphism, replace every oriented
edge by the path spelling the image of its involution, and
read the chord coordinates. This gives the matrices of
$a,b,c$, without a guessed representation or an unproved
descent. The certificate verifies every vertex action, every
edge endpoint, every cycle boundary, and the full intertwining
of edge chains and chord coordinates. It then verifies all
three braid/commutation relations and $(abc)^4=1$ on
$V_D$. The five $X_i$ have exact matrix orders
\[
                         6,6,6,4,4.
 \tag{FP.44.3}
\]
The construction is \texttt{check\_ferrari\_rank129.py}; its
output and the complete matrices and graph are retained in
\texttt{check\_ferrari\_rank129.out} and
\texttt{ferrari\_rank129.npz}.

\begin{statement}{A.131}{Proved}{the harmonic second derived image is not nilpotent}
\end{statement}
In the actual quotient (A.130), set $x=\eta_D(X_1)$,
$y=\eta_D(X_2)$ and
\[
 c=[x,y],\qquad w=[c,xcx^{-1}],\qquad
 u=w^6,\qquad v=(xy^{-1})u(xy^{-1})^{-1}.
 \tag{FP.45.1}
\]
The exact matrix orders are
\[
                  |w|=12,\qquad |u|=|v|=2,
                              \qquad |uv|=6.
 \tag{FP.45.2}
\]
Both $u,v$ belong to the second derived subgroup of
$\langle x,y\rangle$, because $w$ is a commutator of two
elements of its first derived subgroup. In a finite nilpotent
group the two-elements form the unique Sylow two-subgroup;
the product of two involutions must therefore have two-power
order. Equation (FP.45.2) is an exact contradiction.
Thus the pair-second-derived nilpotence test that excluded
all preceding Ferrari truncations fails for this target.

The source words, orders and matrix products are checked in
\path{check_rank129_pair_gate.py}. An independent audit
uses the complete free word, integer-64 arithmetic and the
known identities $x^{-1}=x^5$, $y^{-1}=y^5$, with no imported
inverse routine or order-search implementation. Its source
is \texttt{agent\_lifting/code/check\_rank129\_gate\_audit.py}.
This proves a property of an actual finite quotient. It is
neither a nonzero value of $\Delta$ on GT nor evidence that
an arbitrary finite parameter assignment lifts.

\begin{statement}{A.132}{Proved}{the obstruction is localized in two-branch boundary/proper coupling}
\end{statement}
The deck group $C'/D=(C_2)^5$ acts semisimply on $V_D$.
The trivial character has multiplicity five, and each of its
$31$ nontrivial characters has multiplicity four. The exact
simultaneous eigenspaces are computed from deck translations
on the same Schreier graph, with no new homological marking.

There are $64$ cusp loops in the $D$ cover. Writing
$w_0=t_1t_0t_2$, they are the cycles $q w_0^4q^{-1}$,
one for each of the $64$ cosets modulo the order-four image
of $w_0$ in $Q$. Their span has rank $63$ and is invariant
under $a,b,c$. The quotient, the proper homology, has rank
$66$. On each nontrivial deck character the possible
dimensions are
\[
\begin{array}{c|c|c|c}
 \text{branch points among the four cusps}&
 \text{number of characters}&\text{boundary}&\text{proper}\\\hline
 0&3&4&0\\
 2&24&2&2\\
 4&4&0&4.
\end{array}
 \tag{FP.46.1}
\]
These counts are verified by evaluation on the four actual
cusp words; equivalently they follow from Riemann--Hurwitz
for the corresponding double covers. The full $L$ character
orbits have sizes $24,3,3,1$ among the nontrivial characters:
the final singleton is one of the fully branched covers.

The involution witness (A.131) is supported entirely on the
$24$ two-branch characters. It vanishes on both fully branched
pieces and on the unramified pieces. On the proper rank-$66$
quotient the product of its two involutions has order two;
on the open homology it has order six. Thus the new
non-nilpotence uses the extension of proper homology by
boundary homology and is lost by passing to the compact
Jacobian representation alone.

In the exact binary character basis retained by the code, the
two harmonic orbits
\[
          (1,7,18,20),\qquad (8,14,27,29)
 \tag{FP.46.2}
\]
each have dimension $16$, with boundary and proper dimensions
eight each. On either orbit the two involutions in (A.131)
have product of exact order three. The first orbit's literal
harmonic image is an abelian extension with
\[
 |Q_{\rm lin}|=4608,\quad
 |Q_{\rm lin}'|=384,\quad |Q_{\rm lin}''|=64,
 \quad |Q_{\rm lin}'''|=4,
 \qquad U\simeq\mathbb F_3^4,
\]
so its full order is $4608\cdot3^4=373248$.
Here $Q_{\rm lin}$ retains the separate boundary and proper
blocks, while $U$ is the exact image of the off-diagonal
extension module. Schreier loops in the finite base compute
the entire four-dimensional kernel, not a chosen subspace.

The exact decomposition, branch tests, cusp spans and small
module are respectively
\texttt{check\_rank129\_deck\_blocks.py},
\texttt{check\_rank129\_prym.py},
\texttt{check\_rank129\_boundary.py}, and
\texttt{check\_rank129\_harmonic\_module.py}, with their
outputs and complete matrix data retained alongside them.
These marked data identify a new coupling; they do not yet
identify its character with $\Delta$.

\begin{statement}{A.133}{Proved}{Proved necessary matrix identity, subsequently proved automatic on gentle elements}
\end{statement}
Fix character $1$ in (FP.46.2). All five evaluation-pair
permutation images on its character orbit are abelian.
Consequently every $f\in\widehat F'$ preserves this individual
four-dimensional character space in each of the five factors.
For $\Phi=(1,f)\in\mathcal H$, one has the stronger elementary
fact $f\in\gamma_3(\widehat F)$: in the class-two quotient
write $f=[x,y]^e$; the multiplier-one three-cycle gives
$3e=0$ in $\widehat{\mathbb Z}$, hence $e=0$.

The literal boundary pair images have lower-central order
series $8,2,1$ or $4,1$. Thus all five boundary $f$-blocks
are the identity for every actual gentle element. The proper
pair images for $(X_3,X_4)$ and $(X_5,X_1)$ have order $24$
with central derived subgroup of order two, and that for
$(X_4,X_5)$ has order eight with the same property. Their
proper $f$-blocks are also the identity. Write the five
four-dimensional blocks, in pentagon order, as
\[
 \begin{pmatrix}I&C_{34}\\0&I\end{pmatrix},\quad
 \begin{pmatrix}I&C_{51}\\0&I\end{pmatrix},\quad
 \begin{pmatrix}I&C_{23}\\0&P_{23}\end{pmatrix},\quad
 \begin{pmatrix}I&C_{45}\\0&I\end{pmatrix},\quad
 \begin{pmatrix}I&C_{12}\\0&P_{12}\end{pmatrix}.
\]
All entries are over $\mathbb F_3$; each $C_{ij}$ and $P_{ij}$
is $2\times2$ and is the evaluation of the \emph{same} source
word $f$ in its specified pair. The full pentagon implies
the exact equations
\[
 \boxed{\quad P_{23}P_{12}=I,\qquad
 C_{34}+C_{51}+(C_{23}+C_{45})P_{12}+C_{12}=0.\quad}
 \tag{FP.47.1}
\]
This follows by literal block multiplication. The two harmonic
proper pair groups have derived series
$2304,192,64,4,1$; their values must not be replaced by
independent parameters. All the group orders and class-two
assertions are checked in
\texttt{check\_rank129\_pentagon\_blocks.py} and its output.

Equation (FP.47.1) is a literal boundary/proper matrix identity
in an actual five-point representation. The subsequent exact
normalization in (A.144)--(A.145) proves that it holds for
every multiplier-one gentle element. It therefore supplies
no new comparison between $\kappa_d$ and $\epsilon_2(W)$.
The separate affine extension in (A.146) retains information
that this homology action kills. No implication from a
surviving finite matrix assignment to a genuine GT lift is
asserted.

\subsection{The explicit Prym geometry and the exact limit of the linear cover}


\begin{statement}{A.134}{Proved}{character multiplicities and compact dimensions}
\citep[Sections 2.9 and 2.15]{LitNakamuraTsunogaiYasuda2010}
\end{statement}
Put $D=C''(C')^2$. Since $C'$ is free of rank five,
$C'/D\simeq(C_2)^5$, $[C:D]=256$, and $D$ is free
of rank $1+32(5-1)=129$. Geometrically, A.117 identifies
$C'$ with the fundamental group of
$X=E\setminus E[2]$. Its maximal mod-two abelian cover
$Y\to X$ has degree thirty-two. Every puncture inertia
has order two in the deck group $A=(C_2)^5$, so the
compactification of $Y$ has sixty-four punctures removed.
The Euler characteristic is $32(-4)=-128$, giving genus
thirty-three. Thus its compact and boundary homology
dimensions are sixty-six and sixty-three respectively.

Over $\mathbb F_3$, the $A$-action is semisimple. The
trivial character has homology dimension five; each
of the thirty-one nontrivial characters has dimension
four, by the rank-one local-system Euler characteristic
on $X$. If $r_\chi$ is the number of the four puncture
inertias on which $\chi$ is trivial, then its boundary
multiplicity is $r_\chi$ for nontrivial $\chi$ and three
for the trivial character. Indeed the boundary module
is the sum of the four permutation modules induced from
the inertia groups, modulo the one global boundary
relation. Hence the nontrivial compact multiplicity is
$4-r_\chi$.

The four boundary signs have product one. Restriction
to the three-dimensional boundary subspace has a
two-dimensional kernel, coming from elliptic homology.
Consequently there are exactly
\[
\begin{array}{c|c|c|c}
\text{nontrivial character type}&\text{number}&
 \text{proper dimension}&\text{boundary dimension}\\ \hline
\text{unramified at all four punctures}&3&0&4\\
\text{ramified at exactly two punctures}&24&2&2\\
\text{ramified at all four punctures}&4&4&0.
\end{array}
 \tag{RP.1a}
\]
Including the trivial character gives
$2+24\cdot2+4\cdot4=66$ and
$3+3\cdot4+24\cdot2=63$. These are geometric
character counts, not assumptions about the finite
image of the braid action.

\begin{statement}{A.135}{Proved}{the two-branch covers and their actual Prym curves}
\end{statement}
Let the two selected branch points on $E$ be $Q_1,Q_2$,
both of order two. Translating by $-Q_1$ makes the
branch divisor $O+T$, where $T=Q_2-Q_1\ne O$ is of
order two. The four square roots of this degree-two
line bundle are
\[
                \mathcal O_E(P),\qquad 2P=T.
 \tag{RP.2a}
\]
They give precisely the four double covers with this
fixed branch divisor. Over an algebraically closed
field of characteristic different from two, there is
no additional scalar-twist distinction. A choice of
$P$ is part of the marking and must not be omitted.

Translate once more by $-P$. The branch points become
$-P,P$, and write the translated elliptic curve as
\[
 E:y^2=\prod_{i=1}^3(x-e_i),\qquad a=x(P).
\]
The corresponding genus-two curve has the explicit model
\[
 C_a:y^2=\prod_{i=1}^3(w^2+a-e_i),\qquad x=w^2+a.
 \tag{RP.2b}
\]
The degree-two map to $E$ is the quotient by
$\sigma(w,y)=(-w,y)$. Its two finite fixed points
are exactly the two points over $x=a$. Composing
the quotient with translation by $P$ returns the
original marked branch divisor $O+T$.

The hyperelliptic involution is
$\iota(w,y)=(w,-y)$. The other elliptic quotient,
by $\tau=\iota\sigma$, is
\[
 E'_a:v^2=u\prod_{i=1}^3(u+a-e_i),\qquad
                       u=w^2,\quad v=wy.
 \tag{RP.2c}
\]
It is ramified at the two points at infinity of $C_a$.
Its pullback map $\operatorname{Jac}(E'_a)\to
\operatorname{Jac}(C_a)$ is injective. For a degree-two
ramified map of curves, a nontrivial element of this
kernel would be a line bundle of order two whose
associated connected etale double cover is dominated
by the given double cover; equality of the degrees
would make the original map etale, a contradiction.
The image is anti-invariant under $\sigma$, since
$\tau$ fixes it and $\iota$ acts as minus one on the
Jacobian. It therefore equals the connected Prym
kernel, by dimension. In particular
\[
           \operatorname{Prym}(C_a/E)\simeq E'_a
 \tag{RP.2d}
\]
as elliptic curves. The polarization inherited from
the genus-two Jacobian is twice the principal
elliptic polarization; it does not change this
identification of the underlying elliptic curve.

\begin{statement}{A.136}{Proved}{the signed quadratic Legendre transformation}
\end{statement}
Normalize the chosen two-torsion point as $T=(0,0)$ on
\[
       E_\lambda:y^2=x(x-1)(x-\lambda),\qquad
       \lambda\ne0,1.
\]
For a point $P$ with $2P=T$, its $x$-coordinate satisfies
$a^2=\lambda$. This follows directly from the tangent
line through $P$ and $T$: substituting a line $y=mx$
into the cubic shows that its nonzero intersection
roots have product $\lambda$, and tangency makes both
equal to $a$. Conversely the two points
\[
        P=(a,\pm i\,a(a-1)),\qquad a^2=\lambda,
 \tag{RP.3a}
\]
have $2P=T$. Thus the four choices consist of the two
values of $a$ and the two signs in (RP.3a).

In (RP.2c), put $x=u+a$. Its four branch points on
the $x$-line are $a,0,1,a^2$. The Mobius transformation
and rescaling
\[
       X=\frac{(1-a)x}{x-a},\qquad
       Y=\frac{v(X+a-1)^2}{a(a-1)}
 \tag{RP.3b}
\]
give exactly
\[
       \boxed{\ E'_a:Y^2=a(a-1)X(X-1)(X+a).\ }
 \tag{RP.3c}
\]
For verification, $x=aX/(X+a-1)$ and
\[
 (x-a)x(x-1)(x-a^2)
 =\frac{a^3(a-1)^3 X(X-1)(X+a)}{(X+a-1)^4}.
\]
Thus the new Legendre parameter is $\mu=-a$, together
with the quadratic twist $a(a-1)$. The unsigned formula
$\mu=\pm\sqrt\lambda$ alone loses a nontrivial central
monodromy character.

The complete geometric diagram is over
$S=\mathbb P^1_a\setminus\{0,1,-1,\infty\}$:
the original elliptic homology is pulled back by
$\lambda=a^2$, while the Prym homology is pulled back
by $\mu=-a$ and tensored with the rank-one quadratic
local system of $\sqrt{a(a-1)}$. Its signs around
$0,1,-1,\infty$ are respectively $-1,-1,+1,+1$.
Equivalently, if $T_0,T_1,T_\infty$ are the usual
local Legendre monodromy classes, then the local
Prym classes are
\[
\begin{array}{c|cccc}
 a&0&-1&1&\infty\\ \hline
 H_1(E'_a)&-T_0&T_1&-I&T_\infty.
\end{array}
 \tag{RP.3d}
\]
Here $T_0,T_1$ are nontrivial unipotents and
$T_\infty$ has trace $-2$. The table gives conjugacy
classes of positively oriented local loops. It does
not choose four independent conjugating paths or
identify them with a pre-existing Schreier basis.
For any chosen basepoint and paths, (RP.3b)--(RP.3c)
give an exact simultaneous identification of local
systems, up to the single resulting basis change.

The puncture $a=1$, which the unsigned map $\mu=-a$
would fill, has monodromy $-I$ on the Prym. This is
the precise extra sign retained by the two-branch
blocks. The comparison is not a generic isogeny
between $E_{a^2}$ and $E'_{a}$: at $a=1$ the former
has potentially multiplicative reduction, whereas
the latter is potentially good with the displayed
finite monodromy. An isogeny would preserve that
distinction.

\begin{statement}{A.137}{Proved}{the marked punctures and the intrinsic rank-three extension}
\end{statement}
The open genus-two cover corresponding to a two-branch
character is obtained by removing the preimages of all
four original points of $E[2]$, not only the two branch
points. After the translation used in A.135, these four
points are the four solutions of $2R=T$: they are
$\pm P$, with $x=a$, and $\pm R$, with $x=-a$.
Thus the six removed points on the model (RP.2b) are
\[
 (0,\pm i a(a-1)),\qquad
 (\varepsilon\sqrt{-2a},\delta i a(a+1))
       \quad(\varepsilon,\delta\in\{\pm1\}).
 \tag{RP.5a}
\]
All these points are distinct over
$S=\mathbb P^1\setminus\{0,1,-1,\infty\}$.
Write $C_a^\circ$ for this punctured curve.

The four-dimensional deck-character block is
$H_1(C_a^\circ,\mathbb F_3)^-$ for
$\sigma(w,y)=(-w,y)$. Let $d_+$ and $d_-$ be the
differences of small positively oriented loops about
the two $w$-signs in (RP.5a), with the $y$-sign fixed.
They form its two-dimensional boundary subspace.
The hyperelliptic involution $\iota(w,y)=(w,-y)$
interchanges $d_+$ and $d_-$. It acts as minus one
on all proper genus-two homology. Consequently
\[
 H_1(C_a^\circ,\mathbb F_3)^-
       =L_{\rm bdry}\oplus V_3,
 \tag{RP.5b}
\]
where $L_{\rm bdry}$ is the $+1$ eigenspace of
$\iota$, generated by $d_++d_-$, and $V_3$ is its
$-1$ eigenspace. The latter has proper dimension
two and boundary dimension one, the boundary line
being generated by $d_+-d_-$. The hyperelliptic
involution is canonical for a smooth genus-two
curve; thus this splitting is respected by
parallel transport in the stabilizer of the
chosen double-cover character.

Under the quotient by $\tau=\iota\sigma$, the
removed points map to three points on $E'_a$:
the chosen origin $O$ and a pair $Q_a,-Q_a$.
Using the Legendre coordinates (RP.3b),
\[
 X(Q_a)=\frac{1-a}{2},\qquad
 Y(Q_a)=\frac{a^2-1}{4}\sqrt{2a}.
 \tag{RP.5c}
\]
The other choice of the square root gives $-Q_a$.
The involution induced by $\sigma$ is $[-1]$ on
$E'_a$. Pullback under the degree-two quotient,
whose degree is invertible modulo three, identifies
\[
 V_3\simeq H_1(E'_a\setminus\{O,Q_a,-Q_a\},
                         \mathbb F_3)^{[-1]=-1}.
 \tag{RP.5d}
\]
This remains true though the quotient ramifies at
two interior points: homology of a finite quotient
in coefficients of characteristic prime to the
group order is the invariant homology. On the
boundary, pullback sends the difference of the
$Q_a,-Q_a$ loops to a nonzero scalar multiple of
$d_+-d_-$.

The boundary line in $V_3$ is the quadratic local
system $\chi_{2a}$: a loop exchanging the two
square roots in (RP.5c) exchanges $Q_a,-Q_a$ and
negates their boundary difference. Its signs
around $0,1,-1,\infty$ are $-1,+1,+1,-1$.
Tensoring the whole extension by $\chi_{2a}$
makes this boundary line constant. The corresponding
proper elliptic curve and marked rational section are
\[
 \boxed{\quad
 E''_a:Z^2=\frac{a-1}{2}X(X-1)(X+a),\qquad
 Q''_a=\left(\frac{1-a}{2},\frac{a^2-1}{4}\right).
 \quad}
 \tag{RP.5e}
\]
Over the quadratic cover adjoining $\sqrt{2a}$,
this is the change $Z=Y/\sqrt{2a}$. Thus it is an
exact quadratic-twist identification, not a choice
of unrelated isomorphic fibres.

Dually, the boundary/proper extension is the
relative homology extension of an arc from
$-Q''_a$ to $Q''_a$. If a path and homology basis
are chosen at one basepoint, the difference
between an arc and its continuation is an
integral elliptic cycle; these differences give
the extension cocycle. Its Abel--Jacobi section
is $Q''_a-(-Q''_a)=2Q''_a$. This is the precise
marked-point extension meant here by its Kummer
class. Neither its splitting nor the value of
its reduction modulo three is inferred merely
from the proper elliptic representation.

On the original elliptic base, the hyperelliptic
involution in this construction is $R\mapsto T-R$.
It is therefore the lift of the involution
$[-1]+T$ that fixes the chosen square-root line
bundle $\mathcal O(P)$. In the free-product
description this is one of the three marked
involutions $t_i$; multiplying the lift by the
double-cover deck involution exchanges its
$1$- and $3$-dimensional eigenspace labels but
does not change the splitting. The independent
finite certificate
\texttt{check\_rank129\_open\_prym.py} verifies
this $1+3$ decomposition for each of the twenty-four
two-branch characters and verifies that the
$a,b,c$ actions preserve it between characters.
An explicit choice of character and paths is
still needed to equate an individual matrix
entry with the cocycle of $2Q''_a$.

\begin{statement}{A.138}{Proved}{the section is generically non-torsion}
\end{statement}
Equation (RP.5e) has the integral Weierstrass model
\[
 \mathcal E_a:
 y^2=x^3+2(a-1)^2x^2-4a(a-1)^2x,
 \qquad
 Q_a^{\rm int}=(-(a-1)^2,(a-1)^2(a+1)),
 \tag{RP.6a}
\]
obtained by $x=2(a-1)X$, $y=4(a-1)Z$.
At $a=2$ this specializes to
\[
 E:y^2=x^3+2x^2-8x,\qquad Q=(-1,3),\qquad
                    2Q=(9/4,15/8).
\]
Its discriminant is $2^{12}3^2$, so five and seven
are good primes. Direct enumeration and the group
law give
\[
\begin{array}{c|cc}
 p&\#E(\mathbb F_p)&\operatorname{ord}(\bar Q)\\ \hline
 5&8&4\\
 7&12&6.
\end{array}
 \tag{RP.6b}
\]
If $Q$ were torsion, its order would be divisible
by three from its reduction at seven. Its
three-primary part injects at the good prime
five, a contradiction because that reduced
elliptic curve has order eight. Prime-to-$p$
torsion injectivity here follows directly from
the formal-group filtration of the good
reduction kernel, whose finite successive
quotients are $p$-groups.

Thus $Q$ and $2Q$ are non-torsion. A generic
torsion section would remain torsion at this
smooth specialization, so (RP.5e) is generically
non-torsion as well. The exact arithmetic is
recorded independently in
\texttt{check\_prym\_section.py} and its output.
This argument does not assert that the section
is geometrically indivisible by three; that
would be a separate monodromy or height claim.

\begin{statement}{A.139}{Proved}{exact CM specializations and their comparison scope}
\end{statement}
The new proper elliptic family has invariant
\[
             j(E''_a)=
       \frac{256(a^2+a+1)^3}{a^2(a+1)^2}=j(-a).
\]
If $a^2+a+1=0$, then (RP.6a), after $X=x-2a$,
becomes
\[
                 y^2=X^3+8,\qquad Q=(a,3).
 \tag{RP.7a}
\]
The point is the order-three CM translate of
$(1,3)$. If $a=-2$, after $X=x+6$ the family becomes
\[
                 y^2=X^3-36X,\qquad Q=(-3,-9).
 \tag{RP.7b}
\]
Both displayed points are non-torsion. For the
first curve, the reduction of $(1,3)$ has order
six at five while the whole reduced curve has
order sixteen at thirteen. For the second,
the point has order six at eleven while the
whole reduced curve has order eight at five.
The same prime-to-$p$ argument proves both
claims; the certificate in A.138 checks them.

Relative to the literal R4 models
$E_3:y^2=x^3+1$ and $E_4:y^2=x^3-x$,
(RP.7a) is the quadratic twist by two and
(RP.7b) is the quadratic twist by six.
The exact isomorphisms over the corresponding
quadratic fields are
\[
 (X,y)\longmapsto\left(\frac X2,
                            \frac y{2\sqrt2}\right)
          \quad\hbox{to }E_3,
 \qquad
 (X,y)\longmapsto\left(\frac X6,
                            \frac y{6\sqrt6}\right)
          \quad\hbox{to }E_4.
 \tag{RP.7c}
\]
They carry the two marked points to
\[
 \left(\frac a2,\frac3{2\sqrt2}\right)\in E_3,
 \qquad
 \left(-\frac12,-\frac3{2\sqrt6}\right)\in E_4.
\]
The Tate-module representations before these
isomorphisms differ from those of the R4
curves by the quadratic characters of two
and six. Both characters restrict trivially
to $G_{\mathbb Q^{\rm ab}}$, since that group
fixes the quadratic fields. Thus the proper
torsion representations agree there after
the explicit identifications (RP.7c).
The displayed marked points remain
non-torsion; their division-point extensions
and path cocycles are additional data.
Triviality of the quadratic twists on the
arithmetic subgroup is not a construction
of a common marked path for all of
$\widehat{GT}$.

The corresponding original base curves
$E_{a^2}$ have respectively $j=2048/3$ and
$j=35152/9$. In particular the two elliptic
slots in this construction must be kept
distinct. The CM property of the Prym at
these values does not replace the non-torsion
marked section by the old torsion-frame data.

The characteristic rank-129 cover, its
Prym section, and its finite braid action
are genuine. A comparison of its section's
path cocycle with $\kappa_d$ would still
require a marked full-$\widehat{GT}$ theorem;
the arithmetic specialization does not
provide one. Moreover A.140 proves that
the linear pentagon in this representation
is automatic on the gentle group. Thus
it cannot itself supply the missing theorem,
even though its harmonic pair image has
non-nilpotent second derived subgroup.

\begin{statement}{A.140}{Proved}{the entire linear cover is blind to the gentle pentagon}
\end{statement}
The rank-129 construction passes the pair-second-derived
nilpotence test but ultimately contributes no new pentagon
equation. The following source identities explain why.
Write $u=\xi_1$, $v=\xi_2$, $s=aba$, and
$X_4=s^2=z_3$. The exact Ferrari Nielsen formulas give
\[
 c=ua,\qquad a(u)=u,\qquad
 s(u)=v^{-1},\qquad s(v)=vuv^{-1},\qquad
 cbc=uv^{-1}s.
\]
Consequently
\[
 X_3X_1^{-1}=u^2\in C',\qquad
 X_5X_4^{-1}=uv^{-1}u^{-1}v^{-1}\in C'.
 \tag{RP.8a}
\]
The second membership also follows from its exponent sums
$(0,-2)$ and $C'=K'K^2$. Since $X_4$ is central in
the actual three-strand section, it commutes with
$X_1,X_2$. Thus in $L/C'$ the five inputs are
$(X_1,X_2,X_1,X_4,X_4)$, with $X_4$ central among
them. Every antisymmetric commutator word therefore
has pentagon one in this quotient, and
\[
                         \mathcal P(f)\in C'.
 \tag{RP.8b}
\]

Consider the twenty-four rank-three blocks of A.137.
On the block for a character $\chi:C'/D\to\{\pm1\}$,
every $g\in C'$ acts as the scalar $\chi(g)$, both
on the full block and its boundary line. Tensoring
the block with the dual of that line cancels this
scalar. On the direct sum, with its genuine
permutation of the characters, this is a representation
of $L$ which kills $C'$, hence factors through $L/C'$.
Its pentagon is therefore identically one by (RP.8a).

The line normalization does not preserve an arbitrary
commutator-word value; this qualification is necessary.
It does preserve every actual multiplier-one gentle
$f$. Indeed such an $f$ belongs to $\gamma_3(F)$:
in the class-two free quotient write $f=[x,y]^e$;
the three-cycle relation gives $3e=0$ in
$\widehat{\mathbb Z}$, and hence $e=0$.
The exact signed boundary-line pair images on all
twenty-four characters have nilpotency class at most
two. The two harmonic pair images have lower-central
orders $8,2,1$; the other three have orders $16,2,1$.
Thus every actual $f$ has boundary-line value one
in all five factors. The normalized and original
rank-three pentagon values coincide.

These boundary assertions are finite signed-permutation
identities, not estimates. They are checked on the
whole twenty-four-line representation, including all
its pair orbits. Reduction modulo three is faithful
on signed permutation matrices, so the same identities
hold in the integral boundary lattice and after any
odd coefficient change. An independent audit also
checks directly that each pair commutator commutes
with both generators on all twenty-four lines.

It follows that $\mathcal P(f)$ acts trivially on
each of the twenty-four rank-three blocks. By
(RP.8b), that action is the scalar
$\chi(\mathcal P(f))$. The twenty-four two-branch
characters span $(C'/D)^*$: differences of two
characters with the same boundary restriction
give the two unramified handle characters, and
the six weight-two boundary patterns span the
three-dimensional boundary dual. Therefore
\[
                \boxed{\ \mathcal P(f)\in D
                    =C''(C')^2\quad
                    \text{for every genuine gentle }(1,f).\ }
 \tag{RP.8c}
\]
Every element of $D$ acts trivially by conjugation
on $D^{\rm ab}$. Hence the pentagon is one on
$H_1(D,R)$ for every coefficient ring $R$, not
only on the tested $\mathbb F_3$ quotient.
The apparent boundary/proper equation in A.133
is therefore an identity on the actual gentle
locus. The non-nilpotent harmonic image of
A.131 remains correct, but is insufficient
to make the five-point product detect a new
condition.

This conclusion must be kept separate from
the affine extension
$C/(D'D^3)$. The linear conjugation action
on $H_1(D,\mathbb F_3)$ kills the entire
normal subgroup $D$; the affine extension
retains its class in $D/(D'D^3)$. Equation
(RP.8c) does not kill that class and therefore
does not settle the pentagon in the affine
extension.

\subsection{The actual affine quotient and the translation discarded by homology}


\begin{statement}{A.141}{Proved}{the actual quotient as a path-chain extension}
\end{statement}
Let $C=C_2\amalg C_2\amalg C_2$, generated by
$(t_0,t_1,t_2)$, and put
\[
 D=C''(C')^2,\qquad Q=C/D,\qquad |Q|=256.
\]
The Cayley graph of $Q$ with these three involutions
has $256$ vertices and $384$ unoriented edges.
Orient each edge arbitrarily, with a negative sign
for its reversed orientation. Over $\mathbb F_3$,
write $C_1$ for the edge-chain space and
$I_Q=\ker(\mathbb F_3[Q]\to\mathbb F_3)$ for the
augmentation ideal. The boundary map gives an exact
sequence of $Q$-modules
\[
 0\longrightarrow M=H_1(D,\mathbb F_3)
 \longrightarrow C_1\stackrel{\partial}{\longrightarrow}
 I_Q\longrightarrow0,
 \qquad \dim(M,C_1,I_Q)=(129,384,255).
 \tag{FA.1a}
\]
The dimension follows either from the graph or the
index-$32$ Schreier formula inside the free rank-five
group $C'$. We first use the discrete free product and its
finite covering graph, and then complete the resulting
free finite-index subgroup. Completion preserves its
finite-coefficient abelianization. Thus this graph
computes the stated homology of the profinite subgroup
$D$ with these finite coefficients.

Let $q\cdot c$ denote deck translation of an edge chain.
Inside $C_1\rtimes Q$, define
\[
 \mathcal E=\{(c,q):\partial c=q-1\},\qquad
 (c,q)(d,r)=(c+q\cdot d,qr).
 \tag{FA.1b}
\]
This is a subgroup, since
$(q-1)+q(r-1)=qr-1$. Sending $t_i$ to its one-edge
path from $1$ to $t_i$ defines a homomorphism
$C\to\mathcal E$: the relation $t_i^2=1$ holds
because the translated second edge is the negative
of the first. The image of a word is exactly its
path chain and endpoint.

On $D$ the endpoint is one and the map is the usual
homology map onto $M$, with kernel $D'D^3$. For each
$q\in Q$, choose a path from $1$ to $q$. Its chain,
plus arbitrary cycles from $M$, gives every chain
with boundary $q-1$. Thus the homomorphism is onto
and
\[
       \boxed{\ C/(D'D^3)\simeq\mathcal E,\qquad
                    |\mathcal E|=256\cdot3^{129}.\ }
 \tag{FA.1c}
\]
This is an actual finite group quotient with its
source marking, not a linear approximation to it.

Since $|Q|$ is invertible in $\mathbb F_3$, averaging
an arbitrary linear splitting of $\partial$ gives
a $Q$-equivariant splitting $s:I_Q\to C_1$.
The maps $q\mapsto(s(q-1),q)$ form a group section
of (FA.1b). Consequently $\mathcal E\simeq M\rtimes Q$,
although the path-chain model avoids choosing that
splitting and retains the exact based marking.

\begin{statement}{A.142}{Proved}{centrelessness and the information lost by the linear action}
\end{statement}
The edge module is
\[
 C_1=\bigoplus_{i=0}^2
       \operatorname{Ind}_{\langle t_i\rangle}^{Q}
                      \operatorname{sgn}.
\]
It has no $Q$-invariants: on a fixed edge, the
stabilizing involution reverses orientation, and
$-1\ne1$ in $\mathbb F_3$. Hence
\[
                              M^Q=0.
 \tag{FA.2a}
\]
The action of $Q$ on $M$ is faithful. To prove this
without a numerical faithfulness test, use the same
integral graph cycle lattice $M_{\mathbb Z}$, whose
profinite completion is $D^{\rm ab}$. For $1\ne q\in Q$, the
vertex action has trace zero. An edge fixed by $q$
is reversed, so its trace contribution is minus one.
The homology trace is therefore
\[
          \operatorname{tr}(q\mid M_{\mathbb Z}\otimes\mathbb Q)
                       =1-\#\{q\text{-fixed edges}\}\leq1,
\]
which is strictly below its dimension $129$.
If $q$ acted as the identity modulo three, its
integral action on $\mathbb Z_3$ homology would lie
in the pro-three kernel of
$\operatorname{GL}_{129}(\mathbb Z_3)\to
\operatorname{GL}_{129}(\mathbb F_3)$. But its order
is a power of two, so that action would be the
identity, contradicting the trace. This proves
faithfulness over $\mathbb F_3$.

A central element of $\mathcal E$ consequently
has trivial $Q$ coordinate, since it centralizes
$M$; its remaining coordinate belongs to $M^Q$.
Thus
\[
                         Z(\mathcal E)=1.
 \tag{FA.2b}
\]
The source action of $L$ on $C$ preserves $D'D^3$
and therefore induces actual automorphisms of
$\mathcal E$. On chains it is obtained by replacing
each oriented edge by the literal path spelling
the image of its involution. It is semilinear
with respect to the induced permutation of $Q$
and intertwines the boundary maps. Its restriction
to $M$ is the old rank-129 homology action.

For $m\in M$, however, inner conjugation in the
full path-chain group is
\[
       (c,q)\longmapsto(c+(1-q)m,q).
 \tag{FA.2c}
\]
It is the identity on cycles, whose endpoint is
one, but need not be the identity on chains.
It is the identity everywhere if and only if
$m\in M^Q$, hence if and only if $m=0$.
In any splitting of (FA.1a), this is an upper
triangular shear from the augmentation-ideal
part to the homology part. This is exactly
the information discarded by retaining only
the linear action on $M$.

\begin{statement}{A.143}{Proved}{Proved formulation; the new obstruction is not evaluated here}
\end{statement}
For a genuine gentle multiplier-one pair, A.140
proves $\mathcal P(f)\in D$. Its image in the
new quotient is therefore the definite vector
\[
           m_f=[\mathcal P(f)]\in D/(D'D^3)=M.
 \tag{FA.3a}
\]
The full edge-chain action of its pentagon is
the shear (FA.2c) for this vector. Thus
\[
 \mathcal P(f)=1\text{ in }C/(D'D^3)
 \quad\Longleftrightarrow\quad m_f=0
 \quad\Longleftrightarrow\quad
 \mathcal P(f)\text{ acts trivially on }C_1.
 \tag{FA.3b}
\]
The old rank-129 action is always trivial on
this element and cannot decide any of these
equivalent conditions. The exact sequence
(FA.1a) gives a concrete finite target in
which the remaining vector can be computed.
No calculation in this paragraph proves
$m_f=0$ for all gentle elements, identifies
a coordinate of $m_f$ with $\Delta$, or
promotes a finite solution to full GT.

Geometrically, this extension retains based
paths between the $256$ deck sheets, in
addition to their closed cycles. The genus-two
Prym calculations describe constituents of
the cycle module $M$; they do not supply
these based paths or their extension class.
That distinction is why their linear
boundary/proper coupling can be nontrivial
while the linear pentagon remains automatic.

\subsection{Exact normalization, universal linear cancellation and retained affine translations}


\begin{statement}{A.144}{Proved}{the exact open-Prym normalization and its automatic pentagon}
\end{statement}
Each of the $24$ two-branch four-dimensional character
spaces in (A.132) has a canonical splitting of dimensions
$1+3$. Exactly one of the three deck involutions $t_i$
stabilizes that character with eigenspace dimensions one
and three. The smaller eigenspace is a boundary line; the
larger contains one boundary line and the proper
two-dimensional homology. Computing these deck involutions
on the original graph proves that every $a,b,c$ transition
preserves the two summands, including their character
permutations. Thus the rank-$96$ piece splits as a
rank-$24$ boundary representation plus a rank-$72$ open-Prym
representation.

In the latter, each character transition has a $3\times3$
block $\bigl(\begin{smallmatrix}\alpha&v\\0&M\end{smallmatrix}\bigr)$,
with $\alpha=\pm1$. Tensoring by the inverse boundary line
divides the entire block by $\alpha$. This is a genuine
representation operation on the character bundle, not an
independent rescaling of selected generator matrices. A
single consistent change of the proper basis on each of the
$24$ characters then makes every proper transition have
determinant one. Its existence and consistency around all
generator loops are checked exactly. The resulting
homomorphism is a block-permutation representation with
blocks in $\operatorname{ASL}_2(\mathbb F_3)$, in the
row-vector convention
\[
                  \begin{pmatrix}1&v\\0&M\end{pmatrix},
                         \qquad \det M=1.
\]
Its full $L$ permutation image has order $192$; the five
pure $X_i$ have an elementary abelian permutation image of
order eight. All defining relations of $L$ hold after the
tensor operation and the fixed basis change.

For the untwisted boundary line on \emph{all} $24$ characters,
the five pair groups have lower-central order series, in
pentagon order,
\[
        (16,2,1),\ (16,2,1),\ (8,2,1),\
        (16,2,1),\ (8,2,1).
\]
Since every actual gentle $f$ belongs to $\gamma_3F$ by
(A.133), all its five boundary-line evaluations are one.
The tensor normalization therefore preserves those actual
$f$-values. This preservation is not claimed for arbitrary
commutator words.

The normalized matrices satisfy the exact marked identities
\[
 X_3=X_1,\qquad X_5=X_4,\qquad
 [X_4,X_1]=[X_4,X_2]=1.
 \tag{FP.48.1}
\]
They also follow from source identities. Put
$u=\xi_1$, $v=\xi_2$, $s=aba$. Then
$c=ua$, $a(u)=u$, $s(u)=v^{-1}$ and $s(v)=vuv^{-1}$, whence
\[
 X_3X_1^{-1}=u^2\in C',\qquad
 cbc=uv^{-1}s,\qquad
 X_5X_4^{-1}=uv^{-1}u^{-1}v^{-1}\in C'.
\]
Every $C'$ deck element acts by the same scalar on an
entire character block and its boundary line, so the tensor
operation kills $C'$ itself. Also $X_4$ is the literal
three-strand centre and commutes with $X_1,X_2$.
Equation (FP.48.1) makes the pentagon cancel by commuting
arguments and $f(y,x)=f(x,y)^{-1}$. Thus the apparently new
row-vector equation from the open-Prym coupling is automatic
on all actual gentle elements.

The exact involution eigenspaces and invariant splitting are
\texttt{check\_rank129\_open\_prym.py}; the determinant gauge,
all $24$ boundary-line tests and the marked identities are
\texttt{check\_open\_prym\_affine.py}. Complete matrices are
\texttt{ferrari\_open\_prym72.npz} and
\texttt{ferrari\_open\_prym\_affine72.npz}.

\begin{statement}{A.145}{Proved}{the entire linear rank-129 pentagon is automatic}
\end{statement}
The preceding argument also settles the other character
pieces without a separate parameter enumeration. In $L/C'$,
all $X_i$ centralize $C/C'$, while $X_4,X_5$ lie in this
abelian normal subgroup and $X_3$ differs from $X_1$ by
one of its elements. Hence the pentagon of every commutator
two-cycle word lies in $C'$.

For an actual multiplier-one gentle word, (A.144) shows
that its pentagon has scalar value one on all $24$ two-branch
characters of $C'/D$. These characters span the full
five-dimensional dual of $C'/D$; this is checked in the exact
binary character basis. Therefore
\[
        \boxed{\quad P(f)\in D=C''(C')^2
                   \quad\text{for every }(1,f)\in\mathcal H.\quad}
 \tag{FP.49.1}
\]
Conjugation by an element of $D$ is trivial on $D^{\rm ab}$.
Consequently the pentagon is one in the rank-$129$ homology
\emph{action}, over every coefficient ring, including
$\widehat{\mathbb Z}$. This establishes universal blindness
of that linear test despite the non-nilpotent pair subgroup
proved in (A.131). The intrinsic membership statement
(FP.49.1) must be distinguished from membership in
$D'D^3$, which is the next genuine affine test.

\begin{statement}{A.146}{Proved}{Proved affine construction and nontrivial Brunnian check; the actual-lift obstruction is proved in FP.52}
\end{statement}
The exact graph construction also gives the faithful
relative translation coordinates omitted by its homology
action. Let $C_1$ be the $384$-dimensional edge-chain space
of (A.130), $M=V_D$ its cycle space, and $p_q$ the chosen
tree path from the identity to $q\in Q$. Over $\mathbb F_3$ put
\[
 u=-\frac1{256}\sum_{q\in Q}p_q,
 \qquad \partial u=e_1-\frac1{256}\sum_{q\in Q}e_q.
 \tag{FP.50.1}
\]
The Fox path-chain model, conjugated by $u$, identifies
\[
          C/(D'D^3)\simeq M\rtimes Q.
\]
For a literal graph chain automorphism $T_g$, define
\[
                 r_g=u-T_gu\in M,
                 \qquad M_g=T_g|_M.
 \tag{FP.50.2}
\]
The boundary of $r_g$ is zero because $g$ fixes the identity
vertex and permutes all vertices. The exact cocycle identity
$r_{gh}=r_g+M_gr_h$ proves that
\[
             g\longmapsto(r_g,M_g)
       \quad\hbox{defines } L\longrightarrow
             \mathbb F_3^{129}\rtimes\operatorname{GL}_{129}(\mathbb F_3).
 \tag{FP.50.3}
\]
All braid and full-twist relations are checked in this
affine group, including the translations. The five $X_i$
have orders $6,6,6,12,12$, as opposed to the linear orders
in (A.130).

For an inner element $d\in D$ with homology class $m\in M$,
one has $M_d=I$ and $r_d=m$. Indeed $M^Q=0$, since averaging
identifies the invariant homology with $H_1(C,\mathbb F_3)=0$;
the three generators of $C$ are involutions. Averaging the
exact inner path formula gives $T_du=u-m$.
Combining this with (A.145), the affine pentagon is therefore
precisely the class
\[
                    [P(f)]\in D/(D'D^3).
 \tag{FP.50.4}
\]
It vanishes if and only if $P(f)\in D'D^3$.
The full forgetful base can be retained in a joint quotient.

This affine target has a concrete nonzero Brunnian check:
$[[X_1,X_2],[X_3,X_4]]$ has identity linear part and a
nonzero translation. The full vector is retained in the
certificate. There is also a word in $F''$ satisfying the
two-cycle identity whose pentagon has a nonzero translation;
that test word is explicitly \emph{not} asserted to satisfy
the gentle three-cycle or define a GT element. These checks
show that the affine construction retains actual source
information killed by the linear representation. They do
not yet produce a gentle obstruction character, a genuine
GT counterexample, or a proof that the affine equation
forces $\Delta=0$.

The chain maps, gauge and affine generators are saved in
\texttt{ferrari\_rank129\_affine.npz}; the construction and
all affine relator checks are part of
\texttt{check\_ferrari\_rank129.py}. The exact Brunnian and
two-cycle probes are
\texttt{check\_rank129\_affine\_defect.py} and its output.
The additional literal-word certificate
\texttt{check\_rank129\_inner\_translation.py} takes the
first chord generator of $D$, substitutes its full conjugated
involution words on every graph edge, and verifies identity
linear action and translation exactly equal to the first
homology basis vector. This independently excludes a mere
coboundary of the $129$-dimensional linear representation.
The structural quotient and faithful-shear proof are also
recorded as A.141--A.143 in the literature audit's
\texttt{Ferrari\_affine\_D.tex}.

\subsection{A simultaneous actual witness and automorphism for the affine quotient}


\begin{statement}{A.147}{Proved}{an actual mixed-sign witness passes the full affine rank--129 source quotient}
\end{statement}
There is one actual $k=(1,f_k)\in\mathcal H$, retaining all
properties of A.94, A.114 and A.128, such that
\[
 \boxed{\quad \Delta(k)=1,\quad D(k)=1,\quad
  \mathcal P(f_k)=1\ \text{in}\
  L/(D'D^3).\quad}
\]
In particular this is a condition in the actual relative source
quotient with finite kernel
$\overline C\simeq\mathbb F_3^{129}\rtimes Q$, of order
$2^8 3^{129}$.  It is stronger than satisfying the linear
homology test.  It does not assert the full five-point pentagon.

Let $e=e_3\in\widehat{\mathbb Z}$ have $3$-adic component
$1$ and all other components $0$, and write
\[
 A_e(u,v)=[[u^e,v^e],[v^e,((uv)^{-1})^e]],
 \qquad
 S=A_e(y,z)A_e(x,y)A_e(z,x),\quad z=(xy)^{-1}.
\]
Each of the three factors lies in $F''$ and in every full
orbifold core used in A.94 and A.128.  The same statement
holds after every peripheral permutation.  Its signed
level-twelve value is $[5I]$, so
\[
 S\in F''\cap\operatorname{core}_B\ker\varphi,
 \qquad q_{12}(S)=[5I]^3=[5I].
\]
The single profinite word $S$ also belongs to all added
evaluation kernels away from the exceptional prime three
in A.128 and to the nonharmonic Burau kernels in A.114,
because the corresponding statements hold for every
peripheral conjugate of $A_e$.

Consider any of the five actual affine evaluation maps,
and any of the six outer peripheral permutations of its
source.  Let $H_i$ be its finite affine pair image and
$s_i$ the value of $S$.  The exact same-word certificate
\texttt{check\_rank129\_affine\_seed.py} proves
\[
 s_i^6=1,\qquad s_i^3\in Z(H_i),\qquad
 s_i^4\in\ker\bigl(H_i\longrightarrow
       \mathrm{GL}(B_{63})\times\mathrm{GL}(M/B_{63})\bigr).
\]
Here $B_{63}$ is the literal invariant boundary homology
subspace, not a chosen abstract summand.  All peripheral
orders are computed independently in the affine target;
they divide $12$, so the evaluated exponent $e$ is $4$.
The values $s_i$ have order $1$ on the three nonharmonic
pairs and order $6$ on both harmonic pairs.  The last two
affine $X_i$ have order $12$, although their linear
matrices have order $4$; their translations are included
in every certificate calculation.

To see why these identities are sufficient, the kernel
of the affine graded action consists of matrices of the
form
\[
 \begin{pmatrix}
  I_{63}&C&t_B\\
  0&I_{66}&t_P\\
  0&0&1
 \end{pmatrix}.
\]
It is a finite $3$-group, of nilpotence class at most two
and exponent dividing three.  Its intersection $U_i$
with $H_i$ is normal in $H_i$.  Since $s_i^3$ is central
of order at most two, the normal subgroup
\[
             K_i=U_i\langle s_i^3\rangle
\]
is nilpotent.  It contains $s_i=s_i^3s_i^4$.
Equivalently, every tested value of the one actual word
$S$ has nilpotent normal closure.  No independent choice
of a boundary value, a proper value, or a frame is made.
For the subsequent core construction, use one common subgroup
for all six representatives in each fixed pair image $H_j$:
\[
 K_j^*=U_j\bigl\langle s_{j,\pi}^{\,3}:
                       \pi\in S_3\bigr\rangle.
\]
All these cubes centralize the same two original pair
generators, as checked by the certificate.  Their subgroup
is a finite central two-group; thus $K_j^*$ is normal and
nilpotent and contains every $s_{j,\pi}$.

Let $K$ be the intersection of the $B$-normal cores of
the five inverse images of these common $K_j^*$.
It is open and $B$-normal and contains $S$: the six
outer peripheral representatives have just been checked,
and further inner conjugation is harmless by normality
of $K_j^*$ in $H_j$.  Intersect it further with
$U_0$, an arbitrary open $B$-normal $M_0\supseteq F''$,
an arbitrary finite orbifold condition
$\operatorname{core}_B\varphi^{-1}J$, and any finite list
of the evaluation kernels from A.114 and A.128.  The
resulting subgroup $N$ still contains $S$.  Its
level-twelve image is a two-group containing $[5I]$.
Exactly the coprime $r$-centralizer construction and R2
normalization used in A.94 therefore give an actual
normalized gentle pair with
\[
 f,\beta\in R_{\rm nil}(N),\quad A\in N,
 \qquad q_{12}(A)=[5I],\quad q_{12}(\beta)=1,
 \quad\rho=0.
\]
Because the image of $N$ under each affine evaluation
is contained in the finite nilpotent group $K_j^*$, its
nilpotent residual evaluates to $1$.  Thus all five
actual $f$-factors have trivial affine image.  The cusp
normalization remains inside $R_{\rm nil}(N)$, so it
preserves these conclusions and the retained frames.

The compactness argument of A.94 over all $M_0,J$ and
all finite lists of the earlier conditions now applies
without change, using this one uniform seed $S$.
It gives one $k$ with all the asserted simultaneous
properties.  Finally $\mathcal P(f_k)$ lies in $C$ by
the forgetful fold, and its affine image is $1$.
A.142 and A.146 show that the affine action is faithful on
$\overline C$, so its image in the actual source
$L/(D'D^3)$ is $1$.
This completes the actual lifting proof.

\begin{statement}{A.148}{Proved}{the same witness acts on the actual affine source quotient and fixes its Ferrari kernel}
\end{statement}
Set
\[
 \Gamma_D=L/(D'D^3),
 \qquad \overline C=C/(D'D^3).
\]
The witness in A.147 defines an actual continuous automorphism
of $\Gamma_D$ which fixes $\overline C$ pointwise and induces
its genuine three-strand gentle automorphism on the full
forgetful quotient $\widehat B$.

Here kernel preservation is proved, not inferred from the
vanishing of a pentagon product.  The affine representation
$\rho$ is faithful on all of $\overline C$, not merely on
$M$: its linear deck matrix recovers the $Q$-coordinate,
and, once that coordinate is trivial, A.146 recovers the
$M$-coordinate as the translation.  Thus if $\rho(g)=1$
for $g\in\Gamma_D$, then for every $c\in\overline C$
the commutator $[g,c]$ belongs to $\overline C$ and has
trivial affine image, hence is $1$.

Put
\[
 f_0=f(a^2,b^2),\qquad e=f((aba)^2,c^2),
 \qquad
 a'=a,\quad b'=f_0^{-1}bf_0,\quad c'=e^{-1}ce.
\]
These are the exact multiplier-one braid formulas in the
fixed marking.  A.147 and antisymmetry give
$\rho(f_0)=\rho(e)=1$, so both conjugators centralize
$\overline C$.  Hence $a',b',c'$ have precisely the old
conjugation actions on that kernel.  Under the fold, their
images are the images of $a,b,c$ under the existing genuine
three-strand automorphism $\Phi_3$: the folded value of $e$
is $f(1,q(c)^2)=1$.  In fact $e\in\overline C$ and
centrelessness give $e=1$, although this extra observation
is not needed.

Each of the four defining relators
$aba=bab$, $bcb=cbc$, $ac=ca$, and $(abc)^4=1$,
evaluated on the new generators, has trivial fold and
trivial conjugation action on $\overline C$.  It therefore
belongs to $Z(\overline C)=1$.  The generator assignment
defines a continuous homomorphism
$\widetilde\Phi:L\to\Gamma_D$.
For $w\in C$, its image and its original quotient image
belong to $\overline C$ and have the same inner action;
centrelessness makes them equal.  Thus $D'D^3$ is killed,
and the map descends to an endomorphism of $\Gamma_D$
fixing $\overline C$ pointwise.  An endomorphism of an
extension which is identity on the kernel and an
automorphism on the quotient is bijective: injectivity
follows by taking the quotient first, and surjectivity by
lifting a quotient preimage and correcting in the kernel.
Compactness makes the continuous bijection a topological
automorphism.  Consequently
\[
 \boxed{\quad \Phi_D\in\operatorname{Aut}(\Gamma_D),
   \qquad \Phi_D|_{\overline C}=\operatorname{id},
   \qquad\Delta(k)=1.\quad}
\]
This is an actual relative source action.  It makes no
claim that $k$ belongs to the full profinite GT group.

\subsection{The affine cocycle cannot be removed by a change of origin}


\begin{statement}{A.149}{Proved}{no common affine fixed point, even after restriction to commutators}
\end{statement}
The translations in (A.146) cannot be eliminated by a common
change of origin. The exact systems
$(I-M_g)h=r_g$ for the indicated two generators have the
following ranks over $\mathbb F_3$:
\[
\begin{array}{c|c|c}
 \text{generators}&\text{coefficient rank}&\text{augmented rank}\\\hline
 a,b&125&126\\
 X_1,X_2&111&112\\
 X_2,X_3&111&112\\
 X_3,X_4&117&118\\
 X_4,X_5&120&121.
\end{array}
\]
There is also no common fixed point for any of these pairs'
derived subgroups. For the harmonic pairs, the equations for
the basic commutator and a finite set of its generator
conjugates already have ranks $90,91$; for the last two
pairs they have ranks $64,65$.
These are ranks of explicit inconsistent subsystems, not
claims about the rank of an unenumerated infinite system.
The certificate transforms an equation $Eh=b$ under an
affine generator $(v,M)$ to
$EM^{-1}h=b+EM^{-1}v$ and closes its row space under both
generators and their inverses until contradiction.
Thus even a formula $r_{f(g,h)}=(I-M_{f(g,h)})h_0$ valid
for every commutator word cannot remove the new translations.
This finite cohomological obstruction is compatible with the
actual mixed-parity witness in (A.147); it does not force
$\Delta$ to vanish. Exact systems and solutions or rank
contradictions are in
\texttt{check\_rank129\_fixed\_points.py} and its output.

\subsection{A genuine gentle action on the finite two-kernel quotient}


\begin{statement}{A.150}{Proved}{the finite two-group inputs and their central ambiguity}
\end{statement}
Retain $D=C''(C')^2$, $Q=C/D$ and $\Gamma_D=L/D$.
In the exact cube-chain model of $Q$, the images in
$\operatorname{Aut}(Q)$ of the five pairs have nilpotency
class at most two. The two harmonic pair groups have
order eight and central derived subgroup of order two;
the other three pairs commute. In particular every
$f\in\gamma_3(F)$ evaluates to one in these five
automorphism-pair groups. This is checked directly
from the literal involution-word actions, independently
in \texttt{check\_gentle\_Q\_action.py}: each pair
commutator commutes with both generators. The input
is the actual $256$-element quotient, not a parameter
assignment for $f$.

The following central facts can also be seen directly
in the cube. Write $S_i$ for the sum, over $\mathbb F_2$,
of the four edges in direction $i$. The cycle space
$C'/D$ has dimension five. The action of
$C/C'=(C_2)^3$ on it is faithful: a nonzero translation
in direction $i$ exchanges the two disjoint face
cycles normal to that direction. Therefore a central
element of $Q$ has trivial image in $C/C'$. The
translation-invariant edge chains are the span of
$S_0,S_1,S_2$, and each $S_i$ has boundary the sum
of all eight vertices. Thus
\[
 Z(Q)=\{c_0S_0+c_1S_1+c_2S_2:c_0+c_1+c_2=0\}
           \simeq(C_2)^2\subset Q'.
 \tag{GD.1a}
\]
The exact path substitutions for $a$ interchange
$S_0,S_1$ and fix $S_2$; those for $b$ interchange
$S_0,S_2$ and fix $S_1$. A common invariant vector
in the displayed even-sum plane is therefore zero.
An element of $Q$ fixed by $L$ first centralizes
$C\subset L$, so belongs to this plane. Consequently
\[
                              Q^L=1.
 \tag{GD.1b}
\]
Both assertions are also verified in the independent
finite certificate. In its twelve-edge binary
encoding the three nonzero centre vectors are
$255,3855,4080$.

\begin{statement}{A.151}{Proved}{all relators, kernel stability, and the composition law on $L/D$}
\end{statement}
Let $\Phi=(1,f)$ be a genuine gentle element. Its
two- and three-cycle relations imply
$f\in\gamma_3(F)$, as proved in A.140. Write
\[
 g=f(X_1,X_2),\quad j=f(X_2,X_3),\quad
 h=f(X_4,X_3),\quad k=f(X_5,X_1),\quad l=f(X_4,X_5).
\]
All words in this paragraph are evaluated in the
actual quotient $\Gamma_D$. Define generator images
\[
              a'=a,\qquad b'=g^{-1}bg,
                              \qquad c'=h^{-1}ch.
 \tag{GD.2a}
\]
The first braid relation for $a',b'$ is the genuine
three-strand gentle relation; its proof uses only
the two- and three-cycle equations and the source
free-group automorphism. It is the three-strand
formula recorded in FP.26.4, in the same marking.
For clarity about the central lift, the gentle
extension on the central quotient fixes the
projective braid relation; its residual central
error upstairs has exponent sum zero and hence
is one. The image of $z_3=(ab)^3$ centralizes the
image of the free pure kernel, which is still
that kernel by the genuine free-group
automorphism. It therefore belongs to the
three-strand centre. If it is $z_3^t$, exponent
sum gives $6t=6$ in $\widehat{\mathbb Z}$,
so $t=1$. Thus $z_3$ is fixed, without a
five-point action assumption.

The commutation $[a',c']=1$ holds because $a$
commutes with both arguments $X_4=z_3$ and
$X_3=c^2$ of $h$. For the remaining braid relation,
A.140 gives the actual quotient identity
\[
                   h^{-1}kjlg=1,
                   \qquad h=kjlg.
 \tag{GD.2b}
\]
Here $k$ commutes with $c$, since both $X_5=(bc)^3$
and $a^2$ commute with $c$. Likewise $l$ commutes
with $b$, since both three-strand centres commute
with $b$. Conjugating $b',c'$ by $g$ therefore
gives
\[
       b,\qquad (hg^{-1})^{-1}c(hg^{-1})
                  =l^{-1}j^{-1}cjl.
\]
The genuine three-strand gentle relation for
$b,j^{-1}cj$, followed by conjugation by $l$,
proves the desired braid relation. Thus the
assignment already defines a homomorphism
$\widehat B_4\to\Gamma_D$.

Its image of the central full twist $z_4=(abc)^4$
has trivial fold in $\widehat B$, because the
folded generator images are the genuine gentle
three-strand images of $a,b,a$. Hence this full
twist image belongs to $Q$. By A.150, $g$ and $h$
centralize $Q$, so the new generators have exactly
the old conjugation actions on $Q$. The full
twist image centralizes all new generators,
and therefore belongs to $Q^L=1$. This verifies
the fourth defining relator, yielding a genuine
continuous homomorphism
\[
                    \widetilde\Phi:L\longrightarrow\Gamma_D.
\]

Its fold is the original gentle automorphism of
$\widehat B$. Thus $\widetilde\Phi(C)\subset Q$.
For $w\in C$, its new and old images induce the
same inner action on $Q$, because the generator
actions agree. Consequently
\[
       \delta_\Phi(w)=\widetilde\Phi(w)\bar w^{-1}\in Z(Q).
 \tag{GD.2c}
\]
This is a homomorphism $C\to Z(Q)$: the differences
are central, so the multiplicative error for a
product is the product of the errors. It kills
$C'$, and therefore kills $D\subset C'$.
Since $\bar D=1$, the map $\widetilde\Phi$
kills $D$. This proves the required kernel
stability before defining
\[
                    \Phi_D:\Gamma_D\longrightarrow\Gamma_D.
\]

The maps obey the genuine gentle composition law.
Use the convention $\Phi_{12}=\Phi_1\circ\Phi_2$,
so on the free source
$f_{12}=f_1\Phi_1(f_2)$. On the first pair
$(a^2,b^2)$, the map $\Phi_{1,D}$ fixes the first
argument and conjugates the second by
$g_1=f_1(a^2,b^2)$. Thus
$g_{12}=g_1\Phi_{1,D}(g_2)$. The same calculation
on $(z_3,c^2)$, using $\Phi_{1,D}(z_3)=z_3$,
gives $h_{12}=h_1\Phi_{1,D}(h_2)$.
Substitution in (GD.2a) proves
\[
              \Phi_{12,D}=\Phi_{1,D}\circ\Phi_{2,D}.
\]
The identity pair gives the identity map, and
the inverse gentle pair supplies the inverse.
We have therefore proved an actual continuous
action
\[
                 \boxed{\ \mathcal H\longrightarrow
                             \operatorname{Aut}(L/D).\ }
 \tag{GD.2d}
\]
All defining relators and the kernel have been
checked; this is not an action inferred solely
from the vanishing of a pentagon product.

\begin{statement}{A.152}{Proved}{the entire gentle group fixes $Q$ pointwise, with no inferred action on $M$}
\end{statement}
The correction (GD.2c) factors through
$C^{\rm ab}=(\mathbb F_2)^3$ and is $L$-equivariant,
because the new and old $L$ conjugation actions
on $Q$ agree. The action on $Q'$ is pointwise
the identity. Since $Z(Q)\subset Q'$, composition
adds these corrections. Their possible values form
\[
       \operatorname{Hom}_L((\mathbb F_2)^3,Z(Q))
                              \simeq\mathbb F_2.
 \tag{GD.3a}
\]
Indeed the first module is the three-letter
permutation module, and $Z(Q)$ is its even-sum
two-dimensional summand with the same $S_3$ action.
The one-dimensional Hom-space is also checked
by enumerating its sixty-four possible linear
maps in the exact certificate. This is only
an intermediate bound: the actual correction
always vanishes, as follows.

In $\Gamma_D$, let $J=\overline{\langle z_3,c^2\rangle}$.
Its fold is procyclic, so $J'\subset Q$.
Its conjugation action on $Q$ is abelian
by A.150, so in fact $J'\subset Z(Q)$.
Both generators act trivially on $Z(Q)$:
$z_3\in Q$, and $c^2$ has trivial
three-letter permutation action on the
even-sum plane (GD.1a). Thus $J$ is
nilpotent of class at most two. Since
$f\in\gamma_3(F)$, this proves
\[
                        h=f(z_3,c^2)=1
                         \quad\text{in }\Gamma_D.
\]
The same cyclic-fold and central-action
argument applies to the other two
nonharmonic pairs, giving $k=l=1$.
In particular (GD.2a) fixes $a$ and $c$,
and therefore fixes $u=ca^{-1}=t_1t_0$.
It follows that
$\delta_\Phi(t_0)=\delta_\Phi(t_1)=z$
for some $z\in Z(Q)$. Equivariance under
$a$, which interchanges $t_0,t_1$ in
the abelianization, gives $a(z)=z$.
Equivariance under $b$, which fixes
$t_1$, gives $b(z)=z$. Equation (GD.1b)
forces $z=0$. Finally $b$ interchanges
$t_0,t_2$ in the abelianization, so
$\delta_\Phi(t_2)=0$ as well. Therefore
\[
          \boxed{\ \Phi_D|_{Q}=\operatorname{id}_Q
                     \quad\text{for every }\Phi\in\mathcal H.\ }
 \tag{GD.3b}
\]
There is no nontrivial central-correction
character produced by this quotient.
The restriction is trivial on the entire
gentle group, including the A.129 witness
with $\Delta=1$.

An action on $L/D$ does not, however, determine
an action on $M=H_1(D,\mathbb F_3)$, because
$D$ has already been killed. A lift through
the free-product cover $C\to Q$ requires
peripheral conjugators, and changing those
lifts can change their action on $M$.
The argument above does not construct such
a lift for all gentle elements, nor show
that $k_{\rm rel}$ acts trivially on $M$.
Consequently it supplies no crossed-cocycle
composition law for
$m_f=[\mathcal P(f)]\in D/(D'D^3)$.
Such a law requires an additional specified
action on this module and a compatibility
proof. In particular no three-power index
or odd-power witness argument is inferred
from an unconstructed affine action.

\subsection{One actual witness for every coefficient level of the derived cover}


\begin{statement}{A.153}{Proved}{the whole adelic $D$-abelian family is simultaneously insufficient}
\end{statement}
The witness may in fact be chosen to satisfy
\[
 \boxed{\quad \Delta(k)=1,\quad D(k)=1,\quad
       \mathcal P(f_k)=1\ \text{in}\ L/D'.\quad}
\]
All properties of A.94, A.114 and A.128 are retained.
Here $D=C''(C')^2$ is the one fixed characteristic
subgroup of index $256$ in $C$.  This theorem retains
the full profinite abelianization
$D/D'\simeq\widehat{\mathbb Z}^{129}$, not just its
modulo-three quotient.  It does not replace $D/D'$ by
the nonabelian pro-three completion of $D$.

\emph{The uniform three-adic argument.}
Let $W=C'/D\simeq(C_2)^5$ and put
\[
 E_0=\frac1{32}\sum_{w\in W}\rho(w)
 \quad\text{on}\quad M_3=H_1(D,\mathbb Z_3).
\]
Normality and characteristicity make $E_0$ commute
with the literal $L$-linear action.  It gives an
$L$-stable decomposition
\[
 M_3=M_0\oplus M_\perp,
 \qquad \operatorname{rank}M_0=5,
 \quad\operatorname{rank}M_\perp=124.
\]
Consequently the affine action maps injectively to
the product of its affine actions on the two summands.
The modulo-three certificate for the \emph{same} word
$S$ proves, for all five pairs and all six peripheral
representatives, that its two-primary part acts as the
identity on $M_\perp/3M_\perp$, including its translation.
Equivalently its complementary affine image modulo three
has three-power order.  The kernel of
\[
 \operatorname{Aff}(M_\perp)\longrightarrow
 \operatorname{Aff}(M_\perp/3M_\perp)
\]
is pro-three: successive congruence kernels are additive
groups over $\mathbb F_3$.  Hence that one element's
closed cyclic complementary image is pro-three, and
its two-primary component is exactly $1$ over
$\mathbb Z_3$, at every congruence level.

The rank-five factor is the earlier one, including the
affine cocycle.  To prove this precisely, use the
covering of Schreier graphs for $D\subset C'$ and its
chain pushforward $p_*$.  It is equivariant for the
actual Nielsen word action.  Transfer identifies
$p_*:M_0\to H_1(C',\mathbb Z_3)$ with an isomorphism,
since $32$ is a unit.  For the graph gauges of A.146,
\[
 \partial(p_*u_D)=\partial u_{C'},\qquad
 \delta=p_*u_D-u_{C'}\in H_1(C',\mathbb Z_3).
\]
Thus, for every literal source automorphism $\alpha$,
\[
 p_*t_{\alpha,D}
   =t_{\alpha,C'}+(I-T_{\alpha,C'})\delta.
\]
The two affine actions are therefore conjugate by one
fixed translation.  This is an exact comparison over
$\mathbb Z[1/2]$, not an identification inferred from
modulo-three matrices.

On that rank-five factor, the harmonic three-strand
group has the pro-three-by-central-involution property
proved in A.125 and A.128.  Its linear image modulo
three is $S_3\times\mathrm{SL}_2(\mathbb F_3)$, whose
second derived subgroup is the central involution;
the genuine central full twist supplies a central
two-primary lift, and the remaining kernel is
pro-three.  Since every peripheral substitution of
$S$ belongs to $F''$, its rank-five two-primary
component is central in the pair image.  On the
complementary factor that component is $1$ by the
preceding paragraph.  Hence the two-primary component
of every tested value of $S$ centralizes the \emph{whole}
rank--129 harmonic affine pair over $\mathbb Z_3$.
For each of the three nonharmonic pairs the entire
modulo-three seed value is $1$, so its whole affine
image is pro-three and its two-primary component is
$1$.  This accounts for all five pairs and all thirty
peripheral cases.

Its three-primary component lies in the inverse image
of the modulo-three graded kernel from A.147.  That
inverse image is pro-three, since the modulo-three
graded kernel is a three-group and the congruence
kernel is pro-three.  Therefore, at every modulus
$3^n$, all six values of $S$ in a fixed pair lie in
one normal nilpotent subgroup: take this normal
three-group times the central two-group generated by
their two-primary components.  This establishes the
same R2 input uniformly for every $n$.  Exact identity
of the three-primary component on the graded modules
\emph{modulo $3^n$} is neither claimed nor needed.

\emph{The other primes and the actual source comparison.}
The exact integer chain construction, in the affine
gauge with denominator $256$, proves
\[
 (X_i^{m_i}-I)^2=0,\qquad
 m_i=2\ (i=1,2,3),\qquad m_i=4\ (i=4,5)
\]
as identities of augmented $130$-dimensional matrices
over $\mathbb Z[1/2]$.  Thus at every odd prime
$\ell\ne3$ the procyclic image of $X_i$ has primes
only in $\{2,\ell\}$.  The exponent $e_3$ kills each
of the five $X_i$.  The word $A_{e_3}$, and therefore
$S$, is consequently in every required pair kernel
and every peripheral core at those primes.

For odd $\ell$, the affine construction is faithful
on
\[
 C/(D'D^{\ell^n})\simeq
       (\mathbb Z/\ell^n)^{129}\rtimes Q.
\]
Indeed the integral deck representation is faithful
because its modulo-three representation is faithful;
its finite two-group cannot acquire a kernel on
reduction at an odd prime, since a congruence kernel
is pro-$\ell$.  The deck averaging operator is zero
over $\mathbb Z[1/2]$: the rational graph cycle lattice
has no $Q$-invariants by A.142.  Thus there are no invariant
vectors at any odd modulus, and the affine argument
of A.146 applies verbatim.  It converts factorwise
trivial affine evaluations into the pentagon equation
in the actual quotient $L/(D'D^{\ell^n})$.

At the prime two, $C'/(D'D^{2^n})$ is a finite
two-group, so
$O^2(C')\subseteq D'D^{2^n}$.  The entire source
quotient $L/(D'D^{2^n})$ is therefore already covered
by the retained A.128 equation in $L/O^2(C')$.
No division by $256$ or centreless assertion about a
finite two-group is used at this prime.

Use the one seed $S$ in the R2 finite-intersection
construction of A.147, now retaining any finite list
of odd-prime levels and their common normal nilpotent
pair-image subgroups.  Every source normal core
contains $S$, and its signed value remains $[5I]$.
The same normalization and compactness proof gives
one actual $k$ satisfying every listed condition at
every prime and level, together with all previous
conditions.  Finally
\[
 \bigcap_{\ell\ {
m prime},\ n\ge1}
         D'D^{\ell^n}=D',
\]
because $D/D'\simeq\widehat{\mathbb Z}^{129}$ is
separated by these finite quotients.  This proves the
boxed equation in $L/D'$.

The independent checks at $9$ and $27$ in
\texttt{check\_rank129\_3adic\_seed.py} compute actual
affine peripheral orders and CRT primary components.
They confirm all thirty normal-nilpotent-image gates,
with harmonic seed orders $18$ and $54$, respectively.
They are consistency checks of the uniform proof,
not its premise.  The all-core complementary-affine
identity is explicitly asserted in
\texttt{check\_rank129\_affine\_seed.py}; the exact
integral peripheral and transfer comparisons are
given by the corresponding integral graph certificates.

\begin{statement}{A.154}{Proved}{an actual automorphism of $L/D'$ fixing its full Ferrari kernel}
\end{statement}
The same $k$ defines a continuous automorphism
\[
 \boxed{\quad
 \Phi_{D^{\rm ab}}\in\operatorname{Aut}(L/D'),
 \qquad
 \Phi_{D^{\rm ab}}|_{C/D'}=\operatorname{id},
 \qquad\Delta(k)=1.\quad}
\]
The kernel $C/D'$ is centreless.  If a central element
has deck coordinate $q\in Q$, its action on
$D/D'$ is trivial; faithfulness modulo three gives
$q=1$.  Its remaining coordinate belongs to
$(\widehat{\mathbb Z}^{129})^Q=0$.  The integral sum of the $256$ deck matrices is zero,
since its rational image is the invariant subspace,
which vanishes by A.142. Thus an invariant vector $v$
satisfies $256v=0$. The module $\widehat{\mathbb Z}^{129}$
is torsion-free, so $v=0$, including at the prime two.

For the witness above, both braid conjugators
$f(a^2,b^2)$ and $f((aba)^2,c^2)$ centralize
$C/D'$.  Their commutators with that kernel vanish
in every odd-prime quotient by the faithful affine
calculations, and in every two-primary quotient by
the retained A.128 action; the intersection identity
in A.153 then gives the claim in $C/D'$ itself.
The exact generator-and-relator proof of A.148 now
applies with $C/D'$ in place of $\overline C$.
It proves the defining relators, fixes the kernel
pointwise, kills $D'$, and yields a bijective
continuous endomorphism inducing the genuine
three-strand automorphism on the full base.
This proves an action on the actual source quotient,
without a claim about the full profinite GT group.

\Needspace{10\baselineskip}
\subsection{A faithful Fox target and exact source comparison for the fourth-power cover}


\begin{statement}{A.155}{Proved}{the marked covering chains and their character components}
\end{statement}
Retain the literal Ferrari subgroup $C=C_2\amalg C_2\amalg C_2$
of $L$ and set
\[
 F_5=C',\qquad D_4=F_5'F_5^4=C''(C')^4,\qquad
 A=F_5/D_4\simeq(\mathbb Z/4)^5,\qquad k=\mathbb F_9.
 \tag{FF.1a}
\]
All subgroups in this statement are closed. The group $F_5$ is free
profinite of rank five by A.115. Choose an ordered free basis
$(g_1,\ldots,g_5)$ from the cube Schreier construction, and put
$a_j=\overline g_j\in A$. There are $1024$ deck sheets, so
$[C:D_4]=8\cdot1024=8192$, and $D_4$ is free of rank
$1+1024(5-1)=4097$.

The covering of the five-petal bouquet corresponding to $D_4$
has the exact chain sequence
\[
 0\longrightarrow M=H_1(D_4,k)
 \longrightarrow V=k[A]^5
 \stackrel{\partial}{\longrightarrow}I_A
 \longrightarrow0,\qquad
 \partial=(a_1-1,\ldots,a_5-1),
 \tag{FF.1b}
\]
where $I_A$ is the augmentation ideal of $k[A]$. Thus
$\dim_k(M,V,I_A)=(4097,5120,1023)$.
One can first construct this finite graph for the discrete free
group on the chosen basis and then pass to profinite completion.
Its finite-index subgroup is free, and completion preserves its
finite-coefficient abelianization. This proves that (FF.1b)
computes the indicated continuous profinite homology.

Fix the following based path convention. Its column chain
$d(w)\in V$, together with its endpoint $\bar w\in A$, satisfies
\[
 d(uv)=d(u)+\bar u\,d(v),\qquad
 d(g_j)=e_j,\qquad \partial d(w)=\bar w-1.
 \tag{FF.1c}
\]
These equations also give a direct finite construction for every
profinite word. First put $V_0=\mathbb F_3[A]^5$ and use the
homomorphism from the free profinite group to the finite group
$V_0\rtimes A$ which sends $g_j$ to $(e_j,a_j)$.
Its image is precisely
\[
 \{(v,a)\in V_0\rtimes A:\partial v=a-1\}.
 \tag{FF.1d}
\]
Indeed a path to any specified sheet exists, and its chain can
be changed by every cycle in the graph. On $D_4$, the chain
map is the ordinary mod-three homology map. Its kernel is
exactly $D_4'D_4^3$, so (FF.1d) is the actual finite group
$F_5/(D_4'D_4^3)$. Extending the coefficients injectively from
$\mathbb F_3$ to $k$ gives (FF.1c) with values in $V$.
It does not assert that the original group supplies arbitrary
$k$-cycles.

The field $k$ contains all fourth roots of unity, and $1024$
is invertible in it. Consequently the Fourier decomposition is
exact. For a character $\chi:A\to k^\times$, write
\[
 V_\chi=k^5,\qquad
 \partial_\chi=(\chi(a_1)-1,\ldots,\chi(a_5)-1),\qquad
 M_\chi=\ker\partial_\chi.
 \tag{FF.1e}
\]
For each of the $1023$ nontrivial characters, the row
$\partial_\chi$ is nonzero and there is an exact sequence
\[
 0\longrightarrow M_\chi\simeq k^4
 \longrightarrow V_\chi\simeq k^5
 \stackrel{\partial_\chi}{\longrightarrow}k
 \longrightarrow0.
 \tag{FF.1f}
\]
For the trivial character, the boundary map is zero and
$M_1=V_1\simeq k^5$. Thus
$4097=5+1023\cdot4$. The fifth chain coordinate in
(FF.1f) is part of the target and must be retained.

\begin{statement}{A.156}{Proved}{exact semilinear transport and the inner-action formula}
\end{statement}
An element $\ell\in L$ acts on $F_5$ by its literal conjugation
automorphism $\phi$. Since $D_4$ is characteristic in $F_5$,
there is an induced automorphism $\bar\phi$ of $A$. Define
the marked Fox Jacobian by its columns:
\[
 J_\phi=(d(\phi(g_1))\ \cdots\ d(\phi(g_5)))
       \in\operatorname{Mat}_5(k[A]).
 \tag{FF.2a}
\]
The based chain substitution is the semilinear map
\[
 T_\phi(v)=J_\phi\,\bar\phi(v),\qquad
 T_\phi(av)=\bar\phi(a)T_\phi(v).
 \tag{FF.2b}
\]
Equation (FF.1c), first on the basis and then on products and
inverses, proves
$d(\phi(w))=J_\phi\bar\phi(d(w))$. Consequently
\[
 \partial J_\phi=\bar\phi(\partial),\qquad
 J_{\phi\circ\psi}=J_\phi\bar\phi(J_\psi),\qquad
 T_{\phi\circ\psi}=T_\phi T_\psi.
 \tag{FF.2c}
\]
The inverse automorphism supplies the inverse chain map.
These are genuine actions of $L$ on the finite chain spaces,
obtained from its source conjugation action. All defining
relators of $L$ therefore hold; there is no assumption that
a gentle element acts on a larger five-point group.
Continuity follows directly from the finite construction
(FF.1c), since each Jacobian column and induced deck action
depends continuously on the images of the finitely many
chosen generators.

For clarity, the exact character transport rule is as follows.
Let
\[
 e_\chi=\frac1{|A|}\sum_{a\in A}\chi(a)^{-1}a.
\]
Then $\bar\phi(e_\chi)=e_{\chi'}$, where
\[
 \boxed{\quad \chi'=\chi\circ\bar\phi^{-1},\qquad
 T_\phi:V_\chi\longrightarrow V_{\chi'},\quad
 v\longmapsto J_\phi(\chi')v.\quad}
 \tag{FF.2d}
\]
The Jacobian is evaluated at the \emph{target} character.
In particular boundary compatibility is the exact row identity
\[
                \partial_{\chi'}J_\phi(\chi')=\partial_\chi.
 \tag{FF.2e}
\]
The literal positive and inverse Nielsen words and Fox Laurent
coefficients are stored in
\path{agent_lifting/records/f5_nielsen_fox.json}.
The independent certificate
\path{agent_lifting/code/check_f5_fox_characters.py} checks (FF.2e) for
all $1024$ characters and all six marked generator maps,
as well as both orders of their inverse products. Its saved
output is \path{agent_lifting/records/check_f5_fox_characters.out}.

No commutativity of the pure-input character permutations is
used here. In fact those permutations need not commute modulo
four. An individual value $f(X_i,X_j)$ must be multiplied with
its actual block permutation, unless its action on $A$ has
separately been proved trivial. Restriction to one character
without that proof would not define the required word
evaluation. The direct sum over all nontrivial characters is
$L$-stable and retains every transition in (FF.2d).

For later faithfulness, consider an actual element $c\in F_5$.
Conjugation by $c$ acts trivially on the abelian deck group $A$.
Using (FF.1c) gives, for each basis element,
\[
 d(cg_jc^{-1})
 =d(c)+\bar c e_j-\bar g_jd(c)
 =\bar c e_j+(1-a_j)d(c).
\]
Therefore its full chain action is
\[
 \boxed{\quad T_{\operatorname{Inn}(c)}
          =\bar c I-d(c)\partial,\qquad
 T_{c,\chi}=\chi(\bar c)I-d(c)_\chi\partial_\chi.\quad}
 \tag{FF.2f}
\]
In particular its action on $M_\chi$ is the scalar
$\chi(\bar c)$. If $c\in D_4$, this becomes
$I-d(c)_\chi\partial_\chi$. It is the identity on
$M_\chi$ for every such $c$, whereas on the full $V_\chi$
it is the identity if and only if $d(c)_\chi=0$, for
$\chi\ne1$. The latter equivalence follows because
$\partial_\chi$ is a nonzero row over a field.

\begin{statement}{A.157}{Proved}{the joint target detects the actual fourth-power affine quotient}
\end{statement}
Put
\[
 E_0=C/(F_5'F_5^3)=C/(C''(C')^3),\qquad
 W=\bigoplus_{\chi\ne1}V_\chi.
\]
The old affine target is the actual conjugation action on
$E_0$, not merely its rank-five homology action. Together
with (A.156), it gives a continuous finite representation
\[
 \rho_{\rm joint}:L\longrightarrow
       \operatorname{Aut}(E_0)\times\operatorname{GL}_k(W).
 \tag{FF.3a}
\]
Its kernel on the full Ferrari subgroup $C$ is exactly
\[
 \boxed{\quad C\cap\ker\rho_{\rm joint}=D_4'D_4^3.\quad}
 \tag{FF.3b}
\]
Thus the image of $C$ in this joint target is a faithful copy
of the actual group $C/(D_4'D_4^3)$.
All the quotients here are legitimate: $F_5$ is characteristic
in the normal subgroup $C$, and $D_4'D_4^3$ is characteristic
in the characteristic subgroup $D_4$ of $F_5$, hence normal
in $L$.

Here is a complete proof. First restrict to $F_5$. Suppose
$c\in F_5$ is killed by the
joint target. For each nontrivial character, its restriction
to the four-dimensional space $M_\chi$ is multiplication
by $\chi(\bar c)$, by (FF.2f). The identity action therefore
implies $\chi(\bar c)=1$ for every nontrivial $\chi$.
These characters separate $A$, so $c\in D_4$.
The full five-dimensional blocks now imply
\[
                         d(c)_\chi=0\qquad(\chi\ne1).
 \tag{FF.3c}
\]

It remains to detect the trivial character. Write
$M_0=F_5/(F_5'F_5^3)\simeq\mathbb F_3^5$ and
$A_0=C/F_5\simeq(C_2)^3$. The cube calculation A.116
decomposes $M_0\otimes k$ into the nontrivial characters
$110,101,011,111,111$ of $A_0$. In particular
\[
                              M_0^{A_0}=0.
 \tag{FF.3d}
\]
If $c\in F_5$ acts trivially by conjugation on $E_0$, its
image in the normal abelian subgroup $M_0$ must be fixed
by $A_0$, hence is zero. Thus the old affine component
detects exactly the ordinary mod-three abelianization
of $c$ in $F_5$. That abelianization is
\[
             d(c)_1=(\text{exponent sums of }c\bmod3)
                         \in k^5.
 \tag{FF.3e}
\]
It follows that this last Fourier component also vanishes.
Fourier inversion in $k[A]$ now gives $d(c)=0$.
Since $c\in D_4$, (FF.1c) identifies its kernel as
$D_4'D_4^3$. Conversely, an element of $D_4'D_4^3$
has endpoint one and every chain component zero, and
also lies in $F_5'F_5^3$; it is therefore killed by both
factors in (FF.3a). This proves the asserted kernel equality
on $F_5$.

To extend it to $C$, the same old affine group is sufficient.
The four characters $110,101,011,111$ span $A_0^*$, so the
$A_0$-action on $M_0$ is faithful. A central element of $E_0$
first has trivial image in $A_0$, because it acts trivially
on $M_0$; it then belongs to $M_0^{A_0}=0$. Thus
$Z(E_0)=1$. If $c\in C$ is killed by the old conjugation
target, its image in $E_0$ is central and hence one. Therefore
$c\in F_5'F_5^3\subset F_5$, where the kernel equality is
already proved. This establishes (FF.3b) on all of $C$.

The trivial component can equivalently be described by
transfer. Let $p_*$ be the map on homology induced by
the $1024$-sheeted covering, and let $\operatorname{tr}$
be transfer. Then
\[
 p_*\operatorname{tr}=1024\,\mathrm{id},\qquad
 \operatorname{tr}p_*=\sum_{a\in A}a.
 \tag{FF.3f}
\]
Because $1024$ is a unit in characteristic three,
$p_*:H_1(D_4,k)^A\to H_1(F_5,k)$ is an isomorphism.
This is precisely the rank-five Fourier component in
(FF.3e). The full chain action at the trivial character
alone cannot detect it: (FF.2f) there is always the
identity. The old affine factor in (FF.3a) supplies
exactly this missing information.

Finally let $f$ be a genuine gentle multiplier-one word,
and use the literal five pentagon inputs $X_i$. Its
pentagon defect $\mathcal P(f)$ belongs to $C'$.
For completeness, this weaker conclusion of A.140
already follows in $L/C'$: every pure $X_i$ centralizes
$C/C'$, the images of $X_4,X_5$ belong to that abelian
kernel, and $X_3X_1^{-1}$ also belongs to it.
The three nonharmonic commutator-word evaluations are
therefore one, while the other two cancel by the
two-cycle relation. Consequently (FF.3b) gives the
exact equivalence
\[
 \boxed{\quad
 \rho_{\rm joint}(\mathcal P(f))=1
 \quad\Longleftrightarrow\quad
 \mathcal P(f)\in D_4'D_4^3
 \quad\Longleftrightarrow\quad
 \mathcal P(f)=1\text{ in }L/(D_4'D_4^3).\quad}
 \tag{FF.3g}
\]
The stronger previous result $\mathcal P(f)\in
C''(C')^2$ is compatible with this statement but is
not needed for its proof. The equivalence concerns
the product with its actual semilinear transport;
it does not assert that its individual factors are
blockwise or trivial. It establishes a faithful
test in an actual finite relative quotient, without
asserting that the test holds for every gentle
element or proves the full GT sign comparison.

\begin{statement}{A.158}{Proved}{a conditional transfer to a genuine relative source automorphism}
\end{statement}
Let
\[
 N=D_4'D_4^3,\qquad \Gamma=L/N,\qquad K=C/N,
 \qquad \pi:\Gamma\longrightarrow\widehat B.
 \tag{FF.4a}
\]
For a genuine gentle multiplier-one pair $(1,f)$, put
\[
 g=f(X_1,X_2)=f(a^2,b^2),\qquad
 h=f(X_4,X_3)=f((ab)^3,c^2).
\]
Assume exactly the following two hypotheses in the full joint
chain-and-affine target of A.157:
\[
             \rho_{\rm joint}(g)=1,
             \qquad \rho_{\rm joint}(h)=1.
 \tag{FF.4b}
\]
Then the literal Drinfeld generator assignment
\[
        a\longmapsto a,\qquad
        b\longmapsto g^{-1}bg,\qquad
        c\longmapsto h^{-1}ch
 \tag{FF.4c}
\]
defines a genuine continuous automorphism of $\Gamma$, fixes
$K$ pointwise, and induces the given gentle automorphism on
$\widehat B$. This is a conditional assertion about an actual
relative quotient; no claim is made that (FF.4b) holds for a
particular proposed seed.

To prove it, let $q:L\to\Gamma$ be the quotient map. By A.157,
$N\subset\ker\rho_{\rm joint}$, so the joint representation
factors as $\rho_{\rm joint}=\theta q$, where
$\theta:\Gamma\to\operatorname{im}\rho_{\rm joint}$ is
faithful on $K$. Define (FF.4c) first as a map from the free
profinite group on the three braid generators into $\Gamma$.
By (FF.4b), its composite with $\theta$ agrees on every
generator with the original quotient map.

Its fold is also prescribed. In $\widehat B$, the first
argument of $h$ is one, and $f\in\widehat F_2'$, so $h$
has trivial fold. The three folded images are therefore
$a,\,\bar g^{-1}b\bar g,\,a$. These are the images of
$a,b,a$ under the genuine three-strand gentle automorphism,
whose construction uses only the gentle two- and three-cycle
relations and its free-group automorphism, as in B.35 and
A.151. In particular the folded map is that automorphism
composed with the original fold.

Consider each defining relator of $L$:
\[
             aba=bab,\qquad bcb=cbc,\qquad
             [a,c]=1,\qquad (abc)^4=1.
\]
Its new image has trivial fold and therefore belongs to $K$.
Its new $\theta$-value is also trivial, because the new and
old generator values under $\theta$ agree. Faithfulness of
$\theta$ on $K$ makes the new relator image one. This proves
all four relators, including the central full-twist relation,
and hence constructs a continuous homomorphism
$\widetilde\Phi:L\to\Gamma$.

For every $w\in C$, both $\widetilde\Phi(w)$ and $q(w)$
belong to $K$, and their $\theta$-values agree. The same
faithfulness gives
\[
                         \widetilde\Phi(w)=q(w)
                                  \qquad(w\in C).
 \tag{FF.4d}
\]
In particular $\widetilde\Phi$ kills $N$. Kernel stability
has thus been proved before descending to a homomorphism
$\Phi_\Gamma:\Gamma\to\Gamma$; no lift of an automorphism
of a deck group to its free covering group has been assumed.
The descended map is the identity on $K$ and the given
automorphism on $\widehat B$. It is injective because an
element of its kernel first has trivial fold and then lies
in the kernel of the identity on $K$. It is surjective by
lifting a desired folded preimage and correcting its residual
error by an element of $K$. Its inverse is continuous because
$\Gamma$ is compact Hausdorff. This proves the assertion.

The two hypotheses (FF.4b) concern full word evaluations in
the semilinear target. Triviality only on the four-dimensional
cycle blocks, a solution of the pentagon equation alone, or
nilpotence in a discarded linear quotient does not supply
these hypotheses. In particular this proposition does not
assert a full GT action, a new character comparison, or an
action on an additional homology cover whose kernel stability
has not been established.

\subsection{The complete fourth-power computation and the failure of the previous cyclic seed}


\begin{statement}{A.159}{Proved}{an actual full-chain fourth-power cover target}
\end{statement}
Put $D_4=C''(C')^4$. Then $C'/D_4\simeq(\mathbb Z/4)^5$,
$[C:D_4]=8192$, and $D_4$ is free of rank $4097$.
Use the five cube-Schreier generators $q_1,\ldots,q_5$ of $C'$,
and $\mathbb F_9=\mathbb F_3[i]$, $i^2=-1$.
For every row character $\chi\in(\mathbb Z/4)^5$, retain all five
Fox chain coordinates, with boundary
$d_\chi=(i^{\chi_1}-1,\ldots,i^{\chi_5}-1)$.
An automorphism $g$ with integral abelianization $T_g$ sends source
$\chi$ to target $\eta=\chi T_g^{-1}$ with matrix $J_g(\eta)$.
The identities $d_\eta J_g(\eta)=d_\chi$ and the Fox chain rule
give an actual monomial representation of $L$.
All braid, distant, and full-twist relations hold both in the
literal free substitutions and in the evaluated matrices.

For nontrivial $\chi$, the cycle subspace has dimension four.
The fifth chain coordinate is essential: for $h\in D_4$,
\[
                  J_{\operatorname{Ad}(h)}(\chi)
                    =I-v_h(\chi)d_\chi.
 \tag{FP.54.1}
\]
Thus the full chain action detects the corresponding translation,
while restriction to cycles kills it. Adding the earlier rank-five
affine factor detects the trivial character. The faithful source
construction is A.155--A.157 in
\path{agent_literature/Ferrari_fourth_power_Fox.tex}.
No unproved stability under a GT action is used.

The pure character actions need not commute. In pentagon order
$34,51,23,45,12$, their pair groups have orders $8,8,32,2,32$.
The three nonharmonic groups are abelian; both harmonic groups
have lower central and derived orders $32,2,1$.
The full-chain $X_i$ have exact orders $12,12,12,24,24$.
The four cusp words in the fixed basis are
\[
 q_1q_4^{-1}q_2^{-1},\quad q_1q_3q_5q_1^{-1},\quad
 q_2q_4q_5^{-1}q_2^{-1},\quad
 q_1q_5^{-1}q_3^{-1}q_2q_1^{-1}q_5q_1^{-1}.
\]
Their Fox cycles give the equivariant boundary/proper filtration.
Its total dimensions are $1023$ and $3074$, adding to $4097$.
Certificates are \texttt{check\_ferrari\_d4\_fox.py},
\path{check_d4_seed_normal_closure.py}, and the literal input
\path{agent_lifting/records/f5_nielsen_fox.json}.

\begin{statement}{A.160}{Refuted}{the fixed cyclic seed has nilpotent normal image at the fourth-power stage}
\end{statement}
The seed $S$ of (A.147) does \emph{not} have nilpotent normal image
in (A.159). This concerns that exact controlled source word;
it neither excludes another R2 seed nor settles the main question.
For $x=X_1$, $y=X_2$, set
\[
              u=S(x,y)^3,\qquad v=(yx)u(yx)^{-1}.
 \tag{FP.55.1}
\]
All peripheral orders divide $24$, so $e_3$ is represented by
$16$ in this literal evaluation. The seed has order six; $u,v$
are involutions fixing all character indices. At
$\chi=(1,0,1,1,0)$ the cusp exponents are $(0,1,1,2)$.
The invariant boundary/cycle/chain filtration has dimensions
$1,4,5$. In its proper three-dimensional quotient the matrices
of $u,v$ over $\mathbb F_3[i]$, $i^2=-1$, are
\[
 U=\begin{pmatrix}-1&1&i\\0&-1+i&-1\\0&-1+i&1-i\end{pmatrix},\qquad
 V=\begin{pmatrix}1-i&-i&-1-i\\i&-i&i\\-1-i&-i&1-i\end{pmatrix}.
 \tag{FP.55.2}
\]
They satisfy $U^2=V^2=I$, and $UV$ has exact order three.
The cycle-space and full-chain products have exact order nine.
The obstruction therefore already survives a proper quotient;
it is not an artifact of the fifth chain direction.

The minimal harmonic character orbit containing this block has
sixteen indices
\[
81,86,91,92,386,391,392,397,755,756,761,766,800,805,810,815,
\]
with index $\sum_j4^{j-1}\chi_j$.
The normal closure of $S(x,y)$ contains both involutions in
(FP.55.1). In a finite nilpotent group all two-elements belong
to its Sylow two-subgroup, and their product remains a two-element.
The proper order-three image contradicts this. The product in
the whole $1024$-block target has order eighteen, giving the same
contradiction without passing to a block subquotient.

\path{check_d4_seed_obstruction.py} independently reconstructs
the same source word and saves the exact five-by-five matrices,
filtration basis, and the three order calculations.
The minimal orbit and proper blocks are in
\texttt{ferrari\_d4\_obstruction\_orbit.npz} and
\texttt{ferrari\_d4\_proper16.npz}.
The bounded conjugacy exploration in the auxiliary normal-closure
script is not used as a proof; the two explicit involutions suffice.

\subsection{The remaining primitive translation module}


\begin{statement}{A.161}{Proved}{the primitive-character translation left after the previous affine stage}
\end{statement}
Put $D_2=F_5'F_5^2=C''(C')^2$. For every actual word
$w\in D_2'$, the full joint target has the following exact
reduction. First $w\in F_5'\subset D_4$, so its deck image
in $A$ and its old affine image in $E_0$ are trivial. Moreover
\[
                    d(w)_\chi=0\qquad(\chi^2=1).
 \tag{FF.5a}
\]
There are exactly $32$ such characters, including the trivial
one. To prove (FF.5a), observe that each of them kills $D_2$.
The crossed identity (FF.1c), specialized at such a character,
therefore makes $d_\chi:D_2\to k^5$ an ordinary additive
homomorphism. It kills the closed derived subgroup $D_2'$.

There remain exactly $4^5-2^5=992$ primitive order-four
characters. On each of them write $m_\chi=d(w)_\chi$. The
actual full-chain matrix is
\[
 T_{w,\chi}=I-m_\chi\partial_\chi,\qquad
 \partial_\chi m_\chi=0,\qquad
 \operatorname{rank}(T_{w,\chi}-I)\leq1,\qquad
 (T_{w,\chi}-I)^2=0.
 \tag{FF.5b}
\]
These assertions follow directly from (FF.2f) and
$\partial d(w)=0$. The matrix determines the translation:
for any $j$ with $\chi(a_j)\ne1$,
\[
             m_\chi=
 \frac{e_j-T_{w,\chi}e_j}{\chi(a_j)-1}.
 \tag{FF.5c}
\]
Thus $w=1$ in $C/(D_4'D_4^3)$ if and only if all these
$992$ vectors vanish. No further affine factor is needed
for this restricted class of words. Since the original
chains are defined over $\mathbb F_3$, conjugate characters
satisfy $m_{\chi^3}=m_\chi^3$ coefficientwise. The $496$
conjugate pairs each supply four coordinates over
$\mathbb F_9$, hence at most $3968$ coordinates over
$\mathbb F_3$ in total.

There is also an exact structural meaning to these counts.
In $H_1(D_4,k)$, the characters of order at most two have
combined dimension $5+31\cdot4=129$. They form the
invariant subspace for
\[
                    B=D_2/D_4=2A\simeq(C_2)^5.
\]
Transfer for the covering of index $32$ identifies this
subspace with $H_1(D_2,k)$. Its complementary primitive
part has dimension $4097-129=3968$.
In fact the image of $D_2'$ in $D_4/(D_4'D_4^3)$ is
\emph{exactly} this primitive part over $\mathbb F_3$.
For an elementary verification, consider the finite extension
\[
 1\longrightarrow M_3=D_4/(D_4'D_4^3)
 \longrightarrow D_2/(D_4'D_4^3)
 \longrightarrow B\longrightarrow1.
\]
It splits because $|B|=32$ is invertible on the elementary
abelian three-group $M_3$. Explicitly, for a section with
additive factor set $\alpha(g,h)$, put
$b(g)=|B|^{-1}\sum_h\alpha(g,h)$. Summing the cocycle
identity over its final argument gives
$\alpha(g,h)=b(g)+g b(h)-b(gh)$, and replacing the section
by $(-b(g))$ times that section kills its factor set.
The split extension is therefore $M_3\rtimes B$. As $B$
is abelian, its derived subgroup is precisely $[M_3,B]$,
the sum of the nontrivial $B$-character components. Those
are exactly the primitive $A$-character components above.
The image of $D_2'$ is that derived subgroup, proving the
claim. This surjectivity concerns arbitrary words in
$D_2'$; it does not assert that every such vector arises
as the pentagon defect of a genuine gentle element.

In particular, when the actual witness retained in A.153
has $\mathcal P(f)\in D_2'$, all lower-order components
are already zero, and (FF.5b)--(FF.5c) isolate its remaining
fourth-power affine pentagon test. They do not apply to an
individual factor in a proposed seed's normal closure merely
because that factor was used to construct a pentagon.
The separately computed product of two seed-conjugate
involutions with order nine on full chains, and a size-five
unipotent Jordan block, is a factor-level normal-closure
obstruction. It is not the shear (FF.5b) of an actual
pentagon known to lie in $D_2'$. The latter must have the
stated rank-one, square-zero form on every primitive block.

\subsection{An independent simple-group certificate and local seed correction}


\begin{statement}{A.162}{Proved}{the relevant proper stabilizer image is simple}
\end{statement}
The orbit of this character under $\langle x,y\rangle$ is
\[
\begin{split}
\mathcal O=\{&81,86,91,92,386,391,392,397,\\
             &755,756,761,766,800,805,810,815\},
\end{split}
\]
where index $\sum_j4^{j-1}\chi_j$ uses the fixed cube
Schreier basis of $C'$.  Each character has chain filtration
with dimensions $1,3,1$: boundary, proper homology, and
augmentation quotient.  The four cusp phases at the
initial character are $(0,1,1,2)$, so exactly one cusp
contributes to its boundary line.

Let $H$ be the harmonic full-chain pair image and let
$G$ be the proper three-dimensional image of the
stabilizer of this character.  Exact Schreier generators
for $G$ are supplied by the source words
\[
              x^{-2}y^{-1}x^2y,
              \qquad x^{-3}yxy.
\]
Their matrices generate a group of order $6048$.
The proper image of $S$ has a conjugacy orbit of size
$63$ and normally generates this entire group.
Its derived subgroup also has order $6048$.

Simplicity is certified directly, rather than inferred
from the order or an unproved name for the group.
The enumeration has fourteen conjugacy classes, of sizes
\[
 1,56,63,63,63,378,504,504,504,672,756,756,864,864.
\]
For every nonidentity class, its normal closure has
order $6048$.  Consequently $G$ is nonabelian simple.
The exact script
\texttt{check\_d4\_proper\_stabilizer.py} checks these
claims and records both the source Schreier words and
all class-normal-closure orders.

\begin{statement}{A.163}{Proved}{the local obstruction can be removed by a new actual core seed}
\end{statement}
Let
\[
            \mathcal N=\overline{\langle\!\langle S
                         \rangle\!\rangle}_{\widehat B}
                         \subseteq\widehat F.
\]
This is the closed $B$-normal closure of the one genuine
profinite word $S$.  It is contained in $F''$, in every
original full orbifold core, and in every earlier
$B$-normal seed condition retained in A.153.  Moreover
\[
                     q_{12}(\mathcal N)=\langle[5I]\rangle,
\]
because $q_{12}(S)=[5I]$ is central.

This distinction between $G$ and the image of
$\mathcal N$ is essential.  Each of the six peripheral
values of $S$ fixes all characters, its boundary line,
and its augmentation quotient on the whole orbit
$\mathcal O$.  These statements are explicitly
verified in the certificate.  Therefore the entire
image of $\mathcal N$ has these properties: the six
outer representatives are included, and further
$F$-conjugation preserves the corresponding kernels.
In particular the proper block at a fixed character
is a genuine homomorphism on $\mathcal N$.

For the initial character its image is normal in $G$,
because $\mathcal N$ is normal in $F$ and all the
character-stabilizer elements are represented by words
in $F$.  It contains the nontrivial proper image of $S$;
by simplicity its image is all of $G$.  Transitivity
of $H$ on $\mathcal O$ gives the same assertion in
each of the sixteen transported proper groups.
Precomposing with any of the six peripheral source
automorphisms leaves $\mathcal N$ unchanged, so the
assertion remains valid for every required source core.

Let $P$ be the simultaneous proper image of
$\mathcal N$ in these sixteen characters and six
peripheral evaluations.  It is a subdirect product
of nonabelian simple groups isomorphic to $G$.
Such a subgroup is itself a product of diagonal
simple factors, hence is perfect.  One elementary
proof proceeds inductively by Goursat's lemma:
adjoining a simple factor either adds it independently
or identifies it with one existing simple factor.
A surjection from a product of nonabelian simple
groups onto a simple group factors through one
factor, since the other normal images commute.

It follows that the kernel of
$\mathcal N\to P$ still maps onto
$\langle[5I]\rangle$.  Otherwise the nontrivial
map to this cyclic group of order two would factor
through the perfect group $P$, which is impossible.
Thus there exists an actual controlled seed
\[
 A\in\mathcal N,\qquad P(A)=1,
 \qquad q_{12}(A)=[5I].
\]
This is an actual source-element existence statement.
For a concrete algebraic construction, express the
finite value $P(S)$ as a finite product of commutators
$\prod_i[P(u_i),P(v_i)]$, with $u_i,v_i\in\mathcal N$,
and set
\[
                A=S\left(\prod_i[u_i,v_i]\right)^{-1}.
\]
Every factor is an actual element of the normal closure
of $S$.  Its finite signed value is the claimed one
because the image of $\mathcal N$ under $q_{12}$ is
abelian.  No separately prescribed factor values are
asserted to lift to a gentle element.

The full-chain image of $A$ on this orbit is now
identity on all three graded pieces.  Hence it lies
in the normal unipotent three-group of matrices
\[
 \begin{pmatrix}
  1&*&*\\0&I_3&*\\0&0&1
 \end{pmatrix},
\]
with the actual character permutations also trivial.
Take the $B$-normal core of the inverse image of this
one common normal three-group, retaining all six
source representatives.  It contains $A$, so the
same coprime frame lift and R2 normalization used in
A.94 apply.  Compactness retains all earlier A.153
conditions because $A\in\mathcal N$ belongs to every
earlier normal seed condition containing $S$.
The resulting genuine gentle witness can therefore
be chosen to have $\Delta=1$ and trivial harmonic
$f$-evaluation on this entire full-chain orbit and
its required source cores, while retaining the prior
actual-source results.

\emph{Scope.}  This proves that the displayed
fourth-power orbit does not by itself obstruct the
controlled R2 method.  The orbit is a harmonic-pair
orbit, not the whole $L$-character space.  The result
does not assert that the other fourth-power proper
components have simple images, that one replacement
seed kills all of them, or that the full fourth-power
affine pentagon or the full GT pentagon holds.

\subsection{A simultaneous corrected seed for all primitive fourth-power characters}


\begin{statement}{A.164}{Proved}{all active primitive proper images are simple, and the source seed can be corrected}
\end{statement}
Let $N=\langle\!\langle S\rangle\!\rangle_{\widehat B}$ be the
closed normal subgroup generated by the same actual source seed as
in (A.147). In each harmonic evaluation, each of its six peripheral
seed markings has precisely forty-eight nontrivial primitive proper
blocks, all of dimension three, together with the old proper
two-dimensional block at the trivial character. The complete
transported primitive support consists of forty harmonic character
orbits: thirty-two of size eight and eight of size sixteen, totaling
$384$ characters. No other primitive proper dimension is active.

For each of these forty orbits, the local proper image of $N$ has
order $6048$ and is nonabelian simple. This is an exact assertion
about the simultaneous-source normal subgroup, not merely the
normal closure of one independently chosen matrix. The computation
uses all six peripheral images, transports every character to one
orbit representative, and closes their image under the exact
Schreier stabilizer generators. The stabilizer groups have orders
$6048$, $12096$, or $24192$; in every case the indicated normal
image has order $6048$.

Simplicity is independently certified rather than inferred from
order or perfectness. The reference group has fourteen conjugacy
classes; the normal closure of every nonidentity class is the
whole group. For each of the forty normal images, the certificate
finds its invariant nondegenerate Hermitian form over $\mathbb F_9$
and an explicit orthonormal basis. In these bases every generator
lies in the same reference group. Equal orders give equality,
so all forty groups are explicitly conjugate to the certified
centreless simple group.

The full-chain assertions needed to use this classification are
also checked. All thirty pair/core values fix character indices,
boundary blocks, and augmentation quotients. The three nonharmonic
pair evaluations of $S$ are one. Moreover, conjugation by $abc$
sends the harmonic pair $(X_1,X_2)$ exactly to $(X_2,X_3)$.
Thus the preceding classification covers both harmonic pairs and
every source core.

Let $E$ be the simultaneous image of $N$ on all primitive
five-dimensional chain blocks of all thirty marked maps: the five
evaluation pairs and the six $\widehat B/F$ peripheral markings.
This marking enlargement is essential for retaining a source-normal
kernel. Because $N$ is $\widehat B$-normal, each marked map has
the same image on $N$ as its unmarked map, so the preceding simple
image classification applies to every projection.
Its kernel upon passing to the proper blocks is a finite three-group:
it is a subgroup of the group preserving the boundary/cycle/chain
filtrations and inducing the identity on all graded pieces.
Its proper image is a subdirect product of centreless nonabelian
simple groups, hence a direct product of diagonal simple factors
and in particular perfect. Therefore $E^{\rm ab}$ is a three-group,
and $E$ has no quotient of order two.

On the other hand $q_{12}(N)=\langle5I\rangle\simeq C_2$.
Goursat's lemma applied to the actual joint image
\[
                      N\longrightarrow E\times\langle5I\rangle
\]
now gives the full product: a proper subdirect image would furnish
a common nontrivial quotient, necessarily $C_2$, which $E$ lacks.
Consequently there is an actual source element $A\in N$ such that
\[
 q_{12}(A)=5I,\qquad
 A\text{ has identity image on every primitive full-chain block}.
 \tag{FP.56.1}
\]
This correction remains in $F''$ and in every original normal source
core containing $S$. It also retains every previous nilpotent-image
condition for $N$. If a literal finite expression is desired, it
can be chosen from the dense algebraic normal closure of $S$,
since (FP.56.1) imposes only a finite quotient condition.
No lift of an arbitrary finite $f$ assignment is being inferred.
The accepted R2 construction may therefore be applied with this
corrected source element; the obstruction in (A.160) rules out the
old seed itself, but does not rule out the fourth-power target.

Complete classifications, Hermitian conjugacies, and the exact
source-core checks are respectively
\path{classify_d4_proper_seed_images.py},
\path{check_d4_proper_simple_conjugacies.py}, and
\path{check_d4_global_correction_inputs.py}, with their
JSON certificates and captured outputs. The reference simplicity
calculation is independently duplicated in the lifting audit's
\path{check_d4_proper_stabilizer.py}.

\subsection{One actual witness and source automorphism for the full fourth-power abelian family}


\begin{statement}{A.165}{Proved}{uniform odd-coefficient faithfulness and the exact older-cover factor}
\end{statement}
Let $\ell$ be any odd prime and $n\geq1$. Choose the unramified
coefficient ring $\mathcal O_\ell$ containing all fourth roots
of unity: one can take $\mathbb Z_\ell$ when $\ell\equiv1$
modulo four, and $\mathbb Z_\ell[i]$ when $\ell\equiv3$
modulo four. Put $R=\mathcal O_\ell/\ell^n$ and
\[
 N_{\ell,n}=D_4'D_4^{\ell^n},\qquad
 E_{0,\ell,n}=C/(F_5'F_5^{\ell^n}).
\]
The full Fox representation over $R$, together with the old
affine conjugation action on $E_{0,\ell,n}$, satisfies
\[
 C\cap\ker\rho_{\ell,n}=N_{\ell,n},\qquad
 Z\bigl(C/N_{\ell,n}\bigr)=1.
 \tag{FF.6a}
\]
This includes prime powers, not just residue fields.

For faithfulness, repeat the path construction first over
$\mathbb Z/\ell^n$ and then extend scalars to $R$. The
extension is injective and faithfully flat. Fourier
idempotents exist because $1024$ is a unit. For a nontrivial
character, at least one entry of $\partial_\chi$ is a unit:
distinct fourth roots remain distinct modulo every odd
prime. Thus $\ker\partial_\chi$ is free of rank four,
and $m\partial_\chi=0$ implies $m=0$ over $R$ as well.
The old cube representation has the characters
$110,101,011,111,111$ over $\mathbb Z[1/2]$; this follows
directly by decomposing its signed edge modules with the
eight character idempotents. It has no invariants and its
$(C_2)^3$-action is faithful at every odd modulus. Hence
$E_{0,\ell,n}$ is centreless by the argument in A.157.
Every step of that proof consequently applies and proves
the kernel equality in (FF.6a).

For the centre assertion, a central element first maps to
one in $E_{0,\ell,n}$, so it has a representative in $F_5$.
Its action on the nontrivial Fourier cycle blocks is the
scalar $\chi(\bar c)$, and centrality forces all these
scalars to be one. Its deck coordinate in $A$ is therefore
one, so the element belongs to
$M_{\ell,n}=D_4/N_{\ell,n}$. An invariant of this module
under $F_5/D_4=A$ belongs to the trivial Fourier component.
Transfer identifies that component with
$H_1(F_5,\mathbb Z/\ell^n)$, equivariantly for $C/F_5$.
The old cube representation has no $C/F_5$-invariants,
so $M_{\ell,n}^{C/D_4}=0$. This proves centrelessness.
Scalar extension was faithful, so the conclusion holds in
the actual group over $\mathbb Z/\ell^n$.

There is an exact comparison for the $32$ characters of
order at most two. Write $A_2=F_5/D_2\simeq(C_2)^5$.
The group-ring map $R[A]\to R[A_2]$ sends the full
marked Fox Jacobians for the fourth-power cover to the
full Jacobians for the second-power cover. Under Fourier
transform this is precisely the direct sum of the
order-at-most-two blocks, including their augmentation
coordinates. Its based path map has kernel
\[
                         D_2'D_2^{\ell^n}.
 \tag{FF.6b}
\]
Consequently, if a literal source automorphism acts as the
identity on $C/(D_2'D_2^{\ell^n})$, it acts as the identity
on every one of these full chain blocks at that modulus.
Indeed for each free generator $g_j$, its image differs
from $g_j$ by an element of $D_2'D_2^{\ell^n}$,
which has zero path chain modulo (FF.6b); its deck action
on $A_2$ is also the identity. Both its Jacobian and its
semilinear coefficient action are therefore the identity.
In particular the low-order full-chain representation of
$L$ factors through its actual conjugation action on
$C/(D_2'D_2^{\ell^n})$, and hence through the faithful
older affine action used in A.153. The old affine factor
$E_{0,\ell,n}$ also factors
through this same action, since
$D_2'D_2^{\ell^n}\subset F_5'F_5^{\ell^n}$.
This is an exact comparison of based paths and source
kernels, not merely a transfer isomorphism of cycles.
It is therefore applicable to the normal nilpotent source
conditions retained in A.153, as well as to the pointwise
kernel-fixing conclusion of A.154.

\begin{statement}{A.166}{Proved}{one actual witness passes the full fourth-power abelian family}
\end{statement}
There is one genuine gentle multiplier-one element $k$ such
that
\[
 \boxed{\quad \Delta(k)=1,\qquad D(k)=1,\qquad
       \mathcal P(f_k)=1\ \text{in}\ L/D_4'.\quad}
 \tag{FF.7a}
\]
It simultaneously retains all the earlier conclusions of
A.94, A.114, A.128--A.129, and A.153--A.154. The new quotient
retains $D_4/D_4'\simeq\widehat{\mathbb Z}^{4097}$ in its
entirety. It does not replace this abelianization by the
nonabelian pro-three completion of $D_4$.

We give the source and compactness argument, including the
exact finite inputs. Use the same actual three-supported
cyclic seed $S$ as A.153, and let
\[
         \mathcal N_0=\overline{\langle\!\langle S
                              \rangle\!\rangle}_{\widehat B}.
\]
It lies in the original orbifold cores, in $\widehat F_2''$,
and in every previously retained normal source condition
that contains $S$. Its signed image is
$q_{12}(\mathcal N_0)=\langle[5I]\rangle\simeq C_2$.
The global fourth-power certificate of A.164, with the
independent simple-group proof in A.162, includes all five
pentagon pairs and all six peripheral markings, hence all
thirty source maps. It checks that the images of $S$ fix
character indices, act identically on boundary and
augmentation, and have nontrivial primitive proper support
only in $40$ character orbits for one harmonic pair: $32$
of size eight and eight of size sixteen. The other harmonic
pair is its explicitly checked $abc$-conjugate; both pairs
and all six markings are included in the simultaneous target.
Every active proper normal image
is a certified centreless nonabelian simple group of order
$6048$. The reference group is proved simple by enumeration
of its fourteen conjugacy classes and their normal closures;
all forty images are explicitly conjugated to that reference.
These precise inputs are saved in
\path{agent_five/d4_global_correction_inputs.json},
\path{agent_five/d4_proper_simple_conjugacies.json}, and
\path{agent_lifting/records/d4_proper_stabilizer.json}.
Perfection or group order alone is not substituted for
simplicity in this argument.

Let $E$ be the simultaneous image of $\mathcal N_0$ on
\emph{all} primitive full chain blocks, under these thirty
source maps. Since $\mathcal N_0$ is $B$-normal, the six
markings do not enlarge its possible source values; they
are nevertheless retained explicitly to make all subsequent
kernel conditions $B$-normal. The simultaneous proper image
is a subdirect product of centreless nonabelian simple groups,
so it is a product of diagonal simple factors and is perfect.
To justify this last implication, apply Goursat inductively:
a new simple factor is either independent or identified with
one existing simple factor. A quotient of a product of
nonabelian simple groups onto a simple group factors through
one factor, since the different normal images commute.

The kernel from $E$ to this simultaneous proper image is a
three-group. Every element in the kernel fixes the character
indices and acts identically on all three graded terms of
the boundary/cycle/chain filtration. In an adapted basis it
is a block upper-unitriangular matrix over $\mathbb F_9$;
a product of these finite unipotent groups is a three-group.
This also includes primitive blocks where the proper image
is trivial. It follows that $E^{\rm ab}$ is a three-group,
and hence that $E$ has no nontrivial quotient $C_2$.

Let $\tau:\mathcal N_0\to E$ denote this actual simultaneous
source homomorphism. Its kernel still maps onto
$\langle[5I]\rangle$ under $q_{12}$. Otherwise the nontrivial
quadratic quotient of $\mathcal N_0$ would factor through
$E$, which is impossible. Choose an actual element
\[
       A\in\ker\tau\subset\mathcal N_0,
                    \qquad q_{12}(A)=[5I].
 \tag{FF.7b}
\]
Thus $A$ kills the \emph{whole} primitive full-chain image
modulo three, including its shears, and inherits all the
older normal source conditions. This is an existence
argument in the genuine compact source group; no independent
assignment of pentagon factors is being lifted.

At the prime three, take the common inverse image of the
identity in the primitive chain representation modulo three,
including its six-marking $B$-normal core. It contains $A$.
Its image at every higher three-adic level is contained in
one normal pro-three congruence kernel: successive matrix
congruence kernels are additive groups in characteristic
three. This is a statement about a source subgroup and its
normal image, not the insufficient assertion that the single
element $A$ has three-power order. The low-order blocks and
the old affine factor are quotients of the old $C/D_2'$ action
by A.165. Their common normal images are the pronilpotent
ones already established in A.153. The full new image embeds
in the product of that old pronilpotent image and the new
pro-three congruence image. A closed subgroup of a product
of pronilpotent groups is pronilpotent. At each finite
level, therefore, the old and new requirements give one
common normal nilpotent pair-image subgroup whose source
core contains $A$.

For every odd prime $\ell\ne3$, the exact Gaussian-integer
Fox calculation gives
\[
       (X_i^{m_i}-I)^2=0,\qquad
       m_i=4\ (i=1,2,3),\qquad m_i=8\ (i=4,5),
 \tag{FF.7c}
\]
including the actual character permutations. Every third
peripheral $(X_iX_j)^{-1}$ has the same property with a
power of two at most eight. These identities are checked
on all $1024$ full five-dimensional blocks over
$\mathbb Z[i]$, together with all $L$ relators, inverses,
and boundary identities. The exact input is
\path{agent_lifting/records/d4_integral_peripherals.json}
and its constructing certificate
\path{agent_lifting/code/check_d4_integral_peripherals.py}.
At modulus $\ell^n$, (FF.7c) bounds each peripheral order
by $m_i\ell^n$, since $(I+N)^{\ell^n}=I+\ell^nN$ when
$N^2=0$. The three-supported exponent $e_3$ therefore
kills every peripheral. All the values of $S$, and hence
of its $B$-normal closure $\mathcal N_0$, are trivial in
these full chain targets and all their marked cores.
In particular $A$ satisfies every such condition.

At the prime two there is no Fourier division by a power
of two. The group
$C'/(D_4'D_4^{2^n})$ is a finite two-group, so
\[
                    O^2(C')\subset D_4'D_4^{2^n}.
 \tag{FF.7d}
\]
The retained A.128--A.129 source conditions in $L/O^2(C')$
therefore imply the required equations and factorwise
centralizations at every two-primary level.

For any finite collection of these prime-power conditions,
use $A$ in the same R2 frame lift and normalization as A.94
and A.153. Every prescribed $B$-normal source core contains
$A$, its signed value is still $[5I]$, and the required
normal pair images are nilpotent. At this finite stage,
R2 supplies the prescribed finite source approximations
and target conditions, including the fixed nontrivial
signed value. As in A.94, these are the chosen conditions
$\varphi(f)\in J$ and the required $f$ or frame membership
in the corresponding pronilpotent residuals; one does not
yet assert $D(k)=1$ or all earlier infinite constraints.
The finite target equations for the five $f$-factors are
among these prescribed conditions. They give the
finite-intersection property for the closed conditions
on the compact group of genuine gentle elements.
Compactness, with the original normalization conditions
retained in the same family, then produces one actual
$k$ with $\Delta(k)=1$, $D(k)=1$, every earlier retained
property, and all the new prime-power conditions
simultaneously.

The folded pentagon is one by the two-cycle relation, so
its defect belongs to $C$. At every odd prime the full
joint target kills the defect, and A.165 gives
$\mathcal P(f_k)\in D_4'D_4^{\ell^n}$. The prime-two
containment follows from (FF.7d). Finally
\[
             \bigcap_{\ell\ {\rm prime},\ n\geq1}
                    D_4'D_4^{\ell^n}=D_4',
 \tag{FF.7e}
\]
because $D_4/D_4'\simeq\widehat{\mathbb Z}^{4097}$ is
separated by these finite prime-power quotients. This proves
(FF.7a). The same argument proves factorwise centralization
of $C/D_4'$: at odd primes a factor has identity joint
chain action, so every one of its commutators with $C$ is
killed by the faithful restriction in (FF.6a); at two use
the retained source action, and then use (FF.7e).

\begin{statement}{A.167}{Proved}{a genuine automorphism fixing the full fourth-power Ferrari kernel}
\end{statement}
For the same actual $k$ there is a continuous automorphism
\[
 \boxed{\quad \Phi_4\in\operatorname{Aut}(L/D_4'),
        \qquad \Phi_4|_{C/D_4'}=\mathrm{id},
        \qquad\Delta(k)=1.\quad}
 \tag{FF.8a}
\]
First the profinite kernel $C/D_4'$ is centreless. Let
$Q_4=C/D_4$, of order $8192$, and let $\Lambda$ be the
integral cycle lattice of its Cayley graph on the three
involutions. This graph has $8192$ vertices and $12288$
unoriented edges, so $\operatorname{rank}\Lambda=4097$.
For $q\ne1$ in $Q_4$, the vertex action has trace zero,
and a fixed unoriented edge is reversed. The rational
homology trace is therefore
\[
         \operatorname{tr}(q\mid\Lambda\otimes\mathbb Q)
                     =1-\#\{q\text{-fixed edges}\}\leq1,
\]
strictly below $4097$. This is the homology of an explicit
finite integral graph; no rational homology of a profinite
group is being used. Thus the integral deck action is
faithful. Its reduction modulo any odd prime is faithful
too: a finite two-power-order matrix congruent to the
identity modulo that prime belongs to a pro-odd-prime
congruence kernel and hence is the identity.

The signed edge module has no $Q_4$-invariants over
$\mathbb Q$, because the edge-stabilizing involution
reverses orientation. The rational cycle lattice also
has no invariants. Consequently the integral matrix
$\sum_{q\in Q_4}q$ is zero on $\Lambda$. It follows that
\[
                 (\widehat{\mathbb Z}^{4097})^{Q_4}=0:
\]
an invariant vector would satisfy $8192v=0$, and this
module is torsion-free. A central element of $C/D_4'$
first has trivial deck coordinate by the faithful action
on $\Lambda/3\Lambda$; its remaining coordinate is an
invariant in $D_4/D_4'$, and is therefore zero.

By A.166 both braid conjugators
$g=f_k(a^2,b^2)$ and $h=f_k((ab)^3,c^2)$ centralize
$C/D_4'$. Apply the generator-and-relator argument of
A.158 with the faithful inner action of this centreless
kernel in place of its finite joint representation.
The new generators have the old conjugation actions
on the kernel and the genuine gentle action on the
folded base. Every defining relator has trivial fold
and trivial kernel action, so is one. The resulting
homomorphism fixes the kernel pointwise, kills $D_4'$
before descending, and is an automorphism by the exact
sequence, since it is the identity on the kernel and
an automorphism on the full profinite base. This proves
(FF.8a). It is an actual relative source action and makes
no assertion that $k$ satisfies the full GT pentagon.
